% 本文件由 LaTeX 排版 Agent 根据有效 translation.json 自动生成。
% author=Theodoret; work_id=2702; title=教会史

\chapter{教会史}

% structure: p0001 source=h1 target=omitted reason=page-title-already-rendered-by-parent

卷一\newline{}卷二\newline{}卷三\newline{}卷四\newline{}卷五\par

\SourceRecord{Translated by Blomfield Jackson.；Nicene and Post-Nicene Fathers, Second Series；Vol. 3.；Edited by Philip Schaff and Henry Wace.；Buffalo, NY: Christian Literature Publishing Co.,；1892.；<http://www.newadvent.org/fathers/2702.htm>.}

% structure: p0001 source=h1 target=section reason=page-title

\section{\WorkTitle{教会史（第一卷）}}

% structure: p0002 source=h2 target=subsection reason=relative-heading-rank; base=h2

\subsection{序言——本史之目的。}

当画家在画板和墙壁上描绘古代历史事件时，他们既悦人眼目，又使往事的记忆多年保持鲜活。历史学家以书籍代替画板，以生动的描述代替颜料，从而使往事的记忆更加牢固且持久，因为画家的技艺会被时间毁坏。因此，我也将尝试以文字记录迄今被忽略的教会史事件，认为若不努力而为，任凭遗忘剥夺崇高事迹和有益故事的应有声誉，实为不当。正因如此，朋友们也频频敦促我承担此作。然而，当我将自己的能力与这项事业的规模相比时，我不禁畏缩不前。但信赖一切美善之赐予者的慷慨，我甘愿承担一项超乎己力的任务。\par

巴勒斯坦的\Person{欧瑟伯}已撰写了从神圣宗徒时代到\OriginalTerm{天主所爱之君}{the prince beloved of God}君士坦丁统治时期的教会史。我将从他终止之处开始我的历史。\par

% structure: p0005 source=h2 target=subsection reason=relative-heading-rank; base=h2

\subsection{第一章 亚略异端之起源。}

在邪恶不虔的暴君马克森提、马克西米努斯和李锡尼被推翻后，那些毁灭者如飓风般激起的惊涛归于沉寂；旋风得以遏制，教会从此开始享有平静的安宁。这是为她而建立的，由一位堪受一切赞美的君主\Person{君士坦丁}所建立，他的蒙召，如同神圣宗徒一样，不是出于人，也不是藉着人，而是来自天上。他颁布法律，禁止向偶像献祭，命令建造教堂。他任命基督徒为各省总督，命令尊敬司铎，并以死刑威胁那些胆敢侮辱他们的人。有些人重建了被毁的教堂，另一些人则建造了更宽敞宏伟的新堂。因此，对我们而言，一切尽是欢乐与喜悦，而我们的敌人则被忧郁和绝望所淹没。偶像的庙宇被关闭，但在教堂中频繁举行集会，庆祝节庆。然而，充满嫉妒和邪恶的魔鬼，人类的毁灭者，无法忍受看到教会乘着顺风航行，便激起邪恶计谋的图谋，急于要颠覆由宇宙的创造者和主宰者掌舵的船只。当他开始察觉希腊人的错误已显明，魔鬼的各种诡计已被识破，大多数人崇拜造物主，而不是像从前那样崇拜受造物时，他不敢向我们天主和救主公开宣战；但找到了一些虽以基督徒之名而受尊崇，却仍受野心和虚荣奴役的人，他使他们成为执行其阴谋的合适工具，并藉着他们诱惑其他人重新陷入旧日的错误——不是通过以前设立受造物崇拜的方式，而是使万有的创造者和制造者被降低到与受造物同等的地位。我现在将叙述他在何处以及如何撒下这些\OriginalTerm{莠子}{tares}。\par

\Place{亚历山大里亚}是一座巨大而人口众多的城市，它不仅负责领导埃及，还领导周边地区，即\Place{特拜}和\Place{利比亚}。在上述不虔暴君的统治期间，信德的英勇斗士\Person{伯多禄}获得殉道者的冠冕之后，亚历山大里亚的教会由\Person{阿基拉斯}短暂管理。\Person{亚历山大}接替他，证明自己是福音教义的崇高捍卫者。那时，\Person{亚略}已被列入司铎团名册，并受委托解释圣经，当他看到大司祭职的舵柄被托付给亚历山大时，成了嫉妒攻击的猎物。被这种激情刺伤，他寻找争论和冲突的机会；虽然他看到亚历山大无可指责的行为使他无法提出任何指控，但嫉妒不让他安宁。真理之敌在他身上找到了一个工具，用以搅动和激荡教会愤怒的水流，并说服他反对亚历山大的教义。当宗主教顺从圣经，教导圣子与圣父具有同等尊荣，并与生养祂的天主同一实体时，亚略直接反对真理，断言天主子仅仅是一个受造物或被造物，并添加了那句著名\OriginalTerm{格言}{dictum}：\Quote{曾有一段时期祂不存在；}以及从他的著作中可以了解的其他意见。他坚持教导这些虚假教义，不仅在教会内，也在一般会议和集会中；他甚至挨家挨户地走，试图使人成为他错误的奴隶。亚历山大，强烈依恋宗徒的教义，起初试图用劝诫和忠告来说服他认识错误；但当他看到他疯狂妄为并公开宣告其不敬时，便将他从司铎团中罢免，因为他听见天主的法律大声宣告：\Quote{\Emph{如果你的右眼使你跌倒，就把它剜出来，丢掉它。}}\par

% structure: p0008 source=h2 target=subsection reason=relative-heading-rank; base=h2

\subsection{第二章 主要主教列表}

在这一时期，罗马教会由\Person{西尔维斯特}执掌缰绳。他教座的前任是\Person{米蒂亚德}，后者是那位在迫害期间如此卓著地显扬自己的\Person{马尔切利诺}的继承人。\par

在安提约基雅，\Person{提拉诺}去世后，当教会开始恢复和平时，\Person{维塔利斯}接受了首席权柄，并修复了被暴君们毁坏的所谓\Quote{\Place{帕拉亚}}教堂。\Person{斐洛恭尼}继任，完成了修复工作中所有欠缺的部分；他在李锡尼时期已因其宗教热忱而闻名。\par

在\Person{赫尔蒙}的管理之后，耶路撒冷教会的治理被托付给\Person{玛加略}，其人如其名，他的心智以各种美德为装饰。\par

在同一时期，才华卓著的\Person{亚历山大}治理着君士坦丁堡教会。\par

正是在那时，亚历山大的主教亚历山大察觉阿里乌被权欲所奴役，正聚集那些被他亵渎教义所掳掠的人，并举行私人会议，他致信各大教会的主教，说明其异端。为使我的史书真实性无可置疑，我将在叙述中插入他写给他同名人的信函，信中清晰叙述了我所提及的所有事实。我也将附上阿里乌的信，连同其他为完整叙述所必需的信件，使它们既能证明我著作的真实性，又能使事件经过更加清晰。\par

以下是亚历山大的亚历山大写给与他同名的主教的书信。\par

% structure: p0015 source=h2 target=subsection reason=relative-heading-rank; base=h2

\subsection{第三章 亚历山大的主教亚历山大致君士坦丁堡主教亚历山大的信函}

致他最可敬且同心同德的弟兄亚历山大，亚历山大在主内问安。\par

被贪婪和野心所驱使，心怀恶意者总是图谋破坏最重要教区的福祉。他们以各种借口攻击教会的信仰；被在他们内运行的魔鬼所疯狂，他们随意背离所有虔诚，并践踏对天主审判的畏惧。我自身深受其苦，认为有必要通知你的虔敬，使你能提防他们，免得他们或他们的任何同伙胆敢进入你的教区（因为这些骗子善于欺骗），或传播虚假而似是而非的信函，意图蒙骗那已献身于纯朴而无玷信仰的人。\par

阿里乌和阿基拉斯近日结党，效法科鲁图的野心，并远超过他。科鲁图确实试图以他对他们的指控作为自己有害行为的借口。但他们见他以基督牟利，便拒绝继续服从教会；反而为自己建造洞穴，如强盗一般，如今时常在其中聚集，昼夜在那里散布诽谤，反对基督和我们。他们辱骂一切虔诚的教义，并以犹太人的方式组织帮派对抗基督，否认祂的天主性，宣称祂与其他人同等。他们挑出一切涉及\OriginalTerm{救恩的经世}{dispensation of salvation}和祂为我们而受辱的经文，试图从中收集他们自己的不虔断言，同时回避一切宣告祂永恒的神性和祂与圣父所共有的永不停息荣耀的经文。他们坚持希腊人和犹太人关于耶稣基督所持的不虔教义，因此用尽一切手段追求他们的赞赏。外人嘲笑我们的一切，他们都殷勤实践。他们每日煽动对我们的迫害和叛乱。一方面，他们在法庭上指控我们，怂恿某些被他们诱入错谬的无原则女人作为证人。另一方面，他们允许年轻女子在街头游荡，使基督教蒙受耻辱。不仅如此，他们还胆敢撕裂基督无缝的长衣，那长衣连士兵都不敢分开。\par

当这些与他们的生活方式相符的行动，以及那长期隐藏的邪恶计划终于被我们知晓后，我们一致将他们从敬拜基督天主性的教会中驱逐出去。于是他们到处奔走，联合起来反对我们，甚至向我们那些与我们同心合意的同工司铎们求助，借口寻求与他们的和平与合一，但实际上企图用花言巧语将他们中的一些人拉入他们的谬误中。他们请这些人写一封冗长的信，然后向那些被他们欺骗的人宣读内容，目的是使他们不反悔，反而因看到主教们赞同并支持他们的观点而坚持其邪恶。他们不承认那导致他们被我们驱逐的邪恶教义和行为，而是要么不加评论地传授这些教义，要么通过诡辩和伪造来继续欺骗。如此，他们用动听而卑躬屈膝的言辞掩盖其毁灭性的教义，诱捕粗心大意的人，并抓住一切机会诽谤我们的信仰。因此，许多人被引诱签署了他们的信，并与他们共融，我认为我们同工司铎们的这种举动是极应受谴责的；因为他们不仅违背了规则，甚至还助长了他们反对基督的恶魔行为。为此，亲爱的弟兄们，我立即激励自己，将某些人的不信告知你们，他们说：\Quote{曾经有一个时期，天主子并不存在；} 以及 \Quote{那先前不存在者，后来出现了；当他在某个时候出现时，他就变得像其他每一个人一样。} 他们说，天主从无中创造了万有，并将天主子也包括在有理性的和无理性的受造物之中。与这一教义一致，他们必然断言他本性上是可变的，既能行德行，也能行恶行；因此，通过他由无中被造的假设，他们推翻了圣经的见证，圣经宣告圣言的不变性以及圣言的智慧的神性，而这圣言和智慧就是基督。\Quote{我们也能像他一样成为天主的子女，} 这些可咒骂的人说，\Quote{因为经上记载：\BibleQuote{我养育了子女，使他们成长}。} 当人们向他们提出这段经文的续文，即 \BibleQuote{但他们背叛了我}，并指出这些话与救主不变的性体不合时，他们抛弃了所有敬畏，断言天主预知并预见他的儿子不会背叛他，因此他选择他优先于其他一切。他们还断言，他不是因为他本性上有什么优于其他天主子嗣而被拣选；因为，他们说，没有人生来就是天主的儿子，也没有人与天主有特殊的关系。他们声称，他之所以被拣选，是因为尽管他本性可变，但他那勤勉的品格没有堕落。好像，即使\Person{保禄}和\Person{伯多禄}作同样的努力，他们的儿子身份也不会与他的有任何不同。\par

为确立此疯狂教义，他们侮辱圣经，引用圣咏中关于基督的话：\BibleQuote{\Emph{你爱正义，恨恶不法，因此你的天主以喜乐油膏你，胜过你的同伴。}}\BibleRef{咏 45:8}如今，天主子并非从无中受造，也从未有一刻不存在，这是\Person{若望}宗徒明确教导的，他称祂为\BibleQuote{\Emph{在父怀里的独生子}}\BibleRef{若 1:18}。这位神圣导师想要表明圣父与圣子不可分离，故说\Quote{子在父怀里}。此外，同一\Person{若望}证实天主的圣言不属于从无中受造之物，因为他说\BibleQuote{\Emph{万物是藉着祂造的}}\BibleRef{若 1:3}，并且他还以下话语宣告其独立位格：\BibleQuote{\Emph{在起初已有圣言，圣言与天主同在，圣言就是天主……万物是藉着祂造的；凡受造的，没有一样不是藉着祂造的}}\BibleRef{若 1:1-3}。既然如此，万物都是藉着祂造的，那赋予万物存在者，祂自己又怎会在某一时刻不存在呢？圣言，这创造之能，决不可被定义为与受造物同性质，因为祂确实在起初就存在，万物都是藉着祂造的，并由祂从无中呼唤而存在。\Quote{那存在者}的性质必然与从无中受造之物相反且本质不同。这也表明圣父与圣子之间没有分离，甚至思想中也不能构想分离；而世界从无中受造这一事实，涉及后来其本质的新起源，万物皆由圣父藉着圣子获得了这种存在起源。最虔诚的宗徒\Person{若望}，觉察到\Quote{存在}一词用于天主的圣言时，远超受造物的理解，因此不敢谈论祂的受生或受造，也不敢用等同的措辞称呼造物主与受造物。并非天主子是未受生的，因为唯有圣父是未受生的；而是独生天主的不可言喻的位格超越福音作者甚至天使最敏锐的构思。因此，我认为那些胆敢违命探究此主题的人不应被认为是虔诚的，经上说：\BibleQuote{\Emph{不要寻求过于你难懂的事，也不要查考高于你的事}}\BibleRef{德 3:21}。因为如果许多其他远为低下的事，人心尚不能认知，因而无法达成，如\Person{保禄}所说：\BibleQuote{\Emph{眼所未见，耳所未闻，人心所未想到的，就是天主为爱祂的人所预备的}}\BibleRef{格前 2:9}，天主也对\Person{亚巴郎}说星星不能被他数算\BibleRef{创 15:5}，同样经上说：\BibleQuote{\Emph{谁能数算海边的沙粒或雨滴呢？}}\BibleRef{德 1:2}。那么，除非疯狂，谁还敢探究天主圣言的本性呢？预言之神说：\BibleQuote{\Emph{谁能述说祂的世代？}}\BibleRef{依 53:8}。因此，我们的救主仁慈地对待那些作全世界的柱石的人，想要解除他们追求此知识的负担，便告诉他们这超越他们的自然理解，唯有圣父能洞察此最神圣的奥秘：\BibleQuote{\Emph{除了圣父，没有人认识子；除了子，也没有人认识圣父}}\BibleRef{玛 11:27}。我想，圣父关于同一主题说：\Quote{我的奥秘是为我和我的人。}\par

但那种想象天主子是从无中产生、且祂的发出是在时间中的疯狂愚钝，从\InnerQuote{无}这个词就明显可见，尽管愚昧之人无法看出自己言论的愚昧。因为\InnerQuote{祂不存在}这个短语必然指向时间或某种时代间隔。倘若万物都是藉着祂造的，那么显然，每一个时代、时间、一切时间间隔以及\InnerQuote{不存在}所归属的那个\InnerQuote{那时}，都是由祂造的。说创造了所有时间、时代和季节（这些与\InnerQuote{不存在}相混淆的那一位）曾经不存在，岂不荒谬？因为断言任何受造物的原因可以后于其结果，是极端无知，且完全违背理性。他们说圣子尚未由圣父所生的那段间隔，按照他们的意见，先于天主的上智，而万物是由上智创造的。因此他们与圣经相矛盾，圣经宣告祂是\BibleQuote{\Emph{一切受造物的首生者}}\BibleRef{哥 1:15}。与此教义一致，\Person{保禄}以其惯常的大声音论及祂说：\BibleQuote{\Emph{立祂为万有继承者，并藉着祂创造了诸世界}}\BibleRef{希 1:2}；\BibleQuote{\Emph{因为万有都是藉着祂造的：无论是在天上的、地上的、可见的、不可见的，或是宝座、宰制者、统治者、掌权者，都是藉着祂并为了祂而造的；祂先万有而在}}\BibleRef{哥 1:16-17}。既然\InnerQuote{从无中}这一假设明显不敬，那么结论就是圣父永远是父亲。祂是父亲，因为有圣子持续同在，祂因此被称为父亲。圣子永远与祂同在，圣父便永远完美，毫无缺欠，因为祂不是在时间中、或任何时间间隔内、也不是从先前不存在的东西中生了祂的独生子。\par

难道说天主的智慧曾有不存在的时刻，这不是不虔敬吗？谁说：\BibleQuote{我在祂面前如同受抚养的，我日日是祂的喜乐}？或者说，天主的德能、祂的圣言、或借以认识圣子、指明圣父的任何事物，曾一度不存在，这是有缺陷的吗？断言圣父荣耀的光辉\BibleQuote{曾不存在}，也就毁灭了那光辉所从出的原光；如果曾有一个时刻天主的肖像不存在，那么显然祂所肖像的那一位也不是常存的；不但如此，因着天主位格的精确影像不存在，那永恒拥有这精确影像的那一位也被取消了。由此可见，我们救主的圣子身份与人的儿子身份毫无共同之处。正如已经表明的，祂存在的本性无法用言语表达，并在卓越性上无限超越祂所赋予存在的一切事物，同样，祂因本性分享圣父的天主性的圣子身份，与那些因祂的安排而被收为义子的人的儿子身份有着不可言喻的区别。祂本性不变、完美、自足，而人却易于变化，需要祂的帮助。天主的智慧还能有什么进步？至高的真理或天主圣言还能给自己增添什么？生命或真光还能以任何方式改善吗？假设智慧可能变得愚钝，岂不是更违背本性吗？天主的德能可以跟软弱结合吗？理性本身能被不合理性蒙蔽，或黑暗能与真光相混吗？宗徒不是说：\BibleQuote{光明与黑暗有什么相通？基督与贝里雅耳有什么和谐？}而撒罗满说：\BibleQuote{蛇在磐石上的道}是\BibleQuote{太奇妙}，人的心智无法理解，这\Quote{磐石}按照圣保禄就是基督。然而，人和天使，祂的受造物，已蒙受祂的祝福，使他们能在德行和服从祂的命令中操练自己，从而避免犯罪。正是因此，我们的主因本性是圣父的圣子，受万有朝拜；那些脱去奴役的精神，通过勇敢的行为和在德行上的进步，因着本性是天主圣子的那一位的仁慈，领受了义子的精神的人，也因义子名分成为儿子。\par

祂真实、特有、本性和特殊的圣子身份，由保禄所宣告，他论及天主说：\BibleQuote{祂连自己的儿子也没有顾惜，反而为我们众人把祂交出来}，我们并非因本性是祂的儿子。为了将祂与那些不是\BibleQuote{祂自己的}区别开，他称祂为\BibleQuote{自己的儿子}。福音中也记载：\BibleQuote{这是我的爱子，我所喜悦的}。在圣咏中，救主说：\BibleQuote{上主对我说：你是我的儿子}。通过宣告本性的圣子身份，祂表明除了祂自己以外，没有其他本性的儿子。\par

\BibleQuote{我在母腹中、在晨光之前已生了你}这些话，岂不显然表明了父亲生育的本性圣子身份，这身份不是出于行为的勤勉或道德进步的操练，而是由于本性的单一性？由此可知，圣父独生圣子的圣子身份是不会堕落的；而那些合理受造物的义子身份，他们不是因本性而是祂的儿子，仅仅由于品格的适宜和天主的恩赐，可能会堕落，如经上所载：\BibleQuote{天主的儿子们看见人的女儿，就娶她们为妻}等等。天主借依撒意亚说：\BibleQuote{我抚养儿女，将他们养大，他们竟背叛了我}。\par

亲爱的，我还有许多话要说，但因恐怕进一步劝诫与我同心同德的教师会使你们厌倦，我就略过了。你们既蒙天主教训，并非不知道那刚刚兴起的与教会信仰相违的教训，正是由\Person{厄彼雍}和\Person{阿尔特玛}所传播的，并与安提约基雅主教\Person{撒摩撒达的保禄}的教训相同，后者被全体主教会议绝罚。他的继任者\Person{路济安}，在多年期间退出了与这些主教们的共融。\par

如今在我们当中，出现了\Quote{从无中}的人，他们贪婪地吸吮了这种不虔敬的残渣，同一根苗的旁枝：我指的是\Person{亚略}、\Person{阿基拉斯}和他们整个恶徒团伙。叙利亚的三位主教，不知如何被任命，竟赞同他们，煽动他们造成更致命的狂热。我将他们的判决提交你们的决定。他们牢记一切关于受难、卑辱、空虚自己以及所谓的贫穷，以及救主为我们所取的一切外来名称的篇章，便提出那些章节来否定祂的永恒存在和天主性，却忘记了所有那些宣告祂的光荣、尊贵和与圣父同在的篇章，例如：\BibleQuote{我与父原是一体}。主在这些话中并没有宣称自己是父，也没有将两个本性视为一体；而是说圣父的圣子的本质精确地保存了圣父的相似，祂的本性在一切事上印刻了与圣父相似的印记，是圣父的准确肖像和原型的精确印记。因此，当斐理伯渴望看见圣父，对祂说：\BibleQuote{主，请将父显示给我们}，主极其清楚地对他说：\BibleQuote{看见了我，就是看见了父}，好像圣父在祂肖像那无瑕而生动的镜子中被看见。同样的概念在圣咏中表达，圣人们说：\BibleQuote{在祢的光明中，我们必得见光明}。正是因此，\BibleQuote{凡尊敬圣子的，就是尊敬父}。这是对的，因为人对圣子所说的一切不虔敬的话，也是针对圣父说的。\par

此后，我亲爱的弟兄们，无人能惊讶于我即将详述的那些由他们对我及对我们至虔敬的子民所散布的虚假诽谤。他们不仅向基督的天主性摆阵进攻，且忘恩负义地侮辱我们。他们自认为不屑于与任何古代人物相比较，也不肯与我们从幼年起就熟悉的教师们平起平坐。他们不肯承认任何地方的同工们具有哪怕中等的才智。他们声称唯有他们自己才是智慧者和贫穷者，是教义的发现者，唯有他们才领受了那些从未进入过世上任何其他人头脑的真理。哦，何等邪恶的傲慢！哦，何等极度的愚妄！何等虚妄的夸口，加上疯狂和撒殚般的骄傲，硬化了他们不敬的心！他们不羞愧于反对古圣经的虔敬明净，众同工关于基督的一致虔敬也未能削弱他们的胆量。连魔鬼也不会容忍如此不敬，因为它们也避免说出反对天主子的话。\par

那么，这就是我必须凭我所拥有的能力向那些以粗劣资源向基督扬尘并试图诽谤我们对祂的尊崇者提出的问题。这些编造愚蠢故事的人声称，我们这些拒绝他们关于基督自无有受造的不敬且不合圣经的亵渎的人，教导存在两位无始者。因为这些未经良好教导的人争辩说，必须坚持二者之一：要么相信祂出自无有，要么存在两位无始者。由于他们的无知和缺乏神学实践，他们未能意识到非受造的圣父与祂从无有创造出的理性或非理性受造物之间有多么巨大的距离；而祂——天主圣言——的独生子性体，圣父藉祂从无有创造了宇宙，正如处在二者之间，是由自有永有的圣父所生，正如主亲自作证说：\Quote{凡爱圣父的，也爱由祂所生的圣子。}\par

我们相信，正如教会所教导的，有一位独一的无始圣父，祂的存在没有原因，不变不改，始终处于同一存在状态，既无增进也无减少；祂颁赐了法律、先知和福音；是圣祖、宗徒和诸圣的主；并相信一位主耶稣基督，天主的独生子，不是从无有而生，而是从存在的圣父所生；但不像物质身体那样，藉分离或流溢，如\Person{撒伯里乌}和\Person{瓦伦廷}所教导的；而是以一种不可言说、无法解释的方式，正如我们上面引用的那句话：\Quote{谁能述说祂的世代？}因为没有任何有死的心智能够理解祂位格的性体，正如圣父本身不能被理解，因为理性受造物的本性无法把握祂由圣父所生的方式。\par

但那些被真理之神引导的人无需从我这里学习这些事，因为救主许久以前所说的话语仍然在我们耳中回响：\Quote{除了圣子，没有人认识圣父；除了圣父，也没有人认识圣子。}我们已得知圣子是不变不改、自足完美、如同圣父，只缺少祂的\Quote{无始}。祂是祂圣父精确且完全相似的肖像。因为显然，肖像完全包含了使更大相似性存在的每一样东西，正如主教导我们说：\Quote{我的父比我大。}与此一致，我们相信圣子始终由圣父存在；因为祂是\Emph{祂荣耀的光辉}，是\Emph{祂圣父本体的真像}。但不要让任何人因\Quote{始终}一词而以为圣子是无始的，如同某些心思被蒙蔽者所想的那样；因为说祂是、祂始终存在、且在万世之前，并不是说祂是无始的。\par

人的心智不可能发明一个术语来表达\Quote{无始}的含义。我相信你持此意见；而且我确实确信你的正统观点，即这些术语中没有任何一个以任何方式表示无始。因为所有术语似乎仅仅表示时间的延伸，并不足以表达独生子的天主性及，可谓原始的存在。这些术语被那些竭力阐明奥秘的圣人们使用，他们请求听者原谅，并为他们失败给出合理借口，说\Quote{就我们的领悟所及}。但如果那些声称\Quote{所知道的有限}已为他们\Quote{归于无有}的人期望从人口中得出超越人能力的东西，那么显然，\Quote{是}、\Quote{始终}和\Quote{在万世之前}这些术语远未达到这一期望。但无论它们意味着什么，都与\Quote{无始}不同。因此，圣父作为无始者必须保留祂独特的尊威，无人被称为祂存在的原因；同样，圣子也必须受到应得的尊荣，祂由圣父的生出没有开始；我们必须向祂献上敬拜，正如我们已经说过的，只以虔敬和宗教的方式将\Quote{是}、\Quote{始终}和\Quote{在万世之前}归于祂；然而不拒绝祂的天主性，而将祂与圣父在万事上完全相似归给祂，同时将圣父独有的\Quote{无始}荣耀归于圣父，正如救主亲自说：\Quote{我的父比我大。}\par

此外，除了这关于圣父和圣子的虔敬信仰，我们按圣经的教导，宣认一位圣神，祂感动了旧约的圣人们，以及那被称为新约的众神圣导师。我们相信唯一至公教会，即使全世界合谋攻击她，也不能毁灭她；她战胜了所有异端者的不虔敬攻击；因为我们因她主人的话而勇敢，祂说：\BibleQuote{\Emph{放心，我已战胜了世界}}。此后，我们接受死人复活的道理，我们的主耶稣基督成为这复活的初果；祂带着一个真实而非貌似身体，出自天主之母玛利亚；在时期圆满时寄居在人类中间，为罪过的赦免：祂被钉死，然而祂的天主性却毫无减损。祂从死者中复活，被接升天，坐在高天至尊的右边。\par

在这封信中，我只部分提到这些事，认为正如我所说的，详细讨论每一条会令人厌烦，因为它们已为你们的虔诚勤勉所熟知。我们所教导的就是这些，所宣讲的也是这些；这些就是宗徒教会的信理，我们为此准备赴死，毫不在意那些强迫我们背弃这些信理的人；因为我们绝不会放弃对这些信理的希望，即使他们试图用酷刑来强迫我们。\par

阿黎乌和阿基拉斯及其同伙已被逐出教会，因为他们已经背离了我们虔敬的教义：按照蒙福的保禄所说，他说：\BibleQuote{\Emph{若有人向你们宣讲与你们所领受的不同福音，即使他假装是从天上来的天使，也当受诅咒}}以及\BibleQuote{\Emph{若有人教导不同，并不同意我们主耶稣基督的健全话语，以及合乎虔敬的道理，他是骄傲的，一无所知}}，等等。因此，既然他们已被弟兄们谴责，你们任何人都不要接待他们，也不要理会他们所说或所写的。他们是骗子，散布谎言，从不持守真理。他们到各城去，无非是以友谊为借口，以和平的名义送信，并借虚伪和谄媚换取回信，来欺骗\BibleQuote{\Emph{那被罪过缠累的愚笨妇女}}。我恳求你们，亲爱的弟兄们，要避开那些胆敢如此反对基督的人，他们公开嘲笑基督宗教，急切地想在全然和平的时期在审判台前炫耀，企图激发对我们的迫害，并削弱了基督生成的不可言喻的奥秘。你们要一致联合起来反对他们，正如我们的一些同工已经做的，他们义愤填膺，写信给我反对他们，并签署了我们的信条。\par

我已将这些信件通过我的儿子、执事阿皮约送给你；这些信件出自全埃及和底比斯的牧者们，也出自利比亚、彭塔波利、叙利亚、吕基亚、旁非利亚、亚细亚、卡帕多西亚以及其他邻近地区的牧者们。我相信你们也会效法他们的榜样。我曾多次善意地尝试挽回那些被引入歧途的人，但没有什么方法比我们同工之间联合的显现更能有效地恢复被他们欺骗的平信徒并引导他们悔改。你们要彼此问安，也与你们那里的弟兄们一同问安。我祈祷你们在主内刚强，我亲爱的弟兄们，并愿我从你们对基督的爱中获得果实。\par

\Quote{以下是已被弃绝为异端者的名字：在司铎中：阿黎乌；在执事中：阿基拉斯、欧佐依、埃塔勒斯、路齐乌斯、萨尔马特斯、尤利乌斯、梅纳斯、另一阿黎乌和赫拉狄乌斯。}\par

亚历山大以同样的语气写信给安提约基雅主教菲洛哥尼乌、当时领导贝罗亚教会的欧斯塔提乌，以及所有捍卫宗徒道理的人。但阿黎乌不能忍受保持沉默，而是写信给所有他认为意见相同的人。他给尼苛米底亚主教欧瑟伯的信清楚证明，神圣的亚历山大关于他没有写任何虚假之事。我将在下面插入他的信，以便使那些牵连在其不虔敬中的人的名字广为知晓。\par

% structure: p0038 source=h2 target=subsection reason=relative-heading-rank; base=h2

\subsection{第四章 阿黎乌致尼苛米底亚主教欧瑟伯的信}

致他非常亲爱的主人，天主的人，忠信而正统的欧瑟伯，阿黎乌——被亚历山大主教不义迫害，因着那你们也捍卫的无往不胜的真理——在主内致候。\par

\Quote{\Person{阿摩尼乌斯}，我父亲，即将出发前往\Place{尼科美底亚}，我认为必须藉他向你致意，并告知你因天主和祂的\Person{基督}而对弟兄们怀有的自然亲情，即这位主教大大地耗损并迫害我们，不遗余力地攻击我们。他将我们作为无神论者逐出城，因为我们不同意他所公开宣讲的：天主常存，圣子常存；正如父，子也是如此；子与天主同自存，非受生；祂是永远的；天主在思想或任何间隔上都不先于子；常存的天主，常存的子；祂从非受生者受生；子出自天主本身。\Place{凯撒勒雅}的主教\Person{欧瑟伯}、你的兄弟，\Person{戴奥多都斯}、\Person{保利诺}、\Person{亚大纳削}、\Person{额我略}、\Person{艾提乌斯}，以及所有东方的主教，都被定罪，因为他们说天主的存在先于祂的圣子；除了\Person{斐洛格尼乌斯}、\Person{赫拉尼库斯}和\Person{玛加略}，他们是无学问之人，并信奉了异端意见。他们中有的说子是涌出，有的说子是产生，有的说子也是非受生的。这些是不敬之谈，即使异端者威吓我们一千次死亡，我们也不能听。但我们说、相信、并已教导且仍在教导：子不是非受生的，也绝不是非受生者的任何部分；祂并非由任何质料得到自存；而是藉祂自己的旨意和谋略，在时间之前、万世之前，作为完美天主、独生且不变者而自存，并且在祂被生、或受造、或被预定、或被建立之前，祂并不存在。因为祂不是非受生的。我们被迫害，因为我们说子有开始，而天主没有开始。这就是我们被迫害的原因，也同样因为，我们说祂出自于无有。我们这样说，是因为祂既不是天主的任何部分，也不是任何本质存在。为此我们被迫害；其余的你都知道。我在主内向你告别，记念我们的苦难，我的\OriginalTerm{路西安派}{Lucianist}同道，真诚的\Person{欧瑟伯}。}\par

在这封信中被提名的那些人，\Person{欧瑟伯}是\Place{凯撒勒雅}的主教，\Person{戴奥多都斯}是\Place{劳狄刻雅}的主教，\Person{保利诺}是\Place{提洛}的主教，\Person{亚大纳削}是\Place{阿纳匝尔布}的主教，\Person{额我略}是\Place{贝里图斯}的主教，\Person{艾提乌斯}是\Place{吕达}的主教。\Place{吕达}现在称为\Place{狄奥斯波利斯}。\Person{亚略}以自己与这些人同心合意而自豪。他称\Person{斐洛格尼乌斯}（\Place{安提约基雅}主教）、\Person{赫拉尼库斯}（\Place{特里波利斯}主教）和\Person{玛加略}（\Place{耶路撒冷}主教）为对手。他散布对他们的诽谤，因为他们说圣子是永恒的、存在于万世之前，与圣父同享光荣、同一性体。\par

\Person{欧瑟伯}收到这封信后，也呕吐出自己的不敬，并写信给\Place{提洛}的首领\Person{保利诺}，内容如下。\par

% structure: p0043 source=h2 target=subsection reason=relative-heading-rank; base=h2

\subsection{第五章 \Place{尼科美底亚}主教\Person{欧瑟伯}致\Place{提洛}主教\Person{保利诺}的信}

我主\Person{保利诺}，\Person{欧瑟伯}在主内问安。\par

我主\Person{欧瑟伯}在真理事业上的热心，以及你对此事的沉默，都未能逃过我们的耳目。因此，若我们因我主\Person{欧瑟伯}的热心而欢喜，则另一方面我们为你而忧伤，因为如此之人的沉默似乎是对我们事业的挫败。因此，既然智者不应与他人意见不同并沉默对待真理，我劝你，激发你内心的智慧之灵去写作，并最终开始做对你自己和他人有益的事，特别是当你同意依照圣经写作，并追随其言语和旨意的轨迹时。\par

我们从未听说过有两个非受生者，也没有听说过一个被分为二，也没有学到或相信它曾经历任何有形体的改变；但我们断定非受生者是一个，并由祂而真实存在的那一位也是唯一的，不过并非由祂的性体所造，也丝毫不分有非受生者的本性或性体，在本性和能力上完全相异，并按完美的相似性情和能力而受造，与造物主相似。我们相信祂的起源方式不仅无法用言语表达，而且连思想也不能，不仅对人无法理解，对高于人的一切存有也同样无法理解。我们提出这些见解，并非出于自己的想象，而是从圣经中推论出来，从那里我们得知圣子是受造、被建立、并在与造物主相同的性体和相同的不变且不可言喻的本性中受生；因此主说：\BibleQuote{\Emph{祂在造化的起点创造了我；我从永世被立；山岳尚未出现之前，我已受生。}}\par

如果祂是从祂、或出自祂，作为祂的一部分，或由祂性体流出，就不能说祂是被创造或建立的；这一点，我主，你当然不是不知道。因为出自非受生者的一切，不能说被祂或另一个所创造或建立，因为它自起初就是非受生的。但如果称祂为受生者的事实提供了某种根据，使人相信祂是出于圣父的性体而存在，并且由圣父有了性体的相似，我们回答说，圣经并非只说祂是受生的，也这样称呼那些在本性上与祂完全不同的受造者。因为论到人，经上说：\BibleQuote{\Emph{我养育儿女，将他们养大，他们竟悖逆我。}}另一处又说：\BibleQuote{\Emph{你离弃了生你的天主。}}又说：\BibleQuote{\Emph{谁生出了露珠？}}这句话并不暗示露珠分有天主的本性，而是单纯说明万物都是按祂的旨意受造的。确实，没有一样东西是出于祂的性体，但一切存在都由祂的旨意而被召叫存在。祂是天主；万物都是按照祂的相似，并在祂圣言未来的相似中受造，由祂的自由意志受造。万物都藉祂由天主造成。一切都是出于天主。\par

\Quote{当您收到我的信，并按照天主赐予您的知识和恩宠加以修订后，我恳请您尽快写信给我主亚历山大。我相信，如果您写信给他，必能说服他赞同您的意见。请向主内的所有弟兄问安。我主，愿您蒙受天主的恩宠，并引导您为我们祈祷。}\par

他们就这样彼此通信，互相提供反对真理的武器。当这亵渎的教义在埃及和东方的教会中传播时，每个城市、每个村庄都因神学教义而引发争论和冲突。平民百姓旁观并评判双方的说法，有人支持一方，有人支持另一方。这些确实是值得流泪的悲剧场景。因为不像以往教会受外人和敌人攻击，如今同国的人，同住一屋、同桌共餐，不用矛而用舌头互相争斗。更可悲的是，这些互相攻击的人原是彼此互为肢体，属于同一身体。\par

% structure: p0050 source=h2 target=subsection reason=relative-heading-rank; base=h2

\subsection{第六章 尼西亚大公会议}

这位拥有最深智慧的皇帝听闻这些事后，首先试图堵塞其源头。于是他派遣一位以机智闻名的使者带着信件前往亚历山大，试图平息争端，并期望使争论者和解。但他的希望落空后，便着手召集著名的尼西亚大公会议，并承诺主教及其随行人员的驴、骡和马匹所需费用均由公家提供。当所有能够忍受旅途劳顿的人都到达尼西亚后，皇帝本人也前往那里，既希望见到众多主教，也渴望在他们中间维持团结。他立即安排慷慨供应他们的一切所需。共有三百一十八位主教聚集。罗马主教因年事已高未能出席，但他派遣了两位司铎代表他参加大公会议，并授权他们同意所议决的事项。\par

此时，许多人蒙受丰盛的神恩，许多人像圣宗徒一样，在身上带着主耶稣基督的印记。安提约基雅（米格多尼亚的一座城，叙利亚和亚述人称之为尼西比斯）的主教雅各伯使死者复活并恢复生命，还行了其他许多奇事，无需在此历史中再详细叙述，因为我在题为\WorkTitle{Philotheus}的著作中已经记述了。新凯撒利亚（位于幼发拉底河畔的要塞）的主教保禄曾遭受里基尼乌斯疯狂的迫害，因烙铁烫伤而失去了双手的功用，神经收缩而坏死。有些人被挖出右眼，有些人失去了右臂。埃及的帕弗努提乌斯也在其中。简言之，这次大公会议看起来像一支殉道者集结的军队。然而，这神圣而著名的集会也并非完全没有反对因素；有一些人，虽然少到容易数清，像危险的浅滩一样表面光鲜，实际虽未公开却在暗中支持亚略的亵渎教义。\par

当所有人聚集后，皇帝下令在宫殿中准备一间大厅供他们使用，里面摆放了足够的长凳和座位；这样安排以符合应有的尊严后，他请主教们进入并讨论所提议的议题。皇帝带着少数随从最后进入房间；他身材高大，容貌俊美令人赞叹，更令人惊叹的是他面容上谦逊的神情。在集会中间为他放置了一张矮凳，但他直到征得主教们的许可后才坐下。然后全体神圣集会围坐于他四周。首先立即站起的是伟大的安提约基雅主教欧斯塔修斯，他在前面提到的菲洛戈尼乌斯安息主怀后，被主教、司铎和热爱基督的平信徒一致推举，不得已继任。他用颂词的花环为皇帝加冕，并赞扬他在处理教会事务上所表现的勤奋关注。\par

那位卓越的皇帝接着劝勉主教们要同心合意、和睦相处；他让他们记起已故暴君的残酷，并提醒他们天主在他的治下并藉着他所赐予的可敬和平。他指出，这何等可怕，是的，非常可怕，正当他们的敌人已被消灭、无人敢反对他们之时，他们却彼此攻击，让看热闹的敌人发笑，尤其是他们正在辩论神圣之事，而关于这些事他们已有圣神所默感的文字教导。\Quote{至于福音（他继续说道），以及古代先知的著作和神谕，清楚地教导我们当如何相信有关天主性之事。因此，让一切争辩都丢弃吧；让我们在受天主默感的话语中寻求争议问题的解答。}他就像一位慈爱的儿子对父亲那样，向主教们作了这些和类似的劝勉，竭力使他们教义合一。会议的大多数成员被他的论据说服，彼此之间建立了和睦，并接受了纯正的教义。然而，有少数人反对这些教义，站在亚略一边，前面已经提到过；其中有厄弗所主教\Person{默诺番图}、史基托波利斯主教\Person{帕特罗斐禄}、尼西亚主教\Person{德奥尼}、以及尼若尼亚主教\Person{纳尔奇苏}，尼若尼亚是第二基里基雅的一个城镇，现在称为伊雷诺波利斯；还有马尔玛里加主教\Person{德奥纳}和埃及托勒迈依主教\Person{塞孔多}。他们起草了一份信仰声明，提交给会议。宣读后，它立即被撕毁，被宣布为伪造和虚假。对他们掀起的反对是如此猛烈，加诸他们背叛宗教的谴责如此之多，以致除了\Person{塞孔多}和\Person{德奥纳}外，他们全都站起来，带头公开弃绝亚略。这个不虔敬的人就这样被逐出教会后，大家一致同意起草了一份信经，至今仍被接受；一经签署，会议即告解散。\par

% structure: p0055 source=h2 target=subsection reason=relative-heading-rank; base=h2

\subsection{第七章 从\Person{欧斯塔提}及\Person{亚大纳削}的着作驳斥\OriginalTerm{亚略主义}{Arianism}}

然而，上述主教们并非真诚同意，只是表面上同意。后来他们阴谋反对那些热心宗教的首要人物，以及后者有关他们的著作，就表明了这一点。例如，已经提到的著名安提约基雅主教\Person{欧斯塔提}，在解释\BibleBook{}{箴言}中的经文\BibleQuote{\Emph{主在祂道路的起始、在祂太初的一切作为之前，就造了我}}时，撰写文章反对他们，并驳斥了他们的亵渎言论。\par

\Quote{我现在要继续叙述这些不同的事件是如何发生的。在尼西亚召开了一次大公会议，约二百七十位主教聚集。人数如此之多，我无法说出确切数字，其实我也未费大力气查明这一点。当他们开始探究信仰的本质时，欧瑟伯的声明被提出，其中毫不掩饰地显明了他的亵渎。当众宣读时，因其背离信仰，使听众深为悲伤，并使作者蒙受不可挽回的羞辱。欧瑟伯派被明确定罪后，那篇亵渎的著作当众被撕毁，他们中有些人合谋，以维持和平为借口，迫使所有最能言善辩的人沉默。亚略疯狂者害怕被如此众多的主教会议逐出教会，便急忙声明弃绝和谴责已被定罪的教义，并一致签署了信经。这样，他们通过最可耻的欺骗保住了主教席位——尽管他们本应被撤职——继续有时秘密、有时公开地支持被定罪的教义，用各种论据阴谋反对真理。他们全心全意地要建立这些稗子般的教义，躲避聪明人的审视，回避明察秋毫者，攻击虔诚的宣讲者。但我们不相信这些无信仰者能以此战胜天主。因为即使他们\BibleQuote{\Emph{束腰}}，也\BibleQuote{\Emph{终必粉碎}}，正如\Person{依撒意亚}的庄重预言。}\par

这是伟大的\Person{欧斯塔提}的话。他的战友、真理的捍卫者\Person{亚大纳削}，继承了著名的\Person{亚历山大}的主教职位，在写给非洲人的信函中补充了以下内容。\par

主教们聚集在会议中，渴望驳斥亚略派所捏造的邪恶主张——即圣子是从无有中被创造的，祂是受造物和受造存在，曾有一段时间祂不存在，并且祂在本性上可变化——并一致赞同提出以下符合圣经的声明：圣子在本性上是天主的独生子、圣言、德能和父的唯一智慧；祂正如\Person{若望}所说，是\BibleQuote{\Emph{真天主}}，也如\Person{保禄}所写，是\BibleQuote{\Emph{天主光荣的辉耀，和父本体的印记}}。但欧瑟伯的追随者被自己邪恶的教义所牵引，开始彼此说：\Quote{让我们同意，因为我们也出于天主；\BibleQuote{只有一位天主，万物都出于祂}；\BibleQuote{旧事已过，看，一切都成了新的，一切都出于天主}。他们也特别注重《\WorkTitle{黑马牧人书}》中所包含的内容：\InnerQuote{首先相信只有一位天主，祂创造了并形成了万物，并使他们从无有中成为存在。}}\par

但主教们看穿了他们的恶计和不虔敬的诡计，并对\Quote{出于天主}一词给出了更清晰的阐明，写道：圣子出于天主的本体；这样，那些丝毫不从自身获得存在的受造物，虽被称为\Quote{出于天主}，唯独圣子被称为出于圣父的本体；这是独生子、父真圣言所特有的。这就是主教们为什么写道：祂出于圣父的本体。\par

但是，当那些看来人数不多的亚略派信徒，再次被主教们询问是否承认\Quote{圣子不是受造物，而是能力、独一的智慧、圣父永恒不变的肖像，并且祂是真天主}时，欧瑟伯派信徒彼此示意，表明这些声明同样适用于我们。因为据说，我们是\BibleQuote{天主的肖像和光荣}，且\BibleQuote{我们时常活着的人}；此外，他们说，有许多能力；因为经上记载：\BibleQuote{天主的一切能力都离开了埃及地}。蝗虫和毛虫也被称为\BibleQuote{大能力}。别处又写道：\BibleQuote{万军的天主与我们同在，雅各伯的天主是我们的助佑}。此外，我们不仅仅是天主的子民，更因为圣子称我们为\BibleQuote{兄弟}。\InnerQuote{祂是真天主}的声明并不令我们为难，因为祂存在后，便是真实的。\par

\Quote{这就是亚略派腐败的见解；但对此，主教们察觉了他们在这件事上的诡诈，便从圣经中收集了那些论及基督是光荣、泉源、河流、及那实体的肖像的段落，并引用了如下经文：\InnerQuote{\BibleQuote{在你的光明中，我们才能看见光明}}；以及\InnerQuote{\BibleQuote{我与父原是一体}}。然后，他们更加清楚地简短声明：圣子与圣父同一实体；因为这确实是所引段落的意义。亚略派抱怨这些确切词语不见于圣经，但他们自己的做法却表明这抱怨是无根据的，因为他们自己的不敬断言并非取自圣经；因为经上并未写道：圣子出自虚无，且有一段时期祂不存在；然而他们却抱怨因一些虽非实际见于圣经、却符合真宗教的表述而被定罪。反之，他们自己，仿佛从粪堆中拾得言语，说出了真正属乎尘土之物。而主教们则并非自创言辞，而是从教父那里领受见证，并据此书写。事实上，早在一百三十年前，罗马大城和我们本城的古老主教，曾定罪那些断言圣子是受造物、且与圣父不同体的人。凯撒勒雅主教欧瑟伯知晓这些事实；他一度支持亚略异端，但后来在尼西亚大公会议的信经上签了字。他写信给他教区的民众，主张\OriginalTerm{同体}{consubstantial}一词是\InnerQuote{杰出的主教和博学的作家用来表达圣父和圣子天主性的术语}。}\par

这些人出于对多数人的畏惧，隐藏了自己的不健全想法，同意了会议的决定，从而招致了先知对他们的谴责，因为万有的天主向他们呼喊：\Quote{\BibleQuote{这民族用嘴唇尊敬我，他们的心却远离我。}} 然而，特奥纳斯和塞孔都斯不愿采取这种途径，被一致同意绝罚，视他们为看重亚略亵渎言语胜于福音道理的人。主教们随后回到会议，制定了二十条法规以规范教会纪律。\par

% structure: p0064 source=h2 target=subsection reason=relative-heading-rank; base=h2

\subsection{第八章  关于埃及人梅里提乌斯的事实，从他产生了梅里提乌斯分裂，此分裂存留至今。——关于他的教务会议书信}

在梅里提乌斯被祝圣为主教之后（这离亚略争议不久），他被至圣的亚历山大里亚主教伯多禄（他也获得了殉道的荣冠）判定犯有某些罪行。被伯多禄罢免后，他不服从罢免，反而在特拜德及埃及邻近地区引起骚乱和纷扰，反叛亚历山大里亚的首席权。会议写了一封信给亚历山大里亚教会，说明针对他的反叛行为所议决的事项。信的内容如下：\par

\Emph{教务会议书信。}\par

致因天主的恩宠伟大而圣洁的亚历山大里亚教会，并致在埃及、利比亚及彭塔波利斯的可爱兄弟，在伟大而圣洁的尼西亚大公会议聚集的主教们，在主内致候。\par

蒙天主的恩宠及最虔诚的皇帝君士坦丁（他从不同省份和城市召集我们），伟大的圣洁尼西亚大公会议得以召开，我们判断有必要从整个神圣会议发出一封信，也告知你们什么问题曾被提出和讨论，以及什么已被议决确立。\par

首先，在我们最虔诚的皇帝君士坦丁面前调查了亚略的不敬道理；他的不敬以及他提出的亵渎言语和观点，一致被绝罚，他声称天主圣子出自不存在，在受生之前祂不存在，有一段时期祂不存在，并且祂可以按照自己的自由意志，能够处于美德或邪恶。圣洁的会议绝罚了所有这些断言，甚至拒绝倾听如此不敬且愚蠢的意见以及如此亵渎的表述。关于他的最终决定你们已知晓或即将听到；但我们此刻不再提及，以免我们似乎践踏一个已承受其罪罚的人。他的不敬如此有影响力，以至于使马马里卡主教特奥纳斯和托勒迈斯主教塞孔都斯陷入他的毁灭，而他们也分担了他的刑罚。\par

但埃及藉天主的恩宠，从这些虚假而亵渎的意见中，以及从那些敢于在向来和平的人民中挑起不和与分裂的人手中被解救出来后，却仍留下关于梅勒提乌斯的鲁莽行为以及由他所祝圣之人的问题。亲爱的弟兄们，我们现在告知你们公会议关于此事的法令。神圣的公会议决定，对梅勒提乌斯应予以宽大处理，尽管严格来说他连最小的让步也不配。他被允许留在本城，但被剥夺了所有权力，无论是提名权还是祝圣权，也不得为这些目的出现在任何省或城市中，只保留其职位的空名。凡由他手中领受祝圣的人，必须服从一次更虔诚的重行祝圣，并被接纳进入共融，条件是他们保留职务，但在每个教区和教会中，地位低于我们备受尊敬的同事亚历山大在他们之前所祝圣的人。这样，他们就无权凭己意选择或提名他人担任圣职，实际上也无权在没有属于亚历山大的天主教会和宗徒教会主教们的同意下做任何事。而那些藉天主的恩宠并因你们的祈祷，未在分裂中被查出，并在天主教会和宗徒教会中保持无瑕的人，则有权选举和提名配得圣职的人，并被允许做一切符合法律和教会权威的事。如果教会现任职务者中有人去世，那么这些新近被接纳的人，只要他们显得配得，并得到人民的选举，且选举得到亚历山大的天主教会主教的确认和批准，就可以晋升到死者的荣誉地位。同样的特权也已授予所有其他人。但是，关于梅勒提乌斯，则因他先前的不服从以及他性格的鲁莽和急躁而作出了例外；因为如果他获得丝毫权力，他可能会滥用它，再次引起混乱。这些就是有关埃及和神圣的亚历山大教会的主要事项。至于其他制定的教规或确定的信理，你们将从我们最尊敬的同事和弟兄亚历山大那里听到，他会给你们更准确的信息，因为他亲自指导并参与了所发生的一切。\par

\Quote{我们也要告诉你们一个好消息：按照你们的祈祷，最神圣的逾越节庆的庆祝已得到一致修正，因此我们东方的弟兄们，以前没有与罗马的弟兄们以及你们和所有从起初就这样做的人同时守节，今后将与你们一同庆祝。那么，为我们事业的成功、普遍的和平与和谐、以及一切异端的根除而喜乐，并以更大的尊敬和更热烈的爱接待我们的同工和你们的主教亚历山大，他的到来给我们带来了喜乐，并且在他年事已高之际，为了恢复你们之间的和平而承受了如此多的辛劳。请为我们众人祈祷，使那些已正确决定的事能藉我们的主耶稣基督坚定不移，正如我们相信，是按照天主父在圣神内的美意而成就的，愿光荣归于祂，世世无穷。阿们。}\par

尽管那神圣的主教会议努力用这剂药方治疗梅勒提乌斯派的疾病，但他迷妄的痕迹甚至存留至今；因为有些地区的隐修士团体拒绝遵行纯正的道理，并遵守某些徒劳的纪律要点，与犹太人和撒玛黎雅人的迷妄看法一致。\par

% structure: p0073 source=h2 target=subsection reason=relative-heading-rank; base=h2

\subsection{第九章 君士坦丁皇帝关于大公会议事项写给未与会主教的信函}

伟大的皇帝也将公会议事项的叙述写信给那些未能与会的主教。我认为将这份书信收入我的著作是值得的，因为它清楚地证明了作者的虔诚。\par

\Signature{君士坦丁·奥古斯都致各教会。}\par

鉴于此刻所见的公共福祉是神圣恩宠大能的结果，我首先渴望的是，天主教会中有福的成员能持守同一信德、真诚的爱德和对全能天主的同一宗教形式。但是，既然没有比将一切关于我们至圣宗教的事提交给全体或大多数主教的审查更坚定或更有效的措施来确保这一目标，我尽可能多地召集了他们，并作为你们中的一员坐在他们中间；因为我不愿否认那真理，它是我最大喜乐的泉源，即我是你们的同仆。每一个问题都得到了应有的讨论，直到那取悦全知的天主并有利于合一的道理变得清晰，以致在信德上不应再留下分裂或争论的余地。\par

当时正在讨论最神圣的复活节的纪念事宜，大家一致决定，最好在各地都在同一日庆祝。还有什么比这个使我们获得永生希望的庆节，更公平、更合宜，以致所有人都应以同样的秩序和精确性，在明确的基础上谨慎庆祝呢？首先，我们宣布，庆祝这一神圣庆节时遵循犹太人的习俗是不适当的，因为他们的手已被罪恶玷污，这些可怜人的心灵必然蒙蔽。通过摒弃他们的习俗，我们确立并传给后代一个更为合理的习俗，这一习俗自我们主受难之日以来就一直被遵守。因此，让我们与作为我们仇敌的犹太人毫无共同之处。因为我们从救主那里领受了另一种方式。我们的神圣宗教教导了一条更好、更合法的行为准则。让我们同心合意地行走其中，我所敬重的弟兄们，小心避免与那邪恶之道有任何接触。他们夸口说，没有他们的指示，我们就无法恰当地纪念这个庆节。这是极其荒谬的。因为那些使主致死之后，丧失心智，不受健全理性支配，而是被无节制的激情，随其与生俱来的疯狂所驱使的人，又怎能对任何问题持有正确看法呢？因此，他们如此远离真理，尽可能偏离正确的修正，以致在同一年庆祝第二次逾越节。我们有什么理由跟随那些如此坦白承认不健全且深陷错误的人呢？因为我们绝不能容忍一年庆祝两次逾越节。但即使所有这些事实都不存在，你们自己的睿智也会促使你们以勤勉和祈祷警醒，免得你们纯洁的心灵似乎与如此彻底堕落的人民的习俗有份。还必须记住，在如此重要的一点，即庆祝如此神圣的一个庆节上，分歧是错误的。我们的救主已指定一日来纪念我们的得救，即他至圣的受难。他希望他的天主教会是唯一的，其成员虽然分散在世界各地，却由同一精神，即天主旨意所滋养。请你们虔诚的睿智思考，这是多么邪恶和不恰当：一些人奉献给斋戒的日子，却被他人用在欢宴上；并且在复活节之后，一些人沉浸在节日的喜庆和放松中，而另一些人则投身于规定的斋戒。这种不当必须纠正，所有这些纪念方式的差异都应统一为一种形式，这是天主上智的旨意，我相信你们都会明白。因此，必须纠正这种不规则，以便我们不再与那些弑亲者和杀害我们主的人有任何共同之处。在世界西方、南方和北方地区的所有教会，以及一些东方教会中，遵守着一种有序而卓越的纪念形式；这种形式受到普遍赞扬，我保证你们也愿意同样采纳它，从而欣然接受在\Place{罗马城}、整个\Place{意大利}、在整个非洲、\Place{埃及}、\Place{西班牙}、\Place{高卢}、\Place{不列颠}、\Place{利比亚}、\Place{希腊}、亚洲教区、\Place{本都}和\Place{西里西亚}一致通过的规则，你们不仅要考虑上述地方的教会在数量上更多，而且考虑到，所有人在那个准确推理似乎要求、并且与犹太人的伪誓毫无共同之处的做法上一致同意，这是最虔诚的。\par

\Quote{简言之，总结以上所述，所有人的判断是，神圣的复活节应在同一日举行；因为在如此神圣的事上，存在任何习俗的差异都是不合适的，而采取与错误和罪恶毫无关联的意见更好。有鉴于此，要欣然接受这天上的恩赐和明确的神圣命令；因为主教神圣会议中所行的一切都应归诸天主旨意。因此，当你们向所有我们亲爱的弟兄们传达这封信的内容后，要视自己有义务接受上述决定，并安排此圣日的定期遵守，以便根据我长久以来的愿望，当我与你们面对面相见时，能够与你们在同一日庆祝这神圣的庆节，并与你们一同喜乐，见证魔鬼的残酷被我们的努力、藉着天主的恩宠摧毁，而我们的信德、平安与和谐充满世界。愿天主保佑你们，亲爱的弟兄们。}\par

% structure: p0079 source=h2 target=subsection reason=relative-heading-rank; base=h2

\subsection{第十章 皇帝供应教会每日所需，兼述他的其他德行}

皇帝就这样写信给缺席者。对于出席大公会议的三百一十八人，他表现出极大的仁慈，以十分温和的态度对他们讲话，并赠送他们礼物。他命令预备了许多坐垫供他们休息，并在一场宴会上款待他们全体。最值得尊敬的人，他邀请到自己的餐桌，其余的人则分配到其他桌。他注意到有些人右眼被挖出，得知这种伤残是为了宗教而经受的，就把嘴唇放在伤口上，相信能从亲吻中汲取祝福。宴会结束后，他又向他们赠送了其他礼物。然后他写信给各省总督，指示在每个城市中向贞女、寡妇以及献身于神圣服务的人发放供应金；他衡量年津贴的数额更多是出于自己慷慨的冲动，而非按他们的需要。这笔金额的三分之一至今仍在分发。不敬的\Person{尤利安}扣留了全部。他的继任者赐予了现在分发的金额，当时盛行的饥荒减少了国家的资源。如果过去的养老金是现在的三倍，那么皇帝的慷慨大方由此事实便显而易见。\par

我认为不宜沉默地略过以下情况。一些好争吵的人写了针对某些主教的控告，并将诉状呈递给皇帝。此事发生在和好建立之前，他收下了这些名单，把它们卷成包并用他的戒指封印，命令妥善保管。在和好实现之后，他取出这些文件，当着他们的面烧毁，同时发誓说他没有读过其中一字。他说司铎的罪行不应让群众知晓，以免成为冒犯的机会，使他们毫无畏惧地犯罪。据报道他还补充说，如果他当场发现一位主教在犯奸淫，他会用皇袍盖住那不法之事，以免任何人目击那场景而因此受害。他如此劝诫所有司铎，同时也赐予他们荣誉，然后激励他们各自回到自己的羊群。\par

% structure: p0082 source=h2 target=subsection reason=relative-heading-rank; base=h2

\subsection{第十一章}

我在此要插入一封\Person{欧瑟伯}（\Place{凯撒勒雅}主教）关于信仰的书信，因为它描述了亚略派信徒的厚颜无耻：他们不仅藐视我们的教父，还拒绝他们自己的（教父）。这封信包含了他们疯狂的确凿证据。他们确实尊敬\Person{欧瑟伯}，因为他接受了他们的见解，但公开反驳他的著作。他写这封信给某些亚略派信徒，他们似乎在指控他背叛。这封信本身解释了作者的目的。\par

\Emph{\Person{欧瑟伯}（\Place{凯撒勒雅}主教）的书信，他是当大公会议在\Place{尼西亚}召开时从那里写下的。}\par

你们大概已从其他渠道得知了在\Place{尼西亚}普世大公会议上关于教会信仰的决定，因为伟大事件的声望通常超过对它们的精确记述；但为了避免不严格符合真相的谣言传到你们那里，我认为有必要首先寄给你们我们最初提出的信仰纲领，其次是第二个，即对我们条文增加内容后发表的。以下是我们提出的纲领，它在最虔诚的皇帝面前被宣读，并宣布以正确恰当的语言表达。\par

\Emph{我们所提出的信仰。}\par

\Quote{正如我们在最初的教理讲授中、在我们受洗时，从我们之前的诸位主教所领受的，并如我们从圣经所学，并且无论是作为长老还是作为主教，我们惯常相信和教导的；我们现在也如此相信并这样宣认我们的信仰。如下：——}\par

\Quote{我们信唯一的天主，全能的圣父，造物主，有形无形的万有；并信唯一的主耶稣基督，天主的圣言，出自天主的天主，出自光明的光明，出自生命的生活，独生子，一切受造物的首生者，在万世之前由圣父所生；万物藉祂而造；祂为了我们的救恩降生成人，并居住人间。祂受难，第三日复活，并升到圣父那里；祂将光荣地再来，审判生者与死者。我们也信唯一的圣神。}\par

\Quote{我们相信每一位的存在和持续存在；圣父是真圣父，圣子是真圣子，圣神是真圣神；正如我们的主派遣祂的门徒去宣讲福音时所说：\BibleQuote{你们要去使万民成为门徒，因父及子及圣神之名给他们授洗。}我们坚决确认我们持守这一信仰，我们一直持守它，并将持守它直至死亡，谴责所有不敬神的异端。我们在全能天主和我们的主耶稣基督面前作证，自从我们认识自己以来，我们一直是从心中、从灵魂中如此思想；我们有办法向你们证明，并说服你们，我们过去一直如此相信并宣讲。}\par

当我们提出这个纲领后，完全没有反驳的余地；我们亲爱的皇帝本人首先证明它是最正统的，并且他同意其中的意见；他劝勉其他人签署它，并接受它包含的全部道理，只增加了一个词——\OriginalTerm{同体}{consubstantial}。他解释说，这个词不暗示任何身体条件或变化，因为圣子不是通过分裂或切断从圣父获得祂的存在，因为一个非物质、理智、无形体的本性不会受到任何身体条件或变化的影响。这些事必须被理解为具有神圣而奥妙的含义。我们最智慧、最虔诚的皇帝如此推理。增加\OriginalTerm{同体}{consubstantial}一词促使了以下纲领的制定：——\par

\Emph{大公会议公布的信经。}\par

\Quote{我们信唯一的天主，全能的圣父，有形无形万物的创造者。并信唯一的主耶稣基督，天主子，圣父所生；独生子，即出自圣父的本质，天主出自天主，光明出自光明，真天主出自真天主，受生而非受造，与圣父同体；万物藉祂而造成，无论是天上的还是地上的。祂为了我们人类，并为我们的救恩，从天降下，降生成人，成为人；祂受难，第三日复活；祂升了天，将要审判生者与死者。我们也信圣神。圣而公之宗徒教会绝罚所有说\InnerQuote{圣子曾有一段时间不存在}、\InnerQuote{祂在受生之前不存在}、\InnerQuote{祂是从无中受造}、或\InnerQuote{祂与圣父有不同的本质和不同的实体}、以及\InnerQuote{祂是可变或可改变的}之人。}\par

他们制定这一信条时，我们并未不加审查地放过其中提到圣子是出自圣父的本质且\OriginalTerm{与圣父同质}{consubstantial}的那段话。由此产生了问题和争论，这些术语的含义得到了仔细的检验。于是他们被引导承认，\Quote{与圣父同质}一词意味着圣子出自圣父，但并非作为圣父的一部分。我们认为接受这一意见是恰当的；因为正确的教义教导说，圣子出自圣父，而非圣父本质的一部分。出于对和平的热爱，并且为了不偏离真信仰，我们也接受这一观点，并不拒绝\Quote{与圣父同质}这一术语。出于同样的原因，我们接受了\Quote{受生而非受造}的表达；因为他们主张，\Quote{受造}一词普遍适用于所有藉圣子受造的事物，圣子与这些事物毫无相似之处；因此，祂不是一个受造物，如同祂所造之物，而是拥有超越所有受造物的本质。圣经教导祂是由圣父所生，其生育方式对一切受造物而言是不可理解和无法言喻的。同样，\Quote{与圣父同质}一词，经过探究，也不是按照肉体关系或与可朽之物的相似性被接受的。因为还表明它并不意味本质的分裂、切割，或圣父权能的任何改变、变化或减少，所有这些都与未生圣父的本性格格不入。结论是，\Emph{与圣父同质}这一表达意味着天主子在任何方面都不像祂所造的受造物；而是唯独与生祂的圣父在各方面完全相似：因为祂不出自任何本质或实体，惟独出自圣父。鉴于对教义作出了这样的解释，我们认为赞同它是恰当的，尤其是因为我们知道，在古代，一些博学且著名的主教和作家曾就圣父与圣子的天主性使用过\Quote{与圣父同质}这一术语。\par

这些是我必须就所公布的信德信条通报的情况。我们全体同意，并非未经审查，而是在我们最亲爱的皇帝面前，将提交给我们的观点进行了彻底的审查后，出于上述原因，我们都同意了。我们也允许接受他们附加在信德信条之后的绝罚，因为它禁止使用非圣经的词语；教会的几乎一切混乱和麻烦皆由此而生。既然受默感的圣经没有使用\Quote{从无中}或\Quote{有一个时候祂不存在}的词语，也没有任何其他同类说法，那么断言或教导这些事情似乎不合理。因此，我们判断同意这一意见是妥当的；因为我们以往也从未习惯于使用这样的词语。此外，谴责\Quote{祂在受生之前不存在}的主张似乎并不矛盾，因为所有人都同意：祂按照肉体受生之前，已是天主子。在此，我们的皇帝，最蒙天主所爱，开始论及祂的神性起源和祂在万世之前的存在。祂在实际上受生之前，就已潜在地存在于圣父内，无需生育；圣父始终是圣父，正如祂始终是君王和救主，并潜在地是万有，且从未经历任何存在或行动的变化。\par

\Quote{亲爱的弟兄们，我们认为有必要将这些情况传达给你们，为向你们表明，我们在必须裁决的所有问题上进行了怎样的审查和探究；同时也为证明，我们曾如何坚定地抵抗，甚至直到最后一刻，当表达不当的教义冒犯我们时；而在另一时刻，当那些条款在彻底而诚实地探究其含义后，显得与我们已公布的信德信条中所承认的完全一致，且无任何可反对之处时，我们又无争辩地接受了它们。}\par

% structure: p0096 source=h2 target=subsection reason=relative-heading-rank; base=h2

\subsection{第十二章 驳斥当代亚略派人的亵渎，出自凯撒勒雅主教欧瑟伯的著作}

欧瑟伯清楚证明，上述术语\OriginalTerm{与圣父同质}{consubstantial}并非新造，亦非会议教父的发明；而是从一开始就由父传子。他声称当时全体一致接受了所颁布的信经；并在另一部著作中再次证实了这一点，其中他高度赞扬伟大的君士坦丁的行为。他写着如下：—\par

\Quote{皇帝用拉丁语发表了这番讲话，由翻译译成希腊语，然后他允许议会领袖自由发言。有些人立即开始对邻人提出控诉，而另一些人则采取相互指责和辱骂。每一方都有许多话要说，起初辩论变得非常激烈。皇帝耐心而专注地倾听所有提出的意见，并对各方依次提出的论点给予充分注意。他冷静地努力调和冲突各方；用希腊语温和地对他们讲话，他对这种语言并不陌生，声音甜美而温和。有些人他通过论证说服，另一些人他让他们感到羞愧；他赞扬那些发言良好的人，并激励所有人达成一致；直到最后，他在所有争议点上使他们达到思想和意见的一致，以致他们都同意持守同一信仰，并在同一天庆祝救恩的节日。所决定的事被写成文件，并由所有主教签署。}\par

随后，作者如此继续叙述：——\par

\Quote{当事情这样安排好后，皇帝允许他们返回各自的教区。他们带着极大的喜悦返回，此后一直保持同一种意见，即在皇帝面前所同意的，尽管曾经广泛分离，现在却像同一个身体联合在一起。\Person{君士坦丁}为自己努力的成功而欢欣，通过信件将这些幸福的结果告知远处的人。他下令慷慨地分发大量金钱给乡村和城市的居民，以便以公共庆典庆祝他统治的二十周年。}\par

虽然亚略派不敬地否认其他教父的陈述，但他们应当相信这位教父所写的，他们惯于钦佩这位教父。因此，他们应当接受他对所有主教签署信德声明之一致性的见证。但是，既然他们驳斥自己领袖的意见，他们应当了解\Person{亚略}死亡的极其丑恶和可怕的方式，并全力逃避他作为创始者的不敬教义。因为他的死亡方式很可能不为所有人所知，我将在此叙述它。\par

% structure: p0102 source=h2 target=subsection reason=relative-heading-rank; base=h2

\subsection{第十三章  \WorkTitle{亚他那修关于亚略死亡的书信}摘录}

\Person{亚略}在\Place{亚历山大里亚}逗留了很长时间之后，试图骚乱地再次强行进入教会的集会，声称放弃他的不敬，并承诺接受由教父们起草的信德声明。但未能在神圣的\Person{亚历山大}以及继承\Person{亚历山大}宗主教职位和虔敬的\Person{亚他那修}那里获得信任，他在\Place{尼苛米底亚}主教\Person{优西比乌}的帮助和鼓励下，前往\Place{君士坦丁堡}。他随后进行的阴谋，以及公义的审判者对他们所施的惩罚，都由杰出的\Person{亚他那修}在他写给\Person{阿彼翁}的\WorkTitle{信}中最好地叙述了。因此，我现在将这一段插入我的著作中。他写道：——\par

\Quote{他去世时，我并不在君士坦丁堡；但司铎\Person{玛卡里乌斯}在那里，我从他那里得知了全部情况。\Person{君士坦丁}皇帝受\Person{优西比乌}及其同党怂恿，派人去接\Person{亚略}。\Person{亚略}到达后，皇帝问他是否持守\OriginalTerm{天主教会}{Catholic Church}的信仰。\Person{亚略}于是发誓说他的信仰是正统的，并提交了一份书面的信仰摘要；但他隐瞒了被\Person{亚历山大}主教逐出教会的原因，并狡猾地利用了圣经的言辞。因此，当他起誓声明他并没有坚持\Person{亚历山大}驱逐他的那些错误时，\Person{君士坦丁}就让他走了，说：\InnerQuote{如果你的信仰是正统的，你发誓发得好；但如果你的信仰是不虔敬的，而你却发了誓，就让天主从天上来审判你。}他离开皇帝后，\Person{优西比乌}的同党照例嚣张，要带他进入教堂；但可敬的\Person{亚历山大}，君士坦丁堡主教，拒绝允许，声称异端的发明者不应被接纳进入共融。最后，\Person{优西比乌}的同党发出威胁：\InnerQuote{正如我们不顾你的意愿，成功说服皇帝接来\Person{亚略}，现在，即使你禁止，明天\Person{亚略}也要和我们在这教堂里共融。}他们是在星期六说这话的。\Person{亚历山大}主教对所闻深感悲伤，走进教堂，倾吐哀伤，向天主举手恳求，俯伏在至圣所的地面上祈祷。\Person{玛卡里乌斯}和他一起进去，与他一同祈祷，并听到了他的祈祷。他祈求两件事中的一件。\InnerQuote{如果\Person{亚略}，}他说，\InnerQuote{明天要被接入教会，就请让你的仆人离去，不要使虔敬者与不虔敬者一同灭亡。如果你要宽恕你的教会——我知道你宽恕她——就垂看\Person{优西比乌}追随者的话，不要将你的产业交给毁灭和羞辱。除掉\Person{亚略}，免得他进入教会时，异端似乎也随之而入，不虔敬从此被当作虔敬。}祈祷完毕，主教深怀忧虑地离开了教堂，随后发生了一场可怕而异常的灾祸。\Person{优西比乌}的同党大肆威胁，而主教则求助于祈祷。\Person{亚略}因受同党保护而胆大妄为，说了许多轻浮愚蠢的话，这时他突然因生理需要而退下，立刻，正如经上所记载的，\BibleQuote{他俯首坠下，腹部崩裂}，就断了气，同时失去了共融和生命。这就是\Person{亚略}的结局。\Person{优西比乌}的同党羞惭满面，埋葬了他们所信奉之人的尸体。可敬的\Person{亚历山大}完成了庆典，与教会一同在虔敬和正统中欢庆，与众弟兄一同祈祷，大大光荣了天主。这并不是因为他因\Person{亚略}的死而喜乐——绝对不可——因为\BibleQuote{人人只有一死}；而是因为这件事明显超越了任何人的审判。因为主亲自审判了\Person{优西比乌}同党的威胁和\Person{亚历山大}的祈祷，谴责了亚略异端，表明它不配被接纳进入教会的共融；从而向所有人显明，即使它得到皇帝和所有人的支持，它却为真理本身所定断。}\par

这些是\Person{亚略}从他亲手播下的有害种子所收获的初果，也是日后惩罚他的前奏。他的不虔敬受到了刑罚的定断。\par

现在我要将叙述转向皇帝的虔敬。他致信罗马帝国所有臣民，劝勉他们弃绝从前的错误，接受我们救主的教义，并试图引导他们走向这一真理。他激励各城的主教建造教堂，不仅通过书信鼓励他们，还赠予大量金钱，承担所有建造费用。他自己的这封信函说明了这一点，信的内容如下：\par

% structure: p0107 source=h2 target=subsection reason=relative-heading-rank; base=h2

\subsection{第14章 君士坦丁皇帝关于建造教堂的书信}

\Person{君士坦丁·奥古斯都}，伟大而常胜者，致\Person{优西比乌}。\par

\Quote{我深知并确信，我亲爱的兄弟，由于我们救主基督的仆人至今遭受邪恶的阴谋和暴虐的迫害，所有教堂的建筑或因疏忽而完全毁坏，或因对即将来临的邪恶的恐惧而低于应有的尊严。但现在自由已经恢复，那巨龙，靠着天主眷顾并借助我们，被逐出了帝国的政权，我想神圣的大能已为众人所知，那些因恐惧、不信或堕落而活在错误中的人，现在一旦认识那真正存在者，将被引向真实而正确的生活方式。因此，请你在你所管辖的教堂修复上殷勤努力，并劝勉其他地方的主要主教、司铎和执事同样热心从事这一工作；以便所有现存的教堂得以修复或扩建，并在需要之处建造新的教堂。你及通过你干预的其他人，可以向行政长官和省政府申请一切所需之物；因为他们已接到书面命令，要热忱服从你的圣德所命令的一切。愿天主保佑你，亲爱的兄弟。}\par

皇帝就这样写信给各省的主教，论及建造教堂的事。从他给巴勒斯坦的\Person{优西比乌}的信中，可以轻易了解他是如何设法获取圣经抄本的。\par

% structure: p0111 source=h2 target=subsection reason=relative-heading-rank; base=h2

\subsection{第15章 君士坦丁关于预备圣经抄本的书信}

\Person{君士坦丁·奥古斯都}，伟大而常胜者，致\Person{优西比乌}。\par

\Quote{在以我们名字命名的城市里，通过天主救主的眷顾，有大批人已与圣教会联合。由于那里一切都在迅速进步，我们认为极其重要的是再增建一些教堂。请欣然采纳我们决定并认为适宜告知你睿智的程序，即：你应让人用优质皮纸抄写五十卷书，字迹清晰，便于阅读和携带；这些必须由熟练的、精通其艺的书法家抄写。我指的是圣经的抄本，如你所知，教会会众最需要拥有并使用它们。朕已致函\OriginalTerm{教区总管}{catholicus}，令他务必为这项工程提供一切所需。你应负责确保这些抄本在短时间内完成。完成后，你凭此函可调派两辆公车运送至我们处；这样精美的抄本便可轻易呈交我们检视。请指派你教会的一名执事负责此事；当他到达我们这里时，将得到朕恩慈的证明。愿主保佑你，亲爱的兄弟。}\par

以上所说已足以表明，不，清楚证明皇帝在宗教事务上表现出何等热忱。然而，我还要补充他对我们救主之墓的崇高作为。因为他得知，偶像崇拜者狂怒地将泥土堆在主的坟墓上，急于抹去对祂救恩的一切纪念，并在其上建造了一座献给放纵情欲的女神的神庙，以嘲弄童贞女的诞生。皇帝下令拆毁那可憎的神龛，将被可憎祭品污染的土壤运走，抛到城外远处，并在原址上建造一座宏伟美丽的新圣殿。这一切在他写给耶路撒冷主教玛加略的信中清楚说明，我们已提过此人曾是尼西亚大公会议的成员，并与他的弟兄们一同抵制了亚略的亵渎。以下是这封信。\par

% structure: p0115 source=h2 target=subsection reason=relative-heading-rank; base=h2

\subsection{第十六章 君士坦丁皇帝致耶路撒冷主教玛加略关于建造圣堂的书信}

\Signature{君士坦丁，胜利且伟大的，致玛加略。}\par

我们救主的恩宠如此奇妙，没有任何言语足以表达当前的奇迹。祂至圣苦难的纪念物，竟能在大地之下隐藏如此漫长的岁月，直到那全人类的公敌死去之时，才注定要向祂的解放了的仆人们显现，这事实超越了其他一切可赞叹的事。如果将全世界所有智者聚集一处，试图相称地表述它，他们也无法接近其无限的距离；因为这一奇迹远超人类一切信德的能力，正如天上事物按其本性比人类更强。因此，我首要且唯一的目的就是，正如藉着新的奇迹，对真理的信德日渐坚定，我们所有人的心灵也当更热切地、明智地、热忱地同心合意地致力于神圣法律。我想我的意图现已大致为人所知，我愿你特别确信，我最热切的愿望是用美丽的建筑装饰那被祝圣的地点，我奉天主的命令已将它从污秽偶像的重压下解放出来。因为从一开始，祂就宣布它为圣地，并且在祂揭示出我们主受苦的证明和纪念之时，使其变得更加神圣。\par

因此，我信赖你的睿智，请采取一切必要措施，不仅使\OriginalTerm{巴西利卡}{basilica}本身超越其他所有，而且使其所有安排都无与伦比地优于世上各城中最美丽的建筑。我们已委托我们的朋友德拉基利亚努斯，他担任本省最尊贵\OriginalTerm{总督}{præfect}之职，负责墙体的建造和装饰的监督。他接到我们的命令，雇佣工匠和技师，并提供你认为建筑所需的一切。你视察工程后，请函告我们，你认为哪些柱子或大理石最具装饰性，以便你告知我们所需的一切，可从世界各地运往那里。因为那最奇妙的地方，应当按其尊严加以装饰。我希望从你那里得知，你是否认为巴西利卡的拱顶应加镶板，或以其他方式装饰；因为如果要加镶板，也可以镀金。你的圣德必须尽快向上述官员表明需要哪些工匠、技师以及多少金额；并请及时告知我，不仅关于大理石和柱子，还有镶板天花板，如果你判定这是最美丽的建造方式。愿主保佑你，亲爱的兄弟。\par

% structure: p0119 source=h2 target=subsection reason=relative-heading-rank; base=h2

\subsection{第十七章 海伦娜，君士坦丁皇帝的母亲——她在建造圣堂上的热忱}

这些信件的携带者是一位毫不逊色于皇帝母亲的显赫人物，甚至就是那位因子女而荣耀、其虔诚为众人称颂的她；她生下了那伟大的光明之星，并以虔诚养育了他。她并未因年事极高而畏惧旅途劳顿，而是在她去世前不久（她死于八十岁高龄）承担了这次旅行。\par

当皇后看见救主受难的地方，她立即下令摧毁那里建造的偶像庙宇，并移走其上所立的泥土。当那隐藏已久的坟墓被发现时，在主墓穴附近看到三个十字架。大家都确信其中一个是吾主耶稣基督的十字架，另外两个是与祂同钉的盗贼的。然而他们无法辨别哪一个曾靠近主的身体，哪一个曾接受祂宝血的倾流。但城中明智而圣善的会长玛加略以如下方式解决了这个疑问。他让一位长期患病的高贵妇人，在热切祈祷中依次被每个十字架触摸，从而辨认出救主十字架中的德能。因为当这十字架靠近妇人时，立刻驱逐了顽疾，使她痊愈。\par

皇帝的母亲得知她的愿望实现后，下令将一部分钉子嵌入御盔，以便她儿子的头部免受敌人箭矢的伤害。另一部分钉子被制成马的辔头，不仅为确保皇帝的安全，也为应验古代预言；因为早在很久以前，匝加利亚先知曾预言说：\Quote{\Emph{在马的辔头上必有\Quote{归于全能的上主为圣}。}}\par

她将我们救主的十字架一部分运至宫殿。剩余部分以银匣包裹，委托给城中主教保管，并嘱咐他仔细保护，以便完好无损地传于后世。随后她四处召集工匠和材料，建造了最宽敞、最宏伟的教堂。描述它们的美与宏伟实无必要；因为所有虔诚的人，如果我可以这样说，都急切前往那里瞻仰建筑的壮丽。\par

这位著名而可敬的皇后还行了另一件值得纪念的事。她召集了所有发愿终身守贞的妇女，让她们坐在榻上，自己亲手执行仆役的职责，侍奉她们食物、递杯、斟酒，并取来盆罐倒水为他们洗手。\par

完成了这些和其他可赞颂的行为后，皇后回到儿子身边，不久后，她欢欣地进入了另一个更美好的生命，临别时给予儿子许多虔诚的劝告和热切的祝福。她去世后，人们向她表示了纪念的荣誉，这荣誉是她坚定而热切地事奉天主所应得的。\par

% structure: p0126 source=h2 target=subsection reason=relative-heading-rank; base=h2

\subsection{第十八章  尼苛默底亚主教欧瑟伯的非法调动}

亚略派的人并未停止他们的邪恶阴谋。他们签署信经只是为了披上羊皮，而行狼的勾当。拜占庭（当时这座城市尚未被称为君士坦丁堡）的圣亚历山大，曾以其祈祷刺透亚略的心，此时已蒙召进入更美好的生命。而传播不敬的欧瑟伯，几乎不理会他不久前与其他主教一起同意的定义，立即离开尼苛默底亚，夺取了君士坦丁堡的主教座，直接违反了禁止主教和司铎从一城调往另一城的教规。但那些疯狂到否认天主独生子神性的人，同样违反其他法律，并不足为奇。这并非他首次如此革新；因为他最初受托管理贝黎图斯教区，从那裡跳到了尼苛默底亚。为此，他因显明的不敬被主教会议驱逐，尼该亚主教德奥格尼也一样。此事在君士坦丁皇帝的信函中再次提及；我将在下面插入他写给尼苛默底亚人的信的结尾。\par

% structure: p0128 source=h2 target=subsection reason=relative-heading-rank; base=h2

\subsection{第十九章  君士坦丁皇帝针对欧瑟伯和德奥格尼致尼苛默底亚人的信函}

是谁将这些学说教导给无辜的群众？显然是欧瑟伯，他是暴君暴行的合作者。因为他曾是暴君的受造之物，这已清楚地表明；而且事实上，从主教的被杀以及这些殉道者是真正的主教这一点可以证明。对基督徒的无情迫害大声宣告了这一事实。\par

我在此不提及那些针对我的侮辱，反对派的大部分阴谋正是通过这些侮辱实现的。但他甚至派遣间谍监视我，几乎不阻止为暴君征募军队。不要任何人想象我举出无法证明的事。我拥有明确的证据；因为我已逮捕了他所属的主教和司铎。但我略过所有这些事实。我只是为了让他们对自己的行为感到羞愧而提及，并非出于任何怨恨。\par

有一件事我害怕，有一件事让我焦虑，那就是看见你们被指控为同谋；因为你们受欧瑟伯学说的影响，从而偏离了真理。但如果你们得到一位持纯正忠诚信德的主教，并仰望天主，你们的治愈将是迅速的。这完全取决于你们自己；无疑，你们早就如此行了，若不是上述欧瑟伯在当权者的强烈支持下到来，颠覆了一切纪律。\par

既然有必要再谈一谈\Person{欧瑟伯}，敬请各位耐心记得，在\Place{尼西亞}城曾召开过一次会议，我出于良心，遵命出席，我受此感动，唯一动机是渴望在众人中达成一致，且首先证明并消除源于\Place{亚历山大}城的\Person{阿黎乌}之愚昧、并立即由\Person{欧瑟伯}荒谬有害的阴谋所加强的祸害。但是，亲爱的、备受尊敬的弟兄们，你们不知道\Person{欧瑟伯}，尽管被自己良心的见证所定罪，却多么急切且可耻地坚持支持那些已被普遍谴责的虚假教义。他秘密派人向我为他自己求情，并亲自恳求我帮助阻止他被逐出主教之职，尽管他的罪行已被完全揭露。\OriginalTerm{天主}{God}，我相信祂将继续向你们和我施恩，是我所言真理的见证。我当时自己也受了\Person{欧瑟伯}的迷惑和欺骗，你们将清楚知道。他在一切事上都按自己的意愿行事，内心充满了各种隐秘的邪恶。\par

略去其余恶行的叙述，你们应当知道他与他的愚昧同谋\Person{德敖尼}最近所犯的罪行。我曾下令逮捕\Place{亚历山大}城某些背弃我们信仰的人，他们点燃了纷争的火把。但是这些善人们，即主教们，我本因会议的宽仁保留了他们的补赎，他们不仅收容了那些人，还参与了他们的恶行。因此我决定惩罚这些忘恩负义之人，逮捕他们并将他们放逐到遥远的地区。\par

\Quote{如今你们的责任是怀着那显然你们一直持有、且应当持守的信德仰望天主。这样，我们便有理由为任命纯洁、正统且仁爱的主教而喜乐。若有人提及那些破坏者，或胆敢赞扬他们，就叫他明白，他的狂妄必因我作为天主仆人所受的权柄而被压制。愿天主保全你们，亲爱的弟兄们！}\par

上述主教随即被撤职并放逐。\Person{安菲翁}被委托管理\Place{尼科美底亚}教会，\Person{克雷斯图斯}管理\Place{尼西亞}教会。但被放逐的主教们使用他们惯用的伎俩，滥用了皇帝的仁慈，重开先前的争端，并恢复了他们以前的权势。\par

% structure: p0136 source=h2 target=subsection reason=relative-heading-rank; base=h2

\subsection{第二十章 \Person{欧瑟伯}及其追随者对\Place{安提约基雅}圣\Person{欧斯塔修斯}主教的诡诈阴谋}

\Person{欧瑟伯}，如我所说，强行夺取了\Place{君士坦丁堡}教区。如此，他在该城获得了巨大权势，频繁拜访皇帝并与他亲密交往，他获得自信并密谋对抗那些最先支持真理的人。他起初假装想去\Place{耶路撒冷}参观那里建造的著名建筑；皇帝被他谄媚所欺，允许他以最高荣誉出发，为他提供车辆及其他随行装备和扈从。\Place{尼西亞}主教\Person{德敖尼}，如我们之前所说，是他邪恶阴谋的同谋，与他同行。当他们抵达\Place{安提约基雅}时，他们戴上友谊的面具，受到最大的尊敬。信德的伟大斗士\Person{欧斯塔修斯}以兄弟般的仁慈对待他们。当他们抵达圣地时，与那些与自己意见相同的人会面，即\Place{凯撒勒雅}主教\Person{欧瑟伯}、\Place{西托波利斯}主教\Person{帕特洛非路}、\Place{吕达}主教\Person{阿埃提乌}、\Place{劳迪刻雅}主教\Person{德奥多托}，以及其他怀有\Person{阿黎乌}思想的人；他们将已策划的阴谋告知他们，并同他们一起前往\Place{安提约基雅}。他们旅程的借口是向\Person{欧瑟伯}表示应有的敬意；但真正动机是对宗教的战争。他们贿赂了一个以美貌为交易的卑贱女人，让她出卖她的舌头，然后前往会议，当所有旁观者被命令退去后，他们带进那可怜的女人。她怀中抱着一个婴儿，大声且无耻地断言\Person{欧斯塔修斯}是这婴儿的父亲。\Person{欧斯塔修斯}自知无辜，问她能否提出任何证人证明她所说的话。她回答说不能：然而这些公正的法官竟允许她宣誓，尽管法律上说：\BibleQuote{\Emph{凭两三个证人的口，才能确定事实}}；宗徒也说：\BibleQuote{\Emph{控告长老，除非有两三个证人，不要受理}}。但他们藐视这些神圣的法律，在没有证人的情况下受理了对这位伟人的控告。当那女人再次宣誓说\Person{欧斯塔修斯}是婴儿的父亲时，这些热爱真理的法官判定他为奸淫者。其他坚持教义的主教们，对这些阴谋一无所知，公开反对判决，并建议\Person{欧斯塔修斯}不要服从；阴谋的发起者立即前往皇帝那里，竭力说服他控告是真实的，撤职判决是公正的；他们成功地使这位虔诚与贞洁的斗士以奸淫者和暴君的罪名被放逐。他被押解穿越\Place{色雷斯}，前往\Place{伊利里库姆}的一座城市。\par

% structure: p0138 source=h2 target=subsection reason=relative-heading-rank; base=h2

\subsection{第二十一章 圣\Person{欧斯塔修斯}被放逐后，\Place{安提约基雅}任命了持有异端意见的主教}

\Person{欧拉理乌斯}首先被祝圣，以接替\Person{欧斯大削}。但\Person{欧拉理乌斯}在升任后只活了很短一段时期，便有意将巴勒斯坦的\Person{欧瑟伯}调任到这个主教职位。然而\Person{欧瑟伯}拒绝了任命，皇帝也禁止将此职授予他。接下来又推举了\Person{欧弗洛尼乌斯}，但他也只过了一年零几个月就去世了，于是主教座便授予了\Person{弗拉基卢斯}。所有这些主教都暗中依附亚略异端。因此，那些珍视真宗教的人士，无论是神职人员还是平信徒，大多数都离开了教会，自行集会。他们被称为欧斯大削派，因为他们是在\Person{欧斯大削}被放逐之后才开始集会的。上述那可怜妇人不久后患了一场严重而持久的疾病，于是她承认了自己参与的骗局，并将整个阴谋公之于众，不仅对两三个司铎，而且对很多司铎。她供认自己是被收买才提出这虚假无耻的指控，但她的誓言并非完全虚假，因为那孩子的生父是一个名叫\Person{欧斯大削}的铜匠。这就是这个极好的派别在\Place{安提约基雅}所犯下的一些罪行。\par

% structure: p0140 source=h2 target=subsection reason=relative-heading-rank; base=h2

\subsection{第二十二章 印度人的归化。}

此时，认识天主的真光首次照耀\Place{印度}。皇帝的勇敢和虔诚已传遍天下；蛮族因经验而学会选择和平而非战争，能够毫无恐惧地彼此交往。因此，许多人开始长途旅行：有些出于探索的渴望，有些则基于商业冒险精神。大约此时，一位\Place{提洛}人，精通希腊哲学，渴望深入印度内陆，便带着他的两个年轻外甥出发了。当他实现了自己的愿望后，便登船返回本国。船因需要补充淡水而被迫靠岸，蛮族攻击了它，淹死了一些船员，将其他人俘虏。那位舅父也在被杀者之中，两个少年被带到国王面前。一个名叫\Person{厄德西乌斯}，另一个名叫\Person{弗鲁孟提乌斯}。该国国王随着时间的推移，看出他们的才智，便提拔他们管理自己的家室。若有人怀疑此事的真实性，请回忆\Person{若瑟}在\Place{埃及}王国的历史，以及\Person{达尼尔}的历史，还有那三位真理的斗士，他们从俘虏变为\Place{巴比伦}的王子。国王去世了，但这些年轻人仍留在他儿子身边，并被提升到更大的权柄。由于他们在真宗教中长大，他们劝勉那些来到该国的商人，按照罗马人的习惯聚会，参与神圣礼仪。过了相当长的时间，他们请求国王奖赏他们的服务，允许他们返回本国。他们获得了许可，安全抵达罗马领土。\Person{厄德西乌斯}向\Place{提洛}方向前行，但\Person{弗鲁孟提乌斯}，其宗教热忱超过了对亲人的自然情感，前往\Place{亚历山卓}，告诉该城的主教，印度人深切渴望获得属灵之光。那时，\Person{亚大纳削}执掌那教会的船舵；他听了这故事，然后\Quote{谁，}他说，\Quote{能比你本人更好地驱散无知的迷雾，将神圣宣讲的真光引入这个民族呢？} 说完这番话，他便授予他主教尊位，并派他去从事那个民族的精神耕耘。这位新祝圣的主教离开了这个国家，不惧浩渺的海洋，回到了他未开垦的工作之地。在那里，有天主的恩宠与他一同劳作，他愉快而成功地扮演了农夫的角色，以宗徒奇迹的作为俘获那些试图反驳他话语的人，并由此通过这些奇迹证实他的教导，他继续每天俘获许多灵魂。\par

% structure: p0142 source=h2 target=subsection reason=relative-heading-rank; base=h2

\subsection{第二十三章 伊比利亚人的归化。}

福鲁门提乌斯就这样引导印度人认识天主。大约在同一时期，伊比利亚被一个被掳的妇人引上了真理之路。她恒心祈祷，不让自己享用比铺在地上的麻袋更柔软的床铺，并视禁食为最大的享受。这种刻苦得到了与宗徒相似的恩赐。那些蛮族人不懂医药，患病时习惯互相走访，询问那些曾患同样疾病而痊愈的人是用什么方法治好的。按照这一习俗，一位母亲带着生病的孩子来到这位可敬的妇人面前，询问她是否知道治疗此病的药方。妇人接过孩子，放在自己的床上，向世界的创造主祈祷，求他垂怜孩子，治愈疾病。主垂听了她的祈祷，使孩子痊愈。这位非凡的妇人从此声名大噪；王后患了重病，听说了她的事迹，便派人去请她。被掳的妇人自视甚微，不愿接受王后的邀请。但王后因急需所迫，不顾王室的尊严，亲自跑到被掳的妇人那里。妇人让王后躺在她简陋的床上，再次用祈祷这一有效的药方治疗她的疾病。王后痊愈了，献上黄金、白银、短衣、长袍以及她认为值得拥有的礼物和君王慷慨应赐之物作为酬谢。圣女告诉她，她不需要这些，但她认为最大的报酬是王后认识真宗教。然后，她尽其所能地解释了神圣的教义，并劝告王后为使她痊愈的基督建造一座圣堂。王后回到宫中，因突然痊愈而令她的丈夫惊叹不已；她向他讲述了被掳的妇人所敬拜的天主的德能，并恳求他承认独一的天主，为他建造一座圣堂，并引导全国敬拜他。国王因王后所受的奇迹大为喜悦，但不同意建造圣堂。不久之后，他出去打猎，慈爱的主像对待保禄一样使他成为猎物；突然的黑暗笼罩了他，使他不能移动一步；而与他一同打猎的人却享受着惯常的阳光，只有他被盲目的锁链束缚。在困惑中，他找到了逃脱的方法，因为记起自己先前的不信，他呼求被掳妇人的天主的帮助，黑暗立刻消散了。于是他来到那位奇妙的被掳妇人面前，请她指示他应当如何建造圣堂。那曾用建筑技能充满贝匝肋耳的主，仁慈地使这位妇人能够设计圣堂的蓝图。妇人着手设计，人们开始挖土建造。当建筑完工，屋顶盖好，除了司铎之外一切齐备时，这位可敬的妇人又想方设法获得了他们。因为她劝国王派遣使团前往罗马皇帝那里，请求派遣宗教导师。国王于是为此派遣了使团。君士坦丁皇帝对宗教事业热诚，得知使团的意图后，欣然欢迎使节，并选了一位富有信德、智慧和德行的主教，赠送许多礼物，派他到伊比利亚人那里，使他们认识真天主。他不仅满足了伊比利亚人的请求，还主动承担了保护波斯境内基督徒的责任；因为他得知他们受到外教人的迫害，而他们的国王本身也是错误的奴隶，正策划种种诡计要消灭他们，便写信给他，恳求他本人接受基督宗教，并尊敬其信奉者。君士坦丁本人写的信将更清楚地表明他在此事上的热诚。\par

% structure: p0144 source=h2 target=subsection reason=relative-heading-rank; base=h2

\subsection{第二十四章 君士坦丁皇帝致波斯王沙普尔论基督徒的书信}

在保护神圣信仰时，我享受真理之光；借着追随真理之光，我得以更圆满地认识信仰。因此，正如事实所证明，我承认那最神圣的敬拜教导了对至圣天主的认识。我宣认这一敬拜。以这位天主的德能为盟友，我从大洋的最远边界开始，一个接一个地使世界的每一部分都燃起希望。如今，所有一度被众多暴君奴役、因日常苦难而疲惫不堪、几乎灭绝的民族，因接受国家的保护而重新燃起了生命。\par

我所敬畏的天主，正是我的精锐部队肩负其徽号的那一位，他们为正义的事业而前行，并以辉煌的胜利回报我。我承认，我以永恒的记忆敬畏这位天主。他以纯洁无玷的心灵瞻仰那居于至高天的主。我屈膝呼求他，避免一切可憎的血，一切不雅和不祥的气味，一切咒术之火，以及一切因不法而可耻的错误曾毁灭整个民族并将其投入地狱的污染。\par

天主不容许那些祂在慈祥眷顾中赐予人类以满足需要的恩赐，被人任意滥用。祂只要求人有一颗纯洁的心灵和无瑕的灵魂，并以此衡量他们的德行和虔诚的行为。祂喜悦温良和谦逊，喜爱温良的人，憎恶挑起争端的人；祂喜爱信德，惩罚不信；祂折断一切骄傲的权势，惩罚傲慢者的无礼。祂彻底推翻自高自大的人，按他们的功过赏报谦卑忍耐的人。祂重视正义的统治，以祂的帮助加强它，并以和平的祝福护卫君王的谋略。\par

我知道自己并没有错误，我的兄弟，当我承认这位天主是万有的统治者和圣父时，这是一个真理，而我之前许多坐在皇位上的人却被错误所迷惑，试图否认它。但他们的结局如此可怕，成为全人类的可怕警告，以阻止其他人犯下类似的罪行。在这些人中，我认为那一个人就是被天主的愤怒如同雷霆一般驱赶到你们国家的人，并且他在你们自己的土地上留下了臭名昭著的耻辱纪念碑。实际上，看来这是安排得妥妥当当的，使我们生活的时代因对这类人公开而明显的惩罚而显得突出。我亲眼目睹了那些以非法敕令迫害天主之民的人的结局。因此，我更特别感谢天主，现在藉祂特殊的天主眷顾，为遵守祂法律的人恢复了和平，他们在其中高举和喜乐。\par

\Quote{我预期未来的幸福和安全，只要天主以祂的良善使所有人在那唯一纯正的真实宗教中合而为一。因此，你完全可以理解当我听说\Place{波斯}最优秀的省份中这类人极其丰富时，我是多么喜乐；我指的是基督徒，因为我说的正是他们。那么，一切对你和他们都顺利，因为万有的主将对你仁慈和恩惠。既然你如此强大而虔诚，我将基督徒托付给你照顾，并把他们留在你的保护之中。我恳求你，以符合你良善的慈爱对待他们。你在这件事上的忠信将给你和我们带来难以言喻的福益。}\par

这位卓越的皇帝对所有接受真信仰的人如此关心，以至于他不仅照顾自己的臣民，也照顾其他君主的臣民。因此，他蒙受了天主的特别保护，以至于尽管他掌握整个欧洲、非洲以及亚洲大部分地区的缰绳，他的臣民都乐于服从他的统治。外国民族顺从于他的统治，有些是自愿归顺，有些是在战争中被征服。凯旋纪念碑到处树立，皇帝被称为胜利者。\par

然而，君士坦丁的赞美已被许多其他作者传扬。我们必须继续我们历史的线索。这位应得最高名声的皇帝，将全部心思投入到配得上宗徒的事情上，而那些被接纳进入司铎尊位的人不仅忽视了建立教会，反而竭力将其从根基上拔除。他们捏造各种虚假的指控，反对那些按照宗徒所传教义治理教会的人，并尽力废黜和驱逐他们。他们的嫉妒并不满足于他们针对\Person{欧斯达提乌}捏造的臭名昭著的谎言，而是采用一切手段推翻另一个伟大的信仰堡垒。这些悲惨的事件我现在将尽可能简明地叙述。\par

% structure: p0152 source=h2 target=subsection reason=relative-heading-rank; base=h2

\subsection{第二十五章 对抗\Person{圣亚大纳削}的阴谋}

那位可敬的主教\Person{亚历山大}，曾成功抵制了\Person{亚略}的亵渎言论，在\Place{尼西亚}大公会议五个月后去世，\Person{亚大纳削}继任了\Place{亚历山太}教会的主教职位。他自幼接受神圣研究训练，在担任的每一个教会职务中都赢得了普遍的钦佩。在大公会议上，他如此捍卫宗徒的教义，以至于他赢得了所有真理斗士的赞同，而对手则学会了将这位对手视为私人仇敌和公敌。他当时作为\Person{亚历山大}的随从出席大公会议，年纪很轻，尽管他是首席执事。\par

那些否认天主的独生子的人，听说\Place{亚历山太}教会的舵柄已交到他的手中，他们由于经验知道他对真理的热忱，认为他的统治将导致他们权威的毁灭。因此，他们采取了以下针对他的阴谋。为了掩饰嫌疑，他们贿赂了一些\Person{默利提乌斯}的追随者——这些人在尼西亚大公会议后被罢免，但仍坚持在\Place{特拜得}和\Place{埃及}邻近地区煽动骚乱——并说服他们去见皇帝，指控\Person{亚大纳削}在埃及征税，并将收集的金子交给一个准备篡夺皇权的人。皇帝被这个故事欺骗，\Person{亚大纳削}被带到\Place{君士坦丁堡}。他一到达，就证明指控是虚假的，并恢复了他由天主赐予的职务。这由皇帝写给\Place{亚历山太}教会的一封信所证明，我将只抄录其最后一段。\par

\Emph{\Person{君士坦丁}皇帝致\Place{亚历山太}人信函的一部分}\par

\Quote{相信我，我的弟兄们，那些恶人无法对你们的主教造成任何影响。他们肯定没有别的企图，只是浪费我们的时间，并使他们自己在今生没有悔改的余地。因此，你们要帮助自己，爱那赢得你们爱慕的，并尽一切力量驱逐那些想破坏你们团结的人。仰望天主，彼此相爱。我喜乐地欢迎你们的主教\Person{亚大纳削}；我与他交谈，如同与一个我知道是天主之人的交谈。}\par

% structure: p0157 source=h2 target=subsection reason=relative-heading-rank; base=h2

\subsection{第二十六章 另一次对抗\Person{亚大纳削}的阴谋}

然而，诬蔑\Person{亚大纳削}的人并未停止他们的企图。相反，他们竟捏造了一个如此大胆的虚构故事来攻击他，以至超过了古代悲剧或喜剧作家的任何发明。他们再次贿赂同一派别的人，将他们带到皇帝面前，大声指控那位德行之魁犯了许多可憎的罪行。该派的首领是\Person{尤西比乌}、\Person{特奥格尼斯}，以及\Place{佩林图斯}（现今称为\Place{赫拉克利亚}）的主教\Person{特奥多鲁}。在指控\Person{亚大纳削}犯了他们描述为骇人听闻、不可容忍甚至不可听闻的罪行后，他们说服皇帝在\Place{巴勒斯坦}的\Place{凯撒勒雅}召开一次会议，那里\Person{亚大纳削}有许多敌人，并命令在那里审判他的案件。皇帝完全不知道所设的阴谋，被他们说服，发出了所需的命令。\par

但\Person{圣亚大纳削}深知审判他的人的恶意，拒绝出席会议。这给了反对真理的人一个进一步指控他的借口；他们在皇帝面前控告他傲慢自大。他们的希望并未完全落空；因为皇帝虽然极其宽容，却被他们的陈述激怒，愤然写信给他，命令他前往\Place{提洛}。会议奉命在那里召开，依我之见，是出于怀疑\Person{亚大纳削}因\Place{凯撒勒雅}的主教而畏惧该地。皇帝也写信给会议，其风格符合他虔诚的热忱。他的信如下。\par

% structure: p0160 source=h2 target=subsection reason=relative-heading-rank; base=h2

\subsection{第二十七章 君士坦丁皇帝致提洛会议的书信}

\Person{君士坦丁·奥古斯都}致在\Place{提洛}召开的圣会议。\par

在当今时代的普遍繁荣中，似乎理应让天主教会同样免于困扰，让\Person{基督}的仆人摆脱一切责难。\par

但是某些人，受争论的疯狂欲望所煽动（且不说过着与其身份不相称的生活），正试图将一切陷入混乱。在我看来，这是所有可能灾祸中最大的。因此，我恳请你们，套用一句俗话，火速地、毫不迟延地正式召开一次会议；以便你们能够支持那些需要帮助的人，治愈处境危险的弟兄，使分裂的成员重归一致，趁尚有时日纠正教会的混乱；从而为那些辽阔的地区恢复和谐——唉！这和谐已被少数人的傲慢摧毁了。我相信每个人都会承认，你们所能做的在天主眼中没有任何事情如此悦乐，如此超越我的一切祈祷以及你们自己的祈祷，或如此有助于你们的声誉，如同恢复和平一样。\par

因此，不要延迟；当你们以我们的救主向祂所有仆人要求的那种真诚和忠信（几乎如同我们能听见的话语）聚集时，要以加倍的热忱努力对这些纷争作出适当的了结。\par

我这一方将不遗余力地促进我们宗教的利益。你们来信中建议的一切，我都做了。我已派遣你们指定的那些主教，指示他们前往会议，以便与你们一起商议教会事务。我还派遣了具有执政官品级的\Person{狄奥尼修}，为那些将与你们一同开会的人提供建议，并亲自见证你们的进程，特别是所维持的秩序和规范。如果有人胆敢在此次也违抗我们的命令，拒绝前来会议（虽然我不预见到这一点），将立即派出一名官员，依据皇帝命令将他流放，使他学会不反对皇帝为支持真理而颁布的法令。\par

\Quote{现在留给你们圣德的事，就是凭一致的判断，不偏不倚、不带偏见，依照教会教规作出决定，并为因错误可能造成的过犯设计适当的补救措施；以便教会能摆脱一切责难，我的忧虑得以减轻，纷争者之间得以恢复和平，你们的声望得以增长。愿天主保佑你们，亲爱的弟兄们。}\par

于是主教们前往\Place{提洛}会议。其中也包括那些被指控持有异端教义的人；\Place{加沙}的主教\Person{阿斯克勒帕}便是其中之一。可敬的\Person{亚大纳削}也出席了。我将先叙述指控的悲剧，然后再描述这个著名法庭的诉讼程序。\par

% structure: p0168 source=h2 target=subsection reason=relative-heading-rank; base=h2

\subsection{第二十八章 提洛会议}

\Person{阿尔塞尼}是梅利提安派的主教。他那一派的人将他藏在一个隐蔽的地方，命令他尽可能久地待在那里。然后他们砍下一具尸体的右手，进行防腐处理，放入一个木盒，到处携带，声称那是被\Person{亚大纳削}谋杀的\Person{阿尔塞尼}的手。但无所不见的眼睛没有让\Person{阿尔塞尼}长期隐藏。他先是在\Place{埃及}被看见，然后在\Place{特拜伊德}；后来，天主的眷顾将他引到\Place{提洛}，在那里那出悲剧名声的手被带到了会议面前。\Person{亚大纳削}的朋友们搜寻到他，将他带到一个客栈，强迫他暂时隐藏。清晨，伟大的\Person{亚大纳削}来到了会议。\par

首先带来一个荡妇，她大声无耻地作证说，她曾发誓终身守贞，但\Person{亚大纳削}住在她家时侵犯了她的贞洁。她指控后，被告上前，与他一起的是一位值得称颂的司铎，名叫\Person{弟茂德}。法庭命令\Person{亚大纳削}答辩，但他沉默，仿佛不是\Person{亚大纳削}。然而，\Person{弟茂德}对她这样说：\Quote{女人，我曾与你交谈过吗？或进过你的家吗？}她更加厚颜无耻地回答，在与\Person{弟茂德}争辩时大声尖叫，用手指着他说：\Quote{是你夺走了我的贞洁，是你剥夺了我的贞洁！}还加了其他无耻女人用的不雅言辞。这诬蔑的策划者感到羞愧，所有参与此事的主教都脸红了。\par

这女人被带出法庭，但伟大的\Person{亚大纳削}抗议说，不应把她送走，而应审问她，并查明阴谋的始作俑者。于是，他的指控者大声呼喊说，他还犯了其他更可恶的罪行，他无论如何也无法用任何技巧或诡计洗清；而且，眼睛而非耳朵将根据证据作出判断。说完这些，他们拿出那只著名的盒子，露出那只涂了香料的手。看到这景象，所有旁观者都大声喊叫。有些人相信指控是真的；另一些人则确信是假的，认为\Person{阿尔塞尼}正藏匿在某处。最后，经过一些困难，稍微安静下来后，被告问他的审判者们是否有人认识\Person{阿尔塞尼}。其中几人回答说他很熟悉，\Person{亚大纳削}下令把他带到他面前。然后他又问他们：\Quote{这是真正的\Person{阿尔塞尼}吗？这就是我谋杀的那个人吗？这就是那些人谋杀后砍掉他右手的那个人吗？}当他们承认是同一个人时，\Person{亚大纳削}脱掉他的外衣，露出两只手，右手和左手，说：\Quote{不要有人寻找第三只手，因为人从\OriginalTerm{天主}{Creator}那里只领受了两只手，没有更多。}\par

即使在这明显的证据之后，诬蔑者和参与罪行的审判者，非但没有躲藏起来，或祈求大地裂开吞没他们，反而在集会中喧嚷骚动，宣称\Person{亚大纳削}是巫师，他用魔法咒语迷惑了人们的眼睛。刚才还指控他谋杀的人，现在却竭力要把他撕成碎片并谋杀他。但皇帝委托维持秩序的人救了\Person{亚大纳削}的命，把他拖走，并匆匆送上船。\par

当他出现在皇帝面前时，他描述了所有为了陷害他而编造的戏剧性阴谋。诬蔑者派遣了他们派系的主教到\Place{马勒奥提斯}去，即\Place{尼西亚}主教\Person{德奥格尼}、\Place{佩林图斯}主教\Person{德奥多罗}、\Place{加采东}主教\Person{玛里斯}、\Place{基利家}的\Person{纳尔奇苏}，以及有同样想法的其他人。\Place{马勒奥提斯}是\Place{亚历山大里亚}附近的一个地区，得名于\Place{玛利亚湖}。他们在那里编造了其他谎言，并伪造了审判报告，把已被证明为假的指控与新的指控混杂在一起，仿佛是真的，然后送交皇帝。\par

% structure: p0174 source=h2 target=subsection reason=relative-heading-rank; base=h2

\subsection{第二十九章 耶路撒冷圣堂奉献礼。——圣亚大纳削被流放}

所有出席提洛会议的主教，连同从各处来的其他主教，都奉皇帝之命前往\Place{埃利亚}，去奉献他在那里建造的圣堂。皇帝还派遣了许多性情最和善、以虔诚和忠诚著称的官员，命令他们提供充足的食物，不仅给主教及其随行人员，也给从各处涌到\Place{耶路撒冷}的广大民众。圣祭台用皇室挂饰和嵌有宝石的金器装饰。盛大的庆典结束后，每位主教返回自己的教区。皇帝得知庆典的辉煌壮丽后非常高兴，并祝福一切美善的根源，因为他如此应允了他的祈求。\par

\Person{亚大纳削}抱怨他被不公正定罪，皇帝命令被投诉的主教到庭。他们到达后，不再坚持以前的任何诬蔑，因为他们知道这些诬蔑很容易被驳倒；但他们却让皇帝相信\Person{亚大纳削}曾威胁要阻止粮食出口。皇帝相信了他们的话，把他流放到\Place{高卢}的一个叫\Place{特里尔}的城市。这件事发生在皇帝在位第三十年。\par

% structure: p0177 source=h2 target=subsection reason=relative-heading-rank; base=h2

\subsection{第三十章 真福君士坦丁皇帝的遗嘱}

一年零几个月后，皇帝在\Place{彼提尼亚}的\Place{尼科美底亚}城患病，知道人生无常，他接受了圣洗圣事，他曾打算推迟到在\Place{约旦河}受洗。\par

他将帝国王位的继承者留给了他的三个儿子：\Person{君士坦丁}、\Person{君士坦提乌斯}和最小的\Person{君士坦斯}。\par

他命令伟大的\Person{亚大纳削}返回\Place{亚历山大里亚}，并在\Person{欧瑟伯}面前表达了这一决定，\Person{欧瑟伯}竭力劝阻他。\par

% structure: p0181 source=h2 target=subsection reason=relative-heading-rank; base=h2

\subsection{第三十一章 为君士坦丁辩护}

不应感到惊奇，\Person{君士坦丁}竟如此受欺骗，以致将如此众多伟人流放：因为他相信了那些主教们的断言，他们位高望重，却巧妙地隐藏了恶意。凡熟悉圣经的人都知道，\Person{圣达味}虽为先知，也曾受欺骗；而且不是被司铎，而是被一个仆役、奴隶、无赖所骗。我指的是\Person{漆巴}，他通过谎报\Person{默菲波舍特}来欺骗了君王，从而获得了他的土地。我这样说并非要谴责先知，而是要为皇帝辩护，指出人性的软弱，并教导不应只相信提出控告的人，即使他们看似可信；另一方也应被听取，应留一只耳朵给被控者。\par

% structure: p0183 source=h2 target=subsection reason=relative-heading-rank; base=h2

\subsection{第32章 神圣皇帝君士坦丁的终末}

皇帝如今从他在世上的国度被转移到了一个更好的国度。\par

皇帝的身体被放入金棺，由各省总督、军事指挥官及其他国家官员运往\Place{君士坦丁堡}，全军前后簇拥，无不痛惜哀悼；因为\Person{君士坦丁}对众人如同慈父。皇帝的身体被允许留在皇宫，直到他的儿子们到来，并获得了极高的敬意。但此处无需详述这些细节，因其他作者已有详尽记述。从他们的作品（易于获取）中，可以得知万有主宰如何尊崇祂忠信的仆人。若有人受诱惑而怀疑，让他看看如今在\Person{君士坦丁}的坟墓和雕像旁发生的事，他必承认天主在圣经中所说的真理：\BibleQuote{\Emph{尊崇我的，我必尊崇他们；藐视我的，他们必受轻看。}}\par

\SourceRecord{Translated by Blomfield Jackson.；Nicene and Post-Nicene Fathers, Second Series；Vol. 3.；Edited by Philip Schaff and Henry Wace.；Buffalo, NY: Christian Literature Publishing Co.,；1892.；<http://www.newadvent.org/fathers/27021.htm>.}

% structure: p0001 source=h1 target=section reason=page-title

\section{\WorkTitle{教会史}（第二卷）}

% structure: p0002 source=h2 target=subsection reason=relative-heading-rank; base=h2

\subsection{第一章 \Person{圣亚大纳削}返回}

\Person{圣亚大纳削}返回了\Place{亚历山大}，此前他在\Place{特里尔}停留了两年四个月。\Person{君士坦丁}大帝的长子\Person{小君士坦丁}，其帝国统治范围延伸至\Place{西方高卢}，致函\Place{亚历山大}教会如下。\par

\Emph{君士坦丁大帝之子\Person{小君士坦丁}皇帝致\Place{亚历山大}人书。}\par

\Emph{\Person{君士坦丁}凯撒致\Place{亚历山大}天主教教会人民。}\par

\Quote{我想，你们虔诚的智慧不会不知道，可敬法律的阐释者\Person{亚大纳削}被适时地送往\Place{高卢}，是为了在这些嗜血对手的凶残威胁他神圣头颅之际，使他免受不可弥补的伤害。为了避开这迫在眉睫的危险，他从敌人的牙缝中被救出，留在我管辖的一座城里，在那里他可得到一切必需的供应。然而，他伟大的美德，依靠天主的恩宠，使他蔑视逆境的一切灾祸。我主、我父、有福记忆的\Person{君士坦丁}曾打算恢复他先前的牧职，并把他归还给你们的虔诚；但因皇帝在愿望未实现前被死亡之手止住，我作为他的继承人，认为适合执行这位神圣记忆的君主的意图。你们见到主教本人时，将得知我对他何等尊敬。我受此感召，实因他伟大的美德，以及你们对他归来的殷切渴望。愿天主眷顾看护你们，亲爱的弟兄！}\par

携此信函，\Person{圣亚大纳削}从流放地归来，无论富人穷人、城市居民还是各省百姓，都极其欢欣地欢迎他。只有追随\Person{亚略}疯狂的人对他的归来感到恼火。\Person{欧瑟伯}、\Person{特奥格尼斯}及其同党又故伎重演，试图在年轻皇帝面前中伤他。\par

我现在将叙述\Person{君士坦提乌斯}如何偏离了宗徒的教义。\par

% structure: p0009 source=h2 target=subsection reason=relative-heading-rank; base=h2

\subsection{第二章 \Person{君士坦提乌斯}皇帝偏离真道}

\Person{君士坦提娅}是\Person{李锡尼}的遗孀，\Person{君士坦丁}的同父异母姐妹。她与一位曾浸染\Person{亚略}教义的司铎过从甚密。他并未公开承认自己的谬误，但在与她频繁交谈时，毫不忌讳地宣称\Person{亚略}受到了不公的诽谤。她那不虔敬的丈夫死后，著名的\Person{君士坦丁}竭尽全力安慰她，并竭力使她不致遭受最悲伤的寡居考验。在她临终时，他也陪伴在侧，并给予她一切应有的关心。她便将我提到的这位司铎引荐给皇帝，恳求他接纳并保护他。\Person{君士坦丁}答应了她的请求，不久便履行了诺言。然而，尽管这位司铎享有最充分的言论自由，并得到最尊敬的待遇，他仍不敢暴露他腐化的原则，因为他观察到皇帝对真理的坚定持守。当\Person{君士坦丁}即将被接引到永恒国度时，他立下遗嘱，将其世俗领土分给诸子。他去世时无一子在侧，于是只将遗嘱托付给这位司铎，并要他交给\Person{君士坦提乌斯}，因他离该地较近，将先于其他兄弟到达。司铎执行了这些指示，由此将遗嘱交到\Person{君士坦提乌斯}手中，从而被他认识，并接纳为密友，命其常来见。\Person{君士坦提乌斯}见其心志软弱，有如风中芦苇，便大胆向福音教义宣战。他大声悲叹教会的动荡状态，断言这应归咎于那些在信经中引入不合圣经的\Quote{\OriginalTerm{同质}{consubstantial}}一词之人，并说所有教牧人员与平信徒之间的争端皆由此而起。他诬蔑\Person{亚大纳削}及所有与他意见一致的人，并谋划消灭他们，与他们同谋者还有\Person{欧瑟伯}、\Person{特奥格尼斯}和\Place{培林图斯}主教\Person{狄奥多若}。\par

最后这位主教，其教座通常被称为\Place{赫拉克利亚}，是一位博学之士，曾著有圣经注释。\par

这些主教住在皇帝附近，常去拜见；他们向皇帝保证，\Person{亚大纳削}从流放地归来已造成许多祸患，激起一场风暴，不仅震撼了\Place{埃及}，也波及\Place{巴勒斯坦}、\Place{腓尼基}及邻近地区。\par

% structure: p0013 source=h2 target=subsection reason=relative-heading-rank; base=h2

\subsection{第三章 \Person{圣亚大纳削}第二次流亡——\Person{格里高利}的晋牧与死亡}

主教们用这些及类似的论据攻击软弱的皇帝，说服他把阿塔纳削逐出他的教会。但阿塔纳削及时得知了他们的阴谋，便动身前往西方。欧瑟伯的朋友们向当时罗马的主教儒略呈上了虚假的指控。为遵守教会的法律，儒略传唤了控告者与被控告者到罗马，以便审理此案。阿塔纳削因此前往罗马，但诽谤者拒绝前往，因为他们看出自己的谎言容易被揭穿。但他们见阿塔纳削的羊群没有牧人，便派了一只豺狼而不是牧人去管理他们。这个名叫额我略的人，对羊群的残暴超过了野兽；但六年之后，他被羊群自己消灭了。阿塔纳削前往君士坦斯（君士坦丁，长兄，已阵亡）处，控诉亚略派对他设下的阴谋及他们对信仰的反对。他提醒君士坦斯其父的榜样：他亲自参加了自己召集的伟大而著名的会议，亲临辩论，参与制定法令，并以法律确认它们。皇帝被他父亲的热情所激励，迅速写信给他的兄弟，劝诫他务必完整保存他们所继承的父亲的宗教；\Quote{因为，}他敦促道，\Quote{他因虔诚而壮大了帝国，消灭了罗马的暴君，征服了四方的外邦民族。}君士坦提乌斯因此信而召集东西方的主教前往伊利里库姆的\Place{萨迪卡}，达西亚的首府，以便商议消除教会其他诸多紧迫的困扰。\par

% structure: p0015 source=h2 target=subsection reason=relative-heading-rank; base=h2

\subsection{第四章　君士坦丁堡主教保禄}

\Person{保禄}，君士坦丁堡主教，忠实地持守正统教义，被不健全的亚略派指控煽动叛乱及其他他们通常加于所有宣讲真虔诚者的罪行。人民担心他敌人的阴谋，不允许他去萨迪卡。亚略派利用皇帝的软弱，从皇帝那里获得了放逐保禄的谕令，保禄于是被送往\Place{库库苏斯}，一个小镇，原属卡帕多细亚，今属小亚美尼亚。但这些扰乱公共和平的人并不满足于将可敬的保禄驱逐到荒野。他们派遣了残忍的爪牙以暴力处死他。圣阿塔纳削在其为自己逃亡所写的辩护词中证实了这一事实。他用了如下的话：\Quote{他们追捕君士坦丁堡主教保禄，在卡帕多细亚的库库苏斯城抓住了他，将他勒死，行刑者是\OriginalTerm{总督}{prefect}菲利普，他是他们异端的保护者，是他们最凶残计划的积极执行者。}\par

这就是亚略亵渎之词所导致的谋杀。他们对独生子的疯狂愤怒，伴以对其仆人的残暴行径。\par

% structure: p0018 source=h2 target=subsection reason=relative-heading-rank; base=h2

\subsection{第五章　马其顿纽异端}

亚略派在造成保禄死亡——或者更确切地说，将他送往天国之后，推举马其顿纽接替他的位置。他们认为马其顿纽持相同观点，属于同一派别，因为他像他们一样亵渎圣神。但不久之后，他们也废黜了他，因为他拒绝称圣经所确认为天主圣子的那一位为受造物。在与他分离之后，他成了自己派别的领袖。他教导说，圣子与圣父并非同一实体，但他在每一方面都与圣父相似。他还公开肯定圣神是受造物。这些事件发生在我们叙述后不久。\par

% structure: p0020 source=h2 target=subsection reason=relative-heading-rank; base=h2

\subsection{第六章　在萨迪卡举行的会议}

根据古代记录，有二百五十位主教在\Place{萨迪卡}聚集。伟大的阿塔纳削、前述的加沙主教阿斯克勒帕斯（\Person{阿斯克勒帕斯}），以及安其拉（\Place{安其拉}，加拉提亚首府）主教玛尔凯鲁斯（\Person{玛尔凯鲁斯}，他在尼西亚议会时也担任此主教职），均赴会。诽谤者和亚略派首领们，他们先前审理了阿塔纳削的案件，也出席了。但当他们发现议会的成员坚定持守纯正教义时，竟拒绝进入议会，尽管他们被召请，反而逃跑了，无论控告者还是审判官都是如此。所有这些情况，在议会起草的一封信中得到了更清楚的解释，因此我现在将其插入。\par

\Emph{在萨迪卡聚集的主教们致其他主教的会议信函。}\par

在萨迪卡聚集的圣议会，来自罗马、西班牙、高卢、意大利、坎帕尼亚、卡拉布里亚、阿非利加、撒丁、潘诺尼亚、默西亚、达西亚、达尔达尼亚、小达西亚、马其顿、色萨利、亚该亚、伊庇鲁斯、色雷斯、罗多彼、亚细亚、卡里亚、俾提尼亚、赫勒斯滂、弗里吉亚、彼西底、卡帕多细亚、本都、小弗里吉亚、西里西亚、潘菲利亚、吕底亚、基克拉迪群岛、埃及、底比斯、利比亚、加拉提亚、巴勒斯坦和阿拉伯，致普世教会的主教们，我们在公教会和宗徒教会的同工，以及我们在主内亲爱的弟兄们。愿平安归于你们。\par

阿里乌斯派的疯狂常驱使他们向保存真信仰的天主仆人实施暴力暴行；他们自己引入假道理，迫害那些坚持正统原则的人。他们对信仰的攻击如此猛烈，以至于传到了我们最虔诚的皇帝们的耳中。藉着天主恩宠的合作，皇帝们从不同省份和城市召聚我们，到他们指定在\Place{萨尔迪卡}城举行的神圣大公会议，为的是终结所有纷争，驱逐所有邪恶教义，并在万民中唯独保持基督的宗教。一些东方的主教应我们最虔诚的皇帝们的请求参加了会议，主要是由于针对我们可爱的弟兄和同工（\Person{亚大纳削}，亚历山大里亚主教；\Person{玛尔凯路}，\Place{加拉西亚}\Place{安基拉}主教；\Person{阿斯克肋帕}，\Place{加萨}主教）的传言。也许阿里乌斯派的诽谤已传到你们那里，他们企图以此抢在会议之前，使你们相信他们对无辜者毫无根据的指控，并阻止任何人对他们所坚持的堕落异端产生怀疑。但他们并未被允许长久如此行事。主是教会的保护者；为了教会和我们众人，祂忍受了死亡，并为我们敞开了通往天堂的道路。\par

\Person{欧瑟伯}、\Person{玛里斯}、\Person{狄奥多若}、\Person{特奥尼各}、\Person{乌尔撒基}、\Person{瓦伦斯}、\Person{默诺凡特}和\Person{斯德望}的同党早已写信给\Place{罗马}主教、我们的同工\Person{尤利乌斯}，以反对我们前述的同工：\Place{亚历山大里亚}主教\Person{亚大纳削}、\Place{加拉西亚}\Place{安基拉}主教\Person{玛尔凯路}和\Place{加萨}主教\Person{阿斯克肋帕}。对立派的一些主教也写信给\Person{尤利乌斯}，为\Person{亚大纳削}的无辜作证，并证明\Person{欧瑟伯}追随者所断言的一切不过是谎言和诽谤。阿里乌斯派拒绝听从我们可爱的弟兄和同工\Person{尤利乌斯}的召唤，以及该主教所写的信，清楚证明了他们指控的虚假。因为，如果他们相信他们对我们同工所做和所控告的事是情有可原的，他们就会去\Place{罗马}。但他们在这次伟大而神圣的大公会议中的行为方式，是他们欺骗的明显证据。他们到达\Place{萨尔迪卡}后，发现我们的弟兄\Person{亚大纳削}、\Person{玛尔凯路}、\Person{阿斯克肋帕}等人也在那里；因此他们害怕面对考验，尽管他们不是一次两次，而是多次被召唤。聚集的主教们在那里等候他们，尤其是可敬的\Person{奥西乌斯}，一位因高龄、坚守信仰和为教会劳苦而配得一切尊荣和尊敬的人。所有人都敦促他们加入集会，并利用机会在同工面前证明他们缺席时以口头和书信提出指控的真实性。但他们拒绝听从召唤，正如我们已经说过的那样，因此他们的过激行为证明了他们的陈述虚假，几乎公开宣布了他们所设下的阴谋和诡计。对自己的主张充满信心的人总是准备公开坚持他们。但由于这些控告者不出现证实他们所提出的指控，他们将来藉由惯用伎俩对我们同工提出的任何指控，都将只被视为出于在缺席时诽谤他们的愿望，而没有勇气公开面对他们。\par

他们逃跑了，可爱的弟兄们，不仅因为他们的指控是诽谤，而且因为他们看到有人带着针对他们自己严重而多重的指控来到。锁链和镣铐被展示出来。有些他们曾放逐的人在场；另一些人作为仍被流放者的代表前来。那里站着被他们处死者的亲属和朋友。最严重的是，主教们也出现了，其中一人展示了他们加在他身上的铁链和锁链。其他人作证说死亡紧随他们的虚假指控。因为他们的迷狂已使他们甚至企图杀害一位主教；若不是他从他们手中逃脱，他早已被杀。我们的同工、有福记忆的\Person{特奥杜禄}，带着他们对他的诽谤离世；因为通过诽谤，他被判处死刑。一些人展示了他们所受的刀剑之伤；另一些人作证他们曾遭受饥饿之苦。\par

所有这些供述不是由少数无名之辈作出的，而是由整个教会作出的；这些教会的司铎作证，迫害者曾用刀剑武装士兵对付他们，并用棍棒武装平民；使用了司法威胁，并制造了伪造文件。\Person{特奥尼各}为存心使皇帝对我们同工\Person{亚大纳削}、\Person{玛尔凯路}和\Person{阿斯克肋帕}有偏见而写的信，被曾当过\Person{特奥尼各}执事的人宣读并证实。还证明他们剥光贞女的衣服，焚烧教堂，监禁我们的同工，这一切都是因为阿里乌斯疯狂派的邪恶异端。因为所有拒绝与他们同流合污的人都受到如此对待。\par

意识到犯下所有这些罪行，他们处于极大的困境中。因对其行为感到羞愧，这些行为已无法再隐藏，他们前往\Place{萨尔迪卡}，以为斗胆前往那里会消除所有对他们罪行的怀疑。但当他们看到那些曾被他们诬告的人，以及那些曾受其残酷迫害的人在场，并且也看到几个人带着不可辩驳的指控前来时，他们拒绝进入大公会议。另一方面，我们的同工\Person{亚大纳削}、\Person{玛尔凯路}和\Person{阿斯克肋帕}用尽一切方法劝说他们参加，藉着泪水、紧迫、挑战，承诺不仅要证明他们指控的虚假，还要表明他们严重伤害了自己的教会。但他们被对自己恶行的意识压倒了，以至于逃跑了，通过这次逃跑清楚证明了他们指控的虚假以及他们自己的罪责。\par

然而，虽然他们的诬蔑和背信从一开始就显而易见，如今更被清楚察觉，但我们仍决定按照真理的法则审查此事的情况，免得他们因自己的逃跑而获取重新行骗的借口。\par

在实施这一决议后，我们通过他们的行为证明他们是诬告者，并对我们的同工司铎设下阴谋。\Person{亚尔塞尼}，他们宣称被\Person{亚大纳削}处死，却仍然活着，并在活人之列。仅此事实便足以表明他们的其他指控是虚假的。\par

尽管他们到处散布谣言，说一个圣爵被\Person{亚大纳削}的一位司铎\Person{玛卡里乌斯}打碎，但从\Place{亚历山大里亚}、\Place{玛雷奥提斯}和其他地方来的人都作证说并无此事；\Place{埃及}的主教们也写信给我们同工\Person{尤利乌斯}，声称绝无怀疑曾有这种行为。亚略派声称他们掌握的针对\Person{玛卡里乌斯}的司法事实全由一方编造；在这些文件中，包含了外教人和望教者的证词。其中一位望教者在被询问时回答，他在\Person{玛卡里乌斯}进入教堂时在场。另一个人作证说，他们谈论甚多的\Person{伊斯喀拉斯}当时正卧病在隐修小室。由此可见，当时不可能举行圣事，因为望教者在场，而\Person{伊斯喀拉斯}又缺席；他那时正因病被禁足。那个恶人\Person{伊斯喀拉斯}曾虚假地宣称\Person{亚大纳削}焚烧了一些圣经，并已被定罪，如今却承认当\Person{玛卡里乌斯}来到时他卧病在床；因此他的诬告明显是虚假的。然而他的团体却奖赏了他的诬蔑；他们授予他主教之衔，尽管他连司铎都还不是。有两位从前属于\Person{梅里提乌}、后来被真福\Person{亚历山大}主教接纳、如今与\Person{亚大纳削}同在的司铎前来会议，声明他从未被祝圣为司铎，\Person{梅里提乌}在\Place{玛雷奥提斯}也从未拥有任何教堂或任用任何神职人员。然而，尽管他从未被祝圣为司铎，他们却提升他为主教，为使其头衔能迷惑那些听信他虚假指控的人。\par

我们同工\Person{玛尔塞鲁}的著作也被宣读，清楚显明了\Person{尤西比乌}追随者的两面性；因为\Person{玛尔塞鲁}仅作为疑问提出的问题，他们却指控他当作信德来主张。他在调查前后所作的陈述都被宣读，他的信德被证明是正统的。他并不像他们所说的那样，主张天主的圣言始于因童贞玛利亚的受孕，或其王国将有终结。相反，他写道他的王国无始无终。我们同工\Person{阿斯克勒帕斯}提出了在\Place{安提约基雅}当着原告和\Place{凯撒勒雅}主教\Person{尤西比乌}的面所作的记录，并通过担任审判员的主教们的判决证明了自己的清白。\par

因此，亲爱的弟兄们，他们虽屡次被召却不愿出席会议，并非无缘无故；他们逃跑也并非无缘无故。良心的责备迫使他们逃走，从而同时证明了他们诬蔑的虚妄，以及针对他们提出并证实的指控的真实。除了所有其他控诉理由外，还可补充的是：所有被指控持有亚略异端并因此被开除的人，不仅被他们接纳，而且被提升至最高职位。他们将执事提升为司铎，再从司铎提升为主教；在这一切中，他们被别的动机驱使，唯愿传播和扩散他们的异端，并败坏真信仰。\par

继\Person{尤西比乌}之后，他们的主要领袖是：\Place{赫拉克勒亚}主教\Person{德奥多鲁}、\Place{基里基雅的尼禄尼亚}主教\Person{纳尔基苏斯}、\Place{安提约基雅}主教\Person{斯德望}、\Place{劳迪基雅}主教\Person{乔治}、\Place{巴勒斯坦凯撒勒雅}主教\Person{阿卡基乌}、\Place{亚细亚的厄弗所}主教\Person{梅诺凡图斯}、\Place{默西亚的辛吉杜努姆}主教\Person{乌尔萨基乌}、\Place{潘诺尼亚的穆尔萨}主教\Person{瓦伦斯}。这些主教禁止随他们从东方来的人参加神圣会议，或与天主的教会合一。在前往\Place{萨尔迪卡}的路上，他们在不同地方举行私下集会，并以威胁订立契约，约定一旦到达\Place{萨尔迪卡}，就不参加神圣会议，也不参与其讨论；他们安排一到就仅作形式上的露面，然后立即逃跑。这些事实是由与我们同工的\Place{巴勒斯坦}主教\Person{玛卡里乌斯}和\Place{阿拉伯}主教\Person{阿斯特里乌}告知我们的，他们随同前来\Place{萨尔迪卡}，但拒绝分享他们的不正统。这些主教在神圣会议前控诉他们受到的粗暴对待，以及他们一切交易中缺乏正确原则。他们补充说，在他们中间仍有许多人持正统意见，但这些人被阻止参加会议；有时威胁，有时许诺，被用来使他们留在那一派中。为此原因，他们被迫同住一屋，甚至片刻也不得独处。\par

缄默不语、不加谴责地放过诽谤、监禁、谋杀、鞭打、伪造信件、侮辱、剥光童贞女衣服、放逐、摧毁教堂、纵火行为、将主教从小城调到大教区，及尤其是藉他们的手段兴起反对真信德的不幸的\OriginalTerm{亚略异端}{Arian heresy}，是不对的。为此，我们宣告我们亲爱的弟兄和同仁，\Person{亚他纳削}（\Place{亚历山大里亚}主教）、\Person{玛尔克鲁斯}（\Place{迦拉达}省\Place{安其拉}主教）、\Person{阿斯克肋帕斯}（\Place{加沙}主教），以及与他们一起的所有其他天主的仆人，是清白与纯洁的；我们写信给他们各教区，为使各教会的人知道他们各自主教的清白，只承认他们，等待他们归来。对于像狼一样扑向教会的人，如\Place{亚历山大里亚}的\Person{额我略}、\Place{安其拉}的\Person{巴西略}、\Place{加沙}的\Person{昆提安}，我们命令甚至不可称他们为主教，也不可称他们为基督徒，不可与他们有任何共融，不可收受他们的书信，也不可写信给他们。\par

\Person{狄奥多若}（\Place{欧洲}的\Place{赫拉克莱亚}主教）、\Person{纳尔奇苏斯}（\Place{基里基雅}的\Place{尼罗尼亚斯}主教）、\Person{阿卡基乌斯}（\Place{巴勒斯坦}的\Place{凯撒勒雅}主教）、\Person{斯德望}（\Place{安提约基雅}主教）、\Person{乌尔萨基乌斯}（\Place{莫西亚}的\Place{辛吉杜努姆}主教）、\Person{瓦伦斯}（\Place{潘诺尼亚}的\Place{穆尔萨}主教）、\Person{梅诺方图斯}（\Place{厄弗所}主教）和\Person{乔治}（\Place{劳狄刻雅}主教；因为虽然恐惧阻止他离开东方，但他已被有福的\Place{亚历山大里亚}主教\Person{亚历山大}罢免，并汲取了\OriginalTerm{亚略派}{Arians}的迷妄），由于他们的各种罪行，被神圣会议的一致决定从他们的主教职中逐出。我们规定，他们不仅不得被视为主教，而且应被拒绝与我们共融。因为那些将圣子与圣父的本质和神性分开，并将圣言与圣父疏离的人，应与天主教会分离，并与所有称为基督徒的人疏远。愿他们被你们和所有信徒弃绝，因为他们败坏了真理之言。因为宗徒的诫命吩咐：若有人给你们宣讲另一个福音，与你们所接受的不同，\BibleQuote{他应受诅咒}。命令无人与他们共融；因为\BibleQuote{光明与黑暗不能相通}。远离他们；因为\BibleQuote{基督与\OriginalTerm{贝里雅耳}{Belial}怎能相和？}亲爱的弟兄们，你们要谨慎，不可写信给他们，也不可收受他们的书信。亲爱的弟兄和同仁，你们要努力，如同在精神上与我们一同出席会议，衷心地赞同所颁布的法令，并附上你们的亲笔签名，以保持全世界所有同仁的意见一致。\par

我们宣布以下这些人被绝罚于天主教会：他们声称基督是天主，但不是真天主；祂是子，但不是真子；祂既是\OriginalTerm{受生}{begotten}又是\OriginalTerm{受造}{made}；因为这些人承认，他们通过\Quote{受生}这个术语理解的是受造之物；并且，尽管天主子存在于万世之前，他们却将不存在于时间中、而是在一切时间之前的祂，归给一个开始和一个终结。\par

\Person{瓦伦斯}和\Person{乌尔萨基乌斯}如同毒蛇所生的两条蝮蛇，出自\OriginalTerm{亚略异端}{Arian heresy}。因为他们夸耀地宣称自己是毋庸置疑的基督徒，然而却断言圣言和圣神都被钉死杀害，并且它们死了又复活；并且他们执拗地坚持，如同异端者那样，圣父、圣子、圣神有不同且各别的\OriginalTerm{本质}{essences}。我们被教导并持守大公宗徒的传统、信德和宣认，它们教导：圣父、圣子、圣神有同一本质，异端者称之为\OriginalTerm{实体}{substance}。如果有人问：\Quote{圣子的本质是什么？}我们宣认，那是被承认为唯独圣父的本质；因为圣父从未、也绝不能没有圣子而存在，圣子也不能没有圣父。主张圣父曾没有圣子而存在是极其荒谬的，因为圣子亲自证明了这绝不可能是如此，祂说：\BibleQuote{我在父内，父在我内}；又说：\BibleQuote{我与父原是一体}。我们中无人否认祂是\OriginalTerm{受生}{begotten}的；但我们说祂是在万有之前受生的，无论有形无形；并且祂是总领天使和天使、世界及人类的创造者。经上记载：\BibleQuote{智慧，万物的创造者，教导了我}；又说：\BibleQuote{万物是藉着祂造成的}。\par

如果祂有过开始，祂就不能一直存在，因为永存的圣言没有开始，天主永无终结。我们不说父是子，也不说子是父；而是父是父，父的子是子。我们宣认圣子是圣父的\OriginalTerm{大能}{Power}。我们宣认圣言是天主圣父的圣言，并且除祂以外没有别的。我们相信圣言是真天主，又是智慧和\OriginalTerm{大能}{Power}。我们肯定祂是真子，但不像其他人被称为子的方式：因为他们或因重生而成为神，或因功绩而被称为天主子，并非因与父和子同一本质，如同父和子的情形那样。我们宣认一位\OriginalTerm{独生子}{Only-begotten}和一位\OriginalTerm{首生者}{Firstborn}；但圣言是独生子，祂一直都在父内，并常在父内。我们使用\Quote{首生者}一词与祂的人性有关。但祂在（复活）的新创造中超越（人），因为祂是死者中的首生者。\par

我们宣认天主存在；我们宣认圣父与圣子的天主性是一。没有人否认圣父大于圣子：这不是因为本质不同，也不是因为差异，而仅仅是因为圣父之名大于圣子之名。我们主所说的这句话：\BibleQuote{我与父原是一体}，被那些人解释为指父与子之间的和谐一致；但这是亵渎而乖僻的解释。我们作为天主教徒，一致谴责这种愚昧可悲的意见：因为正如有死的世人之间起了争端又复和好，这些解释者竟说全能天主父与祂的圣子之间也会发生争论与分歧；这种假设完全是荒谬而不能成立的。但我们相信并坚持，这些圣言\BibleQuote{我与父原是一体}所指出的，是圣父与圣子中同一本质的合一。\par

我们也相信圣子与圣父一同为王，祂的国度无始无终，不受时间限制，也永不止息：因为那常存者从未开始存在，也永不会止息。\par

我们相信并接纳圣神——护慰者，就是主所应许并派遣的那一位。我们相信祂是被派遣的。\par

受苦难的并不是圣神，而是祂所披戴的人性；这人性是祂从童贞玛利亚所取的，既是人性便能受苦；因为人是可朽的，而天主是不朽的。我们相信祂第三日复活了，是天主内的人，不是人内的天主；并且祂将那个从罪与腐朽中解救出来的人性作为恩赐献给了祂的父。\par

我们相信，在适当而固定的时刻，祂将要亲自审判万民及其一切行为。\par

\Quote{我们所提及的那些人，其无知和心灵的黑暗如此之深，以至于无法看见真理之光。他们不能理解这些话的含义：\InnerQuote{叫他们也在我们内合而为一}。显然为什么用了\InnerQuote{一}这个词；这是因为宗徒们领受了天主的圣神，然而他们中没有一个人是圣神，也没有一个人是圣言、智慧、大能或独生子。\InnerQuote{如同你}，祂说，\InnerQuote{和我是一体，叫他们也在我们内合而为一}。这些圣言\InnerQuote{叫他们也在我们内合而为一}是精确的：因为主并没有说\InnerQuote{像我同父一样是一体}，而是说\InnerQuote{门徒们彼此联结合一，可以在信德和宣认中合而为一，并因此在天主圣父的恩宠与虔敬中，以及因我们的主耶稣基督的宽仁与慈爱，能够成为一}。}\par

从这封信中可以看出诽谤者的虚伪、先前法官的不公正，以及各项法令的正确性。这些圣教父不仅教导了我们关于天主性体的真理，也教导了降生成人的教义。\par

君士坦斯得知他兄弟的优柔寡断后深感忧虑，对那些策划并诡计多端地利用了这一阴谋的人极为愤怒。他挑选了两位曾参加萨迪卡会议的主教，派他们带着书信去见他的兄弟；他还派遣了一位以虔诚和正直著称的军事指挥官萨利安努斯一同出使。他所交付的书信，与自己相称，不仅包含恳求和劝告，还带有威胁。首先，他要求兄弟注意主教们所说的一切，并查清斯特凡努斯及其同谋的罪行。他还要求他将阿塔纳修斯归还给他的教友；因为控告者的诽谤以及先前法官的不公与恶意已昭然若揭。他补充道，如果兄弟不肯应允他的请求执行这一正义之举，他将亲自前往亚历山大，将阿塔纳修斯归还给渴望他的教友，并驱逐所有反对者。\par

君士坦提乌斯接到这封信时正在安提阿；他同意执行他兄弟所命令的一切。\par

% structure: p0049 source=h2 target=subsection reason=relative-heading-rank; base=h2

\subsection{第七章  关于主教尤弗拉塔斯和文森提乌斯的叙述，以及在安提阿针对他们的阴谋}

真理的一贯反对者对这些事极为不满，以至于策划了一桩臭名昭著的邪恶罪行。\par

两位主教住在山脚附近，而那位军事指挥官则住在另一处的寓所里。\par

此时，\Person{斯德望}执掌\Place{安提约基雅}教会的船舵，几乎使船沉没，因为他利用几个工具进行专横的行为，并借他们的帮助使所有持守正统教义的人陷入重重灾难。这些工具的首领是一个鲁莽轻率的青年，过着极其无耻的生活。他不仅从市场拖走人，殴打侮辱他们，还胆敢闯入私宅，从那里带走无可指摘的男女。但我不愿过于冗长地叙述他的罪行，只讲述他对待主教们的狂妄行为；仅此足以让人了解他对市民施行的非法暴力。他去找城里最下贱的一个妓女，告诉她说有几个陌生人刚到，想与她过夜。他带了十五个同伙，把他们藏在山脚下石墙中，然后去找那妓女。约定信号后，得知同谋者在场，他便走到主教们住宿的客栈院门。一个被收买的仆人开了门。他带那女人进屋，指给她一个主教睡觉的房门，叫她进去。然后他出去叫同伙。他所指的门恰好是年长主教\Person{尤弗拉塔斯}的，他的房间是外面的一间。另一位主教\Person{味增爵}住在里间。那女人进入\Person{尤弗拉塔斯}的房间时，他听到脚步声，当时天黑，便问是谁。她回答，\Person{尤弗拉塔斯}大为惊恐，以为是个魔鬼模仿女人的声音，就呼求\Person{基督}救主相助。与此同时，\Person{奥纳格}——这是那恶党首领之名（这名字对他特别合适，因为他不仅用手，也用脚攻击虔诚之人）——带着他的不法党羽回来，将那些本该审判罪行的人指控为罪犯。吵闹声惊动了所有仆人跑来，\Person{味增爵}也起来了。他们关闭院门，抓住了七个同伙；但\Person{奥纳格}和其余人逃脱了。那女人与被抓的人一起被扣押。天亮时，主教们叫醒了与他们同来的军官，三人一同前往皇宫，控告\Person{斯德望}的狂妄行为，说他的恶行显而易见，无需审判或拷打即可证明。将军大声向皇帝要求，这狂妄行为不应以教务会议处理，而应按普通法律程序，并提议先审讯主教们的随行神职人员，且声明\Person{斯德望}的爪牙也必须受拷打。\Person{斯德望}傲慢地反对，声称神职人员不应受鞭刑。皇帝和主要官员于是决定在皇宫审理此案。首先审问那女人，问她是谁带她去主教们住宿的客栈。她回答说，一个年轻人来找她，告诉她有陌生人想与她作伴；晚上他带她到客栈；他去寻找他的同伙，找到后，带她穿过院门，叫她进入前厅隔壁的房间。她又说，主教问是谁；他惊慌失措，开始祈祷；然后其他人跑来了。\par

% structure: p0053 source=h2 target=subsection reason=relative-heading-rank; base=h2

\subsection{\Emph{第八章 斯德望被罢黜}}

审判官听完这些供词后，命令将被捕者中最年轻的一个带上来。在受鞭刑审问之前，他供出了整个阴谋，并说这是由\Person{奥纳格}策划执行的。\Person{奥纳格}被带进来后，声称他完全是奉\Person{斯德望}之命行事。\Person{斯德望}的罪责就此显明，在场的主教们受命罢黜他，并将他逐出教会。然而，他的被逐并未使教会完全摆脱亚略异端的瘟疫。继他担任首长的\Person{莱昂提乌斯}是\Place{弗吕家}人，性情极其狡猾诡诈，可说如同海中的暗礁。我们将在后面进一步叙述他的事。\par

% structure: p0055 source=h2 target=subsection reason=relative-heading-rank; base=h2

\subsection{\Emph{第九章 圣亚大纳削的第二次返回}}

皇帝\Person{君士坦提乌斯}得知针对主教们的阴谋后，一而再、再而三地写信给伟大的\Person{亚大纳削}，劝勉他从西方回来。我在此插入第二封信，因为它是三封信中最短的一封。\par

\Signature{君士坦提乌斯·奥古斯都，征服者，致亚大纳削。}\par

\Quote{虽然我已在前几封信中通知你，你可以不受骚扰地回到我们的宫廷，以便你能如我热切期望的那样恢复你的主教职位，但我现在再次给你尊下另寄一封信，劝你毫不恐惧或猜疑地立刻乘坐公车回到我们这里，好使你得到你所渴望的一切。}\par

当\Person{亚大纳削}返回时，\Person{君士坦提乌斯}和善地接待了他，命他回到\Place{亚历山大教会}。但宫廷中有一些人受\OriginalTerm{亚略主义}{Arianism}错误的感染，声称\Person{亚大纳削}应当让出一座教堂给那些不愿与他共融的人。此事向皇帝提及，皇帝又向\Person{亚大纳削}提及，他就说，皇帝的命令似乎是公正的；但他也想提出一个请求。皇帝欣然应允只要他提出什么就一定赐予，他就说，在\Place{安提约基}，那些反对与如今拥有教堂的派别共融的人，需要祈祷的殿宇，故将一座天主的殿宇也分配给他们才是公平的。皇帝认为这个请求公正合理；但亚略派的首领们抵制执行，声称双方都不应得到教堂。\Person{君士坦提乌斯}因而对\Person{亚大纳削}极为钦佩，便送他回\Place{亚历山大}。\Person{额我略}已去世，死于\Place{亚历山大}居民之手。百姓为他们的牧者举行盛大庆典；设宴庆祝再度见到他的喜悦，并将赞美归于天主。不久之后，\Person{君士坦斯}辞世。\par

% structure: p0060 source=h2 target=subsection reason=relative-heading-rank; base=h2

\subsection{第十章 \Person{亚大纳削}的第三次流亡与逃亡}

那些对\Person{君士坦提乌斯}的思想取得完全支配、并随意影响他的人，提醒他说，\Person{亚大纳削}曾是他与兄弟之间分歧的起因，这几乎导致自然纽带的破裂，并引发了内战。\Person{君士坦提乌斯}被这些陈述所打动，不仅流放了圣\Person{亚大纳削}，而且判处他死刑；于是他派遣一位军事指挥官\Person{塞巴斯蒂安}，率领一支非常庞大的军队去杀害他，如同对待一名罪犯。一方如何发动攻击，另一方如何逃脱，最好用那位如此受苦并如此奇妙得救者的话语来讲述。\par

于是\Person{亚大纳削}在他的\WorkTitle{为逃亡申辩}中写道：\Quote{但愿他们查考我退避的详情，收集对立派别的证言；因为\OriginalTerm{亚略派}{Arians}曾陪同士兵，既为了激励他们，也为了向不认识我的人指认我。如果他们不为我所叙之事所动，至少让他们羞愧地沉默聆听。那是夜间，一些会众正在守夜，因为预计有一次共融礼。突然一队士兵向他们袭来，由一名将军和五千名武装人员组成，手执出鞘的剑、弓箭和棍棒，正如我已述说过的。将军包围了教堂，以密集队形布阵，防止里面的人出去。我认为在这样的混乱时刻，我不该离开会众，而应首先面对危险；于是我在宝座上坐下，命执事诵读一首圣咏，并命会众回应，\BibleQuote{因为他的仁慈永远常存}。随后我命他们各自回家。但此时将军率兵闯入教堂，包围了圣所，要逮捕我。留下来的圣职人员和平信徒喧嚷地恳请我离开。我坚决拒绝，直到所有其他人退去。我起身，做了祈祷，而后指示所有会众退下。\InnerQuote{我独自面对危险，比你们任何人受伤更好}。当大部分会众离开教堂，其余人正跟随离开时，那些留着的隐修士和部分圣职人员上前把我拉出来。就这样——让真理为我作证——主引领和保护我，我们穿过了士兵中间，他们有的驻扎在圣所周围，有的在教堂里巡逻。于是我不被察觉地走出教堂，并热切感谢天主，我没有抛弃会众，而是在他们安全离开后，得以从那些寻索我性命的人手中逃脱。}\par

% structure: p0063 source=h2 target=subsection reason=relative-heading-rank; base=h2

\subsection{第十一章 \Person{乔治}在\Place{亚历山大}所做的邪恶而大胆的作为}

\Person{亚大纳削}就这样逃脱了对手血迹斑斑的双手，被托付管理羊群的\Person{乔治}，实在也是一匹狼。他对待羊群的残忍超过了狼、熊或豹所能表现的。他迫使发愿永守童贞的年轻妇女不仅否认与\Person{亚大纳削}的共融，而且诅咒教父们的信德。他残忍行动的代理人是指挥军队的军官\Person{塞巴斯蒂安}。他命人在城中心点燃一堆火，将被剥光衣服的贞女靠近火堆，命她们否认信德。虽然她们在信徒与不信者面前构成一幅最令人悲伤和怜悯的景象，但她们认为所有这些羞辱给她们带来了最高的荣耀；她们喜乐地接受了因信德而施与她们的打击。所有这些事实将由她们的牧者更清楚地叙述。\par

在四旬期期间，\Person{格奥尔基乌斯}从\Place{卡帕多细亚}返回，并加重了我们的敌人教给他的恶行。复活节周过后，童贞女被投入监狱，主教们被士兵捆绑拖走，寡妇和孤儿的家园遭洗劫，抢劫和暴力接家挨户进行，基督徒在黑夜中被抓住并从家中拖出。许多房屋被贴了封条。神职人员的弟兄因弟兄的缘故而身处险境。这些暴行极其残忍，但随后发生的更为可怕。在圣神降临节的神圣庆节之后的一周，守斋的人们出来到墓地祈祷，因为他们都拒绝与\Person{格奥尔基乌斯}有任何共融。这个最卑劣的人得知此事后，便煽动军事指挥官\Person{塞巴斯蒂亚努斯}（一个\OriginalTerm{摩尼教徒}{Manichean}）攻击民众；于是，在主日当天，他率领一大队武装士兵，手持出鞘的剑、弓和箭，冲向人群。他发现只有少数基督徒在祈祷，因为大多数人都因时间已晚而离开了。然后，他做出了一个听从这类教师的人所会做的事。他下令点燃大火，把童贞女们带到火旁，试图强迫她们宣称自己信奉\OriginalTerm{亚略信经}{Arian creed}。当他看到她们战胜了，并不理会火时，便下令将她们剥光衣服，鞭打至她们的面孔久久难以辨认。然后他抓住四十个人，对他们施以一种新的酷刑。他下令用带刺的棕榈枝鞭打他们；他们的肉被割裂，以至于有些人因刺深深嵌入而不得不多次接受手术；另一些人无法忍受疼痛而死。所有幸存者，连同那些童贞女，随后被流放到\Place{大绿洲}。他们甚至拒绝将死者的尸体交给亲属埋葬，而是随意抛尸，不加掩埋，藏起来，以显得他们与这些残忍行为无关，对此一无所知。但他们在这愚蠢的期望上受了骗：因为死者的亲友虽为死者的忠信而喜乐，却为尸体的损失而深深悲痛，并广为传播所施暴行的全部情况。\par

\Quote{以下主教被从\Place{埃及}和\Place{利比亚}放逐：\Person{阿摩尼乌斯}、\Person{穆伊乌斯}、\Person{卡尤斯}、\Person{斐洛}、\Person{赫尔墨斯}、\Person{普莱尼乌斯}、\Person{普西诺西里斯}、\Person{尼拉蒙}、\Person{阿加皮乌斯}、\Person{阿纳甘弗斯}、\Person{马尔库斯}、\Person{德拉孔提乌斯}、\Person{阿德尔菲乌斯}、另一位\Person{阿摩尼乌斯}、另一位\Person{马尔库斯}和\Person{雅典诺多鲁斯}；以及司铎\Person{希拉克斯}和\Person{狄奥斯库鲁斯}。所有这些人都被如此残酷地驱逐流放，以致许多人在路上死去，另一些人在流放地死去。迫害者致使三十多位主教丧命。因为，像\Person{阿哈布}一样，他们的心思是要根除真理，倘若可能的话。}\par

\Person{阿塔纳修斯}在致那些遭受如此野蛮对待的童贞女的一封信中，使用了以下言辞：\Quote{你们中不要有人因这些不虔敬的异端者拒绝为你们埋葬、阻止你们的尸体被抬出而悲伤。\OriginalTerm{亚略派}{Arians}的不虔敬已达到如此程度，以至于他们堵塞大门，像许多魔鬼一样坐在坟墓周围，以阻止死者被埋葬。}\par

这些以及许多其他类似的暴行，是\Person{格奥尔基乌斯}在\Place{亚历山大里亚}所行的。\par

圣\Person{阿塔纳修斯}很清楚，没有任何地方可被视为他的安全之所；因为皇帝承诺，无论谁将他活捉，或献上他的首级作为他死亡的证据，都将获得非常丰厚的赏报。\par

% structure: p0070 source=h2 target=subsection reason=relative-heading-rank; base=h2

\subsection{第十二章 米兰公会议}

\Person{君士坦斯}死后，\Person{马格嫩提乌斯}掌握了两方帝国的最高权力；为了镇压他的篡位，\Person{君士坦提乌斯}前往欧洲。但这场战争尽管激烈，并未结束对教会的战争。\Person{君士坦提乌斯}接受了\OriginalTerm{亚略教义}{Arian tenets}，并轻易屈服于他人的影响，他被说服在\Place{意大利}的\Place{米兰}召开一次公会议，首先迫使所有与会主教签署由\Place{提洛}的不公正审判官所颁布的废黜令；然后，既然\Person{阿塔纳修斯}已被逐出教会，再制定另一份\OriginalTerm{信德宣认}{confession of faith}。主教们在接到皇帝的谕函后聚集开会，但他们远未按其指示行事。相反，他们当面告诉皇帝，他所命令的是不义和不虔敬的。因为这勇敢的行为，他们被逐出教会，并被流放到帝国的最偏远边界。\par

可敬的\Person{阿塔纳修斯}在他的\WorkTitle{辩护辞}中这样提到这件事：——\Quote{谁，}他写道，\Quote{能叙述他们犯下的如此暴行？不久前，当各教会享有和平，当人们聚集祈祷时，\Person{利贝里乌斯}（罗马主教）、\Person{保利努斯}（高卢都城主教）、\Person{狄奥尼西乌斯}（意大利都城主教）、\Person{路西斐鲁斯}（撒丁岛诸岛都城主教）和\Person{欧塞比乌斯}（意大利某城主教），他们都是杰出的主教和真理的宣讲者，被逮捕并流放，唯一的原因是他们不能同意\OriginalTerm{亚略异端}{Arian heresy}，也不能签署针对我们捏造的虚假指控。我不必谈论伟大的\Person{奥西乌斯}，那位年迈而忠信的信德宣认者，因为每个人都知道他也被放逐了。在所有主教中，他是最杰出的。有哪次公会议他没有主持，并以其推理的力量说服所有在场者？有哪间教会不仍然保留着他的保护的光荣纪念？有谁曾忧愁地到他那里去，而没有欢喜地离开？有谁曾请求他的帮助，而没有得到他所渴望的一切？然而，他们竟敢攻击这位伟人，仅仅因为他深知他们诽谤的不虔敬，拒绝在那针对我们的狡猾指控上签名。}\par

由上所述，可见亚略派人士对这些圣人的暴行。亚大纳削在同一书中也记述了亚略派首领针对许多其他人所设的众多阴谋：——\Quote{有谁被他们迫害并落入他们手中后，能不遭受他们任意施加的伤害而逃脱？有谁是他们搜寻的对象，被他们找到后，不被处以最痛苦的死亡，或至少被断肢残体？法官所判的刑罚都应归咎于这些异端者；因为法官不过是他们意愿和恶意的工具。何处没有一个地方不留下他们暴行的纪念？若有人与他们的意见不同，他们岂不像依则贝尔一样，诬告并压迫他？何处有教会不因他们对其主教的阴谋而陷入悲伤？安提约基雅为失去忠信正统的欧斯塔削而哀悼。巴兰内为厄乌弗拉提翁哭泣；帕尔托斯和安塔拉杜斯为基玛提乌斯和卡尔特利乌斯哀痛。哈德良堡被召来为挚爱的欧特洛皮乌斯及其继任者路基乌斯哀悼，路基乌斯屡次被锁链捆绑，终因不堪重负而死去。安基拉、贝罗亚和迦萨为马塞鲁斯、基禄斯和阿斯克勒帕斯的缺席而哀伤，他们受尽这个诡诈派别的虐待后被驱逐出境。使者被派去搜寻特拉奇的德奥杜鲁斯和奥林皮乌斯两位主教，以及我和我教区的司铎；若找到我们，我们无疑会被处死。但就在他们策划消灭我们的时候，我们得以逃脱，尽管他们已向总督多纳图斯发送针对奥林皮乌斯的信件，并向非拉格里乌斯发送针对我的信件。}\par

这个不虔敬的派别对最圣洁的基督徒所行的胆大妄为之事即是如此。贺西乌是科尔多瓦的主教，在所有参加尼西亚大公会议的人中最杰出；他也在萨迪卡会议中位居首位。\par

现在我愿在历史中插入一段记载，关于著名的利贝里乌斯为真理辩护向皇帝君士坦提乌斯所作的精彩论述。这些论述被当时一些虔诚者记录下来，以激励他人在神圣之事上同样热忱。利贝里乌斯接替尤利乌斯（西尔维斯特的继任者）治理罗马教会。\par

% structure: p0076 source=h2 target=subsection reason=relative-heading-rank; base=h2

\subsection{第十三章　罗马主教利贝里乌斯与皇帝君士坦提乌斯的会谈}

\Emph{君士坦提乌斯}。——\Quote{我们认为妥当，因你是基督徒又是我城的主教，召你来劝你弃绝与那不虔敬的亚大纳削的一切牵连。因为当议会将他与教会共融分开时，全世界都赞同这一决定。}\par

利贝里乌斯。——\Quote{皇帝啊，教会的判决应依据最严格的公义执行；因此，若蒙陛下喜悦，请召集法庭，若查明亚大纳削该受谴责，则按教会形式对他宣判。因为我们不能未经听审和审讯就定人之罪。}\par

\Emph{君士坦提乌斯}。——\Quote{全世界已谴责了他的不虔敬；但他从起初就嘲笑危险。}\par

利贝里乌斯。——\Quote{那些签署谴责的人并未亲眼目睹任何发生的事；而是受荣耀的欲望和因陛下而遭耻辱的恐惧所驱使。}\par

\Emph{皇帝}。——\Quote{你说的荣耀、恐惧和耻辱是什么意思？}\par

利贝里乌斯。——\Quote{那些不爱天主的荣耀，却更看重陛下赏赐的人，谴责了一个他们既未看见也未审判的人；这与基督徒的原则极不相符。}\par

\Emph{皇帝}。——\Quote{亚大纳削在提洛大公会议上亲自受审，那会议上全世界所有主教都谴责了他。}\par

利贝里乌斯。——\Quote{从未在他面前作出过任何判决。那些在那里集会的人是在他退席后谴责他的。}\par

\Emph{太监欧瑟比}愚蠢地插话。——\Quote{在尼西亚大公会议上已证明他所持的意见与公教信仰完全相悖。}\par

利贝里乌斯。——\Quote{在所有乘船前往马雷奥蒂斯、为起草针对被告的备忘录而被派遣的人中，只有五人宣判了他。被派去的五人中，两人现已去世，即德奥格尼和德奥多鲁斯。其他三人，马里斯、瓦伦斯和乌尔撒基乌斯，仍在世。在萨迪卡会议上，凡为此目的被派往马雷奥蒂斯的人都受到了判决。他们向会议呈递请愿书，请求赦免他们在马雷奥蒂斯起草针对亚大纳削的备忘录，其中包含不实指控和单方面的证词。他们的请愿书仍在我们手中。皇帝啊，我们该支持哪一方？我们该与谁同意并共融？是与那些先谴责亚大纳削、后又为谴责他而请求赦免的人，还是与那些谴责了后者的人？}\par

\Emph{主教埃皮克提图}。——\Quote{皇帝啊，利贝里乌斯并非为信仰辩护，也不是为维护教会的判决，而仅仅是为了在罗马元老面前夸耀自己在辩论中胜过了皇帝。}\par

\Emph{皇帝}（对利贝里乌斯说）。——\Quote{你在宇宙中占多大份额，竟独自一人袒护一个不虔敬之人，破坏帝国和整个世界的和平？}\par

利贝里乌斯。——\Quote{我独自一人并不减弱真理的力量。据古代故事，曾有三人抵抗敕令。}\par

\Emph{太监欧瑟比}。——\Quote{你把我们的皇帝比作拿步高。}\par

\Person{利贝里乌斯}：— \Quote{绝无此事。然而你竟未作任何审讯就鲁莽地定了一个人的罪。我所期望的，首先是签署一份普世性的信仰宣认，以确认在\Place{尼西亚}大公会议上所订立的信经；其次，要召回所有流放的弟兄们，恢复他们各自的主教职位。如果这一切都执行之后，能证明那些如今给教会制造麻烦的人所持的教义符合宗徒的信德，那么我们将全体在\Place{亚历山大里亚}聚集，会见被告、原告及他们的辩护人，在审查案件之后，我们将作出判决。}\par

\Emph{埃皮克提图主教}：— \Quote{不会有足够的驿车来运送如此多的主教。}\par

\Person{利贝里乌斯}：— \Quote{教会的事务没有驿车也能办理。教会能够提供运送各自主教到海边的交通工具。}\par

\Emph{皇帝}：— \Quote{已经通过的判决不应撤销。多数主教的决议应当生效。只有你一个人还和那个不敬神的人保持友谊。}\par

\Person{利贝里乌斯}：— \Quote{哦，皇帝，这是前所未有的事：一位法官竟像对待个人仇敌一样指控不在场的人不敬神。}\par

\Emph{皇帝}：— \Quote{所有人无一例外都受到他的伤害，但没有人像我这样深。他不但不满足于我长兄的死，还不断煽动那位有福记忆的\Person{君士坦斯}对我怀有敌意；但我以极大的克制，容忍了煽动者及其受害者的暴躁。我所取得的所有胜利，甚至包括战胜\Person{马格嫩提乌斯}和\Person{西尔瓦努斯}的胜利，都及不上将这个恶徒逐出教会治理之列。}\par

\Person{利贝里乌斯}：— \Quote{皇帝，请不要利用主教们来发泄你的仇恨和报复；因为他们的手只应该举起行祝福和圣化之事。如果合乎你的意愿，请下令主教们返回自己的住所；如果看来他们与那位今日持守在\Place{尼西亚}所签署的信经信条之真理的人同心，那么就让他们聚集起来，为世界的和平操心，以免一个无辜的人成为遭受谴责的目标。}\par

\Emph{皇帝}：— \Quote{只有一个问题需要处理。我希望你与教会共融，并送你回\Place{罗马}。因此同意和平，签署你的同意书，然后你就返回\Place{罗马}。}\par

\Person{利贝里乌斯}：— \Quote{我已经向那座城里的弟兄们道别了。教会的谕令比居住在\Place{罗马}更为重要。}\par

\Emph{皇帝}：— \Quote{你有三天时间考虑是否签署文件并返回\Place{罗马}；如果不，你必须选择你流放的地点。}\par

\Person{利贝里乌斯}：— \Quote{三天也好，三个月也罢，都不能改变我的想法。请把我送到你愿意的任何地方。}\par

过了两天，皇帝派人召见\Person{利贝里乌斯}，见他心意不变，就下令将他流放到\Place{色雷斯}的\Place{贝洛亚}城。\Person{利贝里乌斯}出发时，皇帝派送五百金币作为他的路费。\Person{利贝里乌斯}对送钱的使者说：\Quote{去，把这些还给皇帝；他需要它们来支付军饷。}皇后也送来了同样数额的钱；他说：\Quote{把它带给皇帝，因为他可能需要它来支付军饷；否则，就把它给\Person{奥克森修斯}和\Person{埃皮克提图}吧，因为他们需要它。}太监\Person{欧塞比乌斯}又给他带来了其他款项，他便这样对他说：\Quote{你把全世界的教会都变成了荒野，现在却向我这个罪犯施舍？走开，先做一个基督徒吧。}三天后，他被流放，没有接受任何给他的东西。\par

% structure: p0103 source=h2 target=subsection reason=relative-heading-rank; base=h2

\subsection{第十四章 论圣\Person{利贝里乌斯}被放逐与归来}

这位真理的胜利斗士，根据皇帝谕令，被遣送到\Place{色雷斯}。此事两年后，\Person{君士坦提乌斯}前往\Place{罗马}。贵妇们敦促她们的丈夫向皇帝请愿，要求将这位牧者归还给他的羊群；她们补充说，如果这一请求不被批准，她们将离弃丈夫，亲自追随她们伟大的牧者。她们的丈夫回答说，他们担心会招致皇帝的愤怒。\Quote{如果我们向他请求，}他们继续说道，\Quote{身为男子，他会认为这是不可饶恕的罪过；但如果你们亲自呈上请愿书，他至少会饶恕你们，并且要么答应你们的请求，要么不伤害你们而让你们离开。}这些高贵妇女采纳了这项建议，以她们惯常的华丽服饰出现在皇帝面前，好让君王从她们的衣着判断其身份，认为她们值得以礼相待、仁慈受接见。她们这样进入御前，恳求他怜悯如此一座大城的处境——它失去了牧者，易受豺狼攻击。皇帝回答说，羊群已有一位能照管它的牧者，城中不需要别的牧者。因为伟大的\Person{利贝里乌斯}被放逐后，他的一位执事，名叫\Person{斐理克斯}，已被任命为主教。他完好地保存了尼西亚信经所陈明的教义，却与那些败坏这信经的人保持共融。因此，城中没有一位市民愿意在他主礼时进入祈祷之所。这些妇女向皇帝陈述了这些事实。她们的劝说成功了；他命令召回被放逐的伟大的\Person{利贝里乌斯}，并由两位主教共同管理教会。皇帝的谕令在竞技场中被宣读，群众高呼，说皇帝的谕令是公正的；观众分为两派，各以其颜色命名，如今每派将有自己的主教。在这样嘲笑皇帝的谕令之后，他们异口同声地喊道：\Quote{惟一的天主，惟一的基督，惟一的主教。}我认为将他们的原话记录下来是恰当的。在这位基督徒群众发出这些虔诚而正义的欢呼后不久，圣\Person{利贝里乌斯}返回，而\Person{斐理克斯}退往另一座城市。\par

为了保持叙述秩序，我将这段记述附在有关米兰主教们行事的叙述之后。现在，我将按时间顺序回到事件的叙述中。\par

% structure: p0106 source=h2 target=subsection reason=relative-heading-rank; base=h2

\subsection{第十五章。\Place{阿里米努姆}会议。}

当所有维护信德的人都被撤走后，那些按照自己意愿塑造皇帝心意的人，自认为他们反对的信德很容易被推翻，而亚流主义得以取而代之，便说服\Person{君士坦提乌斯}在\Place{阿里米努姆}召开东西方主教会议，旨在从信经中删除教父们为抵制\Person{亚流}的诡诈所创制的术语——\OriginalTerm{实体}{substance}和\OriginalTerm{同体}{of one substance}。因为他们认为这些术语造成了教会与教会之间的纷争。当他们在会议中聚集时，亚流派的支持者试图欺骗大多数主教，尤其是那些西方帝国城市的主教，他们为人质朴、不谙世故。他们再三论证，教会的整体绝不能因为两个圣经中找不到的术语而分裂；同时，他们承认描述圣子在一切方面\Quote{\Emph{相似}}于圣父是恰当的，却坚持删除\Quote{\Emph{实体}}一词，认为它不符合圣经。然而，会议看穿了这些观点提出者的动机，因而拒绝他们。正统主教们通过一封信向皇帝陈明他们的心意；他们说，我们是尼西亚会议教父们的儿子和继承人，如果我们胆敢删减他们所签署的任何内容，或增添任何他们妥善制定的东西，我们便表明自己不是真正的儿子，而是我们生身之父的控告者。但他们的信仰宣认的确切措辞将在他们致\Person{君士坦提乌斯}皇帝的信中更准确地给出。\par

% structure: p0108 source=h2 target=subsection reason=relative-heading-rank; base=h2

\subsection{由在\Place{阿里米努姆}召开的会议写给\Person{君士坦提乌斯}皇帝的信。}

我们相信，是应天主的召唤，并出于对陛下虔诚的服从，我们西方教会的主教们在\Place{阿里米农}召开主教会议，为要阐述天主教会的信仰，并揭露其反对者。经过长时间审议，我们明确认为，最好的办法是坚持并守护那从起初确立、经由先知、福音作者和宗徒们宣讲的信仰，借着我们的主耶稣基督——乃是陛下帝国的守卫者、陛下救恩的维护者——一直守护到底。因为在\Place{尼西亚}，在陛下父亲、最荣耀的\Person{君士坦丁}的监督下，已经正确而公正地界定的教义，以及为使人听闻理解而公布宣讲的道理与精神，若加以任何更改，显然是荒谬而不合法的。这信仰注定要对抗并摧毁亚略异端，并且不仅亚略异端，所有其他异端也因它而瓦解。在这信仰上增添任何东西实属危险，减去任何东西则冒巨大风险。若有增添或减损，我们的敌人便可肆意妄为。因此，\Person{乌尔萨西乌斯}和\Person{瓦伦斯}——亚略教义的公开追随者与朋友——已被宣布与我们断绝共融。为了保持他们在共融中的位置，他们请求获准一个\OriginalTerm{补赎的机会}{locus penitentiæ}，并对他们承认有误的所有点予以宽恕，正如他们自己所写的文件所证明，并借此获得了恩惠与宽恕。这些事情发生时，正值在\Place{米兰}召开主教会议，罗马教会的司铎们也出席了。众所周知，\Person{君士坦丁}——虽已去世，但值得纪念——曾以完全精确和审慎的方式阐述了所拟定的信经；如今他领受圣洗后去世，安息于他所应得的平安。我们认为，在他之后我们进行任何创新，并羞辱所有那些制定并表达这一教义的圣宣信者和殉道者，是荒谬的——因为他们的思想始终被教会古老的纽带所束缚。天主借着我们的主耶稣基督，将这一信仰传至陛下统治的时代，靠祂的恩宠，陛下的帝国如此广阔，统治全世界。然而，那些可悲可怜的人，又胆大妄为地宣称自己是他们那邪恶见解的宣讲者，并着手推翻真理的全部力量。因为在陛下命令召开会议时，他们暴露了自己不诚实的意图。他们开始通过恶行与混乱进行创新。他们网罗了某些追随者——一位名叫\Person{格尔马尼乌斯}、\Person{奥克森提乌斯}和\Person{卡尤斯}的人，他们是异端与纷争的煽动者，其教义虽只一种，却超越了一大堆亵渎言论。然而，当他们意识到我们并不赞同他们的想法，也不同情他们邪恶的企图时，他们进入我们的集会，仿佛要提出其他建议，但很短时间就足以暴露他们的真实意图。因此，为了避免教会管理时常陷入同样的困境，也防止它们被扰乱与无序的漩涡所吞噬，我们认为安全之策是保持之前所界定之事固定不变，并将上述之人与我们断绝共融。为此，我们派遣使者到陛下面前，以表明和解释本次会议信中表达的心意。我们首要叮嘱这些使者，须按照古老正确的界定守护真理。他们要向陛下禀明，不像\Person{乌尔萨西乌斯}和\Person{瓦伦斯}那样，说真理被推翻后就会有和平；因为和平的破坏者如何能成为和平的使者？而是，这些变革将会带来纷争与扰乱，既涉及其他城市，也涉及罗马教会。因此，我们恳请陛下以仁慈的耳朵和温和的态度接待我们的使者，不要允许任何新事物轻慢逝者。准许我们持守我们祖先的界定和决议，我们毫不犹豫地声称，他们所做的一切都是本着睿智与智慧，并藉着圣神而行的。现在寻求引入的新事物，正使信徒充满不信，使非信徒充满轻信。\par

我们请求您下令，为远方的、同样受年老与贫困困扰的主教们提供回乡的便利，以免教会长期缺少主教。\par

\Quote{此外我们恳求一件事：既不可从已确立的教义中减去什么，也不可增添什么，而要使一切保持完整，正如它们由陛下父亲的虔诚保存至今日。愿我们不再劳苦，也不再远离自己的教区，而让主教与自己的百姓在和平、祈祷与朝拜中度过时日，为陛下的帝国、健康和和平献上祈祷，愿天主永远赐予您这一切。我们的使者亦将根据圣经训诲陛下，他们携带着主教们的签名与问候。}\par

信函写就，使者派出，但宫廷的高级官员们虽然接到了公文并呈递给他们的君主，却拒绝引见使者，理由是陛下正忙于国事。他们这样做，是希望主教们因延误而恼怒，并急于返回托付给他们的城市，最终被迫自己解散并破除反异端的壁垒。但他们的诡计落空了，因为信仰的高贵捍卫者们又向皇帝发出了第二封信，恳请他接见使者并解散会议。我现将此信附录如下。\par

\Emph{会议致\Person{君士坦提乌斯}的第二封信}\par

致常胜者、虔诚的皇帝\Person{君士坦提乌斯}，在\Place{阿里米农}集会的众主教问安。\par

\Quote{至显赫的主上和皇帝，我们已收到陛下仁慈的信函，通知我们，因国事繁忙，您迄今未能召见我们的使者。您吩咐我们等候他们的回报，以便陛下能就我们眼前的目标以及前任的法令作出决定。但我们斗胆在此信中对陛下重申我们先前提议的要点，因为我们丝毫未退让。我们恳求您仁慈地接纳我们卑谦的信函，在此我们向陛下作出答复，以及我们命使者呈交陛下的要项。陛下与我们都同样清楚，这是何等严重且不合宜的境况——在您最幸福的统治时期，竟有如此多的教会似乎没有主教。因此，我们再次恳求最荣耀的皇帝陛下，若合您心意，请命我们在冬季的严寒之前返回各自的教会，使我们能够与我们的子民，一如既往，为陛下的帝国健康兴旺，向全能天主及祂的圣子、我们的主和救主基督献上最恳切的祈祷。}\par

% structure: p0116 source=h2 target=subsection reason=relative-heading-rank; base=h2

\subsection{Chapter 16. 关于在\Place{色雷斯}{Thrace}的\Place{尼卡}{Nica}召开之会议，以及该处草拟之信经}

这封信之后，他们激怒了皇帝，并争取到多数主教（尽管违背他们的意愿）到色雷斯一个名为尼卡的小镇。他们欺骗了一些单纯的人，恐吓了其他人，以执行他们一贯损害真宗教的计谋，即从信经中删去\Quote{本体}和\Quote{同一本体}这些词，并代之以\Quote{相似}一词。我将他们的公式插入这部历史中，不是因为它措辞恰当，而是因为它证实了亚流派的分裂，因为即便当今的不满者也未接受它。现在，他们传讲\Quote{不相似}以代替\Quote{相似}。\par

\Emph{在色雷斯的尼卡提出的不正统信经}\par

\Quote{我们信唯一真天主，全能圣父，万物由祂而有。并信天主的独生子，祂在万世之前、一切起始之前，由天主所生，万物由祂造成，有形无形；独生者，唯一由圣父所生的独生子，天主出自天主：相似祂的父，正如经上所载，祂的诞生无人知晓，唯有生祂的圣父。我们知道天主的这独生子，由父所遣，为摧毁罪恶与死亡，从天上降下，由圣神和童贞玛利亚按肉身而生，正如经上所载。祂与门徒同行，当救恩计划按父的旨意完成时，被钉、死亡、埋葬，并降至阴间，连阴府自身也战栗。第三日祂从死者中复活，与门徒同行四十日。祂被接升天，坐在圣父的右边，将在复活的最末日，在圣父的光荣中来临，照各人的行为予以赏报。我们信圣神，即天主的独生子耶稣基督，祂既是天主又是主，应许赐予人护慰者，正如经上所载，真理之神。这圣神祂亲自在升天并坐在父右边之后派遣，从那里来审判生者死者。但\Quote{本体}一词，因教父们过于简单插入，且不为百姓所理解，既因未见于圣经而成为绊脚石，我们觉得宜予删除，且今后不许再提及\Quote{本体}，因为圣经从未提及父与子的\Quote{本体}。也不得在圣父、圣子、圣神的位格上称\Quote{一个本质}。我们称子相似于父，正如圣经所称和教导的；但所有异端，无论已被谴责者，或任何反对此公布文件的（若有），皆当绝罚。}\par

这份信经由主教们签署，有些是被吓倒，有些是被哄骗；但那些拒绝服从的人被流放到世界最遥远的地区。\par

% structure: p0121 source=h2 target=subsection reason=relative-heading-rank; base=h2

\subsection{Chapter 17. 关于罗马主教\Person{达玛苏}{Damasus}与西方主教关于\Place{阿里米农}{Ariminum}会议之教务会议法案}

所有真理的捍卫者，特别是西方捍卫者，谴责这一公式，可见于他们写给伊里利亚人的信函。首先署名者是\Person{达玛苏}{Damasus}，他在\Person{利贝里}{Liberius}之后获得\Place{罗马}{Rome}教会之首席，并兼备许多美德。与他一同签署的，有在罗马聚会的\Place{意大利}{Italy}和\Place{迦拉达}{Galatia}（今称\Place{高卢}{Gaul}）的九十位主教。我本会插入他们的名字，但觉得多余。\par

-------------------------\par

在\Place{罗马}{Rome}以神圣会议聚集的主教们——\Person{达玛苏}{Damasus}、\Person{瓦勒里安}{Valerianus}及其他各位，致\Place{伊里利亚}{Illyria}诸主教，我们亲爱的弟兄们，在天主内问安。\par

我们相信我们，天主的司铎，由之教化其余众人是正当的，我们持守并教导百姓那建立在宗徒教导上的圣信经，且绝不偏离教父们的定义。然而，通过\Place{高卢}{Gaul}和\Place{维内提亚}{Venetia}弟兄们的报告，我们得知某些人已陷入异端。\par

主教的职责不仅是预防这种祸害，而且还要抵制任何出现的分歧教导，无论是由于不完整的教导，还是不良注释者的简单读者。他们应决心不滑入危险的道路，而是每当分歧的意见传到他们耳中时，要坚持我们教父的教义。因此，已决定在此事上特别谴责\Place{米兰}{Milan}的\Person{奥森提乌}{Auxentius}。所以，罗马帝国中所有法律的教师都应得到良好训导，不因分歧教义玷污信仰。\par

当异端者的邪恶开始兴盛，亚略异端者的亵渎正如现今一样向前爬行之际，我们的教父，三百一十八位主教，即罗马帝国中最神圣的高级圣职者，在尼西亚集会。他们为对抗魔鬼的武器所设立的屏障，以及用以击退其致命毒药的解毒剂，便是他们的宣认：圣父与圣子同一本体、同一神性、同一德能、同一权能、同一形像，圣神亦同一本质与本体。凡不如此思想者，被视为与我们的共融分离。他们的集议值得一切尊敬，他们的定义是健全的。但某些人意图借后来的其他讨论来败坏并玷污它。然而，从一开始，那些在阿里米努姆被逼迫操纵或创新信仰的主教们，便如此纠正了错误，以至于他们承认自己为相反论据所迷惑，或承认自己没有觉察到与尼西亚教父所传意见有任何矛盾。从聚集在阿里米努姆的主教人数，并不会产生任何偏见，因为众所周知，罗马主教（其意见本应优先于一切而等待）、味增爵（其无瑕的主教职务已持续多年），以及其余诸人，均未给予这类教义以支持。这一点更为重要，因为，正如已经说过的，那些似乎被诱骗而屈服的人，在他们更明智的时刻，自己也表示了不赞成。\par

你的诚挚由此看出，这一唯一信仰，即建于尼西亚、基于宗徒权威的信仰，应当永远稳妥地持守。你看出，我们这些自认公教的东西方主教们，与西方主教们一同以它为荣。我们相信，不久之后，那些持不同见解者，应当毫不迟延地被排除出我们的共融，并被剥夺主教之名，好使他们的羊群摆脱错误，自由地呼吸。因为当他们自己受错误之害时，不能指望他们纠正子民的错误。愿你的尊长意见与所有天主司祭的意见和谐一致。我们相信你在此意见上是坚定稳固的，因此我们理应正确地与你们一同相信。愿你的爱德以答复使我们喜乐。\par

\Quote{亲爱的弟兄们，祝你们平安。}\par

% structure: p0130 source=h2 target=subsection reason=relative-heading-rank; base=h2

\subsection{第十八章 亚历山大里亚主教亚大纳削关于同一大公会议的书信}

伟大的亚大纳削也在他致非洲人的信中，就阿里米努姆议会如此写道。在此情形下，谁能容忍提及阿里米努姆议会或任何其他议会，而不提尼西亚议会呢？谁不憎恶摒弃教父们的言辞，而偏爱那些在阿里米努姆以暴力和党派之争所引入的言辞呢？谁愿与这些人——那些连自己的话都不接受的人——为伍呢？在他们自己的十几次主教会议中，他们已如所述，忽而规定这个，忽而规定那个；而现今，他们自己正公开指责一个接一个的这些会议。他们正遭受古时犹大叛逆者所遭遇的命运。因为正如\BibleBook{耶肋米亚}{先知书}所载：\Quote{\Emph{他们离弃了我这活水的泉源，为自己凿出蓄水池，破裂的蓄水池，不能蓄水}}，同样，这些人反对大公会议，为自己凿出了许多会议，但全都证明是虚妄的，如同\Quote{\Emph{不结粮的嫩芽}}。因此，我们不要接纳那些引证阿里米努姆议会或任何其他议会而非尼西亚议会的人，因为事实上，那些引证阿里米努姆的人似乎也不知道那里所行的事；他们若知道，就会闭口不言。因为你们，亲爱的，已从你们在那次会议中的代表们得知，并因此非常清楚，乌尔撒基、瓦冷、欧多克西、奥克森提，以及同他们一起的德莫斐禄，被要求诅咒亚略异端，却推辞，宁愿作其捍卫者，于是均因提出违反尼西亚法令的提案而被罢免。相反，那些真正是主的仆人并持守正确信仰的主教——大约二百位——声明他们只坚持尼西亚议会，并且拒绝考虑对其法令作任何删减或增添。他们已将此结论通报给君士坦提乌斯，议会是由他下令召集的。\par

另一方面，那些在阿里米努姆被罢免的主教们，却被君士坦提乌斯接纳，并且成功地使那判处他们的二百位主教受到严重侮辱，并威胁不让他们返回教区，还须在色雷斯遭受严酷对待，且正值冬季，以迫使他们接受那些革新者的措施。\par

因此，若我们听到有人援引阿里米努姆，让我们反问：请先给我们看那罢免的判词，以及主教们所起草的文件，其中声明他们不求超越尼西亚教父所拟定的条文，也不诉诸任何其他议会，唯独尼西亚议会。事实上，恰恰是这些事实他们隐瞒，却极力宣扬色雷斯受胁迫的宣认。他们只是表明自己是亚略异端的友人，与纯正信仰格格不入。只要有人愿意将那次大公会议与这些人所援引的其他会议并列比较，他就会看到一边是真宗教，另一边是愚昧与混乱。尼西亚的教父们聚集，并非在被罢免之后，而是在宣认圣子出自圣父的本体之后。这些人曾被罢免一次、两次，又在阿里米努姆第三次被罢免，然后竟敢规定将本体或本质归于天主是错误的。亚略狂热的帮派在西方反对真理教义所编造的诡计与阴谋，竟是如此奇怪、如此繁多。\par

% structure: p0134 source=h2 target=subsection reason=relative-heading-rank; base=h2

\subsection{第十九章 论安提约基雅主教良修的诡诈，以及弗拉维安与狄奥多若的勇敢}

在\Person{安提约基亚}，\Person{普拉西杜斯}由\Person{斯德望}继任，后者被逐出教会。\Person{莱昂提乌斯}随后接受了首席之职，但违反了尼采亚会议的法令，因为他自阉，成为阉人。他这一鲁莽行为的原因，真福\Person{亚大纳削}如此叙述：看来，\Person{莱昂提乌斯}是因一个名叫\Person{欧斯托利亚}的年轻女子而成了诽谤的牺牲品。当他发现自己无法与她同住时，便为她自阉，以便可以自由地与她同居。但他并未消除嫌疑，反而因此被革除司铎职。\Person{亚大纳削}关于他早年生平的其他部分就这样记载了。现在我将简要揭露他的恶行。他虽分享亚略的谬误，却始终试图隐藏自己的不健全。他观察到神职人员和其余百姓分为两派：一派在给予圣子光荣时，使用连接词\Quote{和}；另一派则对圣子使用介词\Quote{藉着}，并对圣神应用\Quote{在}。他自己则沉默地诵读全部光荣颂，站在他附近的人所能听到的只有\Quote{直到永远}。若不是他灵魂中极度的邪恶通过其他方式被暴露，或许会有人说他采用这种策略是出于促进人民和睦的愿望。但当他给真理的捍卫者造成许多损害，并继续全力支持不敬之事的推进者时，他被揭露隐藏了自己的不健全。他既受制于对人民的恐惧，也受制于\Person{君士坦提乌斯}对那些胆敢说圣子与圣父不相似者所发出的严厉威胁。然而，他的真实情感通过其行为被证实。宗徒教义的追随者从未从他那里得到任何祝圣或丝毫鼓励。另一方面，支持亚略迷信的人，既被允许完全自由地表达意见，又时常被接纳到司铎职位。此时，\Person{尤诺米乌斯}的师傅\Person{阿埃提乌斯}，他以其思辨助长了亚略谬误，被接纳为执事。然而，\Person{弗拉维雅努斯}和\Person{狄奥多鲁斯}，他们过着苦修生活，并公开捍卫宗徒法令，公开抗议\Person{莱昂提乌斯}对真宗教的攻击。他们敦促说，一个在邪恶中长大并图谋以不敬虔赢得名声的人，竟被视为配得上执事职，这是教会的耻辱。他们还威胁说，他们将脱离他的共融，旅行到西帝国，向世界揭露他的阴谋。\Person{莱昂提乌斯}此时惊慌，暂停了\Person{阿埃提乌斯}的神圣职务，但继续对他表示特别青睐。\par

那对杰出伴侣\Person{弗拉维雅努斯}和\Person{狄奥多鲁斯}，虽尚未被接纳到司铎职且仍列于平信徒中，但昼夜工作以激发人们对真理的热忱。他们是首先将诗班分为两部分，并教他们以对唱方式咏唱\Person{达味}的圣咏的人。这种做法首先在\Place{安提约基亚}引入，然后向各方传播，渗透到地极。其创始人如今将热爱天主圣言与圣工的人聚集在殉道者的教堂中，与他们一起整夜向天主咏唱圣咏。\par

当\Person{莱昂提乌斯}察觉到这一点时，他认为试图阻止他们并不安全，因为他看到百姓对这些杰出人士极为好感。然而，他在言语上装出礼貌的样子，请求他们在教堂里举行这种敬拜。他们深知他的恶意。尽管如此，他们还是听从了他的命令，并随即召集他们的诗班来到教堂，劝勉他们歌颂善主。然而，没有任何事能促使\Person{莱昂提乌斯}纠正他的邪恶，他反而戴上了公正的面具，隐藏了\Person{斯德望}和\Person{普拉西杜斯}的邪恶。那些接受司铎和执事信仰败坏的人，虽然过着卑鄙不规则的生活，他却将他们列入名册；而其他具有各种美德并坚定遵守宗徒教义的人，他却不予承认。于是，在神职人员中，大多数是沾染异端的人，而大多数平信徒却是信德的捍卫者，甚至专业教师也缺乏勇气揭露他们的亵渎。实际上，\Person{普拉西杜斯}、\Person{斯德望}和\Person{莱昂提乌斯}在\Place{安提约基亚}所行的不敬和邪恶之事如此之多，以至于需要一部专门的历史，并且如此可怕，以致配得上\Person{达味}的哀歌；因为对他们也必说：\BibleQuote{看，祢的仇敌正在喧嚷；恨祢的人都昂首挺胸。他们密谋反抗祢的子民，商议攻击祢的密友。他们说：\InnerQuote{来，我们将他们除灭，使他们不再成为一个民族；愿以色列的名字不再被记念。}}\par

现在我们继续叙述的进程。\par

% structure: p0139 source=h2 target=subsection reason=relative-heading-rank; base=h2

\subsection{第二十章 关于\Person{欧多克修斯}（来自\Place{日耳曼尼西亚}）的创新，以及\Place{安西耳}的\Person{巴西里乌斯}和\Place{塞巴斯特}的\Person{欧斯塔修斯}对他的热忱反对}

\Place{日耳曼尼西亚}是一座位于\Place{基里基雅}、\Place{叙利亚}和\Place{卡帕多细亚}沿海的城市，属于称为\Place{幼发拉底斯亚}的省份。该教会首领\Person{欧多克修斯}一听到\Person{莱昂提乌斯}的死讯，便直奔\Place{安提约基亚}并夺取了主教座，在那里如同野猪一样摧毁了主的葡萄园。他甚至不像\Person{莱昂提乌斯}那样试图隐藏自己的恶行，而是直接攻击宗徒法令，使所有敢于反驳他的人都陷入各种麻烦。此时，\Person{巴西里乌斯}继\Person{马塞鲁斯}之后，执掌\Place{加拉提雅}首府\Place{安西耳}的教会；\Place{塞巴斯特}，\Place{阿尔美尼亚}的首府，则在\Person{欧斯塔修斯}的引导之下。这些主教一听到\Person{欧多克修斯}的邪恶和疯狂，便写信告知\Person{君士坦提乌斯}皇帝的放肆行为。\Person{君士坦提乌斯}此时仍逗留在西方，在僭主们死后，正努力医治他们所造成伤害。两位主教都为皇帝所熟知，并因其高尚品格对皇帝有很大的影响力。\par

% structure: p0141 source=h2 target=subsection reason=relative-heading-rank; base=h2

\subsection{第二十一章 关于第二次\Place{尼采亚}大公会议}

君士坦提乌斯收到这些公文后，写信给安提约基雅人，否认他曾将安提约基雅主教座席托付给欧多克修斯，正如欧多克修斯公开宣布的那样。他命令将欧多克修斯放逐，并因其在\Place{比推尼}的\Place{尼西亚}所采取的行动而受惩罚，他在那里曾命令召集主教会议。欧多克修斯本人说服了掌管皇室事务的官员将尼西亚定为会议地点。但那位至高无上的统治者与治理者，祂知晓未来如同过去，以一场大地震阻止了集会，导致该城大部分被毁，居民大多丧生。得知此事后，聚集的主教们惊慌失措，各自返回自己的教会。但我认为这是天主上智的安排，因为在那座城市中，圣教父们曾定义了宗徒信仰的教义。而在同一座城市，后来聚集的主教们却意图制定相反的教义。名称的相同肯定会为亚略派一伙提供欺骗的手段，并欺骗天真的灵魂。他们本想称会议为\Quote{尼西亚}，并使之与昔日著名的会议混同。但眷顾教会的那位解散了这次会议。\par

% structure: p0143 source=h2 target=subsection reason=relative-heading-rank; base=h2

\subsection{第二十二章 论在\Place{伊苏利亚}的\Place{塞琉西亚}举行的会议}

过了一段时间，在欧多克修斯的控告者的建议下，君士坦提乌斯命令在\Place{塞琉西亚}举行会议。这座\Place{伊苏利亚}城镇位于海边，是该地区的首府。东方的诸位主教，以及来自\Place{亚细亚}的\Place{本都}地区的主教们，被命令在此聚集。\par

\Place{凯撒勒雅}——\Place{巴勒斯坦}的首府——的主教座席，现在由\Person{阿卡西乌斯}担任，他继任了\Person{优西比乌斯}。他曾在\Place{撒尔底加}会议中被谴责，但他对如此众多主教组成的会议表示轻蔑，并拒绝接受他们的不利裁决。在\Place{耶路撒冷}，我常提到的\Person{玛加略}由\Person{玛克西穆斯}继任，此人因其为宗教奋斗而引人注目，因为他被剥夺了右眼，并致残了右臂。\par

当他被接升入不知衰老的生命后，\Person{济利禄}——宗徒法令的热切捍卫者——被尊以主教职。这些人彼此争夺首位，给国家带来了巨大灾难。\Person{阿卡西乌斯}抓住一个小借口，免去了\Person{济利禄}的职务，将他从\Place{耶路撒冷}驱逐出去。但\Person{济利禄}经过\Place{安提约基雅}，发现该城没有牧人，便来到\Place{塔尔索}，与当时该教区的主教——杰出的\Person{西尔瓦努斯}——同住。\Person{阿卡西乌斯}一得知此事，就写信给\Person{西尔瓦努斯}，告知他\Person{济利禄}被免职。然而，\Person{西尔瓦努斯}出于对\Person{济利禄}的尊重，并且也不无对当地居民的疑惧——他们非常喜欢这位陌生人的教导——拒绝禁止他参与教会职务的施行。然而，当他们到达\Place{塞琉西亚}时，\Person{济利禄}加入了\Person{巴西略}、\Person{欧斯塔修}、\Person{西尔瓦努斯}以及其余人在会议中的一派。但是当\Person{阿卡西乌斯}加入聚集的主教们——共一百五十人——时，他拒绝在\Person{济利禄}被逐出他们之前与它（指会议）有任何关联，因为\Person{济利禄}已被剥夺了主教职。有一些渴望和平的人恳求\Person{济利禄}退出，并承诺在决定法令之后会调查他的案件。他不肯让步，于是\Person{阿卡西乌斯}离开他们，走了出去。然后他遇见\Person{欧多克修斯}，消除了他的恐惧，并鼓励他，承诺会做他的朋友和支持者。就这样，他阻止\Person{欧多克修斯}参与会议，并与他一起前往\Place{君士坦丁堡}。\par

% structure: p0147 source=h2 target=subsection reason=relative-heading-rank; base=h2

\subsection{第二十三章 论在\Place{君士坦丁堡}的正统主教们所遭遇的事}

\Person{君士坦提乌斯}从西方返回后，在\Place{君士坦丁堡}逗留了一段时间。在那里，\Person{阿卡西乌斯}在皇帝面前对聚集的主教们提出诸多指控，称他们是一群邪恶之徒，为了教会的毁灭与败坏而召集，从而激起了皇帝的愤怒。\Person{君士坦提乌斯}尤其因对\Person{济利禄}的指控而激动，因为\Person{阿卡西乌斯}说：\Quote{因为那件圣袍，即著名的\Person{君士坦丁}皇帝为荣耀\Place{耶路撒冷}教会而赐给该城主教\Person{玛加略}的，用来在施行神圣洗礼礼仪时穿着的，完全用金线制成的圣袍，已被\Person{济利禄}卖掉了。它已被买去，}他接着说，\Quote{被某个舞台舞者买去；当他穿着它跳舞时，跌倒了并死了。与这样的\Person{济利禄}，}他继续道，\Quote{他们竟自居为全世界其余部分判断和决定。}朝廷中有势力的一方以此为由，劝皇帝不要召集全体会议，因为他们对大多数人的一致感到恐慌，而只召集十位主要领导人物。其中有\Place{亚美尼亚}的\Person{欧斯塔修}、\Place{加拉提亚}的\Person{巴西略}、\Place{塔尔索}的\Person{西尔瓦努斯}和\Place{基齐库斯}的\Person{厄肋乌息}。\par

他们到达后，敦促皇帝断定欧多克修斯犯有亵渎和违法之罪。然而，君士坦提乌斯受对立派别唆使，回答说应先就信仰问题作出决定，然后再调查欧多克修斯的案件。巴西里乌斯仗着以往的交情，大胆向皇帝提出异议，说他攻击宗徒的法令；但君士坦提乌斯对此不悦，命巴西里乌斯住口，说：\Quote{因为教会的骚乱是由于你。} 巴西里乌斯沉默后，尤斯塔修斯插话说：\Quote{陛下，既然您希望对信仰问题作出决定，请考虑欧多克修斯对独生者轻率说出的亵渎之词。} 他一面说一面拿出信仰宣示，其中除许多不敬之词外，还有以下表述：\Quote{不同术语说的事物，其本质也不同：} \Quote{只有一位天主圣父，万物出于祂；只有一位主耶稣基督，万物藉着祂。} 现在术语\InnerQuote{出于祂}不同于术语\InnerQuote{藉着祂}；所以圣子不同于天主圣父。君士坦提乌斯命人宣读这份信仰宣示，并对其中包含的亵渎感到不悦。于是他问欧多克修斯是否是他起草的。欧多克修斯立刻否认作者身份，说这是阿埃提乌斯写的。这位阿埃提乌斯就是利昂提乌斯因惧怕弗拉维亚努斯和狄奥多鲁斯的指控，先前将他从执事职位上贬黜的那一位。他也是格奥尔吉乌斯的支持者，后者是亚历山大里亚人的奸诈敌人，言语不敬，行为邪恶。此时他与尤诺米乌斯和欧多克修斯为伍；因为利昂提乌斯死后，欧多克修斯暴力夺取了安提阿教会的主教宝座，他便与尤诺米乌斯从埃及返回，发现欧多克修斯与自己思想一致，既奢侈享乐又信仰异端，于是选择安提阿作为最适宜的居所，他和尤诺米乌斯都成了欧多克修斯席上的常客。他最大的抱负是做一个成功的食客，整天流连于这个或那个人的餐桌大快朵颐。皇帝早已听说这一切，现在下令将阿埃提乌斯带到他面前。阿埃提乌斯到来后，君士坦提乌斯给他看那件文件，并询问他是否是其语言作者。阿埃提乌斯完全不知发生了什么事，也不了解询问的意图，以为供认会赢得赞赏，便承认自己是那些短语的作者。于是皇帝看出他罪大恶极，立即判处流放，押送到弗里吉亚的一个地方。这样，阿埃提乌斯因亵渎而蒙羞，被逐出皇宫。然后尤斯塔修斯声称欧多克修斯也持有相同观点，因为阿埃提乌斯曾与他同住同食，并根据他的判断起草了这份亵渎性信条。为证明欧多克修斯参与起草文件的论点，他提出一个事实：除了欧多克修斯本人，没有人将这份文件归于阿埃提乌斯。对此，皇帝命令法官不得凭猜测判决，而必须对事实进行精确调查。尤斯塔修斯表示同意，并敦促欧多克修斯通过诅咒阿埃提乌斯的作品来证明他不同意归于他的观点。皇帝非常乐意地接受了这个建议，并相应地下达命令；但欧多克修斯退缩了，用许多手段规避服从。但当皇帝发怒，威胁要将他与阿埃提乌斯一同流放，因为他参与了受到惩罚的亵渎时，他便否认了自己的教义，尽管当时和后来他都坚持这一教义。然而，他又反过来抗议尤斯塔修斯派，说他们有责任谴责\Quote{\OriginalTerm{同质}{Homoüsion}}这个词，因为这个词未见圣经。\par

西尔瓦努斯则指出，他们有责任拒绝并从神圣集会中排除\Quote{出自非存在者}、\Quote{受造物}和\Quote{另一本质}这些短语，这些术语也未见圣经，在先知或宗徒的著作中都找不到。君士坦提乌斯断定这是正确的，并命亚流派人宣布诅咒。起初他们坚持拒绝；但最后，当他们看到皇帝的愤怒时，便勉强同意诅咒西尔瓦努斯摆在他们面前的术语。但他们更加坚决地要求谴责\Quote{\OriginalTerm{同质}{Homoüsion}}。然而，西尔瓦努斯用无可辩驳的逻辑向亚流派人及皇帝提出真理：如果天主圣言不是出自非存在者，祂就不是受造物，也不是另一本质。那么祂与生育祂的天主是同一本质，如天主出自天主、光出自光，并且与生育者具有同一本性。他有力而真实地提出这一论点，但他的听众中无人信服。阿卡西乌斯和欧多克修斯一派喧闹起来；皇帝发怒，威胁要将他们逐出教会。于是埃琉修斯、西尔瓦努斯等人说，惩罚权属于皇帝，但决定敬神或不敬神的问题属于他们，他们抗议说：\Quote{我们不会背叛教父的教义。}\par

君士坦提乌斯本应钦佩他们的智慧和勇气，以及他们对宗徒法令的勇敢捍卫，但他将他们逐出教会，并命另选他人接替。于是欧多克修斯暴力夺取了君士坦丁堡教会；埃琉修斯被逐出基齐库斯后，尤诺米乌斯被任命接替他。\par

% structure: p0152 source=h2 target=subsection reason=relative-heading-rank; base=h2

\subsection{第24章。反对阿埃提乌斯的主教训谕}

这些事过后，皇帝命令以正式公文宣判定\Person{阿埃提乌斯}有罪。从命之下，他的同党便判决了自己的同伙。于是他们致函\Place{亚历山大里亚}主教\Person{乔治}，关于他的信函，我将在本史中留一席之地，以揭露他们的邪恶，因为他们对待友人和仇敌的方式完全一样。\par

全体公会议致\Person{乔治}信函抄本，反对其执事\Person{阿埃提乌斯}，因其不虔的亵渎。\par

致可敬的主、\Place{亚历山大里亚}主教\Person{乔治}，\Place{君士坦丁堡}聚集的圣主教会议致候。\par

因\Person{阿埃提乌斯}被公会议定罪，因其非法且极冒犯的著作，主教们已按教会法规处理了他。他被免除执事职，并被逐出教会；我们的告诫已发出，禁止任何人阅读其非法书信，因其无用且无价值，应予弃置。我们进一步对他附加了绝罚，若他坚持己见，也对其支持者附加绝罚。\par

按理说，所有在公会议中聚集的主教本应憎恶并赞同对这名制造罪行、骚乱、分裂，在全世界引发动荡和教会与教会对立之人的判决。但尽管我们祈祷，且出乎我们一切意料，\Person{塞拉斯}、\Person{斯特凡}、\Person{赫利奥多罗斯}和\Person{德奥斐洛}及其同伙并未与我们一同投票，甚至不肯签署对他的判决，尽管\Person{塞拉斯}指控上述\Person{阿埃提乌斯}另一次狂妄自大的实例，声称他更厚颜无耻地跳出来宣告，天主向宗徒们隐瞒的事如今已启示给他。即便在\Person{塞拉斯}报告了\Person{阿埃提乌斯}这些狂言后，上述主教仍未羞愧，也不能被说服与我们一同投票定他之罪。然而我们以极大的忍耐长时间包容了他们，时而愤怒，时而恳求，时而催促他们与我们联合，使公会议的决定一致；我们长久坚持，希望他们能听取并同意而让步。但当这一切忍耐都不能使他们羞于接受我们对上述罪犯的声明时，我们视教会规则比人情更珍贵，对他们宣布了绝罚令，给予他们六个月的悔改、补赎及表达与公会议联合与和谐意愿的期限。若在规定时间内他们回转，同意与兄弟和好，并同意关于\Person{阿埃提乌斯}的法令，我们决定他们应被接纳入教会，恢复其在公会议中的权威及我们的友爱。但如果他们顽固坚持，偏爱人情胜过教会法规和我们的友爱，那么我们判定他们被解除主教职。若他们遭贬黜，则必须另立主教，使合法教会得享有序并与自身统一，同时各国所有主教以同一心意和同一谋略宣示同一教义，保持爱德的联系。\par

为使您知悉公会议的法令，我们特派此函致阁下，恳求您遵守之，并藉基督的恩宠正确平安地治理您辖下的各教会。\par

% structure: p0159 source=h2 target=subsection reason=relative-heading-rank; base=h2

\subsection{第二十五章　使欧诺米派与亚略派分离的原因}

\Person{欧诺米}在其著作中赞扬\Person{阿埃提乌斯}，称他为天主之人，并以许多溢美之词尊崇他。然而他当时与那些否定\Person{阿埃提乌斯}的一派关系密切，并因他们而得选为主教。\par

如今，\Person{尤多克修斯}和\Person{阿卡基乌斯}的追随者——他们已同意在\Place{色雷斯尼凯}制定的法令（本史前已提及）——任命了其他主教，接替\Person{巴西略}和\Person{埃莱乌西乌斯}党徒的各教会。其余细节，我以为不必详述。我只打算叙述与\Person{欧诺米}有关的事。\par

因为当\Person{欧诺米}在\Person{埃莱乌西乌斯}在世时夺取了\Place{基济科斯}主教区后，\Person{尤多克修斯}敦促他隐藏自己的意见，不向那些寻找借口迫害他的党徒透露。\Person{尤多克修斯}之所以提出此劝告，既因他知道该教区信仰纯正，也因他体验到\Person{君士坦丁}对主张天主独生子为受造物之党徒所显的愤怒。\Person{尤多克修斯}对\Person{欧诺米}说：\Quote{我们先等待时机；一旦到来，我们将宣讲现在所隐瞒的；教导无知的；争取或强迫或惩罚我们的对手。}\Person{欧诺米}听从了这些建议，以晦暗的阴影掩藏其不虔的教义。他的听者中那些受天主圣言培育的人清楚看出，其言辞表面下隐藏着污秽的谬误溃烂。\par

但他们虽心中苦恼，仍以为公开抗议不如谨慎行事，遂戴上异端假面具，前往主教私宅恳求他体恤那些被不同教义所左右之人的苦痛，并清楚阐释真理。欧诺弥受此鼓舞，遂吐露了自己暗中所持的见解。代表团接着指出，阻止整个教区分享真理是不公平的，甚至是完全错误的。藉着这些及类似的说辞，他被诱导在教会的公开聚会中暴露了自己的亵渎。于是，他的反对者们满腔怒火赶赴君士坦丁堡；首先在欧多克修面前控告他，欧多克修拒绝接见他们后，他们便求见皇帝，哀叹他们的主教在他们中间所造成的败坏。他们说：\Quote{欧诺弥的讲道比亚略的亵渎更不虔敬。}君士坦提乌震怒，命欧多克修传唤欧诺弥，若证实其罪，便褫夺其主教之职。欧多克修当然一再受到控告者的催促，却继续拖延行动。于是他们再次大声疾呼向皇帝申诉，说欧多克修未遵行任何一项谕旨，且全然漠视一座重要城市被交付给欧诺弥的亵渎。君士坦提乌便对欧多克修说：你若不去提拿欧诺弥并审讯他，一旦证实对他的指控，便惩罚他，我就将你流放。这一威胁吓住了欧多克修，于是他写信给欧诺弥，让他逃离基齐库，并告诉他只能怪自己未听从给他的暗示。欧诺弥惊恐之下退去，却不堪忍受耻辱，力图将自己的背叛归咎于欧多克修，声称他和艾提乌都受到了残忍对待。从那时起，他建立了自己的派别，所有认同他思想并谴责他背叛的人都脱离了欧多克修，与欧诺弥联合，至今仍以他的名字为名。于是欧诺弥成为异端始祖，凭自己特殊的罪过加深了亚略的亵渎。他建立自己的派别，是因为他受野心奴役，事实清楚证明了这一点。因为艾提乌被定罪流放时，欧诺弥拒绝与他同行，尽管他称其为师傅和天主的人，却仍与欧多克修保持密切联系。\par

但轮到他时，他为他的罪恶付出了代价；他不服从公会议的决定，却开始祝圣主教和司铎，尽管他自己已被剥夺了主教品秩。这些就是君士坦丁堡所发生的事。\par

% structure: p0165 source=h2 target=subsection reason=relative-heading-rank; base=h2

\subsection{第二十六章 论尼西比城之围及雅各伯主教宗徒般的生活}

波斯王沙普尔向罗马人开战时，君士坦提乌集结军队向安提约基雅进军。然而，击退敌人的并非罗马军队，而是罗马军中虔诚者所崇拜的天主。这场胜利如何取得，我接下来便要叙述。\par

尼西比，有时被称为安提约基雅·米格多尼亚，位于波斯与罗马帝国的边界。在尼西比，我方才提到的雅各伯同时是主教、守护者和统帅。他是一位真正闪耀着宗徒品性恩宠的人。他非凡的、令人难忘的奇迹，我已在\WorkTitle{宗教历史}中详尽叙述，此处认为无需赘述。\par

然而，为着眼前的话题，我将记述一事。\Person{雅各伯}治理的城池当时已归罗马人所有，正被波斯军队围困。封锁持续了七十天。\Quote{赫勒波利斯}及许多其他攻城器械被推至城墙前。城邑被栅栏和壕沟围住，但仍坚守未破。流经城中央的\Place{米格多尼乌斯河}，最后波斯人在上游相当远处截断其水流，并加高两岸堤岸以蓄水。当他们见到大量积水并已开始漫过堤坝时，便突然将水如攻城器械般冲向城墙。冲击力极大；壁垒无法承受，崩塌倒下。城墙另一侧，即\Place{米格多尼乌斯河}出口处，遭遇同样命运；它抵挡不住冲击，也被冲毁。\Person{沙普尔}一见此景，便以为能攻占余下的城池，当日他休息等待泥泞干涸、河流可涉。次日他全力进攻，指望从缺口进入城中。但他发现城墙两侧都已重建，他的一切辛劳付诸东流。因为那位圣人通过祈祷，使军队和其余市民充满勇力，他们既重建了城墙，又抵挡了器械，击退了进犯之敌。而这一切，他并未靠近城墙，而是在教堂内恳求万有的主。此外，\Person{沙普尔}不仅对城墙重建之速感到惊异，也为另一景象所震慑。因为他看见城垛上站着一位君王模样的人，浑身闪耀着紫袍和王冠。他以为那是罗马皇帝，便以死威胁侍从，怪他们未报告皇帝驾临；但侍从坚持说他们报告属实，\Person{君士坦提乌斯}在\Place{安提约基雅}，他这才领悟了异象的意义，喊道：\Quote{他们的天主在为罗马人作战。}于是这可怜人在盛怒之下向空中掷出一支标枪，尽管他知道不能击中无体之灵，却无法克制怒气。因此，卓越的\Person{厄弗辣因}（他是叙利亚人中最好的作家）恳求神圣的\Person{雅各伯}登上城墙去察看那些蛮族，并向他们施放诅咒之箭。于是这位神人同意，爬上一座塔楼，但当他看到那数不清的军队时，他没有发出别的诅咒，只求让蚊蚋和蠓虫降临到他们身上，借此微小动物，他们可以认识到罗马保护者的权能。他的祈祷之后，成群的蚊蚋和蠓虫飞来；它们充满大象的鼻腔，以及马匹和其他动物的耳朵鼻孔。这些小生物的攻击难以忍受，它们挣断缰绳，掀翻骑手，使得队伍混乱。波斯人弃营而逃。于是这位可怜的君主通过一次轻微而仁慈的惩罚，认识到了保护虔敬者的天主的大能，便率军班师返国，这次围攻所收获的只是耻辱而非胜利。\par

% structure: p0169 source=h2 target=subsection reason=relative-heading-rank; base=h2

\subsection{第二十七章　\Place{安提约基雅}会议及其反对\Person{圣默肋提乌斯}之事}

此时，\Person{君士坦提乌斯}驻跸于\Place{安提约基雅}。波斯战争已经结束；有一段和平时期，他再次召集主教们，目的是要他们否认\Quote{同质}(\OriginalTerm{同质}{homoousios})和\Quote{异质}(\OriginalTerm{异质}{heteroousios})两种公式。\Person{莱昂提乌斯}死后，\Person{欧多克修斯}夺取了安提约基雅的主教座，但在他被驱逐并非法立足（经过多次主教会议之后）于君士坦丁堡后，安提约基雅教会便没有牧人了。因此，聚会的主教们（从各地大量聚集）宣称，他们的首要任务是给羊群提供一位牧者，然后与他一起商讨信仰问题。恰巧，统治亚美尼亚某城的\Person{梅利提乌斯}正因治下百姓的不顺从而忧心，当时闲居别处。亚略派信徒以为\Person{梅利提乌斯}与他们思想相同，是他们的教义的拥护者。于是他们请求\Person{君士坦提乌斯}将安提约基雅教会的缰绳交到他手中。确实，他们为了建立自己的不虔敬，肆无忌惮地践踏一切法律；非法行为已成为他们亵渎的基础；这并非他们不规则行径的孤例。另一方面，维护宗徒教义的人完全了解伟大的\Person{梅利提乌斯}的纯正，清楚知道他纯洁的品格和丰富的德性，达成一致表决，并采取措施立即将他们的决议写成文书并由全体签署。双方将此文件作为妥协的契约，委托给一位捍卫真理的高贵主教——\Place{萨莫萨塔}的\Person{欧塞比乌斯}保管。当伟大的\Person{梅利提乌斯}奉诏抵达时，所有高级圣职者、教会其他各级人员以及全城居民都出来迎接他。那里还有犹太人和外邦人，都急切地想见到伟大的\Person{梅利提乌斯}。皇帝命\Person{梅利提乌斯}和其他能言善辩的人向群众讲解经文\BibleQuote{上主在造化之初，在太初创造万物之前，就生了我}\BibleRef{箴 8:22}，并命令熟练的文书当场记下每个人的话，以为这样他们的训导会更精确。首先是\Place{劳迪基雅}的\Person{乔治乌斯}发泄了他污秽的异端。之后是\Place{凯撒勒雅}的\Person{阿卡基乌斯}提出了一种折中的教义，虽远离敌人的亵渎，但未能保持宗徒教义的纯正无瑕。然后伟大的\Person{梅利提乌斯}站起来，展示了信仰准则的不可弯曲的线条，因为他像木匠使用规尺一样使用真理，避免了过度和不足。于是众人热烈鼓掌，恳求他给他们一个简短的教导总结。于是他伸出三根手指，收回两根，剩下一根，说出那句令人难忘的话：\Quote{在思想上是三位，但我们论及一位。}\par

对此教导，那些心中藏着亚略瘟疫的人磨砺他们的舌头，开始散布巧妙的诽谤，宣称\Person{梅利提乌斯}是\OriginalTerm{撒伯里乌派}{Sabellian}。他们就这样说服了反复无常的皇帝（他如同著名的\Place{欧里普斯}海峡，轻易地改变流向），使他将\Person{梅利提乌斯}遣返回乡。\par

\Person{欧佐伊乌斯}，一个亚略教义的公开辩护者，立刻被擢升接替其位；此人正是当年身为执事时，伟大的\Person{亚历山大}与\Person{亚略}一同贬黜的那一位。如今，保持纯正的那部分百姓与不纯正的分离，聚集在位于城中称为\Place{帕莱亚}地区的宗徒教会。\par

事实上，自对著名的\Person{欧斯塔提乌斯}发动攻击以来，他们忍受亚略主义的可憎之事已达三十年，期待某种有利的变化。但看到不虔敬愈演愈烈，忠于宗徒教义的人公开受攻击并受秘密阴谋威胁，\Person{梅利提乌斯}被流放，异端的捍卫者\Person{欧佐伊乌斯}被立为主教取代他时，他们想起了对\Person{罗特}说的话：\BibleQuote{逃命吧}；此外，福音的法律明确规定：\BibleQuote{若是你的右眼使你跌倒，剜出它来，丢掉它}。主对脚手也制定了同样的法律，并补充说：\BibleQuote{你失去一个肢体，比全身被投入地狱里，为你更好。}\par

就这样造成了教会的分裂。\par

% structure: p0175 source=h2 target=subsection reason=relative-heading-rank; base=h2

\subsection{第二十八章　关于\Place{萨莫萨塔}的主教\Person{欧塞比乌斯}}

前面提到的那位令人钦佩的\Person{欧塞比乌斯}，他曾受托保管共同决议，当他看到盟约被违反，便返回了自己的教区。这时，一些人对那份书面文件感到不安，说服\Person{君士坦提乌斯}派遣信使去取回它。皇帝便派了一名使用驿站快马、迅速传递消息的官员。他到达后传达了皇帝的口信，但令人钦佩的\Person{欧塞比乌斯}说：\Quote{除非受托保管文件的全体会众指示我，我不能交出它。}这个答复报告给了皇帝。他勃然大怒，又派人去告诉\Person{欧塞比乌斯}，命令他交出文件，并补充说如果他拒绝，就下令砍掉他的右手。但他写这些只是为了吓唬主教，因为送信的使者奉命不要执行威胁。当\Person{欧塞比乌斯}打开信，看到皇帝威胁的惩罚时，他伸出右手和左手，让那人把两只手都砍掉。他说：\Quote{这份明确证明亚略派邪恶的文书，我绝不交出。}\par

当\Person{君士坦提乌斯}得知这勇敢的决心时，他惊愕不已，不禁赞叹不已；因为连敌人也因伟大行为的力量而不得不钦佩对手的成功。\par

此时，\OriginalTerm{君士坦狄乌斯}{Constantius}得知他先前立为\Place{欧罗巴}的\OriginalTerm{凯撒}{Cæsar}的\OriginalTerm{尤利安}{Julian}意图夺取王权，并集结军队反抗其主上。他因此从\Place{叙利亚}出发，死于\Place{基里基亚}{Cilicia}。他也没有得到他父亲为他留下的助佑；因为他没有完整保存他父亲虔诚的遗产，因而痛悔自己信仰的改变。\par

\SourceRecord{Translated by Blomfield Jackson.；Nicene and Post-Nicene Fathers, Second Series；Vol. 3.；Edited by Philip Schaff and Henry Wace.；Buffalo, NY: Christian Literature Publishing Co.,；1892.；<http://www.newadvent.org/fathers/27022.htm>.}

% structure: p0001 source=h1 target=section reason=page-title

\section{\WorkTitle{教会史（第三卷）}}

% structure: p0002 source=h2 target=subsection reason=relative-heading-rank; base=h2

\subsection{第一章  论尤利安执政；他自幼受虔诚教育而堕入不敬；以及他起初隐藏不敬，后表露无遗}

如前所述，\Person{君士坦提乌斯}叹惋悲伤而终，因其偏离了父辈的信德。\Person{尤利安}自欧洲渡往亚洲途中闻其死讯，欣然承袭王权，再无对手。\par

早年，\Person{尤利安}尚为孩童时，与其兄\Person{加卢斯}一同接受了纯洁而虔诚的教导。\par

在其少年及青年初期，他继续持守同一教义。\Person{君士坦提乌斯}因恐亲族觊觎帝位，诛杀之；\Person{尤利安}因惧怕其堂兄，得以列入读经员行列，常在教会集会中向众人朗读圣经。\par

他还建造了一座殉道者堂；但殉道者见其背教，拒受奉献；因地基如建者心意般不稳，圣堂在祝圣前便倒塌了。\Person{尤利安}的童年及青年便是如此。然而当\Person{君士坦提乌斯}因对抗\Person{玛格嫩提乌斯}之战而西行时，他立了\Person{加卢斯}为东方凯撒——加卢斯富有虔诚，且持守到底。此时\Person{尤利安}抛却先前有益之虑，因不义的自信而一心夺取帝国权杖。于是，途经希腊时，他寻访预言者和占卜者，欲知能否得偿所愿。遇一人，许诺预言此事，领他进入一偶像庙宇，引入内殿，召唤欺诳的魔鬼。魔鬼以惯常形貌显现时，恐惧迫使\Person{尤利安}在额上划十字圣号。魔鬼一见主得胜的记号，便想起自己的溃败，立即逃散。术士得知逃跑之因后责备他；\Person{尤利安}承认恐惧，并感叹十字架的大能，因魔鬼不能忍受见其记号而逃跑。术士说：\Quote{善士，切莫如此想；它们并非如你所言因惧怕而逃，而是因厌恶你所行而离去。}如此，他欺骗了这可怜人，给他行秘仪，使他充满可憎之事。\par

如此，权欲使这可怜人丧失了一切真宗教。然而，获得最高权力后，他隐藏不敬良久；尤其担心那些受过真宗教原则教导的军队——先由杰出的\Person{君士坦丁}（他使他们脱离旧错，训导于真理之道）所教导，后由其诸子（他们巩固了父亲所授的教导）所巩固。因为\Person{君士坦提乌斯}虽被其左右之人误导，不接纳\OriginalTerm{同质}{\GreekText{ὁ} \GreekText{μοούσιον}}一词，却真诚接受了其含义——他称天主圣言为真子，在万世之前由父所生，公开反对那些胆敢称祂为受造物者，并绝对禁止偶像崇拜。\par

我还要叙述他另一项高尚事迹，作为其热心神业的确证。在对抗\Person{玛格嫩提乌斯}的战役中，他曾集合全军，劝他们全体一同参与神圣奥迹，\Quote{因为，}他说，\Quote{生命的终结总是不确定的，尤其在战争中，无数箭矢从双方射出，刀剑、战斧及其他兵器袭击人，造成暴死。因此，各人应当穿上那件最需要的宝贵外袍，以备来世之用。若有人此刻不愿穿上此衣，就请离开这里回家去。我不愿与那些在神圣礼仪中无份无关的士兵作战。}\par

% structure: p0009 source=h2 target=subsection reason=relative-heading-rank; base=h2

\subsection{第二章  论主教们的回归与\Person{保利努斯}的祝圣}

\Person{尤利安}对这些情况了如指掌，并未表露灵魂的不敬。为吸引所有主教顺从于其统治，他下令甚至那些被\Person{君士坦提乌斯}逐出教会、流亡帝国最偏远之地者，也可返回各自的教会。于是，此谕令颁布后，神圣的\Person{梅利提乌斯}回到了安提约基雅，而远负盛名的\Person{阿塔纳修}回到了亚历山大里亚。\par

但\Person{尤西比乌斯}、意大利的\Person{希拉利乌斯}以及管辖\Place{撒丁岛}的\Person{路济弗尔}，因\Person{君士坦提乌斯}之命被放逐至埃及边界的\Place{泰拜德}，如今他们与其余意见相同者相遇，并声称各教会应当和谐一致。因他们不仅遭受对手的攻击，且彼此不睦。在\Place{安提约基雅}，教会的健全肢体已分裂为二：一方面，那些从一开始即为可敬的\Person{尤斯塔修斯}而与其余人分离者，自行聚集；另一方面，那些与令人钦佩的\Person{梅利提乌斯}一道远离亚略派者，则在\Quote{所谓的Palæa}举行礼仪。双方均使用同一信经，因二者皆为尼西亚会议所定教义的捍卫者。使他们分离的唯有彼此间的争执和对各自领袖的尊崇；即便其中一人去世，也未止息纷争。\Person{尤斯塔修斯}在\Person{梅利提乌斯}被选之前已去世，而正统派在\Person{梅利提乌斯}被放逐及\Person{尤佐伊乌斯}当选后，便与不虔者的共融分离，自行聚集；称为\OriginalTerm{尤斯塔修斯派}{Eustathians}的那一派不愿与这些联合。\Person{尤西比乌斯}派与\Person{路济弗尔}派者设法寻找使二者合一的方法。因此，\Person{尤西比乌斯}恳求\Person{路济弗尔}前往\Place{亚历山大里亚}，就此事与伟大的\Person{亚大纳削}商议，而他本人则打算承担促成和解的工作。\par

\Person{路济弗尔}却没有前往\Place{亚历山大里亚}，而是去了\Place{安提约基雅}。在那里，他提出了许多促进双方和睦的论据。由司铎\Person{保利努斯}领导的尤斯塔修斯派坚持反对。见此情形，\Person{路济弗尔}采取了不当之举，祝圣\Person{保利努斯}为其主教。\par

\Person{路济弗尔}的此举延长了纷争，持续了八十五年，直至最值得赞扬的\Person{亚历山大}主教任期。\par

\Place{安提约基雅}教会的航舵一交到他手中，他即尝试各种方法，以极大的热忱和精力促进合一，从而将断裂的肢体重新接合于教会整体。然而当时，\Person{路济弗尔}却使争吵加剧，并在\Place{安提约基雅}逗留了相当长的时间；而\Person{尤西比乌斯}抵达当地后，得知错误的医治使疾病极难治愈，便乘船西去。\par

\Person{路济弗尔}返回\Place{撒丁岛}后，对教会教义作了某些增补，接受这些的人便以他命名，并在相当长的时间内被称为\OriginalTerm{路济弗尔派}{Luciferians}。但过了一段时间，这教义之火也熄灭了，被人遗忘。这就是主教们返回后所发生的事。\par

% structure: p0016 source=h2 target=subsection reason=relative-heading-rank; base=h2

\subsection{第三章　论异教徒从\Person{尤里安}获得权力后对基督徒所做之事的数量与性质}

当\Person{尤里安}公开其不虔时，各城充满了纷争。被偶像崇拜的欺骗所奴役的人壮起胆来，打开偶像的庙宇，开始举行那些本应从人类记忆中消失的污秽仪式。他们再次点燃祭坛之火，以牺牲的血污秽土地，以燔祭的烟雾污染空气。他们因所侍奉的魔鬼而狂乱，像狂欢者般在街道上奔跑，以低级的舞台笑料攻击圣人，并以不洁游行的一切侮辱和粗鄙行为对待他们。\par

另一方面，敬虔的追随者无法容忍他们的亵渎，以侮辱回应侮辱，并试图驳斥他们所尊崇的谬误。而行不义的人对此感到气愤，君主所给予的自由鼓励了他们的大胆，便在基督徒中给予致命打击。\par

皇帝本应顾及臣民的和平，但他因极度不虔，自己却以相互的愤怒使人民疯狂。对于残暴者对抗和平者所胆敢行的事，他视而不见，并将重要的民事和军事官职委托给野蛮不虔之人，这些人虽不敢公开强迫真敬虔的热爱者献祭，却以各种侮辱对待他们。此外，伟大的\Person{君士坦丁}所赐予神圣职务的一切殊荣，\Person{尤里安}均予剥夺。\par

要讲述当时偶像欺骗的奴隶所胆敢行的一切事，需要一部单记这些罪行的历史，但我将从大量事例中选取几件。在巴勒斯坦的\Place{阿斯卡隆}和\Place{迦萨}，有司铎级男子和终生守贞的妇女被剖腹，塞满大麦，喂给猪吃。在同属该地区的\Place{塞巴斯特}，\Person{洗者若翰}的棺椁被打开，骸骨被焚烧，灰烬被抛散。\par

谁能不流泪地讲述在腓尼基所行的卑劣之事？在\Place{黎巴嫩}附近的\Place{赫利奥波利斯}，住着一位名叫\Person{基里尔}的执事。在\Person{君士坦丁}统治时期，他因神圣的热忱，打碎了那里许多受人崇拜的偶像。现在，那些名声败坏之人念念不忘此事，不仅杀了他，还剖开他的肚子，吞食了他的肝脏。但他们的罪行不能逃过全视之眼，他们受到了应得的惩罚：所有参与这令人憎恶的邪恶之人牙齿全部同时脱落，舌头也腐烂并脱落；他们还丧失了视力，以所受的苦难宣告了圣善的能力。\par

在邻近的\Place{厄梅撒}城，他们将新建的教堂奉献给了那具有女性形态的\Person{狄奥尼索斯}，并在其中设立了他那荒谬的雌雄同体像。在\Place{色雷斯}著名城市\Place{多里斯托隆}，胜利的斗士\Person{埃米利亚诺斯}被色雷斯全境总督\Person{卡皮托利诺斯}投入熊熊燃烧的火堆。然而，若要真正以戏剧般的庄重叙述\Place{阿勒图萨}主教\Person{马尔库斯}的悲惨命运，则需要\Person{埃斯库罗斯}或\Person{索福克勒斯}的口才。在\Person{君士坦提乌斯}时代，他曾摧毁某座偶像庙宇，并在其原址建造了一座教堂；阿勒图萨人一得知\Person{尤利安}的意图，便公开表露了敌意。起初，遵照福音的训导，\Person{马尔库斯}试图逃脱；但当他得知自己的一些同乡因他而被捕时，他便返回并把自己交给了那些嗜血之人。他们抓住他后，既不怜悯他的高龄，也不尊重他深厚的德行；尽管他以教义和生活之美而著称，但他们首先剥光了他的衣服，击打他，在他每一肢体上施以鞭笞，然后把他扔进污秽的下水道，拖出来后，又交给一群少年，吩咐他们用笔毫不留情地刺他。之后，他们把他放进一个篮子，涂上腌汁和蜂蜜，在盛夏的露天中悬挂起来，招引黄蜂和蜜蜂前来赴宴。他们这样做的目的是强迫他要么恢复他所摧毁的庙宇，要么支付重建的费用。然而，\Person{马尔库斯}忍受了所有这些痛苦的折磨，并声称他不会同意他们的任何要求。他的敌人以为他因贫穷而无力支付，便减免了一半的要求，命他支付其余部分；但\Person{马尔库斯}高悬在空中，被笔刺伤，被黄蜂和蜜蜂吞噬，却不仅没有显出痛苦的样子，反而用不断的嘲弄讥讽他那些不敬的折磨者，说道：\Quote{你们是世俗的，属于大地；我却是崇高的，被高举的。}最后，他们只求他支付一小部分钱；但他却说：\Quote{给一个奥波尔与给全部一样不敬。}于是，他们沮丧地放了他，不禁钦佩他的坚毅，因为他的话给他们上了新的圣洁之课。\par

% structure: p0023 source=h2 target=subsection reason=relative-heading-rank; base=h2

\subsection{第四章　\Person{尤利安}颁布反对基督徒的法律}

\Emph{无数}其他的恶行在那时由恶人向义人施行于陆地和海上，遍及全世界，因为天主的敌人如今毫不掩饰地开始制定反对真宗教的法律。首先，他禁止\Quote{加里利人}的儿子（他试图这样称呼救主的崇拜者）参与诗歌、修辞和哲学的学习，因为他说，用谚语的话来说：\Quote{我们被自己翅膀上的箭射中，}因为我们从自己的书籍中取得武器来与我们作战。\par

此后，他又颁布另一道法令，命令将加里利人从军队中驱逐出去。\par

% structure: p0026 source=h2 target=subsection reason=relative-heading-rank; base=h2

\subsection{第五章　圣\Person{亚大纳削}的第四次流亡与逃亡}

此时，真理的胜利斗士\Person{亚大纳削}又经历了一次危险，因为魔鬼不能容忍他口舌与祈祷的力量，于是武装了他们的臣仆来辱骂他。他们发出许多声音，恳求邪恶的冠军流放\Person{亚大纳削}，并且还加上这一点：如果\Person{亚大纳削}留下，就没有一个异教徒会留下，因为他会把他们都争取到他那一边。\Person{尤利安}被这些恳求所动，不仅判\Person{亚大纳削}流放，而且判处死刑。他的百姓颤抖，但据记载，他预言了这场风暴的迅速消散，因为他说：\Quote{这是一片很快就会消逝的云。}然而，他一得知执行皇帝谕令的人到来，便离开了；他在河岸找到一条船，出发前往\Place{提拜德}。被指定执行他死刑的军官得知他逃走，便拼命追赶；但他的一位朋友赶在前头，告诉他军官正在迅速追来。于是他的一些同伴恳求他逃往旷野，但他命令舵手将船头转向\Place{亚历山大里亚}。于是他们划船迎向追捕者，而执行死刑的人来了，他说：\Quote{亚大纳削离此多远？}\Quote{不远，}\Person{亚大纳削}说，就这样摆脱了他的敌人，而他自己则返回\Place{亚历山大里亚}，在那里躲藏直到\Person{尤利安}统治的剩余时间。\par

% structure: p0028 source=h2 target=subsection reason=relative-heading-rank; base=h2

\subsection{第六章　论\Person{阿波罗}与\Place{达佛涅}，以及圣\Person{巴彼拉}}

\Person{尤利安}想要对波斯人发动战役，便派遣他最信任的官员前往罗马帝国内的所有神谕所，而他自己则作为恳求者前往\Place{达佛涅}的皮提亚神谕所，恳求它向他启示未来。神谕回答说，附近躺着的尸体正在阻碍占卜；必须先将它们移到别处；然后他才会发出预言，因为他说：\Quote{如果这树林不洁净，我就什么也不能说。}那时，那里正安放着胜利的殉道者\Person{巴彼拉}和与他一同光荣殉道的少年们的遗骸，而说谎的先知显然因\Person{巴彼拉}的圣洁影响而无法说出他惯常的谎言。\Person{尤利安}知道这一点，因为他过去的虔诚曾教导他胜利的殉道者的能力，于是他没有移动任何别的尸体，只命令基督的敬拜者迁移胜利的殉道者的遗骸。他们满怀喜乐地走向树林，将棺材放在车上，走在它前面，带领着浩大的人群，唱着\Person{达味}的圣咏，每到间歇便呼喊：\Quote{愿所有崇拜铸像的人都蒙羞。}因为他们明白，迁移殉道者的遗骸意味着恶魔的失败。\par

% structure: p0030 source=h2 target=subsection reason=relative-heading-rank; base=h2

\subsection{第七章　论宣信者\Person{德奥多若}}

\Person{尤利安}无法忍受这些行为带给他的羞辱，次日便下令逮捕合唱游行的首领。当时\Person{萨路斯提乌斯}任城长，是个不义的仆人，但他仍试图劝说君王不要允许那些渴望荣耀的基督徒达成他们的目的。然而当他看到皇帝无法抑制怒火时，便逮捕了一名洋溢着神圣热情的优雅青年——这青年正在广场上行走——将他当众吊在刑架上，用鞭子抽裂他的背，用钩爪般的刑具划伤他的肋部。他从黎明一直做到日暮；然后给他戴上铁链，下令把他关押在狱中。次日早晨，他向\Person{尤利安}报告了所作的事，述说了那青年的坚毅，并补充说，这件事对他们自己而言是失败，对基督徒却是胜利。天主的仇敌被这真理说服，便不再容忍如此对待，下令释放\Person{狄奥多若}——这是那位在真理之战中年轻而光荣的斗士的名字。当被问及在遭受那些最残酷最野蛮的酷刑时是否感到任何疼痛时，他回答说，起初确实感到一点疼痛，但随后有一个人出现在他面前，不断地用凉爽柔软的巾帕擦拭他脸上的汗水，并鼓励他要勇敢。\Quote{因此，}他说，\Quote{当刽子手停止时，我并不高兴，反而烦恼，因为那带给我灵魂慰藉的人也随他们离开了。}但谎言之占卜的恶魔立时增加了殉道者的光荣，暴露了自己的虚谎；因为一道从天而降的雷霆烧毁了整个神龛，把皮提亚的神像本身化为细尘——那神像是木制的，表面镀金。\Person{尤利安}的叔父、东方城长\Person{尤利亚努斯}夜间得知此事，全速骑马来到\Place{达芙涅}，急切地想给他所崇拜的神祇提供援助；但当他看到那所谓的偶像化为粉末时，便鞭打了掌管神庙的官员，因为他猜测火灾是某个基督徒所为。然而他们虽受虐待，却不能忍受说谎，坚持说火不是从下面，而是从上面开始的。此外，一些附近的乡民也前来作证，说他们亲眼看见雷霆从天而降。\par

% structure: p0032 source=h2 target=subsection reason=relative-heading-rank; base=h2

\subsection{第八章　论圣器及供给的没收}

即使恶人得知了这些事，他们仍结阵对抗万有的天主；君王命人将圣器移交国库。他用钉子封住了\Person{君士坦丁}所建大教堂的门，并宣布它对其惯常聚集的朝拜者关闭。此时该教堂由亚略派占据。东方城长\Person{尤利亚努斯}、皇室司库\Person{费利克斯}，以及掌管皇帝私库和财产的官员\Person{埃尔皮狄乌斯}——一位罗马习惯称之为\OriginalTerm{私产伯爵}{Comes privatarum}的官员——一同进入了圣殿。据说\Person{费利克斯}和\Person{埃尔皮狄乌斯}本是基督徒，但为了取悦不敬的君王而背弃了真宗教。\Person{尤利亚努斯}在圣桌上做出了极端猥亵的行为，当\Person{尤佐伊乌斯}试图阻止他时，他打了他的脸，并告诉他——据故事所述——基督徒的命运就是得不到神祇的保护。但\Person{费利克斯}凝视着圣器的华美——这些圣器因\Person{君士坦丁}和\Person{君士坦提乌斯}的慷慨而装点得辉煌——说：\Quote{看哪，玛利亚的儿子是用怎样的器皿服侍的。}但不久之后，他们就为这些疯狂而亵渎的鲁莽行为付出了代价。\par

% structure: p0034 source=h2 target=subsection reason=relative-heading-rank; base=h2

\subsection{第九章　论皇帝叔父\Person{尤利亚努斯}及\Person{费利克斯}的遭遇}

\Person{尤利亚努斯}立即患上一种痛苦的疾病；他的内脏腐烂，再也无法通过正常的排泄器官排出粪便，反而在他亵渎的那一刻，他那污秽的口成了排泄的器官。\par

据说他的妻子是一位信仰显著的妇人，她这样对丈夫说：\Quote{丈夫，你应当赞美我们的救主基督，因为他通过你的惩罚向你显示了他独特的大能。因为如果你没有经历这些天降的灾祸，你永远也不会知道你所攻击的是谁。}于是，因着这些话和他沉重的苦难，这可怜的人明白了自己疾病的原因，恳求皇帝将教堂归还给那些被剥夺的人。然而他未能得到所请，就这样结束了生命。\par

\Person{费利克斯}本人也突然被天降的鞭打击倒，整日整夜从口中呕血，因为他身体的所有血管都将汇聚的血液输向这一器官：当他的血全部流尽时，他死了，被交付于永死。\par

这些人因他们的恶行所受的惩罚就是如此。\par

% structure: p0039 source=h2 target=subsection reason=relative-heading-rank; base=h2

\subsection{第十章　论司铎之子}

一位司铎之子、自幼在邪恶中长大的青年，此时归向了真宗教。有一位以虔诚著称、被列入执事之列的夫人，是他母亲的密友。当他尚在幼年时随母亲去拜访她，她便以慈爱接待他，并敦促他归向真宗教。母亲去世后，这青年常去拜访她，并得享她一如既往的教导。那夫人的劝诫深深打动了他，他便询问自己的老师：如何才能既逃脱父亲的迷信，又有份于她所宣讲的真理？她回答说，他必须离开父亲，而尊敬那创造他父亲及他自己本身的主；他必须寻找另一座城市，在那里隐藏起来，躲避那邪恶皇帝的暴力；并答应为他安排此事。于是青年说：\Quote{从今以后，我要来把我的灵魂交托给你。}不久之后，犹利安来到达芙涅庆祝一个公共节日。同行的有青年的父亲，他既是司铎，又习惯于随侍皇帝；他们的父亲带着青年和他的兄弟，他们被指派负责圣殿的服务，并负有礼仪性地洒在御膳上的职责。达芙涅节期一般持续七天。第一天，青年站在皇帝的卧榻旁，按规定洒在肉食上，将它们彻底玷污。然后他全速跑向安提约基雅，来到那位可敬的夫人那里，说：\Quote{我到你这里来了；我遵守了我的诺言。你当关心每个人的救恩，实现你的承诺。}她立刻起身，将青年领到天主的仆人默肋提乌斯那里，后者吩咐他先在客栈楼上待一段时间。他的父亲在达芙涅到处寻找那孩子后，返回城里，搜遍大街小巷，四处张望，渴望能撞见他的儿子。最后他来到神圣的默肋提乌斯所住的客栈；抬头一看，见他儿子正从格子窗里窥视。他便跑上去，拉他下来，把他带回家。他先用许多鞭子抽打他，然后把烧红的铁签烙他的脚、手和背，然后把他关进卧室，从外面锁上门，返回达芙涅。我亲耳听那人年老时讲述此事，他还补充说，他受到恩宠的感动和充满，便打碎了他父亲所有的偶像，嘲笑它们无能为力。后来，当他想起自己所行的事，害怕父亲回来，便恳求他的主基督欣允他的行为，打破门闩，开启大门。他说：\Quote{因为我正是为了你才如此受苦、如此行事。}他告诉我：\Quote{正当我这样说时，门闩脱落，门扇大开，我跑回我的女导师那里。她给我穿上女装，用她的有篷马车带我到神圣的默肋提乌斯那里。他把我托付给耶路撒冷的主教，当时是济利禄，我们夜间出发前往巴勒斯坦。}犹利安死后，这青年也引导他的父亲走上了真理之路。他把这件事连同其他事情一起告诉了我。这样，这些人就被引导认识了天主，并成为有份于救恩的人。\par

% structure: p0041 source=h2 target=subsection reason=relative-heading-rank; base=h2

\subsection{第11章 论圣殉道者犹文提努斯与玛克西米努斯}

这时，犹利安更加无节制，或者说是更少羞耻地开始武装自己反对真宗教，表面上戴着温和的面具，却一直在准备罗网和陷阱，凡被它们欺骗的人都落入邪恶的毁灭中。他首先用污秽的祭献污染城里和达芙涅的井泉，使每个使用泉水的人都沾染可憎之物。然后他彻底污染了市场上陈列的食物，因为面包、肉、水果、蔬菜以及各种食品都被洒过。那些蒙救主之名所召的人看见所行的事，就叹息哀哭，表示憎恶；然而他们还是吃了，因为他们记起宗徒的法律：\BibleQuote{凡在肉市上出卖的，你们只管吃，不要为良心而查问什么。}军队中有两位军官，是皇帝随从的盾牌手，在一次宴会上，他们以颇为激烈的言辞哀叹所行之事的可憎，并采用了巴比伦光荣青年的可赞美的话语：\BibleQuote{你竟把我们交在一个不虔敬的、比地上所有民族更背教的君王手中。}宾客中有一人告发了此事，皇帝逮捕了这两位可敬的人，并试图通过审问他们来查明他们所说的话。他们欣然接受皇帝的讯问，视之为自由发言的机会，以崇高的热忱回答说：\Quote{大人，我们是在真宗教中长大的；我们服从了最优越的法律，即君士坦丁及其子们的法律；现在我们看见世界充满污染，食物和饮料都被可憎的祭献玷污，我们哀叹。我们在家中哀悼这些事，如今在陛下面前表达我们的悲痛，因为这是您统治中唯一使我们不满的事。}那位被称为最温和、最富哲理的皇帝一听到这话，就摘下了温和的面具，露出了不虔敬的面孔。他命令对他们施加残忍痛苦的鞭打，并剥夺了他们的生命；或者不如说，使他们脱离那悲伤的时期，赐予他们胜利的冠冕？他假装惩罚他们并非因为真宗教（他们正是因此而真正被杀害），而是因为他们的傲慢，因为他宣称他们因侮辱皇帝而受罚，并下令将此消息公开传播，从而吝惜给予这些真理的斗士以殉道者的名号和荣誉。一人名叫犹文提努斯，另一人名叫玛克西米努斯。安提约基雅城尊崇他们为真宗教的捍卫者，将他们安葬在一座宏伟的坟墓中，直到今日，他们仍以一年一度的节日受到尊崇。\par

其他担任公职且有声誉的人同样使用了类似的勇敢言词，赢得了相似的殉道荣冠。\par

% structure: p0044 source=h2 target=subsection reason=relative-heading-rank; base=h2

\subsection{第十二章 论伟大的皇帝\Person{瓦伦提尼安}}

\Person{瓦伦提尼安}不久之后成为皇帝，当时他是一位护民官，指挥驻扎在皇宫的禁卫军。他毫不掩饰自己对真信仰的热忱。有一次，当那位昏聩的皇帝在庄严的游行中进入\Place{福尔图娜神庙}的圣所时，庙门两侧站着庙仆，如他们以为的那样，用洒水刷净化所有进入的人。瓦伦提尼安走在皇帝前面，注意到有一滴水掉在了自己的斗篷上，便用拳头打了侍从一下，\Quote{因为，}他说，\Quote{我不是被净化，而是被玷污了。}为此事他赢得了两个帝国。可咒骂的\Person{尤利安}看到所发生的事后，将他遣送到沙漠中的一座堡垒，命他留驻在那里；但过了一年零几个月，他获得了帝国，作为宣认信仰的赏报，因为公义的审判者不仅在来世尊敬那些关心神圣事物的人，而且有时甚至在今生就为善行赐予报酬，通过祂现在丰盛地赐予的，来加强对将来赏报的希望。\par

但那位暴君又设下另一计谋反对真理。当他按照古老习俗坐上皇座，向他的军队分发黄金时，他违反习俗，让人抬进一个装满热炭的祭台，并将乳香放在桌上，命令每个将要领受黄金的人先向祭台投乳香，然后从他自己的右手领取黄金。大多数人完全没有意识到这个陷阱；但那些事先得到警告的人假装生病，逃过了这个残酷的圈套。另一些人因贪图金钱而轻视自己的救恩，还有一群人因怯懦而放弃了信仰。\par

% structure: p0047 source=h2 target=subsection reason=relative-heading-rank; base=h2

\subsection{第十三章 论其他宣信者}

在这次致命的金钱分发之后，一些领受者正在一同宴饮。其中一人拿起杯子，在它上面划救恩标记之前没有喝。\par

一位宾客为此责备他，说这与刚才所做的事完全矛盾。\Quote{什么，}他说，\Quote{我做了什么矛盾的事？}于是有人提醒他关于祭台和乳香，以及他否认信仰的事；因为这些都是与基督徒的宣认相违背的。众人听到这话，多数宴饮者都呻吟哀叹，扯下大把头发。他们从宴席起身，跑过广场，大声宣称他们是基督徒，被皇帝的诡计欺骗了，他们撤回自己的背教行为，并准备尽力挽回他们无意中遭受的失败。他们喊着这些话跑向皇宫，大声谴责暴君的诡计，恳求将他们投入火焰，好使他们正如被火玷污一样，也藉火得以洁净。这一切话语使那恶人失去理智，一时冲动下令斩首他们；但当他们被带出城外时，一大群人开始跟随着他们，惊叹他们的坚毅，并为他们为真理的勇敢而欢喜。当他们到达通常处决罪犯的地点时，他们中最年长的人恳求刽子手先砍下最年轻者的头，免得他因目睹其他人的被杀而失去勇气。他刚跪在地上，刽子手拔出刀，就有一个人跑来宣布缓刑，还在远处就喊着叫停下行刑。这时最年轻的兵士因免死而痛苦。\Quote{啊，}他说，\Quote{\Person{罗马努斯}（他的名字叫罗马努斯）不配被称为基督的殉道者。}那卑鄙的伎俩者停止行刑，是出于嫉妒：他吝于让信仰的捍卫者获得荣誉。他们的判决减为流放，被驱逐到城墙之外和帝国最偏远的地区。\par

% structure: p0050 source=h2 target=subsection reason=relative-heading-rank; base=h2

\subsection{第十四章 论公爵\Person{阿尔特米}。论女执事\Person{普布利亚}及其神圣勇敢}

\Person{阿尔特米}指挥在埃及的军队。他在\Person{君士坦丁}时代获得这一指挥权，并毁坏了大部分偶像。为此，\Person{尤利安}不仅没收了他的财产，还下令将他斩首。\par

这些以及诸如此类的事，正是那被不虔敬之人描述为最温和、最不易动怒的人所行的。\par

我现在将在我的历史中记录一位极卓越的女子的高尚故事，因为就连妇女，以神圣的热忱武装自己，也藐视了\Person{尤利安}疯狂的愤怒。\par

在那些日子里，有一位名叫普布利亚的妇人，德高望重，因善行而著名。她短暂地受过婚姻的枷锁，并将最美好的果实献给天主，从这片沃土中生出了若望，他长期担任安提约基雅的首席司铎，并多次当选宗徒之座，但屡次拒绝此尊位。她带领一群矢志终身守贞的童女，赞颂创造并救赎她的天主。一日，皇帝经过，因她们视那毁灭者为可鄙可嘲之物，便更加高声奏乐，尤其咏唱那些嘲笑偶像无能的圣咏，并用达味的话说：\Quote{外邦的偶像是金银，是人手所造的。}在描述它们的无觉之后，她们又加上：\Quote{愿制造它们并信赖它们的人成为像它们一样。}尤利安听见，非常恼怒，命令她们在他经过时安静。然而她丝毫不理会他的命令，反而更加用力地歌唱。当皇帝再次经过时，她命她们高唱：\Quote{愿天主兴起，使他的仇敌四散。}于是尤利安怒令将那唱诗班的女教师带到他面前；虽然他看出她年高德劭，却既不顾怜她的白发，也不尊重她高尚的品格，反而命一些侍卫掌掴她的双耳，用暴力使她的面颊红肿。她却视这侮辱为荣耀，回家后照常以她的属灵诗歌攻击他，正如那圣咏的作者和导师驱走折磨撒乌耳的恶灵一样。\par

% structure: p0055 source=h2 target=subsection reason=relative-heading-rank; base=h2

\subsection{第十五章　论犹太人，论他们建殿的企图及从天降下的灾祸}

尤利安，他的灵魂已成为毁灭性恶魔的居所，如狂乱的祭司般前行，始终对真宗教发怒。因此他现在也武装了犹太人反对基督的信徒。他首先召聚一些人，询问他们：既然他们的律法规定他们有献祭的义务，为何他们不献祭？他们回答说他们的崇拜仅限于某一特定地点，于是这位天主的仇敌立刻下令重建被毁的圣殿，自负地以为他能够使主的预言失效，而实际上他正彰显了这预言的真实。犹太人听了他的话非常高兴，将他的命令传遍天下各地的同胞。他们从四面八方急忙赶来，同样贡献金钱和热忱进行这工程；皇帝也尽其所能提供一切，与其说是出于慷慨的骄傲，不如说是出于对真理的敌意。他还派遣了一位合适的官员作为总督去执行他亵渎的命令。据说他们制作了银制的镐、铲和篮子。当他们开始挖掘并运出土时，一大群人整天继续工作，但到了夜间，先前运走的土又自行从谷中返回。他们还摧毁了先前建筑的残余，意图全部重建；但当他们聚集了成千上万斗的白垩和石灰时，突然刮起狂风，暴雨、风暴和旋风将所有东西吹散到四面八方。他们仍然在疯狂中继续，并未被天主的宽容所唤醒。然后发生了大地震，足以使完全不认识天主作为的人心惊胆战；当他们仍未敬畏时，火从挖掘的地基中喷出，烧死了大部分挖掘者，并驱散了其余的人。此外，当一大群人在邻近的建筑物中夜间睡觉时，那房子连同屋顶突然倒塌，压死了所有人。那一夜及次夜，救恩的十字架记号在天上明亮发光，犹太人的衣服上也充满十字架，不是明亮的，而是黑色的。当天主的仇敌看见这些事时，因着从天降下的灾祸而恐惧，便逃散，各自回家，并承认了他们祖先钉死在十字架上的那位的天主性。尤利安听说了这些事，因为每一个人都在传说。但他像法郎一样硬了心肠。\par

% structure: p0057 source=h2 target=subsection reason=relative-heading-rank; base=h2

\subsection{第十六章　论远征波斯}

波斯人一听说君士坦提乌斯去世，便鼓起勇气，宣战并越过罗马帝国的边界。因此尤利安决定集结他的军队，虽然这是一支没有天主护佑的队伍。他先派人到德尔斐、提洛、多多纳及其他神谕所，询问占卜者他是否该出征。他们命他出征并应许胜利。我附上其中一个神谕以证明它们的虚假。其内容如下：\Quote{现在，我们众神都开始在那野兽河畔夺取胜利的战利品，而在他们中我——阿瑞斯，勇敢的战争喧嚣之发起者——将为首。}让那些称皮提亚为智慧之神、缪斯之主的人嘲笑这话语的荒谬吧。我既已发现它的虚假，倒要怜悯那受它欺骗的人。这神谕称底格里斯河为\InnerQuote{野兽}，因为河和动物同名。它发源于亚美尼亚的群山，流经亚述，注入波斯湾。被这些神谕所迷惑，那不幸的人沉迷于得胜的梦想，在与波斯人作战后，又幻想讨伐加里肋亚人——他如此称呼基督徒，以为这样可以辱没他们。然而，他虽有学识，却应该想到，更改名称并不会损害声誉，因为即使苏格拉底被称为克里底亚，毕达哥拉斯被称为法拉里斯，他们也不会因改名而蒙羞——尼柔斯即使被称为忒尔西忒斯，也不会失去自然赋予他的俊美。尤利安知道这些事，却全不放在心上，以为他可以用不恰当的名号来伤害我们。他相信了神谕的谎言，并威胁要在我们的教堂里设立淫欲女神的雕像。\par

% structure: p0059 source=h2 target=subsection reason=relative-heading-rank; base=h2

\subsection{第十七章　论贝洛亚地方元老的直言不讳}

在发出这些威胁之后，他被一个贝洛亚人压了下去。这个人因其在长官中居于首位而闻名，但又因他的热忱而更加闻名。当他看到自己的儿子陷入流行的异教时，将他赶出家门并公开与他断绝关系。那青年来到帝国近旁的城市，向皇帝陈明了自己的意见和父亲的判决。皇帝让他放心，并答应使他与父亲和好。当他到达贝洛亚时，他邀请官员和身居高位的人赴宴。那青年请求者的父亲也在其中，父子二人被命坐在皇帝的卧榻上。宴会中间，尤利安对父亲说：\Quote{在我看来，强迫一个心向别处、不愿改变忠诚的人是不对的。你的儿子不愿追随你的教义。不要强迫他。即使是我，虽然很容易强迫你，也不试图强迫你追随我的教义。}于是，父亲被对神圣真理的信德所激动，使辩论更加尖锐，喊道：\Quote{大人，}他说，\Quote{你说的是这个被天主憎恨、宁要谎言不要真理的可怜虫吗？}\par

尤利安再次戴上温和的面具说：\Quote{别骂了，伙计。}然后，转向那青年，\Quote{我，}他说，\Quote{会照顾你，因为我没能说服你父亲这样做。}我提起这件事，特意指出不仅这位可敬的人的勇敢值得赞扬，而且非常多的人藐视尤利安的统治。\par

% structure: p0062 source=h2 target=subsection reason=relative-heading-rank; base=h2

\subsection{第十八章　论家教的预言}

另一个例子是在安提阿一位优秀的人，负责管教少年们，他比通常的家教受过更好的教育，并且是当时首席教师、著名的诡辩家利巴尼乌斯的密友。\par

当时利巴尼乌斯是异教徒，期待胜利，并牢记尤利安的威胁。有一天，他嘲笑我们的信仰，对家教说：\Quote{木匠的儿子在做什么呢？}家教充满神圣的恩宠，预言了即将发生的事。\Quote{诡辩家，}他说，\Quote{您嘲笑的木匠之子，万物的创造者，正在做棺材。}\par

几天后，那可恶之人的死讯传来。他被抬出来，躺在棺材里。他威胁的夸耀被证明是虚妄的，天主受到了光荣。\par

% structure: p0066 source=h2 target=subsection reason=relative-heading-rank; base=h2

\subsection{第十九章　论圣尤利安隐修士的预言}

一个在肉身中效法无形体者生活的人，即绰号叙利亚语\Quote{萨巴斯}的尤利安，我在\WorkTitle{宗教史}中写了他的生平。当他听说不虔敬的暴君的威胁时，更加热心地向万有的天主献上祈祷。就在尤利安被杀的那一天，他在祈祷时听说了这件事，尽管隐修院距离军队有二十多站之远。据说，当他大声呼求主、恳求仁慈的主时，他突然止住眼泪，进入狂喜的喜乐，面容发光，以此表明他灵魂中的喜乐。他的朋友们看到这一变化，恳求他告诉他们喜乐的原因。\Quote{野猪，}他说，\Quote{主葡萄园的敌人，已因他对主的恶行而受到惩罚；他死了。他的祸害结束了。}众人一听这话，就欢喜跳跃，向天主唱起感恩之歌。从带来皇帝死讯的人那里，他们得知正是在那可怕的人被杀的那天和时辰，这位年迈的圣人知道了并宣布了它。\par

% structure: p0068 source=h2 target=subsection reason=relative-heading-rank; base=h2

\subsection{第二十章　论尤利安皇帝在波斯的死亡}

尤利安的愚蠢更因他的死而显明。他渡过分隔罗马帝国与波斯的河流，将军队带过，然后立即烧毁了船只，使他的士兵不是出于自愿而是被迫服从。最好的将军惯于激励部队，如果看到他们气馁，就鼓励他们、提升他们的希望；但尤利安烧毁了退却的桥梁，断绝了一切美好的希望。他无能的进一步证明是他未能履行四处征集粮秣、为部队提供给养的职责。尤利安既没有下令从罗马运来补给，也没有通过掠夺敌国领土来充足供应。他把有人居住的世界抛在身后，坚持穿越荒野行军。他的士兵没有足够的饮食，没有向导，在荒芜之地迷失行军。因此，他们看到了他们最英明的皇帝的愚蠢。在他们的怨言和牢骚中，他们突然发现那疯狂反抗造物主的他受伤致死。制造战争喧嚣的阿瑞斯从未像他承诺的那样来帮助他；洛西亚斯给了虚假的占卜；那位喜悦雷电的神并没有向致命一击的人掷下闪电；他威胁的夸耀被打翻在地。那施行正义一击的人的名字至今无人知晓。有人说他被一个无形者所伤，有人说是被称为以实玛利人的游牧者之一，有人说是因荒野饥荒无法忍受的骑兵。但无论是人还是天使动用了刀剑，无疑行事者是执行天主旨意的使者。据说，当尤利安受伤后，他用手捧满血，向空中抛洒并喊道：\Quote{你赢了，加利利人。}这样，他同时发出了胜利的承认和亵渎。他是如此疯狂。\par

% structure: p0070 source=h2 target=subsection reason=relative-heading-rank; base=h2

\subsection{第二十一章　论他在哈兰的巫术死后被揭露。他被杀后，他巫术的幻术被揭露。因为哈兰是一座仍然保留着他虚假宗教遗物的城市。}

尤利安从厄得撒向左转，因为那里以真宗教的恩宠为装饰；而他徒然愚昧地行经卡尔雷时，来到那受不敬神之人尊敬的庙宇，与他的同伴行过某些污秽之礼后，锁上并封了门，还安置守卫，命令他们在他返回之前不得让人进入。当他的死讯传来，不虔的统治被虔诚的统治接替，圣所被打开，在里面发现了这位已故皇帝刚毅、智慧与虔诚的证据。因为看见一个女人被她的头发高高吊起，双手伸开。那恶棍剖开了她的肚子，我猜想他是从她的肝脏得知他对波斯人的胜利。\par

这就是在卡尔雷发现的可憎之事。\par

% structure: p0073 source=h2 target=subsection reason=relative-heading-rank; base=h2

\subsection{第二十二章  论在安提约基雅皇宫中发现的人头及那里的公开欢庆}

据说在安提约基雅，皇宫中发现许多装满人头的箱子，还有许多满是尸体的井。这便是邪恶神祇的教诲。\par

安提约基雅人听闻尤利安之死，便沉浸在欢庆与喜乐中；不仅在教堂和殉道者圣所中表现出狂喜，甚至在剧院里也宣告了十字架的胜利，并将尤利安的预言视为笑柄。在此我将记录安提约基雅人那值得称赞的言论，以便保存在后世子孙的记忆中，因为他们异口同声地喊道：\Quote{\Person{马克西穆斯}，你这愚人，你的神谕在哪里？因为天主和他的基督已经得胜了。} 这样说是因为当时有一个名叫\Person{马克西穆斯}的人，自诩为哲学家，实为行巫术者，并自夸能预言未来。但安提约基雅人，从光荣的同伴\Person{伯多禄}和\Person{保禄}那里接受了神圣的教导，对万有的主和救主满怀热忱，坚持诅咒尤利安直到最后。他们的情感被那对象完全知晓，于是他写了一本书反对他们，并称之为\WorkTitle{反胡须者}。\par

这因暴君之死而生的欢庆将结束我历史书的这一卷，因为在我看来，将尤利安不虔的统治与正义的统治连接起来是不合宜的。\par

\SourceRecord{Translated by Blomfield Jackson.；Nicene and Post-Nicene Fathers, Second Series；Vol. 3.；Edited by Philip Schaff and Henry Wace.；Buffalo, NY: Christian Literature Publishing Co.,；1892.；<http://www.newadvent.org/fathers/27023.htm>.}

% structure: p0001 source=h1 target=section reason=page-title

\section{教会史（第四卷）}

% structure: p0002 source=h2 target=subsection reason=relative-heading-rank; base=h2

\subsection{第一章　论\Person{若维安}的统治与孝爱}

儒利安被杀后，将军与督军们召开会议，商议应由何人继承皇位，以拯救作战中的军队，并恢复罗马的国运——此时因已故皇帝的鲁莽，按俗话说，已处于千钧一发的危险之中。但首领们正在商议时，军队集会要求若维安为皇帝，尽管他既非将军亦非次高官阶；但他却是一位异常杰出、因诸多缘故而广为人知的人物。他身材高大，灵魂崇高。在战争与严酷斗争中，他常是首先冲锋者。他勇敢地反对不敬，不畏暴君的权势，其热忱足以将他列于基督的殉道者之中。于是将军们接受了士兵的一致投票，视之为天主的拣选。这位勇敢者被领上前，被安置在一座匆忙搭起的高台上。军队以皇帝的尊号向他欢呼，称他为奥古斯都和凯撒。若维安以其一贯的直率，无畏于在场的指挥官们，也顾不得军队近日的背教行径，卓越地说：\Quote{我是基督徒。我不能统治这样的人。我不能指挥儒利安的军队，它已在邪恶的纪律中受训。这样的人，被剥去了天主眷顾的遮盖，将轻易而可笑地沦为敌人的猎物。}军队听了这话，齐声喊道：\Quote{皇帝啊，不要犹豫；不要以为统率我们是可鄙之事。你将要统治的是在真理的训练中养育的基督徒；我们的老兵曾在君士坦丁本人亲自训导，我们中的年轻人则由君士坦提乌斯教导。这位已故者的帝国仅存续了短短数年，这数年的时间太短，不足以在那些被他迷惑的人身上留下烙印。}\par

% structure: p0004 source=h2 target=subsection reason=relative-heading-rank; base=h2

\subsection{第二章　论\Person{亚大纳削}的返回}

皇帝因这些话而\Emph{欣喜}，便着手为将来筹谋国家的安全，以及如何将军队从征战中无损地带回。他无需过多商议，便立刻收获了由真宗教的种子所结出的果实，因为万有的天主明证了祂自己的眷顾，使一切困难消失。波斯王一得知\Person{若维安}即位，便立即派遣使节商议和平；不仅如此，他还提供军队的补给，并指示在沙漠中为他们设立市场。休战协定为期三十年，军队安全地从战争中返回家园。皇帝踏上自己领土后的第一道法令，是召回被流放的主教，并宣布将教会归还给那些持守尼西亚信经不渝的会众。他还致函著名的这些教义的捍卫者\Person{亚大纳削}，恳求他写一封信，其中包含关于宗教事务的确切教导。\Person{亚大纳削}召集了最有学问的主教与他相会，并回信劝勉皇帝持守尼西亚所传授的信仰，因为这信仰与宗徒的训导一致。我急切地希望造福所有可能遇见这封信的人，故在此将其附上。\par

% structure: p0006 source=h2 target=subsection reason=relative-heading-rank; base=h2

\subsection{第三章　致皇帝\Person{若维安}有关信仰的主教会议信函}

致最虔诚、最仁慈、胜利的\Person{约维安}奥古斯都：\Person{亚大纳削}及集会的其余主教，以埃及至\Place{特拜德}和\Place{利比亚}全体主教的名义。明智地偏好和追求神圣事物，适合一位蒙天主眷爱的君王。愿你这样以真理持守你的心于天主手中，并在平安中统治多年。因陛下近来表示愿从我等学习天主教的信仰，我们感谢了主，并决定首先提醒陛下关于教父们在\Place{尼西亚}所宣认的信仰。有些人蔑视这信仰，并策划了针对我们的许多不同攻击，因我们拒绝服从亚略异端。他们成了天主教会中异端和分裂的创始者。在我们主\Person{耶稣基督}内的真实而虔诚的信仰，已如从圣经所认识和宣读的那样，向所有人显明了。藉此信仰，殉道的圣人们得以成全，如今已离世与主同在。这信仰本应在各处屹立不倒，若非某些异端者的邪恶竟敢尝试篡改它；因为\Person{亚略}及其党羽企图败坏它，并引入不敬虔以毁灭它，声称天主子是出自无有、受造物、被造之物，且可改变。藉此手段，他们欺骗了许多人，以致连那些似乎颇有身份的人也被他们引入歧途。然后，我们圣洁的教父们主动采取行动，如我们所说，在\Place{尼西亚}集会，弃绝了亚略异端，并签署了天主教的信仰，以便通过向全世界宣告真理来扑灭异端的火焰。因此，这信仰在整个教会中被知晓并宣讲。但因有些人愿重新挑起亚略异端，竟有僭妄置尼西亚教父所宣认的信仰于不顾，而另一些人则假装接受它，实则否认它，歪曲\OriginalTerm{同体}{\GreekText{ὁμοούσιον}}的意义，从而亵渎圣神，声称圣神是受造物且是通过圣子的手段而被造的存在者，我们，迫不得已看到这类亵渎对人民造成的伤害，急忙向陛下呈上尼西亚所宣认的信仰，使陛下知道它是如何精确地制定，以及那些教导与之相悖的人的错误何其大。最圣洁的奥古斯都，请知悉，这信仰是从永远所宣讲的信仰，是尼西亚集会的教父们所宣认的信仰。全世界所有教会都同意此信仰，在\Place{西班牙}、\Place{不列颠}、\Place{高卢}、全\Place{意大利}和\Place{坎帕尼亚}、\Place{达尔马提亚}和\Place{密细亚}、\Place{马其顿}、全\Place{希腊}、全\Place{非洲}、\Place{撒丁}、\Place{塞浦路斯}、\Place{克里特}、\Place{旁非利亚}和\Place{伊扫利亚}、\Place{吕西亚}、全\Place{埃及}和\Place{利比亚}、\Place{本都}、\Place{卡帕多细亚}及邻近地区，以及\Place{东方}所有教会，除少数接受亚略主义之外。关于以上所有教会，我们经过检验得知其意见。我们收有信函，并且我们知道，最虔诚的奥古斯都，尽管有少数人反驳这信仰，他们不能损害整个有人居住的世界的决定。\par

在长期受到亚略异端有害影响之后，他们更加激烈地抗拒真宗教。为了禀报陛下，虽然您早已熟悉，我们仍费心附上那三百一十八位主教在\Place{尼西亚}所宣认的信仰。其内容如下。\par

我们信唯一的天主，全能的圣父，造天地及一切有形无形者；又信唯一的主\Person{耶稣基督}，天主子，由圣父所生，即出自圣父的性体，天主出自天主，光明出自光明，真天主出自真天主：生而非受造，与圣父同体，万物（无论在天上或在地上）都藉着祂而受造。祂为了我们人类和我们的救恩，从天降下，降生成人。祂受难，第三日复活。祂升了天，将来要审判生者死者。我们也信圣神。圣而公教会、宗徒的教会弃绝那些说：天主子曾有未存在之时，在祂受生以前祂不存在，祂出自无有，或说祂是不同的本质或不同的本体，或是受造物，或可变化或改变。在此信仰中，最虔诚的奥古斯都，所有人都必须遵守为神圣和宗徒的，也不可有人试图通过说服性的推理和言辞争辩来改变它，如同起初亚略的狂人们所做的那样，他们主张天主子出自无有，曾有未存在之时，祂是受造、被造且可变化的。因此，如我们所述，\Place{尼西亚}会议弃绝了此异端，并宣认了真理的信仰。因为他们并没有简单地说子相似父，以免祂被认为仅仅是相似神，而是真天主出自天主。他们颁布了\Quote{\OriginalTerm{同体}{Homoüsion}}这一术语，因为它专属于一个真正真实父亲所生的真正真实儿子。然而他们并未将圣神与圣父和圣子分开，而是与圣父、圣子一同，在唯一的天主圣三信仰中光荣圣神，因为天主圣三的神性是一。\par

% structure: p0010 source=h2 target=subsection reason=relative-heading-rank; base=h2

\subsection{第四章 恢复给教会的津贴；以及皇帝的驾崩。}

当皇帝收到此信，他先前对神圣事物的认识和倾向得到了确认，于是颁布第二道敕令，命令归还伟大的\Person{君士坦丁}所拨给教会的谷物数量。因为\Person{儒利安}，正如一个与我们主救主开战的人所可预期的，甚至停止了这一给养；并且由于儒利安的邪恶所导致的降临帝国的饥荒，使得无法征收君士坦丁法令规定的贡献，\Person{约维安}命令目前先供应三分之一，并应许饥荒停止后他将全额供应。\par

在以此类谕令标志其统治开端之后，约维安从安提约基雅启程前往博斯普鲁斯；但在比提尼亚与加拉提亚交界处的村庄达达斯塔奈，他去世了。他以最宏伟、最美好的支持和支柱踏上了从此世离去的旅程，但所有经历过他统治之宽仁的人都留在悲痛中。我以为，至高统治者正是为了定我们的不义之罪，既向我们显示美善之事，又将它们收回；藉前者，祂教导我们祂能轻易赐予祂所愿的一切；藉后者，祂定我们为不配拥有它们之罪，并指引我们转向更美好的生命。\par

% structure: p0013 source=h2 target=subsection reason=relative-heading-rank; base=h2

\subsection{第五章 论瓦伦提尼安努斯在位，及其如何立其弟瓦伦斯为共治者}

当军队得知皇帝突然驾崩，他们如丧父般哀悼已故的君王，并立瓦伦提尼安为帝。他正是那位击打神庙军官并被遣送至城堡的人。他不仅以勇气著称，亦以审慎、节制、正义和伟岸的身材闻名。他的性格如此王者风范且宽宏大量，以至于当军队企图为他任命一位同僚共享皇位时，他说出了那句广为流传的著名话语：\Quote{在我成为皇帝之前，士兵们，是你们将帝国的缰绳交给我；如今我已执掌，则是我而非你们为我国家操心。} 全军对他的话惊叹不已，并欣然顺从了其权威的指引。然而，瓦伦提尼安从潘诺尼亚召来他的弟弟，与他分享帝国。但愿他从未这样做过！瓦伦斯当时尚未接受谬误的教义，他被托付治理亚洲和埃及，而瓦伦提尼安则将欧洲归于自己。他前往西部诸省，以宣布真宗教为开端，训导他们于一切公义。当亚略派的米兰主教奥克森提乌斯（他在几次会议上被谴责）去世后，皇帝召集主教们并如此致辞：\Quote{你们既然受圣经养育，深知应如何为被赋予主教职者，以及他应如何不仅以言语，更以生命正确统治其子民；展现自身为各样美德的典范，并使自己的言行为其教导的见证。如今请将这样一个人置于你们的总主教宝座上，使我们这些治理王国的人能真诚地在他面前低头，并欢迎他的责备——因为作为人，我们难免有时失足——如同医生的治疗。}\par

% structure: p0015 source=h2 target=subsection reason=relative-heading-rank; base=h2

\subsection{第六章 论盎博罗削，米兰主教的选举}

皇帝如此说完，议会便恳请他这位睿智虔诚的君王做出选择。他答道：\Quote{这责任对我们来说过于重大。你们这些被神圣恩宠所尊荣、从上方接受了光照的人，会做出更好的选择。} 于是他们离开皇帝御前，开始私下商议。与此同时，米兰民众陷入分裂，一些人渴望立此人，另一些人则渴望立彼人。那些曾受奥克森提乌斯谬误感染的人，倾向于选择观点相同者；而正统派一方则同样急于拥立与他们观念一致的主教。当执掌该区首席文官职位的盎博罗削得知这场争执时，担心发生暴乱，便匆忙赶往教堂；争斗顿时平息。民众齐声呼喊：\Quote{立盎博罗削为我们的牧者！}——尽管此时他尚未领洗。消息传到皇帝耳中，他立即下令为这位可敬之人施以圣洗圣事并授以圣职，因为他知道他的判断如木匠的准绳般正直，他的裁决比天平的横梁更精确。他更从持相反意见者达成的一致中得出结论：这选择是出于天主的旨意。于是盎博罗削接受了圣洗圣事及总主教职务的恩宠。最卓越的皇帝亲临现场，据说向他的主和救主献上了如下赞美诗：\Quote{全能的主和救主，我们感谢祢；我将人的身体托付给此人照管；祢将他们的灵魂托付给他，并已显示我的选择是公义的。}\par

不多日之后，神圣的盎博罗削以最大的自由向皇帝进言，并指出官员的某些行事不当。瓦伦提尼安说，这种自由对他而言并不新鲜，并且他对此十分熟悉，不仅没有反对，反而欣然同意任命他为主教。皇帝继续说：\Quote{继续，按天主的法律所吩咐的，医治我们灵魂的过失。}\par

这就是瓦伦提尼安在米兰的言行。\par

% structure: p0019 source=h2 target=subsection reason=relative-heading-rank; base=h2

\subsection{第七章 瓦伦提尼安努斯与瓦伦斯两位皇帝致亚洲教区关于\OriginalTerm{同质}{Homoüsion}的信函，闻悉在亚洲与弗里几亚有人争论天主旨意}

瓦伦提尼安下令在伊利里亚姆召开一次会议，并将与会主教所批准的谕令送至争论者处。他们决定坚守尼西亚所颁布的信经，皇帝本人也与其弟联名致函，敦促他们遵守这些谕令。\par

这道谕令清楚宣告了皇帝的虔诚，同时也展现了当时瓦伦斯在神圣教义上的纯正。因此我将全文列出。\Quote{伟大的皇帝们，永享尊荣，尊荣常胜，瓦伦提尼安努斯、瓦伦斯和格拉提安，致亚洲、弗里几亚、卡洛弗里几亚·帕卡提亚那的主教们，在主内问安。}\par

在\Place{伊利里库姆}召开了一次大公会议，对救恩的圣言进行了大量讨论之后，三倍真福的主教们宣告：圣父、圣子、圣神的天主圣三具有同一实体。他们钦崇这天主圣三，毫不松懈地履行他们应尽的本分，即对伟大君王的敬礼。朕的旨意是宣讲这天主圣三，使无人能说：\Quote{我们接受统治此世的君王的宗教，却不顾及那将救恩讯息赐予我们者}，因为我们的天主福音含有这判断说：\BibleQuote{我们应把凯撒的归给凯撒，把天主的归给天主。}\par

你们这些主教，你们这些救恩的圣言的捍卫者，你们说什么？如果这是你们的宣言，那么就这样继续彼此相爱，停止辱骂皇权吧。不要再迫害那些勤勉事奉天主的人，因着他们的祈祷，地上的战争得以止息，背教天使的袭击也被击退。这些人藉着祈祷努力击退一切有害的魔鬼，他们知道如何按法律纳税，不抗拒君王的权力，而是以纯洁的心灵遵守天上君王的诫命，并服从我们的法律。但你们已被证明是不顺从的。我们尝试了一切方法，但你们已经放弃了。然而我们愿与你们无干，如同\Person{比拉多}在基督居住在我们中间时受审，他不愿杀死祂，当他们恳求祂死时，他转向东方，讨水洗手，并洗手说：\BibleQuote{我对这义人的血是无罪的。}\par

因此，朕的陛下始终下令：在基督田地里工作的人不得受迫害、压迫或虐待；伟大君王的管家也不得被放逐；免得今天在我们的君主之下你们似乎兴旺富足，然后与你们的恶谋士一同践踏他的盟约，如同在\Person{匝加利亚}的血的事件中那样，但他和他的同党被我们的天上君王\Person{耶稣基督}在祂来临时消灭了，被交付于死亡的审判，他们与助纣为虐的致命恶魔一同被灭。朕已将这些命令口头传达给\Person{阿梅吉丢斯}、\Person{切罗尼乌斯}、\Person{达玛苏}、\Person{拉姆蓬}和\Person{布伦提西乌斯}，并将实际法令也发给你们，使你们知道在可敬的主教会议上所制定的内容。\par

此信后附上主教会议的谕令，其内容简述如下。\par

我们遵照伟大而正统的主教会议，明认圣子与圣父具有同一实体。我们并不像某些人先前那样理解\Quote{同一实体}一词，他们以虚假的附和签署了自己的名字；也不像现今某些人那样，称旧信经的起草者为教父，却使这词的含义失效，他们追随那些声称\Quote{同一实体}意为\Quote{相似}的陈述的作者，其理解是：既然圣子与祂所造的任何受造物都不可比拟，祂唯独与圣父相似。因为那些这样不敬地思考的人将圣子定义为\Quote{圣父的特殊受造物}，但我们，与现今在\Place{罗马}和\Place{高卢}的主教会议一道，坚持圣父、圣子、圣神具有同一个实体，存在于三个位格，即三个完美的本质中。我们根据\WorkTitle{尼西亚信经}的阐述，明认同一实体的天主圣子由\Person{童贞玛利亚}取得肉身，并在人间支搭帐幕，为我们完成了全部救恩计划，包括诞生、受难、复活和升天；祂将在审判之日再次来临，按各人的生活方式报答我们，带着肉身显现，彰显祂的神性权能，祂是带着肉身的天主，而不是带着神性的人。\par

我们谴责那些有不同想法的人，也谴责那些不诚实地谴责以下说法的人，即\Quote{圣子在受生之前并不存在}，而是写道甚至在实际受生之前，祂就潜在地在父内。因为这对所有受造物来说是真实的，它们并不像圣子那样永远与父同在，圣子是由永恒受生而生的。\par

以上是皇帝诏书的简短摘要。我现在附上主教会议的实际公函。\par

% structure: p0029 source=h2 target=subsection reason=relative-heading-rank; base=h2

\subsection{第八章　\Place{伊利里库姆}主教会议关于信仰的公函}

\Place{伊利里库姆}的主教们致\Place{亚细亚}、\Place{夫黎基雅}及\Place{卡洛弗里基亚帕卡提亚纳}教区的天主的教会和主教们，在主内问安。\par

聚会并进行长期查考后，我们关于救恩的圣言宣告：圣父、圣子、圣神的天主圣三具有同一实体。我们认为写信给你们是合适的，并非要以虚妄的争论来讨论关乎敬拜天主圣三的事，而是以谦卑之心，堪当此职。\par

此信由我们亲爱的弟兄、同劳者\Person{埃尔皮迪乌斯}司铎送达。因为并非写在我们的书信中，而是在我们的救主\Person{耶稣基督}的书上写着：\BibleQuote{我是属\Person{保禄}的，我是属\Person{阿颇罗}的，我是属\Person{刻法}的，我是属\Person{基督}的。\Person{保禄}为你们被钉死在十字架上吗？或者你们领受圣洗归于\Person{保禄}的名下吗？}\par

确实，我们的谦卑起初似乎不适合给你们写任何信，因为你们的宣讲给你们管辖下的整个区域带来了极大的恐惧，你们竟将圣神与圣父和圣子分离。因此我们被迫派遣我们的主内同劳者\Person{埃尔皮迪乌斯}到你们那里，以查实你们的宣讲是否确实如此，并携带\Place{罗马}帝国政府的这道公文。\par

凡不将天主圣三视为同一实体者，应受绝罚；若有人被发现与他们共融，也应受绝罚。\par

但对于宣讲天主圣三具有同一实体的人，天国已经预备好了。\par

因此，弟兄们，我们劝勉你们：不可传授别的道理，也不可持守任何其他虚妄的信仰，而要始终、处处宣讲天主圣三具有同一实体，使你们能继承天国。\par

在撰写这一点时，我们也受到提醒，要给你们写这封信，关于我们的同工现在或将来被任命为主教的事：如果在已经担任过公职的主教中有正直的人，就从他们中选；如果没有，就从司铎阶层中选；同样，也要从实际的司铎阶层中任命司铎和执事，使他们在各方面无可指摘，而不是从元老院和军队中选。\par

我们本不愿给你们写一封长信，因为一个全权代表，我们的主内同工\Person{埃尔皮迪乌斯}，被派遣去仔细查问你们的宣讲，是否真的如同我们从我们的主内同工\Person{尤斯塔修斯}那里所听到的那样。\par

最后，如果你们曾经误入歧途，就脱去旧人，穿上新人。同一位弟兄和同工\Person{埃尔皮迪乌斯}将教导你们如何宣讲真信德：与天主圣父同一本体的天主圣三，同圣子及圣神一起，应受钦崇、荣耀和显扬，圣父在圣子内，圣子在圣父内，连同圣神，直到永远。因为既然这事已显明，我们就能够按照以前在\Place{尼西亚}由教父们所确立的信德，明确宣认天主圣三是同一本体。只要宣讲这信德，我们就能避开通灭的魔鬼的陷阱。当魔鬼被消灭时，我们就能在和平中彼此以和平的书信致敬。\par

因此，我们写信给你们，是为使你们知道那些\OriginalTerm{亚略狂人}{Ariomaniacs}的罢黜，他们不承认圣子是出自圣父的本体，也不承认圣神。我们附上他们的名字：\Person{波利克罗尼乌斯}、\Person{泰莱马库斯}、\Person{福斯图斯}、\Person{阿斯克勒皮亚德斯}、\Person{阿曼提乌斯}、\Person{克娄帕特}。\par

\Quote{我们如此写作，是为光荣圣父、圣子及圣神，直到永远，阿们。我们恳求圣父及圣子我们的救主耶稣基督，同圣神，使你们多年安康。}\par

% structure: p0042 source=h2 target=subsection reason=relative-heading-rank; base=h2

\subsection{第九章　论\OriginalTerm{欧迪乌斯派}{Audiani}的异端}

英明的皇帝如此留意宗徒的法令，但\Person{欧迪乌斯}，一个在种族和语言上都是叙利亚的人，那时出现作为新法令的发明者。他早已开始孵化不义，如今暴露了他的真面目。起初，他以荒谬的方式理解经文\Quote{让我们按我们的肖像，按我们的模样造人。}由于未能理解圣经的含义，他认为天主具有人的形状，并推测天主被身体的部分所包裹；因为圣经常用人体部位的名称来描述天主的行为，因为通过这些方式，天主的眷顾更容易被无法感知非物质理念的心智所理解。除了这种不敬，欧迪乌斯还添加了其他类似的谬误。通过折衷的过程，他采纳了\Person{摩尼}的一些教义，并否认宇宙的天主是火或黑暗的创造者。但这些及所有类似的错误都被他党派的支持者所掩盖。\par

他们声称自己与教会的集会分离。但由于他们中有些人索取可诅咒的高利贷，有些人非法与女人同居而无婚姻约束，而那些清白之人则与有罪者自由交往，他们以独自生活为由掩盖他们教义的亵渎。然而，这个借口是无耻的，是法利塞人教导的自然结果，因为法利塞人在质问圣宗徒时，控告灵魂与身体的医师说：\Quote{你的老师为什么同税吏和罪人一起吃饭？}透过先知，天主论及这样的人说：\Quote{他们说：\InnerQuote{你不可靠近我，因为我是圣洁的。}这在我眼中是烟雾中的火。}但现在不是驳斥他们无理错误的时候。因此，我继续叙述余下部分。\par

% structure: p0045 source=h2 target=subsection reason=relative-heading-rank; base=h2

\subsection{第十章　论\OriginalTerm{美撒里派}{Messaliani}的异端}

此时也兴起了美撒里派的异端。将其名译为希腊语的人称他们为\OriginalTerm{欧基特派}{Euchitæ}。\par

他们还有另一个自然产生于他们行为方式的称呼。由于他们受到某个魔鬼的影响，他们以为那是圣神的降临，因此被称为狂热者。\par

完全感染了这种瘟疫的人视手工劳动为不义而逃避；他们沉溺于怠惰，称他们梦中的想象为预言。这异端的领袖是\Person{达多埃斯}、\Person{萨巴斯}、\Person{阿德尔菲乌斯}、\Person{赫尔马斯}和\Person{西梅奥内斯}，以及其他一些人，他们不远离教会的共融，声称基督我们的师傅所说的神圣食物\BibleQuote{谁吃我的肉，并喝我的血，必得永生}既无益也无害。\par

为了掩盖他们的不健全，他们甚至在定罪后还无耻地否认，并弃绝那些意见与他们内心想法一致的人。\par

在这种情况下，\Person{莱托伊乌斯}，\Place{梅利提尼}教会的首领，一个充满神圣热忱的人，看到许多隐修院，或者我宁愿说是强盗窝，深受此病毒害。因此他焚烧了它们，并从羊群中赶走了豺狼。\par

同样，杰出的\Person{安斐洛基乌斯}，受托管理吕考尼亚人的都会，治理全体人民，一得知这场瘟疫侵入他的教区，就使它离开他的边界，并使他牧养的羊群免受感染。\par

\Person{弗拉维安努斯}，那位名闻遐迩的\Place{安提约基雅}人大司祭，得知这些人住在\Place{厄德萨}，并以他们特有的毒害攻击所有与他们接触的人，便派遣了一队隐修士，将他们带到\Place{安提约基雅}，并以下列方式在他们否认自己的异端时定了他们的罪。他说，他们的控告者是在诽谤他们，证人在作伪证；而\Person{阿德尔斐乌斯}，一位年纪很大的人，他以和善的言辞招呼他，并吩咐他坐在自己身边。然后他说：\Quote{我们，可敬的阁下，年事已高，对人性和那敌对我们的魔鬼的诡计有更准确的认识，并凭经验学到了恩宠恩赐的特质。但这些年轻人对这些事没有清楚的认识，不能忍受聆听灵性的教导。因此请告诉我，你是如何理解那敌对之神退去，圣神的恩宠降临的？}老人被这些话打动，倾吐了他所有隐秘的毒液，因为他说，对于领受圣洗圣事的人没有任何益处，只有通过恳切的祈祷才能驱走内住的魔鬼，因为每一个生于世界的人，从其第一祖先那里，就像继承本性一样继承了魔鬼的奴役；但当这些魔鬼被驱逐后，圣神便降临，以可感知和可见的记号显示祂的临在，立刻使身体脱离情欲的冲动，完全除去灵魂向恶的倾向；其结果是，不再需要抑制身体的斋戒，也不再需要约束它并教它如何正当行走的教导或训练。不仅如此，领受此恩赐的人不仅从身体的放纵运动中得解放，还能清楚地预见未来，并以肉眼看见天主圣三。\par

这样，神圣的\Person{弗拉维安努斯}挖到了那污浊的源头，并成功暴露了它的支流。然后他这样对那可怜的老人说：\Quote{你这在邪恶日子里变老的人啊，你自己的口定了你的罪，不是我，你的嘴唇作证反对你。}他们的错谬被这样揭露后，他们被逐出\Place{叙利亚}，退到\Place{旁非利亚}，并用他们那有害的教义充满了那里。\par

% structure: p0054 source=h2 target=subsection reason=relative-heading-rank; base=h2

\subsection{第十一章 \Person{瓦伦斯}如何落入异端}

现在我要继续我的叙述，并描述那场为教会激起许多巨浪的风暴的开端。\Person{瓦伦斯}最初接受皇权时，以忠于宗徒教义而著称。但当哥特人跨越\Place{多瑙河}并蹂躏\Place{色雷斯}时，他决定集结军队向它们进军；因此他决心不穿上神圣恩宠的外衣就不上阵，而要先用圣洗圣事的全副武装保护自己。形成这个决定时，他的行为既好又明智，但他随后的行为暴露出非常软弱的性格，导致了放弃真理。他的命运与我们的第一祖先\Person{亚当}相同；因为他也被妻子的论点说服，失去了自由的地位，不仅成了俘虏，而且成了女人花言巧语的顺从听众。他的妻子早已陷入\Person{亚略}的陷阱，现在她诱捕了丈夫，并说服他同她一起跌入亵渎之坑。他们的领袖和启蒙者是\Person{尤多克修斯}，他仍然掌控着\Place{君士坦丁堡}的船舵，结果船没有被向前驾驶，而是沉入了海底。\par

% structure: p0056 source=h2 target=subsection reason=relative-heading-rank; base=h2

\subsection{第十二章 \Person{瓦伦斯}如何流放德行高尚的主教们}

就在\Person{瓦伦斯}受洗的时候，\Person{尤多克修斯}用誓言束缚了这个不幸的人，要他在自己教义的不敬神中坚持，并将持有相反意见的人逐出每个主教座。因此\Person{瓦伦斯}放弃了宗徒教导，投向了反对派；不久他就履行了誓言的其他部分；因为他从\Place{安提约基雅}驱逐了伟大的\Person{梅里提乌斯}，从\Place{撒摩撒塔}驱逐了神圣的\Person{优西比乌}，并剥夺了\Place{劳迪西亚}她可敬的牧者\Person{佩拉吉乌斯}。\Person{佩拉吉乌斯}在很年轻时就承担了婚姻的轭，在婚礼当日，在他新婚的第一天，他说服他的新娘选择贞洁而非夫妻交合，并教导她接受兄弟情谊取代婚姻结合。因此他将所有荣誉归于节制，并且自身也拥有与她和睦相处的姊妹美德，基于这些原因，他被一致选为主教。然而，就连他生命与言行的光辉也未能震慑真理的敌人。\Person{瓦伦斯}也将他放逐到\Place{阿拉伯}，神圣的\Person{梅里提乌斯}到\Place{亚美尼亚}，而\Person{优西比乌}，那位在宗徒工作中不知疲倦的劳动者，到\Place{色雷斯}。他确实不知疲倦，因为当他得知许多教会现在失去了牧者时，他游走于\Place{叙利亚}、\Place{腓尼基}和\Place{巴勒斯坦}，穿着战袍，头戴冠冕，祝圣司铎和执事，并填补教会的其他秩次；如果偶然遇到与他有相同情怀的主教，他就任命他们去空的教会。\par

% structure: p0058 source=h2 target=subsection reason=relative-heading-rank; base=h2

\subsection{第十三章 论\Place{撒摩撒塔}主教\Person{优西比乌}及其他}

关于\Person{优西比乌}在收到命令他前往\Place{色雷斯}的皇帝敕令后所表现的勇气和谨慎，我认为所有迄今为止不知道的人都应该听一听。\par

这道谕旨的送达者于傍晚抵达目的地，\Person{尤西比乌斯}劝他保持沉默，隐藏此行的原因。主教说：\Quote{因为民众在神圣的热忱中长大，若他们知道你为何而来，会把你淹死，而我将为你的死负责。}说完此话并照常举行了晚祷后，这位老人独自徒步出发，趁着夜色。他将自己的意图告诉了一个家仆，那仆人跟着他，只带着一个坐垫和一本书。当他到达河岸（因为\Place{幼发拉底河}紧贴着城墙流淌）时，他登上一艘船，吩咐船夫划向\Place{宙格玛}。天亮时，主教已抵达\Place{宙格玛}，而\Place{萨莫萨塔}城充满了哭泣和哀号，因为上述家仆将给他的命令报告给\Person{尤西比乌斯}的朋友们，告诉他们他愿与谁同行，以及要携带哪些书。于是全体会众哀哭他们牧者的离去，河上挤满了渡船。\par

他们来到他所在之处，看见自己心爱的牧者，便以哀哭和叹息泪流成河，试图说服他留下，不要将羊群弃于豺狼。但一切徒劳，他给他们读了宗徒的法律，这法律明确命令我们要服从长官和掌权者。他们听后，有人送他金子，有人送银子，有人送衣服，还有人送仆人，仿佛他正要启程前往某个陌生而遥远的异地。主教除了几位更亲密的朋友所赠的一些微薄礼物外，一无所取，然后对全体会众给予训导和祈祷，并劝勉他们勇敢地为宗徒的法令站立。\par

于是他出发前往\Place{多瑙河}，而他的朋友们返回本城，互相鼓励，等待豺狼的袭击。\par

我若认为，若他们的信德之热忱与真诚在我的历史中缺乏纪念便是亏待了他们，那么我现在要继续描述它。\par

亚略派在剥夺了羊群极其优秀的牧者之后，设立了另一位主教取代他；但是城中居民，无论是贫穷放牧还是富裕显赫，无论是仆人、工匠、农夫、园丁，无论男女老少，都没有像往常一样来到教会聚会。新主教独自一人生活；没有一个人看他一眼或与他说一句话。然而有报告说他表现得彬彬有礼、温和节制，以下事例可资证明。有一次他想洗澡，他的仆人便关上了浴池的门，不让想进来的人进入。当他看到门前的人群时，便命令打开门，并指示每个人都可以自由使用浴池。他在里面的厅堂中也表现出同样的行为；因为他看到某些人站在他身旁，而他正在洗澡，便请他们与他共用热水。他们沉默着。他认为他们的犹豫是出于对他的尊敬，便迅速起身走了出去，而这些人实际上认为连水都受到了他异端的污染，因此把水全部倒入下水道，另外为自己准备了新的水。获悉此事后，那位闯入者离开了城市，因为他认为继续住在一个厌恶他、将他视为公敌的城市是愚蠢和荒谬的。\Person{尤诺米乌斯}（他叫这个名字）离开\Place{萨莫萨塔}后，\Person{路济乌斯}，一只不折不扣的狼、羊群的敌人，被任命接替。但是羊群虽然完全没有牧人，却自己牧放了自己，坚持不懈地保持宗徒教义的全部纯洁。新闯入者如何被憎恶，以下叙述将阐明。\par

一些少年在广场上打球玩乐，正值\Person{路济乌斯}路过。球偶然掉落，滚到了驴的腿间。少年们大声喊叫，因为他们认为球被污染了。\Person{路济乌斯}看到这情景，便让随从停下来了解发生了什么事。少年们生了火，把球抛过火焰，以为这样就能净化它。我确实知道这不过是孩童的举动，是古代习俗的残余；但它足以证明该城对亚略派怀有多么深的仇恨。\par

然而\Person{路济乌斯}并不效法\Person{尤诺米乌斯}的温和，反而说服当局流放了许多其他神职，并将神圣教义的最杰出捍卫者遣送到罗马帝国最远的边境：执事\Person{埃沃尔奇乌斯}被送往\Place{奥阿西斯}的一个废弃村庄；\Person{安提约古}，他有幸与伟大的\Person{尤西比乌斯}有亲戚关系（他是其兄弟之子），并因自己的高尚品格而更显卓越，且具有司铎品级，被送往\Place{亚美尼亚}的偏远地区。这位\Person{安提约古}如何勇敢地为神圣法令争辩，将由以下事实表明。当神圣的\Person{尤西比乌斯}经过多次冲突（每一次都获得胜利）而殉道身亡后，照例举行了民众的会议，\Place{佩拉}的主教\Person{约维努斯}也来到，他曾在短时间内与亚略派共融。\Person{安提约古}被一致推选为叔父的继任者。当他被带到圣台前并奉命屈膝时，他转过身，看见\Person{约维努斯}正将右手按在他头上。他拨开那只手，命令他从祝圣者中离开，说他不能忍受一只曾接受以亵渎方式庆祝的奥迹的手。\par

这些事发生在稍后。在我现在所说的时候，他被迁往\Place{亚美尼亚}内陆。\par

神圣的\Person{尤西比乌斯}住在\Place{多瑙河}畔，哥特人正在那里蹂躏\Place{色雷斯}并围攻城市，正如他自己的作品中所描述。\par

% structure: p0069 source=h2 target=subsection reason=relative-heading-rank; base=h2

\subsection{第十四章 论圣\Person{巴尔塞斯}，以及\Place{厄德撒}主教及其同伴的流放}

\Person{巴耳塞斯}的名声如今不仅在他自己的城市\Place{厄德撒}和邻近城镇中伟大，而且在\Place{腓尼基}、\Place{埃及}和\Place{忒拜}，他穿越所有这些地区都享有因其伟大德行而赢得的崇高声誉。\Person{瓦伦斯}将他放逐到\Place{阿剌杜斯}岛，但当皇帝得知无数人群涌向那里，因为\Person{巴耳塞斯}充满宗徒的恩宠，并以言语驱除疾病时，便将他送到\Place{埃及}的\Place{奥克西林库斯}；但那里也因他的名声吸引了所有人，这位堪配天堂的老人被带到该地区蛮族附近的一个偏远城堡，名为\Place{斐诺}。据说在\Place{阿剌杜斯}，他的床铺保存至今，在那里受到极大的尊敬，因为许多病人躺在上面，并凭他们的信德获得痊愈。\par

% structure: p0071 source=h2 target=subsection reason=relative-heading-rank; base=h2

\subsection{第十五章 论发生在\Place{厄德撒}的迫害，以及\Place{厄德撒}的司铎\Person{欧洛吉乌}和\Person{普罗托革内斯}}

如今\Person{瓦伦斯}第二次将羊群剥夺了他们的牧人后，将一只狼安置在他们之上代替他。当这位皇帝到达\Place{厄德撒}时，全体民众已经弃城，聚集在城前。他命令名叫\Person{莫德斯图}的总督，召集他手下负责征收贡赋的部队，带领所有在场的武装力量，用棍棒和棍子击打，必要时使用其他战争武器，驱散聚集的群众。清晨，当总督执行此命令时，在穿过广场的路上，他看到一个妇女怀抱婴儿，急速奔跑。她轻视了军队，强行穿过他们的队列：因为被神圣热忱点燃的灵魂不怕人，并将此类恐怖视为可笑的游戏。总督看见她，明白了所发生的事后，命人将她带到面前，询问她往哪里去。她说：\Quote{我听说正计划攻击主的仆人；我想加入我在信德中的朋友们，好与他们分享你所施加的屠杀。}总督说：\Quote{但那婴儿，你究竟带他去做什么？}她说：\Quote{好叫他与我分享我所渴望的死亡。}\par

当总督从这妇女口中听到这些，并通过她发现了激励全体人民的热情后，他向皇帝报告，并指出预谋屠杀的无用。他说：\Quote{我们只会从这一行为中收获不名誉，而无法扑灭这些人的精神。}然后他不允许群众遭受他们所预期的酷刑，命令他们的领袖——我指的是司铎和执事——带到面前，并向他们提供两个选择：要么诱导羊群与狼共融，要么被流放到某个偏远地区。然后他召集群众到面前，用温和的言辞劝他们服从皇帝的法令，力辩说少数可数的人对抗如此大帝国的君主简直是疯狂。人群默然无语。然后总督转向他们的领袖\Person{欧洛吉乌}，一个优秀的人，说：\Quote{你为什么对我所说的话不予回答？}\Person{欧洛吉乌}说：\Quote{我认为当没有向我提问时，我不必回答。}总督说：\Quote{但我用了许多论点劝你们采取对你们有利的做法。}\Person{欧洛吉乌}反驳说这些请求是向全体群众提出的，他认为自己贸然上前回答是荒谬的；\Quote{但是，}他继续说，\Quote{如果你问我个人意见，我会告诉你。}总督说：\Quote{那么，与皇帝共融。}\Person{欧洛吉乌}带着愉快的讽刺继续说：\Quote{他难道既领受了祭司职分又领受了帝国吗？}总督于是察觉他不是在认真说话，很不高兴，在向老人倾泻斥责之后，补充说：\Quote{我没有这样说，你这傻瓜；我劝你们与皇帝与之共融的人共融。}对此老人回答说他们有一位牧人，并服从他的指示，于是他们中的八十人被捕，被放逐到\Place{色雷斯}。在途中，他们处处受到极其崇高的尊敬，城市和村庄出来迎接他们，尊崇他们如同得胜的运动员。但嫉妒武装了他们的反对者向皇帝报告说，被视为耻辱的事实际上给这些人带来了巨大的荣誉；于是\Person{瓦伦斯}命令将他们分成两人一组，送往不同方向，一些到\Place{色雷斯}，一些到\Place{阿拉伯}最远的地区，另一些到\Place{忒拜}的城镇；据说，本性相连的人被野蛮人拆散，兄弟与兄弟分离。他们的领袖\Person{欧洛吉乌}与次位的\Person{普罗托革内斯}被放逐到\Place{安提诺}。\par

我也不会让这些人的德行湮没无闻。他们发现本市主教与自己心意相同，因此参加了教会的聚会；但看到聚会人数很少，经询问得知本市居民都是外邦人，他们便如所预料地感到忧伤，哀叹他们的不信。但他们认为仅仅忧伤还不够，而是尽其所能致力于使这些人痊愈。神圣的\Person{欧洛吉乌斯}被关在一个小房间里，日以继夜地向宇宙的天主献上祈求；令人钦佩的\Person{普罗托革涅斯}受过良好教育，并熟练于快速书写，他选定了一个合适的地点，开办了一所男童学校，以教师身份设立，不仅教导学生们快速书写的技艺，还教导他们天主的圣言。他教他们\Person{达味}的圣咏，并让他们学习宗徒道理中最重要的事项。其中一个男孩生病了，\Person{普罗托革涅斯}前往他家中，握住病人的手，藉祈祷驱除了疾病。其他男孩的父母听说此事后，将他请到家中，恳求他帮助病人；但他拒绝在病人领受圣洗恩赐之前请求天主驱除疾病。由于渴望孩子们健康，父母们欣然同意，最终获得了灵魂和身体的救恩。每当他劝服任何一个健康的人接受天主的恩宠时，他就领他到\Person{欧洛吉乌斯}那里，敲门恳求他开门，并将主的印记盖在这猎物上。当\Person{欧洛吉乌斯}因祈祷被打断而烦恼时，\Person{普罗托革涅斯}常说，拯救迷路者更为重要。他做这些奇妙的工作，向如此多的人传授神圣知识之光，同时始终将首位让给别人，并将他的战利品带到\Person{欧洛吉乌斯}面前，这令所有看到他的事迹的人钦佩。他们正确地推测，\Person{欧洛吉乌斯}的德行远为更伟大崇高。\par

在风暴平息、完全恢复平静之后，他们奉命返回家乡，全体人民嚎啕痛哭护送他们，尤其受到该教会的主教的护送，他如今失去了他们的牧养。当他们到家时，伟大的\Person{巴尔塞斯}已被接到无痛苦的生命中，神圣的\Person{欧洛吉乌斯}受托执掌他所领导过的教会的舵；卓越的\Person{普罗托革涅斯}被委以\Place{卡雷}的牧职，那是一片荒芜之地，长满了异教的荆棘，需要大量的劳作。但这些事发生在和平重归教会之后。\par

% structure: p0076 source=h2 target=subsection reason=relative-heading-rank; base=h2

\subsection{第16章 关于神圣的\Person{巴西略}，\Place{凯撒勒雅}的主教，以及\Person{瓦伦斯}和长官\Person{莫德斯图斯}对他采取的措施}

几乎可以说，\Person{瓦伦斯}剥夺了每一教会的牧者，并前往\Place{卡帕多细亚}的\Place{凯撒勒雅}，那时正是伟大的\Person{巴西略}（世界之光）的教座所在地。他事先派长官前往，命令他要么说服\Person{巴西略}接受\Person{欧多克修斯}的共融，要么在他拒绝时，以流放惩罚他。由于他事先知道主教的高尚声誉，起初不愿攻击他，因为他担心主教若勇敢地迎击并抵挡他的攻击，会给其他人树立勇气的榜样。这巧妙的计谋如蜘蛛网般无效。因为古代流传的故事对其余的主教们已足够，他们犹如城墙环中的堡垒，使信德之墙岿然不动。\par

然而，长官抵达\Place{凯撒勒雅}后，派人请来伟大的\Person{巴西略}。他恭敬地接待他，用温和有礼的语言劝他顺应时势所需，不要因教义上的一点细微差别而抛弃如此众多的教会。此外，他还许诺给他皇帝的友谊，并指出藉此友谊他可以为许多人带来巨大的利益。神圣的人说：\Quote{这种谈话适用于小男孩，因为他们和类似的人容易吞下这种诱饵。但那些受天主圣言滋养的人，连神圣信经的一个音节也不肯放弃，并且为了信经，若有必要，他们准备接受各种死亡。皇帝的友谊，如果与真宗教结合，我珍视它；否则，我视它为致命之物。}\par

然后长官被激怒，宣称\Person{巴西略}神智不清。但神圣的人说：\Quote{但这个疯狂我祈求常在我身。}主教随即被命令退下，思考应采取的途径，并在次日宣布他的决定。此外，威胁与劝告并举。据说这位杰出主教的答复是：\Quote{明天我仍会是今天这个人；你自己不要改变，尽管执行你的威胁吧。}在这些讨论之后，长官遇见皇帝，报告了谈话，指出主教的德行和他性格中无畏的丈夫气概。皇帝一言不发地走进内宫。在宫中，他看到来自天上的灾祸已降临，因为他的儿子卧病在死亡的门口，他的妻子被多种疾病困扰。于是他认识到这些悲伤的缘由，便恳求那位他曾以惩罚威胁过的神圣之人到他的府上来。他的官员执行了皇帝的旨意，然后伟大的\Person{巴西略}来到了皇宫。\par

\Person{瓦伦斯} 见皇帝之子濒临死亡，便应许他若由虔诚者领受圣洗圣事即可复生，并以此承诺离去。但皇帝如同愚昧的 \Person{黑落德}，记起自己的誓言，命在场的部分亚略派人士为孩童施洗，孩童随即死去。于是 \Person{瓦伦斯} 后悔，意识到遵守誓言何等危险，便来到神圣殿宇，领受了伟大的 \Person{巴西略} 的教导，并在祭台前奉献了惯常的礼品。主教更命他进入神圣帷幕之内，与他同坐，详谈天主圣意，亦聆听其言。\par

当时有一人名为 \Person{德摩斯梯尼}，为御膳房总管，他粗鲁地责骂这位教导天下之人，却犯了文法错误。\Person{巴西略} 微笑说：\Quote{我们在此见到一位无知的德摩斯梯尼。} 德摩斯梯尼动怒出言威胁，他续道：\Quote{你的事是料理羹汤调味；你不可能明白神学，因你的耳朵被堵住了。} 他如此说，皇帝甚为喜悦，便将自己在那里的几块良田赐给他照料下的穷人，因为那些人身患重病，急需照料与医治。\par

伟大的 \Person{巴西略} 就这样避开了皇帝的首次攻击。但当皇帝第二次前来时，他的正确判断受到欺骗他的谋士阻挠，他忘记了先前发生的事，命 \Person{巴西略} 归入敌对派系。劝说无效后，他下令执行放逐谕令。然而当他试图签字时，竟连一字的一划也写不成，因为笔折断了。第二支、第三支笔同样折断，他仍竭力签署那道邪恶谕令，手却颤抖，全身战栗，心灵充满恐惧；他用双手撕碎文件，如此天地的主宰证明了：正是祂准许其余人等经受这些苦难，却使 \Person{巴西略} 比所设的网罗更坚强，并借 \Person{巴西略} 的一切事迹宣告了自己的全能，同时又传扬了善人的勇气。于是 \Person{瓦伦斯} 攻击未果。\par

% structure: p0083 source=h2 target=subsection reason=relative-heading-rank; base=h2

\subsection{第十七章  论伟大的 \Person{亚大纳削} 之死与 \Person{伯多禄} 的当选}

在 \Place{亚历山大里亚}，胜利的 \Person{亚大纳削} 历经了一切战斗——每场均获荣冠——终于从劳苦中得到解脱，进入那无劳苦的生命。于是 \Person{伯多禄}，一位极为卓越的人，继承了教座。他蒙福的前任首先拣选了他，圣职人员和有地位有官职者一致赞同，全体民众也用欢呼表达喜悦。他曾分担 \Person{亚大纳削} 的沉重劳苦，无论在本地还是外地都常伴随其左右，与他一同经历诸多危险。因此邻近的主教们急忙前来会合；那些住在苦修学府中的人也离开那里加入队伍，众人一同恳求选立 \Person{伯多禄} 接续 \Person{亚大纳削} 的宗主教席位。\par

% structure: p0085 source=h2 target=subsection reason=relative-heading-rank; base=h2

\subsection{第十八章  论 \Person{伯多禄} 的废黜与亚略派 \Person{路基约} 的引入}

他们刚将他安置在主教座，该省总督便聚集一群希腊人与犹太人，包围了教堂的墙壁，命 \Person{伯多禄} 出来，并威胁若不从便将放逐他。他这样做的借口是执行皇帝的旨意，使反对者陷入困境，但实际是被自己的不虔之情所驱策。因他沉迷于偶像崇拜，将教会的风暴视为辉煌的节庆。然而可敬的 \Person{伯多禄} 见到这突如其来的冲突，便暗中离开，登上一艘驶往 \Place{罗马} 的船。\par

数日后，\Person{欧佐约} 与 \Person{路基约} 一同从 \Place{安提约基雅} 前来，将众教会交给他。此人正是 \Place{萨莫萨塔} 已尝过其不虔与不法滋味者。但那些在 \Person{亚大纳削} 教导中栽培的民众，如今见所供灵粮迥异，便远离了教会的集会。\par

\Person{路基约} 用偶像崇拜者作随从，鞭打一些人，囚禁另一些人；又驱赶一些人逃亡，粗暴残忍地洗劫他们的家宅。但这一切在可敬的 \Person{伯多禄} 的信函中有更好的叙述。我先举一事例说明 \Person{路基约} 的不虔行为，然后将该信函载入本书。\par

在埃及有一些天使般生活与谈吐的人，为逃避国家的不安，选择在旷野独居。他们使沙质贫瘠的土地结出果实；因为结给天主的果实是何等甘甜美好，乃是他们赖以生活的德行之律。在众多引领这种生活方式的人中，有闻名遐迩的\Person{安当}，他是克修学校中最卓越的大师，使旷野成为其隐修士们的德行训练场。在他完成一切伟大光荣的劳作、抵达不再有忧患之风侵袭的港湾后，他的追随者们便遭到可悲不幸的\Person{路济乌斯}迫害。那些神圣团体的一切领袖——著名的\Person{玛加略}、与他同名的另一位、\Person{依西多禄}及其余的人——都被从山洞中拖出，遣送到一个由不敬神的人居住的岛屿，那里从未有过任何虔敬导师的祝福。当船靠近岛屿海岸时，岛民所崇敬的魔鬼离开了它自古作为居所的偶像，并使司祭的女儿充满狂乱。她受灵感驱使，冲到船夫们靠岸的海滩上。魔鬼以这女孩的舌头为工具，通过她喊出那在斐理伯被邪神之灵附身的妇人所说的话，众人——无论男女——都听见了，说：\Quote{唉，你们基督仆人的能力啊！我们到处被你们从城镇、村庄、山丘、高地、无人居住的荒野赶出；我们原指望在这小岛上远离你们的利箭，但我们的指望落空了。你们被你们的迫害者派到这里来，不是要受他们伤害，而是要赶走我们。我们正要离开这岛，因为我们被你们德行的锐利光芒所伤。}说完这些话及类似的话，他们把少女摔倒在地，然后一起逃走了。但那神圣团体为女孩祈祷，扶起她，把她完好且神智清醒地交还给她的父亲。\par

目睹这奇迹的旁观者扑倒在新来者脚下，恳求参与救恩的途径。他们拆毁了偶像的丛林，受教诲之光明照耀，领受了圣洗圣事的恩宠。这些事传到\Place{亚历山大里亚}，所有民众聚集起来辱骂\Person{路济乌斯}，说若不释放那神圣的圣人团体，天主的义怒必降在他们身上。于是\Person{路济乌斯}怕城中发生骚乱，便允许这些圣隐修士回到他们的洞穴。此例足以显示他不敬神的邪恶。他所敢于犯下的罪恶行径，将由杰出的\Person{伯多禄}的书信更清楚地说明。我犹豫是否全文录下，故只从中摘引若干片段。\par

% structure: p0091 source=h2 target=subsection reason=relative-heading-rank; base=h2

\subsection{第十九章 \Place{亚历山大里亚}在亚略派\Person{路济乌斯}时期事件叙述，取自\Place{亚历山大里亚}主教\Person{伯多禄}的一封信}

该省省长\Person{帕拉迪乌斯}，教派为外教人，惯常在偶像前俯伏，常怀向基督开战之念。在召集上述兵力后，他出发攻击教会，仿佛正奋勇征服一个外敌。随后，众所周知，最骇人听闻的事发生了；光是想到要叙述这事，回忆就使我充满痛苦。我曾泪如雨下，若非藉神圣默想缓解悲伤，我必长久深陷悲恸。人群闯入称为德奥纳的教堂，那里代替圣洁的话语的是偶像的颂赞；那里曾经宣读圣经之处，可听到不雅的鼓掌声与不男不女的下流言语；那里对基督的童贞女施加了言语无法述说的暴行，因为\Quote{连提起来也是可耻的}。一听到这些恶行，有善意者便掩耳，祈求宁可变聋也不听那污秽言语。但愿他们只满足于言语的罪，而不以行为超越言语的邪恶；因为无论侮辱如何严重，凡有基督智慧及其神圣教训内住的人皆能忍受。然而这些恶徒，这愤怒的器皿预备遭毁灭，他们皱起鼻子，从鼻孔中流出——我姑且这样说——如同从泉源涌出的污秽声响，并从基督的圣童贞女身上撕扯衣服，这些女子的言谈举止恰如圣人的写照；他们赤裸裸地拖着她们游街示众，如同出生时一般；他们随意对她们行淫秽戏弄；其行径野蛮残忍。若有人怜恤地干预并劝其慈悲，便被击伤而去。唉！我何其不幸！许多童贞女遭受野蛮蹂躏；许多少女被棍棒击头，倒地无声，甚至不许为她们收尸；其悲痛欲绝的父母至今找不到他们的尸体。但为何要诉诸与更大苦难相比似乎微小的悲哀呢？为何滞留于此而不赶赴更紧迫的事件？当你们听到这些时，我知道你们必惊讶，并必与我们长久无言，惊异于主的仁慈，他没有将所有一切都完全灭绝。就在祭台上，不敬神的人行了那些事，正如经上所载：\Quote{在我们祖先的日子未曾发生过，也未曾听说过。}\par

一个男孩，他否认自己的性别，想要装作女孩，眼睛涂着锑粉（正如经上所载），脸上擦着胭脂，像他们的偶像一样，穿着女装，被放在我们呼求圣神降临的神圣祭台上跳舞、挥手、旋转，仿佛他是在某个下流舞台的前台，而旁人大声嘲笑并粗鲁地发出不雅喊叫。但既然这在他们看来实在算是得体而非不当，他们便继续进行他们视为符合其无耻行径的举动；他们挑出一个以极其卑鄙闻名的人，让他立刻脱掉所有衣服和一切羞耻，将他赤裸裸地（如同他出生时一样）放在教会的宝座上，并称他为反对基督的卑鄙辩护者。然后，他以神圣之词说出无耻的邪恶，以可敬的教义说出放荡的淫秽，以虔诚说出不虔，以节操说出奸淫、通奸、肮脏的情欲、偷窃；教导饕餮和醉酒以及其余一切对人的生命都有益处。事态如此，连我也离开了教会——我怎能留在军队进来、暴民被收买行凶、所有人都争相牟利、外教人的暴民许下大肆诺言的地方呢？——于是，确实，派来了一个继任我的人。他名叫\Person{路济乌斯}，像买这世界的某种尊位一样买下了主教职位，热切要维持一只狼的恶名和品行。没有任何正统主教会议选举他；没有真正神职人员的投票；没有平信徒要求他；正如教会的法律所规定的。\par

\Person{路济乌斯}不能没有排场地进入该城，因此陪伴他的是——不是主教，不是司铎，不是执事，不是众多平信徒；没有隐修士在他前面吟咏圣经中的圣咏；而是\Person{欧佐伊乌斯}，他曾是\Place{亚历山大}城的执事，很久以前在伟大的神圣\Place{尼西亚}大公会议上与\Person{亚略}一同被贬黜，最近又被提升来统治和蹂躏\Place{安提约基雅}教座；还有司库\Person{玛格努斯}，他以各种不虔而臭名昭著，率领着一支庞大的军队。在\Person{尤利安}统治期间，这个\Person{玛格努斯}曾烧毁了\Place{腓尼基}名城\Place{贝里图斯}的教堂；而在可纪念的\Person{约维安}统治期间，他多次恳求皇恩才勉强逃脱斩首，被迫自费将其重建。\par

现在，我恳请你们的热忱兴起为我们辩护。从我写的内容，你们应该能计算出这个\Person{路济乌斯}兴起反对我们时对天主的教会所造成的伤害的性质和程度。他多次被你们的虔诚和各地正统主教拒绝，却占领了一座有正当理由将他视为敌人并如此对待的城市。因为他不仅像圣咏中那位亵渎的愚人那样说\Quote{基督不是真天主。}而且，他自己败坏，也败坏别人，为那些崇拜受造物而非造物主的人不断对救主发出的亵渎而欢喜。这个恶棍的观点与外教人完全相同，他为什么不敢崇拜一个新造的神呢？因为人们当着面用这些话欢迎他：\Quote{欢迎，主教，因为你否认圣子。\OriginalTerm{塞拉比斯}{Serapis}爱你，把你带到了我们这里。}他们就是这样称呼自己本地的偶像。然后，没有片刻耽搁，上述的\Person{玛格努斯}——\Person{路济乌斯}邪恶中不可分离的同伙，残忍的保镖，野蛮的副官——集合了他所管辖的所有群众，逮捕了十九名司铎和执事，其中一些已经八十岁，指控他们涉及某种对罗马法律的卑鄙违反。他组成了一个公开的法庭，并且由于对基督徒维护德行的法律无知，试图强迫他们放弃从宗徒经教父们传给我们的祖传信仰。他甚至声称这将是令最慈悲、最宽厚的\Person{瓦伦斯·奥古斯都}满意的。\Quote{可怜的人啊}，他喊道，\Quote{接受吧，接受亚略派的教义；即使你们用真正的敬拜来敬拜天主，如果你们这样做不是出于自愿而是被迫的，天主也会宽恕你们。非自主的强迫总有辩护，而自由行动则是有责任的，且常伴随着谴责。好好考虑这些论点；自愿地来吧；不要再拖延了；签署现在由\Person{路济乌斯}宣讲的亚略的教义，}（他这样点名介绍他）\Quote{要确信，如果你们服从，你们将从君王那里得到财富和荣誉；如果你们拒绝，你们将被锁链、拉肢架、酷刑、鞭子和残忍的折磨所惩罚；你们将被剥夺财产和产业；你们将被流放，并被判居住在荒凉之地。}\par

\Quote{停止吧，停止试图用这些话来恐吓我们，不要再妄言。我们敬拜的不是一个迟来或新发明的神。如果你们愿意，就在你们怒气的虚妄风暴中向我们吐唾沫，像一阵狂风吹向我们来吧。我们坚持真宗教的教义直至死亡；我们从未认为天主是软弱的、不明智的或不真实的，像这不虔的亚略派所教导的那样，认为天主有时是圣父，有时又不是圣父，使圣子成为一个时间和暂时的存在。因为如果像亚略狂们所说，圣子是受造物，与圣父不是本性上的同体，那么圣父也会因圣子的非存在而归于非存在，因为他们断言圣父有一个时期不是圣父。但如果他永远是圣父，他的后裔是真正出自他，而非通过派生（因为天主是不能受苦的），那么谁说圣子——万物藉着恩宠藉他得以存在——\InnerQuote{有一个时期他并不存在}，岂不是疯狂和愚昧吗？}\par

这些人因背离了我们在世界各地聚集于\Place{尼西亚}的教父们，并弃绝了亚略的虚假教义（而今这教义正由这位后来者辩护），而真正成了无父之人。教父们规定了圣子并非如你们现今强逼我们所说的那样，与圣父有不同的实体，而是同一且相同的。他们虔诚的理智清楚地领悟了这一点，因此，通过对神圣术语的充分对照，他们承认圣子与圣父是\OriginalTerm{同性同体}{consubstantial}。\par

他们提出这些及类似论点后，被囚禁了多日，希望他们能被诱使背离其正确心智；但他们反而像竞技场上最高贵的运动员一样，粉碎了一切恐惧，并不时以父辈的英勇事迹为思想膏油，借神圣思想的帮助，在虔诚中保持了更高贵的恒心，将刑架视为德行的训练场。当他们如此奋战时，正如蒙福的\Person{保禄}所写，成了天使和世人的一场景象，全城的人都跑来观看基督的运动员，他们以坚忍的毅力战胜了施刑法官的鞭打，以忍耐赢得了反对不敬虔的凯旋，展现了对抗亚略派的胜利。因此，他们的凶恶敌人以为能通过威胁和酷刑征服他们，并将他们交给基督的敌人。就这样，这残暴不仁的暴君以残忍的诡计施加酷刑，恶待他们，而全体百姓则站着哀哭，以各种方式表达悲痛。随后，他再次集结他那些在无序中训练的军队，在港口传唤殉道者受审，或更应说，是预先定罪的判决，而偶像崇拜者和犹太人也按他们的方式发出雇来的喊叫，反对他们。由于他们拒绝向亚略狂徒的明显异端屈服，他们被判处（全体百姓含泪站在法庭前）从亚历山大城流放到腓尼基的\Place{赫利奥波利斯}，那个地方的居民都沉迷于偶像，甚至连基督的名字都受不了听到。\par

在他们登船的命令下达后，\Person{玛格努斯}驻守在港口，因为他在公共澡堂附近宣判了他们的刑罚。他拔出出鞘的剑给他们看，以为这样能恐吓那些曾用双刃剑一次又一次击倒敌对魔鬼的人。于是他命他们出海，尽管他们没有带上任何食物，也没有为流放准备任何慰藉。说来奇怪，几乎难以置信，海面泛起泡沫；我想，大海是不愿、或可说是拒绝接纳这些好人于其上，从而分担不义判决的份。现在，连无知的人也看清了法官的残暴意图，真可以说：\Quote{为此，天惊骇不已。}\par

全城哀叹，直至今日仍在悲悼。一些人交替捶胸，发出巨大的响声；另一些人同时举起手和眼向天，作为对所受冤屈的见证，并几乎用话语说：\BibleQuote{天哪，要听；地啊，侧耳而听，}这是何等非法之事。如今一片哭泣哀号；全城回荡着歌唱与叹息，泪水如河从每只眼中流出，几乎要用其潮水淹没大海。前述的\Person{玛格努斯}在港口命令桨手升帆，同时传来少女与主妇、老翁与青年的混合哭喊，众人齐声悲泣哀号，群众的喧嚣淹没了泡沫海水因波涛而起的咆哮。于是殉道者乘船前往\Place{赫利奥波利斯}，那里每个人都沉迷于迷信，魔鬼的享乐之道盛行，而城市四周群山接天，正适合野兽的恐怖巢穴。他们留下的所有朋友，无论是公开在城中央，还是各自私下，都呻吟着吐出悲伤的话语，甚至被市府长官\Person{帕拉狄乌斯}下令禁止哭泣；此人自己恰好是个完全沉迷于迷信的人。许多哀悼者先被逮捕投入监狱，然后被鞭打、用梳理梳撕裂、受酷刑，并因其神圣热忱而成为教会斗士，被遣送往斐尼苏斯和普罗科涅苏斯的矿场。\newline{}\par

其中大多数是隐修士，度着独修的刻苦生活，约有二十三人。不久之后，我们敬爱的罗马主教\Person{达玛稣}所派来给我们送安慰与共融信函的执事，被刽子手双手反绑，如同臭名昭著的罪犯，公开游街示众。在遭受了加于杀人犯身上的酷刑之后，他的脖颈被石头和铅块狠狠地鞭打。他登船启航，额上带着神圣十字架的标记，与其余人一样；无人援助，也无人诱惑他，他被送往\Place{芬内苏斯}的铜矿。在地方官对幼弱孩童施加的酷刑中，有些孩子被当场弃置，不得神圣葬礼；尽管父母、兄弟、亲属，乃至全城都恳求至少能给予这一慰藉。但唉，法官的残忍啊——若那仅会判刑的人还能称为法官的话！为真道奋勇作战的人竟落得比杀人犯更悲惨的下场，尸首无人埋葬。光荣的战士被抛给野兽和猛禽吞噬。那些因良心不安而想向父母表达同情的人，竟如犯法般被斩首。哪一部罗马法律，甚至哪个外邦的习俗，曾因向父母表达同情而施加刑罚？古人中何时有过如此非法之举？希伯来人的男婴确实曾奉法老之命被杀，但法老的谕令出于嫉妒和恐惧。今日的残忍远比法老的时代更甚。若在不义之中可作选择，他们的不公远胜于我们的。若非法之事尚可论好坏，那么他们的过错又强多少呢？实际上，不义始终是不义。\par

我所记述之事难以置信、不近人情、可怕、野蛮、残忍、冷酷。但在这一切之中，亚流派狂热的追随者却像骄傲地欢腾，而全城都在哀悼；因为正如\BibleBook{出谷}{纪}所载：\BibleQuote{没有一个家里没有死人。}\par

那些从不满足于作恶的人策划了新的骚动。他们始终以恶行施展其恶念，将不义的毒液喷向该省的主教们，利用前述司库\Person{玛格诺斯}作为他们不义的工具。\par

有些人被他们交给元老院，有些人则被他们随意诱捕；他们不遗余力，从四面八方搜捕所有人，意图使其陷入不敬；他们四处游走，如同异端的真正父亲魔鬼一样，寻找可吞噬的人。\par

最终，在多次徒劳的努力后，他们将十一位埃及主教流放到\Place{迪奥凯撒肋雅}，那是一座由杀害主的犹太人所居住的城市。这些主教从幼年到老年，在旷野中度着刻苦的独修生活，以理性和纪律克制私欲的倾向，无畏地宣讲虔诚的信德，汲取虔诚的教义，多次战胜魔鬼，常以德行使仇敌羞愧，并以最智慧的论证公开驳斥亚流派异端。然而，他们如地狱一般，不满足于兄弟之死，这些愚妄之徒渴望以恶行博取名声，试图在全世界留下他们残忍的记忆。如今，他们竟激起皇帝的注意，反对安提约基雅的天主教会中的某些神职人员，以及一些前来作证反对其恶行的卓越隐修士。他们使这些人被流放到\Place{本都}的\Place{新凯撒肋雅}，不久之后便因当地土地贫瘠而丧命。此时期上演了如此悲剧，本当缄默遗忘，但在历史中记下一笔，以谴责那些以口舌攻击独生子的人，他们被亵渎的疯狂所感染，不仅箭指宇宙之主，更向祂忠实的仆人发动无休止的战争。\par

% structure: p0106 source=h2 target=subsection reason=relative-heading-rank; base=h2

\subsection{第二十章　论萨拉森人的女王\Person{玛维亚}及隐修士\Person{梅瑟}被祝圣为主教}

此时，依市玛耳人正在罗马边境附近蹂躏乡村。他们由一位王后\Person{玛维亚}率领，她不顾天性赋予的女性身份，展现出男子的精神与勇气。多次交战后，她签订了停战协定，并在接受了神圣知识的光照后，恳求使一位名叫\Person{梅瑟}的人升任她部落的大司祭，此人住在埃及和巴勒斯坦交界处。\Person{瓦伦斯}答应了这一请求，并命令将这位圣人护送到\Place{亚历山大里亚}，那里是附近最便利之地，以接受主教恩宠。当他抵达并看到\Person{路济乌斯}试图给他覆手时，他说：\Quote{愿天主不许我由你的手受祝圣；圣神的恩宠不会因你的呼召而降临。}\Person{路济乌斯}说：\Quote{你凭什么如此猜测？}他回答说：\Quote{我不是在说猜测，而是明确知道；因为你反对宗徒的谕令，言语攻击它们，你无法无天的行为与你的亵渎之言相称。有哪个不敬之人因你而未曾嘲笑教会的聚会？有哪位卓越之人未被流放？哪种野蛮的残暴不被你每日的行为所超越？}如此，这位勇敢者说道；那杀人犯听了，想要杀害他，却害怕再次点燃业已平息的战火。因此，他命令带出\Person{梅瑟}所要求的其他主教。\Person{梅瑟}领受了合乎正道信仰的主教恩宠后，回到请求他的百姓中，以宗徒般的教导和奇迹引导他们走上通往真理的道路。\par

这些就是在天主的眷顾安排下，\Person{路济乌斯}在\Place{亚历山大里亚}所做之事。\par

% structure: p0109 source=h2 target=subsection reason=relative-heading-rank; base=h2

\subsection{第二十一章}

在\Place{君士坦丁堡}，亚略异端者将虔诚的司铎们装满一条船，不加压舱物就驱逐出海，又派他们自己的人乘另一条船，下令焚烧司铎们的船。于是，他们既与大海又与火焰搏斗，最后沉入深海，赢得了殉道者的冠冕。\par

\Person{瓦伦斯}在\Place{安提约基雅}停留了相当长的时间，并给予所有以基督徒名义作掩护之人——外教人、犹太人及其他——完全的自由，任其宣讲与福音相悖的教义。这谬误的奴仆甚至行外教礼仪，于是，那在\Person{犹利安}之后被\Person{若维安}扑灭的欺诈之火，如今因\Person{瓦伦斯}的许可再度燃起。犹太人、狄俄倪索斯和得墨忒耳的礼仪，不再像在虔敬的统治下那样只在隐蔽处举行，而是在广场上由狂欢者狂奔乱舞。\Person{瓦伦斯}唯一敌视的，是那些持守宗徒教义之人。他首先将他们逐出教堂——殊荣的\Person{若维安}曾将新建的教堂也赐予他们。当他们聚集在山崖附近，以赞美诗尊崇他们的主，享受天主的圣言，忍受各种天气的侵袭——时而雨雪严寒，时而酷热——连这点微薄的庇护也被剥夺，军队被派去将他们远远驱散。\par

% structure: p0112 source=h2 target=subsection reason=relative-heading-rank; base=h2

\subsection{第二十二章　\Person{弗拉维亚诺斯}与\Person{狄奥多禄}如何聚集\Place{安提约基雅}的正统教会。}

\Person{弗拉维亚诺斯}与\Person{狄奥多禄}如同防波堤，挡住了涌来的浪潮。他们的牧者\Person{默肋提乌斯}被迫远居他乡。但此二人在抵御豺狼时，以勇气与智谋作战，并悉心照料羊群。如今，他们被从山崖下驱逐，便在邻近的河岸牧养羊群。他们不能像\Place{巴比伦}的俘虏那样，把琴挂在柳树上，却继续在天主统治的所有地方歌颂他们的创造者与施恩者。然而，即便在此处，仇敌也不容虔诚的牧者聚集赞美主的会众。于是，这对卓越的牧人再次将他们的羊群聚集在练兵场，在那里秘密地给他们展示精神食粮。\Person{狄奥多禄}以他的智慧与勇气，如同清澈浩荡的江河，灌溉自己的田地，淹没了对手的亵渎之言，丝毫不顾自己显赫的出身，甘愿为信德受苦。\par

卓越的\Person{弗拉维亚诺斯}，也出身最高贵，却认为虔诚才是唯一的高贵。他如同一位赛场的教练，膏抹伟大的\Person{狄奥多禄}，仿佛他是将参加五场竞赛的运动员。\par

那时，他本人并不在教会礼仪中讲道，却为讲道者提供丰富的论据和圣经思想，使他们能将箭矢射向\Person{亚略}的亵渎，而他则像从箭囊中递出智慧的箭矢。无论在家或在外，他滔滔不绝地讲论，轻易撕破异端者的网罗，指出他们的防御不过是蛛网。在这些竞赛中，他得到那位\Person{阿弗拉特}的帮助——我在\WorkTitle{宗教史}中记述了他的生平——此人宁愿羊群的福祉胜过自己的安息，放弃他那规矩独修的静室，承担起牧人的辛劳。既然在其他著作中已写了这些事，我认为现在不必赘述他所积累的美德财富，但我要叙述他的一件善行范例，因它与本史极为相关。\par

% structure: p0116 source=h2 target=subsection reason=relative-heading-rank; base=h2

\subsection{第二十三章　论圣人隐修士\Person{阿弗拉特}}

在\Place{奥龙特斯}河以北坐落着宫殿。河以南，城墙之上建有一座巨大的两层柱廊，两侧有高塔。在宫殿与河流之间有一条公共道路，从镇上经过这个区域的城门通向郊外的乡村。虔诚的\Person{阿弗拉阿特}曾有一次沿着这条大道前往士兵训练场，以便履行牧养他羊群的职责。皇帝恰巧从宫殿的一个廊台上往下看，看见他披着一件未经鞣制的山羊皮斗篷，虽然年事已高，却走得很快。当有人指出这就是全镇那时都敬重的\Person{阿弗拉阿特}时，皇帝喊道：\Quote{你往哪里去？告诉我们。}他机敏而巧妙地回答说：\Quote{为你的帝国祈祷。}皇帝说：\Quote{你最好待在家里，像隐修士一样独自祈祷。}那位神人回答说：\Quote{是的，我本来应该这样做，并且直到现在我也一直这样做，只要救主的羊群平安无事；但现在它们被严重搅扰，大有被野兽捕食的危险，我不得不竭尽全力拯救这些幼小者。因为，先生，请告诉我，假如我是一个坐在闺房看家的女孩，看见一道火焰落下，我父的家着了火，我该做什么？告诉我：是坐在里面，不管房子着火，等待火焰逼近？还是告别我的闺房，跑来跑去取水，设法扑灭火焰？你当然会说后者，因为一个敏捷而活泼的女孩会这样做。先生，我现在所做的正是这样。你点燃了我们父的家，我们正在四处奔跑，努力扑灭它。}阿弗拉阿特这样说，皇帝威胁了他，但没再多说。一位皇帝的寝宫侍从，比较激烈地威胁了那位神人，遭遇了如下命运。他负责浴室，在这次谈话之后立即下来为皇帝准备浴室。一进去他就失去了理智，在凉水尚未掺入之前就踏进了沸水，结果丧了命。皇帝坐着等他通报浴室已准备好进入，过了相当长的时间后，他派其他官员去报告延迟的原因。他们进去查看了整个房间，发现那位内侍被热力杀死，倒在沸水中死了。这事传到皇帝耳中，他们感受到了\Person{阿弗拉阿特}祈祷的力量。然而他们并没有离开那邪恶的教义，反而像\Person{法郎}一样硬了心；那位执迷不悟的皇帝，虽然知道了圣人的奇迹，却仍坚持他那反对虔诚信仰的疯狂暴怒。\par

% structure: p0118 source=h2 target=subsection reason=relative-heading-rank; base=h2

\subsection{第二十四章 论圣隐修士\Person{尤利安努斯}}

此时，那位我已经提到过的著名\Person{尤利安努斯}，也被迫离开旷野来到\Place{安提约基雅}，因为当谎言的养子、轻易捏造诽谤者——我指的是亚略派人——声称这位伟人是他们一派时，真理之光\Person{弗拉维亚努斯}、\Person{狄奥多若}和\Person{阿弗拉阿特}派遣了一位德行的斗士\Person{阿卡基乌斯}（他后来极其明智地治理了\Place{贝罗厄亚}教会）到著名的\Person{尤利安努斯}那里，恳求他怜悯如此众多的民众，同时揭露敌人的谎言，确认真理的宣告。\Person{尤利安努斯}往返\Place{安提约基雅}以及在那座大城市本身所行的奇迹，我在我的《\WorkTitle{虔诚史}》中已有记述，所有愿意了解的人都容易读到。但我确信，凡是研究过人性的人都不会怀疑，他吸引了全城的人到我们的集会中，因为非凡的事物通常能吸引所有人追随它。他行了大奇事这一事实，甚至由真理的敌人所证实。\par

在此之前，在\Person{君士坦提乌斯}统治时期，伟大的\Person{安多尼}在\Place{亚历山大里亚}也做了同样的事，因为他离开了旷野，在那座城市里四处走动，告诉所有人：\Person{阿塔纳修斯}是真道理的宣讲者，亚略派是真理的敌人。因此，那些虔诚的人知道如何根据每个具体时机调整自己：何时保持安静和休息，何时离开旷野前往城镇。\par

% structure: p0121 source=h2 target=subsection reason=relative-heading-rank; base=h2

\subsection{第二十五章 论此时有哪些其他隐修士出类拔萃}

此时还有其他一些人，放射出独修生活哲学的灿烂光芒。在\Place{卡尔基斯}旷野，\Person{阿维图斯}、\Person{玛尔奇安努斯}和\Person{阿布拉阿梅斯}，以及更多我难以一一列举的人，在可感知的身体中努力过一种超越感知的生活。在\Place{阿帕梅亚}地区，\Person{阿加佩图斯}、\Person{西默盎}、\Person{保禄斯}和其他人收获了最高智慧的果实。\par

在\Place{泽格马}地区，有\Person{普布利乌斯}和\Person{保禄斯}。在\Place{叙利安}地区，著名的\Person{阿克塞普塞玛斯}被关在斗室里六十年，既未被人看见，也未与人交谈。可敬的\Person{则乌玛提乌斯}，虽然失明，却常去坚固羊群，与狼搏斗；于是他们烧了他的斗室，但极其忠信的\Person{特拉雅努斯}将军为他另建了一间，还给了他其他照顾。在\Place{安提约基雅}附近，\Person{玛利亚努斯}、\Person{欧瑟比乌斯}、\Person{阿米安努斯}、\Person{帕拉狄乌斯}、\Person{西默盎}、\Person{阿布拉阿梅斯}等人，使天主的肖像保持完好；但所有这些人的生平我已记叙过。然而，靠近大城的那座山装饰得像一座草原，因为在那里闪耀着：来自加拉提亚的\Person{伯多禄}、与他同名的埃及人、\Person{罗玛努斯}、\Person{塞维鲁斯}、\Person{则诺}、\Person{梅瑟}和\Person{玛尔库斯}，以及许多不为世人所知、却为天主所认识的人。\par

% structure: p0124 source=h2 target=subsection reason=relative-heading-rank; base=h2

\subsection{第二十六章 论\Place{亚历山大里亚}的\Person{狄迪莫斯}和\Place{叙利亚}的\Person{厄弗冷}}

当时在\Place{厄德撒}，可敬的\Person{厄弗冷}大放异彩；在\Place{亚历山大}，\Person{狄迪慕斯}亦然。二人皆撰文驳斥违背真理的教义。\Person{厄弗冷}使用叙利亚语，放射出神性恩宠的光辉。他完全未受异教教育玷污，却能揭露异教谬误的细微之处，并使一切异端诡计的虚弱暴露无遗。\Person{巴尔达桑}之子\Person{哈耳摩尼}曾创作某些歌曲，以旋律之甜美混合其不敬，迷惑听众，引他们走向毁灭。\Person{厄弗冷}采纳了这些歌曲的曲调，却配以虔敬之词，从而给听众既带来极大的愉悦，又赐予治愈的良药。这些歌曲至今仍用于欢庆我们得胜殉道者的节日。\par

然而\Person{狄迪慕斯}自幼失明，却受过诗歌、修辞、算术、几何、天文、\Person{亚里士多德}的逻辑以及\Person{柏拉图}的雄辩术的教育。他仅凭听觉接受所有这些学科的教导——并非因其传递真理，而是因其可能成为真理对抗谎言的武器。对于圣经，他不仅学到了声音，更领悟了含义。因此，在当时，在过苦修生活与追求德行的人们中，这些人卓然超群。\par

% structure: p0127 source=h2 target=subsection reason=relative-heading-rank; base=h2

\subsection{第二十七章  当时在\Place{亚细亚}与\Place{本都}哪些主教著称于世}

主教中有两位\Person{格列高利}，一位是纳齐盎的\Person{格列高利}，另一位是尼撒的\Person{格列高利}；后者是伟大的\Person{巴西略}的弟弟，前者是\Person{巴西略}的朋友与同工。他们在\Place{卡帕多西亚}是虔敬的首要捍卫者；与他们并驾齐驱的还有\Person{伯多禄}，他与\Person{巴西略}和\Person{格列高利}同父母所生，虽未像他们一样受过异邦教育，却同样过着卓越辉煌的生活。\par

在\Place{彼息狄亚}，\Person{奥普提慕斯}；在\Place{吕考尼亚}，\Person{盎斐罗基乌斯}，均为维护祖传信仰而战斗在最前列，击退了敌人的进攻。\par

在西方，\Place{罗马}主教\Person{达玛苏斯}，以及受托治理\Place{米兰}的\Person{盎博罗削}，击退了来自远方的攻击者。与这些人一起，被迫居住在偏远地区的主教们通过著作坚固了朋友，瓦解了敌人——如此，宇宙的治理者赐予了能应付如此巨大风暴的舵手。面对敌人的暴力，祂将军长的美德列阵迎战，并提供了适宜的手段来抵御这些困难时期的纷扰；不仅教会因他们慈爱的主得到这种保护，而且被认为配得另一种引导。\par

% structure: p0131 source=h2 target=subsection reason=relative-heading-rank; base=h2

\subsection{第二十八章  \Person{瓦伦斯}就战争致信伟大的\Person{瓦伦提尼安}，以及后者如何回复}

主激起哥特人作战，并将那只会与虔诚者争斗的人引向\Place{博斯普鲁斯}。那时这个虚妄的人第一次意识到自己的软弱，便派人向他兄弟请求援军。但\Person{瓦伦提尼安}回答说，帮助一个与天主作战的人是邪恶的，更应制止他的鲁莽行为。这使那不幸的人更加执迷不悟，但他并未中止自己的鲁莽行动，仍坚持与真理为敌。\par

% structure: p0133 source=h2 target=subsection reason=relative-heading-rank; base=h2

\subsection{第二十九章  关于伯爵\Person{特伦提乌斯}的虔敬}

\Person{特伦提乌斯}，一位杰出的将领，以虔敬著称，他竖立了胜利的纪念柱，从\Place{亚美尼亚}返回。当\Person{瓦伦斯}命他选择一件赏赐时，他提出的请求正合一位在虔敬中养育之人所应选：他既不要金银，也不要土地、官职或房屋，只求准许将那一个教会赐给那些为宗徒的教导而不惜一切的人。\Person{瓦伦斯}收下请求书，但得知内容后，愤怒地将其撕碎，命\Person{特伦提乌斯}另求他物。然而，伯爵捡起请求书的碎片，说：\Quote{大人，我已得了我的赏报，不再求别的。万物的审判者是审判我意向的审判者。}\par

% structure: p0135 source=h2 target=subsection reason=relative-heading-rank; base=h2

\subsection{第三十章  关于\Person{特拉雅努斯}将军的直率之言}

\Person{瓦伦斯}渡过\Place{博斯普鲁斯}进入\Place{色雷斯}后，先是在\Place{君士坦丁堡}呆了相当长的时间，对战争结局感到恐惧。他曾派遣\Person{特拉雅努斯}统军对抗蛮族。当这位将军战败归来时，皇帝严厉斥责他，指责他软弱怯懦。\Person{特拉雅努斯}勇敢地，如同一个壮士所应做的，回答说：\Quote{我并未战败，大人，是你因与天主作战而放弃了胜利，并将祂的援助转向了蛮族。你攻击祂，祂就站在他们一边，因为胜利属于天主，并临于天主所引领的人。你难道不知道，}他继续说，\Quote{你将谁赶出了他们的教会，你又把这些教会交给了谁来管理？} 像\Person{特拉雅努斯}一样的将领\Person{阿林透斯}和\Person{维克托}证实了他所说的话的真实性，并恳求皇帝不要因基于事实的指责而发怒。\par

% structure: p0137 source=h2 target=subsection reason=relative-heading-rank; base=h2

\subsection{第三十一章  关于\Place{君士坦丁堡}的隐修士\Person{依撒格}与\Place{斯基泰}主教\Person{布雷塔尼奥}}

据传，在\Place{君士坦丁堡}独修的\Person{依撒格}，见\Person{瓦伦斯}率军出征，便大声喊道：\Quote{皇帝啊，你往哪里去？与天主作战，而不是以祂为盟友？正是天主亲自激起了蛮族攻击你，因为你激动了许多口舌亵渎祂，并将祂的敬拜者从他们的圣所中赶出。停止你的征伐，终止战争吧！将卓越的牧者归还给羊群，你便能不费吹灰之力赢得胜利；但若你执意作战而不这样做，\BibleQuote{你将亲身体验到踢刺棒是何等艰难}。你永远回不来，并且会毁灭你的军队。}于是皇帝怒气冲冲地回答：\Quote{我必回来，并要杀死你，以此惩罚你的谎言预言。}但\Person{依撒格}未被此威胁吓倒，喊道：\Quote{若我所说的话证明是假的，你就杀了我。}\par

\Person{布雷塔尼奥}，一位以各种德行著称的人，受托治理\Place{斯基泰}所有城市的教务，他满怀热忱，抗议教义的败坏以及皇帝对圣徒们的无法无天的攻击，用虔诚的\Person{达味}的言语喊道：\BibleQuote{我在君王前也论说你的法度，并不以为羞耻。}\par

% structure: p0140 source=h2 target=subsection reason=relative-heading-rank; base=h2

\subsection{第三十二章 \Person{瓦伦斯}对\Place{哥特人}的征讨及其不虔敬之罚}

\Person{瓦伦斯}却拒绝了这些优秀的顾问，派出军队投入战斗，而他自己则坐在一个小村庄等待胜利。他的军队抵挡不住\Place{蛮族}的进攻，转身逃跑，在逃跑中被一一杀死，\Place{罗马人}拼命奔逃，\Place{蛮族}则全力追赶。当\Person{瓦伦斯}听到失败的消息时，他试图躲藏在他所在的村庄里，但当\Place{蛮族}到来时，他们放火烧了那个地方，连同那不虔敬的敌人一起烧死。就这样，\Person{瓦伦斯}在此生为他所犯的错误付出了代价。\par

% structure: p0142 source=h2 target=subsection reason=relative-heading-rank; base=h2

\subsection{第三十三章 \Place{哥特人}如何被亚略异端所沾染}

对于那些不知道情况的人来说，或许值得解释一下\Place{哥特人}如何染上了亚略异端的瘟疫。在他们渡过\Place{多瑙河}并与\Person{瓦伦斯}讲和之后，臭名昭著的\Person{欧多克修斯}（他当时在场）向皇帝建议，说服\Place{哥特人}接受与他共融。他们确实早已领受了神圣知识的光芒，并在宗徒教义中受培育，但\Person{欧多克修斯}说：\Quote{但现在，}\Quote{意见一致会使和平更加稳固。} \Person{瓦伦斯}赞成这个建议，并向\Place{哥特人}的首领们提出教义上的协议，但他们回答说，他们不愿背弃祖传的教导。在所述时期，他们的主教\Person{乌尔菲拉}被他们绝对服从，他们视他的话为不可违背的法律。\Person{欧多克修斯}部分凭借其口才的魅力，部分凭借其提议中所设的贿赂，成功诱使他（\Person{乌尔菲拉}）说服\Place{蛮族}接受与皇帝共融。于是\Person{乌尔菲拉}以不同派系之间的争吵实际上是个人之间的竞争、并不涉及教义差异为借口，说服了他们。结果是直到今日，\Place{哥特人}断言\OriginalTerm{圣父}{Father}大于\OriginalTerm{圣子}{Son}，但拒绝称\OriginalTerm{圣子}{Son}为\OriginalTerm{受造物}{creature}，尽管他们与那些如此宣称的人共融。然而不能说他们完全放弃了祖传的教导，因为\Person{乌尔菲拉}在努力说服他们与\Person{欧多克修斯}和\Person{瓦伦斯}共融时，否认教义有任何差别，并说分歧只是出于虚妄的争吵。\par

\SourceRecord{Translated by Blomfield Jackson.；Nicene and Post-Nicene Fathers, Second Series；Vol. 3.；Edited by Philip Schaff and Henry Wace.；Buffalo, NY: Christian Literature Publishing Co.,；1892.；<http://www.newadvent.org/fathers/27024.htm>.}

% structure: p0001 source=h1 target=section reason=page-title

\section{教会史（第五卷）}

% structure: p0002 source=h2 target=subsection reason=relative-heading-rank; base=h2

\subsection{第一章 论皇帝\Person{格拉提安}的虔诚}

主天主如何容忍那些向祂发怒的人，并惩罚那些滥用祂耐心的人，这事由\Person{瓦伦斯}的作为和命运清楚教导。因为慈爱的主使用仁慈和公义如同砝码和天平；每当祂看见任何人因罪恶深重而逾越仁慈的界限时，祂便以公正的惩罚阻止他走向更极端的地步。\par

如今，\Person{格拉提安}，\Person{瓦伦提尼安}之子、\Person{瓦伦斯}之侄，取得了整个罗马帝国。他父亲去世后已继承欧洲的权杖，在父亲生前已与他共治。\Person{瓦伦斯}死后无嗣，他又获得了亚洲以及利比亚的部分地区。\par

% structure: p0005 source=h2 target=subsection reason=relative-heading-rank; base=h2

\subsection{第二章 论主教的归回}

皇帝立即明确表明他遵循真宗教，并向万有的主献上他王国的初果，颁布谕令，命令被放逐的牧者返回并回到他们的羊群，并命令将圣殿交付给与\Person{达玛苏}共融的会众。\par

这位\Person{达玛苏}，\Person{利贝流}在罗马教区的继任者，是一位备受赞颂的圣德之人，并自愿在言行上成为宗徒教义的护卫者。为实施其谕令，\Person{格拉提安}派遣了当时非常著名的将军\Person{萨波尔}，命令他将宣讲\Person{亚略}亵渎言论的人如野兽般逐出神圣的羊栈，并使优秀的牧者回归天主的羊群。\par

在每一处，这都毫无争议地执行了，唯独在东方首都\Place{安提约基雅}除外，那里燃起了一场纷争，我将描述之。\par

% structure: p0009 source=h2 target=subsection reason=relative-heading-rank; base=h2

\subsection{第三章 论由\Person{保利诺}引起的分裂、由\Place{劳底嘉}的\Person{阿波利拿里}引起的革新以及\Person{梅利提乌}的哲学}

已经叙述过宗徒教义的护卫者如何分裂为两派：在针对伟大的\Person{尤斯塔提乌}的阴谋之后，一派如何因憎恶\Person{亚略}的可憎之事而自行聚集，以\Person{保利诺}为主教，而另一派在\Person{尤佐伊}被祝圣后，如何与杰出的\Person{梅利提乌}一起与不虔敬者分离，经历了前述的危险，并受到\Person{梅利提乌}给予他们的智慧训导。除此以外，\Place{劳底嘉}的\Person{阿波利拿里}自命为第三派的领袖，虽然他披着虔诚的面具，看似维护宗徒教义，但很快就被发现是公开的敌人。关于天主性，他使用了不正确的论证，并提出了一些等级尊威的观念。他还胆敢使降生成人的奥迹不完全，断言那受托管理身体的理性灵魂未能获得救恩。因为根据他的论证，天主圣言没有取这个灵魂，因此既没有赐予它医治的恩赐，也没有分给它祂的部分尊威。这样一来，有形的肉身被无形的能力所朝拜，而那按天主的肖像造的灵魂却仍然留在罪的耻辱之下。他还口吐许多其他错误，他那失足盲目理智。有时他甚至准备承认肉身取自童贞圣母，有时又主张它与天主圣言一同从天降下，还有时又说祂成了血肉，没有从我们取任何东西。其他无稽之谈和琐碎之事，我认为无需重复，他将其与天主的福音应许混在一起。通过这类论证，他不仅用危险的教义充满自己的朋友，甚至也传给一些我们当中的人。随着时间的推移，当他们看到自己的微不足道，并看见教会的荣光时，除少数人外，所有的人都聚集到教会的共融中。但他们没有完全抛弃先前的谬误，并用它感染了许多端正的人。这就是教会中关于肉身与神性一性、将独生子的受难归于神性以及其他导致平信徒与司铎之间分歧的论点增长的起源。但这些属于较晚的时期。在我所说的那个时候，当将军\Person{萨波尔}到达并出示了皇帝谕令时，\Person{保利诺}声称他赞同\Person{达玛苏}，而\Person{阿波利拿里}隐藏他的谬误，也如此说。神圣的\Person{梅利提乌}则没有表示，容忍了他们的争论。\Person{弗拉维安}以其智慧享负盛誉，那时仍在司铎之列，起初对\Person{保利诺}在官员面前说：\Quote{我亲爱的朋友，如果你接受与\Person{达玛苏}共融，请清楚向我们指出教义如何一致，因为他虽然承认天主圣三的一性，却公开宣讲三个本质。你反而否认本质的三性。那么请向我们说明这些教义如何和谐，并按谕令接受教会的管理。}如此用论证使\Person{保利诺}无言以对后，他转向\Person{阿波利拿里}说：\Quote{朋友，我很惊讶，你竟对真理发动如此猛烈的战争，而你却清楚知道可敬的\Person{达玛苏}如何坚持我们的本性被天主圣言完全取去；但你却固执地持相反意见，因为你剥夺了我们理智的救恩。如果我们对你的这些指控是假的，现在就否认你所创的新奇说法；接受\Person{达玛苏}的教导，并接受圣所的管理。}\par

就这样，\Person{弗拉维安}以他极大的智慧，用他真实的推理阻止了他们的大胆言论。\par

\Person{梅肋提乌斯}，在所有人中最温良，这样和蔼而亲切地对\Person{保理诺}说：\Quote{羊群之主已将这群羊交在我手里；你已接受其余羊群的托管；我们的小羊在真宗教中彼此共融。因此，我亲爱的朋友，让我们合并羊群；让我们停止关于管理它们的争论，一同牧放，共同照料。倘若首座是纷争的起因，我要设法除去那纷争。我将把圣福音放在首座上；我们分坐其两侧；若我先离世，我的朋友，你将独掌羊群的领导权。若你先我而去，我将尽力照料羊群。}这位神圣的\Person{梅肋提乌斯}如此温和而亲切地说话。\Person{保理诺}没有同意。官员对所说的话作出裁决，将各教会交给了伟大的\Person{梅肋提乌斯}。\Person{保理诺}继续领导最初分裂出去的羊群。\par

% structure: p0013 source=h2 target=subsection reason=relative-heading-rank; base=h2

\subsection{第四章 论撒摩沙塔主教欧瑟比}

\Emph{阿波林}在未能获得教会治理权之后，继续公开宣讲他的新奇教义，并自命为异端之首。他大部分时间住在\Place{老底嘉}；但在\Place{安提约基雅}，他已祝圣了\Person{维塔利乌斯}，此人品德卓越，自幼受宗徒教义培育，但后来染上异端。我前述的\Person{狄奥多禄}，在大风暴中挽救了即将沉没的教会之船，已被神圣的\Person{梅肋提乌斯}任命为\Place{塔尔索}主教，并接受了对\Place{基里基雅}人的牧职。\Person{梅肋提乌斯}将\Place{阿帕美亚}教区托付给\Person{若望}，此人出身显赫，但因其自身优秀品质而比祖辈更卓越，他在教导和生活之美上同样出众。在暴风雨时期，他在可敬的\Person{斯德望}支持下，驾驶着同道信友的集会；但后者被神圣的\Person{梅肋提乌斯}调往另一处战场，因为得知\Place{革尔玛尼基雅}已被欧道克修斯派的瘟疫污染，他被派往那里作为医生以阻挡疾病，他既受过完整的异教教育，又受到神圣教义的培育。他没有辜负人们对他的期望，因为他以灵性教导的力量将豺狼变成了绵羊。\par

伟大的\Person{欧瑟比}从流亡归来后，祝圣了\Person{阿卡丘}——他在\Place{贝洛雅}声名显赫，并在\Place{希拉波里}祝圣了\Person{德奥多托}——其苦修生活至今为人传颂。此外，\Person{欧瑟比}被任命为\Place{卡尔基斯}教区主教，\Person{依西多禄}被任命为我们自己的\Place{居鲁士}城主教；两人都是可敬佩的人，在神圣热忱上出类拔萃。\par

据报道，\Person{梅肋提乌斯}也祝圣了\Person{欧洛纠}为\Place{厄德萨}的牧者——虔诚的\Person{巴尔塞斯}已在那里去世；\Person{欧洛纠}是宗徒教义的著名捍卫者，曾与\Person{普罗托革诺}一同被流放到\Place{安提诺}。\Person{欧洛纠}将他在艰苦服务中的同伴\Person{普罗托革诺}托付给\Place{卡雷}的牧职，作为患病城市的治疗医生。\par

最后，神圣的\Person{欧瑟比}祝圣了\Person{玛里斯}为\Place{多利刻}主教——当时那座小城被亚略派的瘟疫感染。伟大的\Person{欧瑟比}来到\Place{多利刻}，意图将这位极其可敬、以各种美德著称的\Person{玛里斯}安置在主教授席上。当他进城时，一名完全被亚略派瘟疫感染的女人从屋顶掉下一片瓦，砸中他的头部，使他受伤，不久后他便离世进入更好的生命。临死时，他嘱咐旁观者不要对肇事妇女施加丝毫惩罚，并以誓言约束他们服从。如此，他效法了自己的主，他被钉十字架的人说：\BibleQuote{父啊，宽赦他们吧，因为他们不知道他们做的是什么。}\par

同样，他也效法了他的同仆\Person{斯德望}的榜样，\Person{斯德望}在石头如暴风雨般砸向他后，大声喊道：\BibleQuote{主，不要将这罪归于他们。}伟大的\Person{欧瑟比}在经过许多不同的斗争后如此去世。他曾在\Place{色雷斯}逃脱了蛮族，但未能逃脱不敬虔异端者的暴力，并借此赢得了殉道者的冠冕。\par

这些事件发生在主教们回归之后。此时，\Person{格拉提安}得知\Place{色雷斯}正被烧死\Person{瓦伦斯}的蛮族蹂躏，于是他离开\Place{意大利}，前往\Place{潘诺尼亚}。\par

% structure: p0020 source=h2 target=subsection reason=relative-heading-rank; base=h2

\subsection{第五章 论德奥多西的征讨}

此时，\Person{德奥多西}因祖先的辉煌及自身的勇气而享有盛誉。正因如此，他时常遭受对手的嫉妒，便住在\Place{西班牙}——他出生并成长之地。皇帝不知采取何种措施，因为蛮族因胜利而自大，既几乎无敌，也看似无敌，于是形成一个想法：任命\Person{德奥多西}为最高统帅将是摆脱困境的出路。因此，他立即从\Place{西班牙}召来\Person{德奥多西}，任命他为全军统帅，并派遣他率领集结的军队。\par

狄奥多西因信德而信心倍增，毅然进军。进入色雷斯后，他看见蛮族前来迎战，便布阵御敌。两军交锋，敌人抵挡不住攻击，阵脚大乱。蛮族溃败逃跑，得胜者全速追击。蛮族遭大量屠杀，不仅被罗马人所杀，也自相残杀。大部分蛮族倒下后，少数逃脱者渡过多瑙河。这位伟大的统帅将所率军队分驻邻近城镇，随即骑马疾驰前往格拉提安皇帝处，亲自报告自己的胜利。就连皇帝本人也对这一事件感到惊愕，他所听到的消息似乎难以置信，而一些心怀嫉妒的人则散布谣言说他已逃跑且失去了军队。他唯一的回应是请反对者派人查明蛮族死者的数目，\Quote{因为，}他说，\Quote{仅从他们的战利品就很容易得知其数目。}此话一出，皇帝便让步，派遣官员去调查并汇报战况。\par

% structure: p0023 source=h2 target=subsection reason=relative-heading-rank; base=h2

\subsection{第六章　狄奥多西的统治与他的梦}

这位伟大的统帅留下，随后看见了一个奇妙的异象，这异象清晰地由宇宙的天主亲自显示给他。在异象中，他仿佛看见安提约基雅教会首领、神圣的默肋提乌斯为他披上皇袍，并将皇冠戴在他头上。他看见异象那夜之后的清晨，他将此事告诉了一位密友，这位朋友指出这梦清晰明了，毫无隐晦或含糊之处。\par

最多不过几天之后，派去调查战事的官员返回，报告说有大批蛮族被射杀。\par

于是皇帝确信他选择狄奥多西担任统帅是明智之举，便任命他为皇帝，并将瓦伦斯管辖的那部分领土交给他治理。\par

随后，格拉提安前往意大利，并派遣狄奥多西前往他所负责的地区。狄奥多西一登上帝位，便首先关注教会的和谐，下令他所辖领域的主教们火速前往君士坦丁堡。帝国的那部分地区是当时唯一受亚略派瘟疫感染的地方，因为西方已幸免于此。这是因为君士坦丁大帝的长子君士坦丁和幼子康斯坦斯完整地保存了他们父亲的信仰，而西方的皇帝瓦伦提尼安也保持了纯正的真宗教。\par

% structure: p0028 source=h2 target=subsection reason=relative-heading-rank; base=h2

\subsection{第七章　亚略派的著名领袖}

帝国东部从多个渠道受到感染。埃及亚历山大城的司铎亚略在此滋生亵渎之论。巴勒斯坦的欧瑟比、帕特罗斐鲁斯和埃提乌斯，腓尼基的保利努斯和格雷戈里乌斯，劳底嘉的狄奥多图斯及其继任者乔治乌斯，以及之后的西里西亚的亚大纳削和纳西苏斯，滋养了这些邪恶播下的种子。比提尼亚的欧瑟比和狄奥格尼斯，厄弗所的梅诺方图斯，佩林图斯的狄奥多鲁斯和加采东的马里斯，以及色雷斯其他一些以恶行闻名的人，长期以来一直浇灌并培育着莠草的庄稼。这些坏农夫得到了君士坦提乌斯的漠不关心和瓦伦斯的恶意的帮助。\par

因此，皇帝只召集了他自己帝国的主教们到君士坦丁堡会合。他们到来，共计一百五十人。狄奥多西禁止任何人告诉他哪一位是伟大的默肋提乌斯，因为他希望凭他的梦认出这位主教。全体主教进入皇宫，然后，没有理会其余所有人，狄奥多西跑到伟大的默肋提乌斯面前，像一个爱父亲的孩子，长时间以孝子般的喜悦注视着他，然后伸出双臂拥抱他，亲吻他的眼睛、嘴唇、胸膛、头以及给他加冕的手。然后他告诉默肋提乌斯他的梦。其余所有主教也都受到礼貌的欢迎，并被教导要以父亲般的身份认真讨论提交给他们的问题。\par

% structure: p0031 source=h2 target=subsection reason=relative-heading-rank; base=h2

\subsection{第八章　在君士坦丁堡召开的公会议}

那时，纳齐盎的新近牧者正住在君士坦丁堡，不断抵制亚略派的亵渎言论，以福音的教导浇灌圣洁的子民，将迷失在羊群之外的游荡者找回，使他们远离毒草。因此，他曾使一小群羊变成一大群。当神圣的默肋提乌斯见到他时，深知制定教规者的目的——为了防止野心勃勃的企图，禁止迁调主教——便确认格雷戈里为君士坦丁堡的主教。不久之后，神圣的默肋提乌斯安息于那无痛苦的生命，所有伟大演说家的葬礼赞词都向他表示敬意。\par

亚历山大城的主教提摩泰，他是继亚大纳削之后担任宗主教之位的伯多禄的继承者，祝圣了马克西穆斯以取代可敬的格雷戈里——一位犬儒学派的哲学家，他刚剃去犬儒式的长发，并被阿波利纳里的浮夸修辞所迷惑。但这种荒谬是与会主教们所不能容忍的——他们是可敬的人，充满神圣的热忱和智慧，如伟大的巴西略的继承者赫拉第乌斯，巴西略的兄弟格雷戈里和伯多禄，来自吕考尼亚的安斐洛基乌斯，来自彼西底的奥普提姆斯，来自西里西亚的狄奥多若。\par

出席公会议的还有劳底嘉的佩拉吉乌斯，厄德萨的欧洛吉乌斯，阿卡基乌斯，我们的依西多洛，耶路撒冷的济利禄，巴勒斯坦凯撒勒雅的格辣西亚，他因学识和圣德而著称，以及许多其他的德性斗士。\par

凡我所提名者，皆与埃及人分离，与伟大的\Person{格列高利}{Gregory}一同举行神圣礼仪。但他本人恳求他们——他们聚在一起原为促进和谐——将一切对个人的过错问题置于彼此和好的目标之下。他说：\Quote{因为，我将脱离许多忧虑，重新过我珍视的宁静生活；而你们，在长期痛苦的战争之后，将获得渴望已久的平安。还有什么比刚刚逃脱敌人武器的人，反而消耗自己的力量互相伤害更荒谬的呢？这样做，我们只会成为对手的笑柄。那么，请找一位有见识、能承担重任并妥善处理事务的贤人，立他为主教。}这些优秀的牧者受此劝勉而动，便委任了一位出身高贵、以各种美德和显赫家世闻名的人，即\Person{涅克塔里}{Nectarius}，作为那座伟大城市的主教。\Person{马克西莫}{Maximus}因参与\Person{阿波利拿里}{Apollinarius}的疯狂，被剥夺主教职权，予以摈弃。接着，他们颁布了关于教会良好治理的法规，并公布了在\Place{尼西亚}{Nicæa}所定信仰的确认书。然后，他们各自返回本国。次年夏季，他们中多数人再次聚集于同一城市，因教会需要再度被召；他们收到了西方主教们的公函，邀请他们前往罗马，那里正召集一次大公会议。但他们恳请免予远行外地，认为这样做无益。不过他们写信指出了教会所遭遇的风暴，并暗示西方主教对此处理疏忽。他们还在信中附上了\OriginalTerm{宗徒道理}{apostolic doctrine}的摘要；但其言辞的大胆与智慧，将由信函本身更清楚地显明。\par

% structure: p0036 source=h2 target=subsection reason=relative-heading-rank; base=h2

\subsection{第九章　从君士坦丁堡公会议发出的公函}

致可敬的诸位主内尊长、我们的兄弟和同工\Person{达玛稣}{Damasus}、\Person{盎博罗削}{Ambrosius}、\Person{布里东}{Britton}、\Person{瓦勒利安}{Valerianus}、\Person{阿朔里}{Ascholius}、\Person{阿内米}{Anemius}、\Person{巴西略}{Basilius}以及在伟大的\Place{罗马}{Rome}城聚集的其他圣主教们：在伟大的\Place{君士坦丁堡}{Constantinople}城聚集的正统主教圣公会议，在主内问安。\par

若要列举我们因亚略派权势而遭受的一切苦难，并试图向尊长们报告，仿佛你们尚不熟悉——这可能显得多余。因为我们并不认为你们的虔诚会如此不重视我们所遭遇的事，以至于需要了解这些必然引起你们同情的事。况且，袭击我们的风暴也并非因其微小而可被忽视。我们的迫害昨日才发生，其声音仍在那些受难者以及因爱心而感同身受者耳中回响。可以说，就在一两日前，有些人在异乡挣脱锁链，经历重重苦难返回自己的教会；另一些人的遗骸——那些流亡中死去的人——被运回家乡；还有一些人，即使在流亡归来后，仍见异端者的激情沸腾，如蒙福的\Person{斯德望}{Stephen}般被他们用石头击毙，在自己家乡遭遇比异乡更悲惨的命运。另有一些人受尽各种酷刑，至今躯体上仍带着伤痕和基督的印记。\par

谁能尽述罚款、褫夺公权、个人财产充公、阴谋、暴行、监禁的事呢？诚然，在我们身上实现了一切数不尽的患难，或许因为我们为自己的罪付出代价，又或许因为仁慈的天主借众多苦难来试炼我们。为此，我们感谢天主，他藉这样的磨难训练了他的仆人，并按其丰厚的怜悯再次使我们得享安息。我们确实需要长久的闲暇、时间和辛劳来重建教会，好使她能如医生医治久病之身，循序渐进地驱除疾病，恢复古代真宗教的健康。诚然，总体而言，我们似乎已从迫害的暴力中得救，并刚刚收复长久以来被\OriginalTerm{异端者}{heretics}掠夺的教会。但豺狼仍烦扰我们——它们虽被逐出羊圈，却在林间各处蹂躏羊群，胆敢举行敌对集会，在百姓中煽动叛乱，凡能损害教会的事无所不为。\par

所以，正如我们已说过的，我们更加必须长久劳苦。然而，既然你们通过我们最虔诚的皇帝的来信，邀请我们加入你们正在\Place{罗马}因天主准许而召集的会议（如同我们是你们自己的肢体），显出了你们的兄弟之爱，目的乃是：既然当时只有我们单独遭受迫害，如今我们的皇帝在真宗教上与我们已经一致，你们就不应再与我们分离统治，而是我们——用宗徒的话说——应与你们一同为王，我们的祈祷是：若有可能，全体一起离开我们的教会，宁可满足我们渴望见到你们的愿望，而不顾及教会的需要。因为谁将给我们鸽子般的翅膀，我们就能飞去安息？但这条道路看来会使正在恢复中的教会毫无保护，而且对我们多数人来说是不可能的，因为按照在\Place{阿奎莱亚}会议之后一年前从你们的圣座寄给最虔诚的皇帝\Person{德奥多西}的信件，我们已前往\Place{君士坦丁堡}，只装备了足以旅行到\Place{君士坦丁堡}的装备，并带来了各省留守主教们对此会议的同意。我们未曾期待任何更长的旅程，在抵达\Place{君士坦丁堡}之前也未曾听到任何消息。除此之外，由于指定时间的限制，既无法为更长的旅程做准备，也无法告知我们共融的主教们并取得他们的同意，因此前往\Place{罗马}的旅程对多数人来说是不可能的。因此，我们在当前情况下采取了次优方案，既为更好地管理教会，也为显明我们对你们的爱，大力恳请我们最可敬、可敬的同僚兼兄弟主教\Person{西里亚库斯}、\Person{欧瑟比乌斯}和\Person{普里斯西安}同意前往你们那里。\par

通过他们，我们愿意表明：我们的意向完全是为和平，以合一为唯一目的，并且我们充满对正信的热忱。因为，我们无论遭受异端者带来的迫害、苦难、皇帝的威胁、君王的暴行，或任何其他考验，都是为了福音信德而承受的。这信德由在\Place{彼提尼雅}的\Place{尼西亚}的三百一十八位教父所批准。这信德应当为你们、为我们、为所有不曲解真信德之言的人所足够；因为它是古老的信德，是我们圣洗的信德，是教导我们相信因圣父、圣子及圣神之名的信德。\par

根据这信仰，圣父、圣子、圣神有一个天主性、能力和本体；尊严相等，威严相等，存在于三个完美的本体和三个完美的位格中。如此，既没有空间容纳\Person{撒柏略}的异端（因混淆本体或消灭个别性），也废除了欧诺米派、亚略派和反圣神派的亵渎，他们将本体、本性、天主性分割，并在未受造、同体、同永的天主圣三上引入一个后来的、受造的、不同本体的本性。此外，我们毫无扭曲地保存了主的降生成人的教义，持守这传承：肉身的安排既非无灵魂、亦非无理智、也非不完全；并且深知天主的圣言在万世之前已是完全的，在末世为了我们的救恩成为完全的人。\par

以上就是我们所无畏而坦诚宣讲的教义概要；关于这些，如果你们愿意阅读\Place{安提约基雅}会议的记录，以及去年在\Place{君士坦丁堡}举行的大公会议所发布的记录，你们将能进一步满意，因为在其中我们更详细地陈述了我们信德的宣认，并附加了对新派人士最近所记载的异端的绝罚。\par

\Quote{至于个别教会的特殊管理，正如你们所知，有一古老惯例，由在\Place{尼西亚}的圣教父们所颁布的规定所确认：在每个省，由本省主教们并经他们同意，连同邻近主教，应按需要举行祝圣。在遵守这些惯例时，请注意其他教会已由我们管理，并且最著名教会的司铎已公开任命。因此，对于\Place{君士坦丁堡}那新建的（若可以这样说）教会——我们最近藉天主的仁慈从异端的亵渎中，如从狮口夺回——我们在大公会议面前，在大家同意下，在最虔诚的皇帝\Person{德奥多西}面前，并经全体圣职人员和全城同意，祝圣了可敬、最虔诚的\Person{奈克塔留}为主教。至于\Place{叙利亚}最古老、真正宗徒的教会——那里首先赐予了基督徒这高贵名称——本省和东方教区的主教们集会，并按教会法规祝圣了可敬、最虔诚的\Person{弗拉维亚努斯}为主教，得全体教会同意，他们仿佛异口同声表达对他们的敬意。这合法的祝圣也获得了大公会议的认可。至于\Place{耶路撒冷}教会——众教会之母——我们宣告可敬、最虔诚的\Person{济利禄}是主教，他不久前由本省主教们按法规祝圣，并在多处与亚略派进行了良好战斗。我们恳求你们的圣座，因着属灵之爱的介入和敬畏天主之情的感动，使人性感情受约束，并将教会的建立看得比个人恩惠或宠爱更重要，从而为这些由我们正确、合法地处理的事而欢喜。这样，既然我们之间在信仰上一致，基督徒的友爱已建立，我们将不再使用宗徒所谴责的那话：\InnerQuote{我是属保禄的，我是属阿颇罗的，我是属刻法的}，而全体都显出是属基督的——他在我们内并未分开——藉天主的恩宠，我们将使教会身体保持不裂，并勇敢地站立在主的审判台前。}\par

这些事是他们针对亚略、艾提乌斯和欧诺米的狂妄而写的；此外也针对撒伯流、佛提努、马塞鲁、萨莫萨塔的保禄和马其顿尼。同样，他们明确谴责了亚波利拿里的新奇说法，其中写道：\Quote{我们保存了主降生成人的道理，持守传统，即那肉身的安排既不是无灵魂的，也不是无理性的，也不是不完全的。}\par

% structure: p0046 source=h2 target=subsection reason=relative-heading-rank; base=h2

\subsection{第十章——罗马主教达玛苏斯针对亚波利拿里和提摩泰的会议书信。}

当最可赞颂的达玛苏斯听到这一异端兴起时，他宣布不仅谴责亚波利拿里，也谴责他的追随者提摩泰。我认为有必要将他就此向东方帝国的主教们传达的书信载入我的历史。\par

\Emph{罗马主教达玛苏斯的书信}\par

最可敬的儿子们：既然你们的爱德向宗座献上了应有的尊敬，你们也应以同样的丰盛为自己领受。因为即使在圣教会中——那圣宗徒曾坐席并教导我们应当如何管理托付给我们的船舵——我们仍然承认自己不堪当这尊荣，但我们正因此尽己所能，竭力追求那真福的荣耀。你们要知道，我们已谴责了不洁的提摩泰，异端者亚波利拿里的弟子，连同他那不敬神的道理，并确信今后他的余毒不再有任何分量。但若那古蛇，虽曾被击打一次两次，仍为自身的毁灭而复苏——他虽存在于教会之外，却从不停止以他的致命毒液试图倾覆某些不忠信的人——你们要像躲避瘟疫一样躲避他，常铭记宗徒的信仰，即教父们在尼西亚所书写的那一信仰；你们当坚定地站在稳固的基础上，在信仰中岿然不动，从今以后不让你们的圣职人员或平信徒听从虚妄的言语和无益的争论，因为我们已给出一个规范，凡自称为基督徒者应持守它，就是宗徒们所传授的规范，如圣保禄所说：\Quote{若有人给你们宣讲福音，与你们所接受的不同，当受弃绝。}因为基督，天主子、我们的主，藉自己的苦难赐给了人类丰富的救恩，为将整个陷于罪恶的人从一切罪恶中释放出来。若有人说基督缺少人性或天主性，他充满魔鬼的灵，并自称是地狱之子。\par

\Quote{那么你们为何再次向我要求谴责提摩泰？看，在宗座的判决下，在亚历山大主教伯多禄面前，他和他的老师亚波利拿里一同被谴责了，他们将在审判之日承受当得的刑罚和折磨。但若他凭着自己的宣信改变了在基督内的真希望，而说服了一些较不稳固的人，似乎怀有某种希望，那么凡故意反抗教会秩序的人，也要同他一起灭亡。愿天主保守你们身心安康，最可敬的儿子们。}\par

在伟大的罗马集会的主教们也写了其他反对其他异端的文字，我认为有必要将它们插入我的历史。\par

% structure: p0052 source=h2 target=subsection reason=relative-heading-rank; base=h2

\subsection{第十一章——教宗达玛苏斯在塞萨洛尼基时写给马其顿的保禄主教的公教信仰宣认}

尼西亚大公会议之后兴起了这一谬误。有些人竟以亵渎的口舌声称圣神是由圣子受造的。我们因此弃绝那些不公开宣认圣神与圣父及圣子是同一性体、同一权能的人。同样，我们弃绝那些追随撒伯流的谬误，说圣父与圣子为同一者的人。我们弃绝亚略和欧诺米，他们以同样不敬神但措辞不同的方式主张圣子和圣神是受造物。我们弃绝马其顿派，他们出自亚略的根，虽改变了名称却未改变不敬神。我们弃绝佛提努，他重振了厄彼昂的异端，宣认我们的主耶稣基督只是玛利亚所生。我们弃绝那些主张有两位圣子的——一位在万世之前，另一位在从玛利亚取得肉体之后。我们也弃绝所有主张天主的圣言在人的肉体中运行时代替了理性灵魂的人。因为这位天主的圣言自己并非代替一个理性的、理智的灵魂而存在于他自己的身体内，而是取了我们无罪且有理性的灵魂并拯救了它。我们也弃绝那些说天主的圣言因伸展和收缩而与圣父分离，并亵渎地主张他没有本质存有或注定要死的人。\par

凡由一教会转到其他教会的人，我们暂视他们与我们的共融隔绝，直到他们返回最初被祝圣的城市为止。\par

若有人在他处之后，有人被祝圣接替他的位置，则那放弃自己城市的人应被视为被剥夺主教职，直到他的继任者安息于主内为止。\par

若有人否认圣父是永恒的，圣子是永恒的，圣神是永恒的，此人应受弃绝。\par

若有人否认圣子是由圣父所生，即出自他的天主性体，此人应受弃绝。\par

若有人否认天主子是真天主，全能、全知，与圣父同等，此人应受弃绝。\par

若有人说天主子在世生活于肉体内时，没有在天上并与圣父同在，此人应受弃绝。\par

若有人声称在十字架的苦难中，天主子是以天主性承受其痛苦，而非以他所取的仆人形像中的理性灵魂和肉体，如圣经所说，此人应受弃绝。\par

若有人否认天主的圣言在肉体中受苦，在肉体中尝了死，是死者中的首生者，因为圣子是生命并赋予生命者，此人应受弃绝。\par

若有人否认他坐在圣父的右边，带着他所取的肉体，并要以这肉体审判生者死者，此人应受弃绝。\par

若有任何人否认圣神确实并绝对出自圣父，并否认圣子出于天主性体、是天主出于天主，愿他受\OriginalTerm{弃绝}{anathema}。\par

若有任何人否认圣神是全能的、全知的和无所不在的，如同圣父的圣子一样，愿他受\OriginalTerm{弃绝}{anathema}。\par

若有任何人说圣神是受造物，或是藉圣子造成的，愿他受\OriginalTerm{弃绝}{anathema}。\par

若有任何人否认圣父藉那降生成人的圣子和圣神创造了所有有形无形的万物，愿他受\OriginalTerm{弃绝}{anathema}。\par

若有任何人否认圣父、圣子及圣神有一个\OriginalTerm{天主性}{Godhead}、一个权能、一个主权、一个光荣、一个主道、一个国度、一个旨意和一个真理，愿他受\OriginalTerm{弃绝}{anathema}。\par

若有任何人否认圣父、圣子及圣神是三个真实的\OriginalTerm{位格}{persons}，永生永存，包含一切有形无形的万物，全能，审判万物，赐予万物生命，创造并保存万物，愿他受\OriginalTerm{弃绝}{anathema}。\par

若有任何人否认圣神当受一切受造物的敬拜，如同圣子和圣父一样，愿他受\OriginalTerm{弃绝}{anathema}。\par

任何人若对圣父和圣子持有正确的想法，却对圣神没有正确的持守，愿他受\OriginalTerm{弃绝}{anathema}，因他是\OriginalTerm{异端者}{heretic}；因为所有对天主圣子和圣神没有正确想法的异端者，都被证明陷于犹太人和外邦人的不信之中。若有人分开\OriginalTerm{天主性}{Godhead}，说圣父是一位天主，圣子是一位天主，圣神是一位天主，并坚持它们被称为诸神而非天主，因我们相信并认识的圣父、圣子及圣神的唯一天主性和主权——即一个天主在三个实质中——或撤去圣子和圣神，以致暗示只有圣父被称为天主并被相信为独一的天主，愿他受\OriginalTerm{弃绝}{anathema}。\par

因为\Quote{诸神}之名是由天主赐予天使和所有圣人的；但论到圣父、圣子及圣神，由于它们唯一而同等的天主性，所称谓并宣布的不是\Quote{诸神}之名，而是\Quote{我们的天主}之名，使我们相信我们是因圣父、圣子及圣神之名受洗，而非因总领天使或天使之名，如同异端者、犹太人或愚昧的外邦人那样。\par

这就是基督徒的救恩：相信\OriginalTerm{天主圣三}{Trinity}，即相信圣父、圣子及圣神，并因这唯一的天主性、权能、神性和实体受洗，而寄望于祂。\par

这些事发生在\Person{格拉提安}在世期间。\par

% structure: p0074 source=h2 target=subsection reason=relative-heading-rank; base=h2

\subsection{第十二章　论\Person{格拉提安}之死与\Person{马克西穆斯}的统治}

\Person{格拉提安}在战争胜利和明智审慎的治理中，因阴谋而结束了自己的生命。他没有留下儿子继承帝国，只有一个与父亲同名的兄弟，即尚是少年的\Person{瓦伦提尼安}。于是\Person{马克西穆斯}因轻视\Person{瓦伦提尼安}的年轻，夺取了西方的皇位。\par

% structure: p0076 source=h2 target=subsection reason=relative-heading-rank; base=h2

\subsection{第十三章　论\Person{瓦伦提尼安}之妻\Person{尤斯蒂娜}及其反对\Person{盎博罗削}的阴谋}

此时，伟大的\Person{瓦伦提尼安}之妻、年轻皇帝之母\Person{尤斯蒂娜}，将她早先领受的\OriginalTerm{亚流}{Arian}教导的种子告知了她的儿子。她深知丈夫信仰的热忱，故在丈夫生前一直竭力隐瞒自己的观点；但看到儿子性情温和顺服，她便鼓起勇气，将那欺骗性的教义提了出来。那少年以为母亲的劝告是明智有益的，因为天性如此布置了诱饵，使他看不到下面致命的鱼钩。他首先就此事与\Person{盎博罗削}沟通，心想若能说服主教，便能轻易使其余的人就范。然而，\Person{盎博罗削}却努力提醒他父亲的虔诚，劝他不可破坏他所继承的遗产。他也向他解释了两种教义的区别：一种是符合主和宗徒的教导，另一种则完全与之相悖，并与精神律法之法规为敌。\par

那年轻人，如同年轻人常有的，再加上他那本身即为欺骗受害者的母亲的挑动，不仅没有接受所提出的论据，反而失去了耐心，一怒之下，准备用军团兵和盾牌手的队伍包围通往教堂的通道。然而，当他得知这位杰出的斗士丝毫不为他的举动所惊——因为\Person{盎博罗削}视它们如同某些人用来吓唬婴儿的鬼怪和妖精——他极其愤怒，公开命令他离开教堂。\Person{盎博罗削}说：\Quote{我不会这样做的。我不愿将羊栈交给豺狼，也不愿将天主的殿宇出卖给亵渎者。你若想杀我，就在这里用你的剑或矛刺入我的身体。我将欢迎这样的死亡。}\par

% structure: p0079 source=h2 target=subsection reason=relative-heading-rank; base=h2

\subsection{第十四章　论僭主\Person{马克西穆斯}给\Person{瓦伦提尼安}的信息}

过了一段时间，\Person{马克西穆斯}得知对真理的大声宣告者所进行的攻击，便派遣文书给\Person{瓦伦提尼安}，责令他停止对真正宗教的战争，并劝诫他不要放弃父亲的信仰。若他的建议被忽视，他还威胁要发动战争，且以实际行动证实了他信中所写——他集结军队，向\Person{瓦伦提尼安}当时所在的\Place{米兰}进军。当\Person{瓦伦提尼安}听闻他临近时，逃往了\Place{伊利里库姆}。他已从经验中学会了听从母亲建议所带来的好处。\par

% structure: p0081 source=h2 target=subsection reason=relative-heading-rank; base=h2

\subsection{第十五章　论皇帝\Person{狄奥多西}就此事所写的信函}

当杰出的皇帝\Person{狄奥多西}听闻皇帝的所作所为以及僭主\Person{马克西穆斯}写给他的信后，他写信给那逃亡的少年，内容如下：\Quote{你不可惊讶为何恐慌临于你而胜利归给你的敌人；因你是在对抗虔诚，而他则在为虔诚而战。你抛弃了它，如今赤身逃跑；他身披它的全副武装，正在制服赤身露体的你，因为那赐予我们真正宗教律法者，总是站在它一边。}\par

\Person{德奥多西}尚在远方时如此写道；及至他听闻\Person{瓦伦提尼安}逃走并前去救援，见他流亡，投奔自己帝国，他首先想到的是救助其灵魂，驱除侵入的不敬神瘟疫，使之归向父辈的真宗教。于是他命他安心，便进军讨伐篡位者。他未流血便将帝国归还少年，并诛杀了\Person{马克西穆斯}。因他若不为盟友之死向凶手复仇，便觉自己理亏，且违背了与\Person{格拉提安}所立之约。\par

% structure: p0084 source=h2 target=subsection reason=relative-heading-rank; base=h2

\subsection{第十六章　论\Place{依科尼雍}主教\Person{安斐洛基}}

皇帝归来时，那位我常提及的可敬的\Person{安斐洛基}前来恳求，将亚略派聚会从各城驱逐。皇帝认为请求过于严厉，予以拒绝。极有智慧的安斐洛基当时默不作声，因他想到一个令人难忘的计策。下次他进宫，看见皇帝身旁站着新近被立为皇帝的儿子\Person{阿卡狄}，便如常向\Person{德奥多西}请安，却未向阿卡狄行礼。皇帝以为此疏忽是因忘记，命安斐洛基上前向儿子请安。他说：\Quote{陛下，我向您致敬已足矣。}德奥多西为失礼而恼怒，说：\Quote{对我儿子的不敬，就是对我的无礼。}此时，极有智慧的安斐洛基才揭示其行为的用意，高声说：\Quote{陛下，您看，您不容忍对您儿子的不敬，对冒犯他者忿怒不已。那么请相信，全宇宙的天主厌恶那些亵渎独生子的人，恨他们如忘恩负义对待救主和恩主者。}\par

于是皇帝明白了主教的意图，对他所言所行都很钦佩。他不再拖延，颁布谕令，禁止异端集会。\par

但要完全逃脱人类共同仇敌的一切罗网，绝非易事。常有这样的人，他避开了淫欲的陷阱，却陷于贪婪的网罗；若他胜过贪念，另一边又有嫉妒的陷坑；即使他安然跳过此坑，又会发现另一边有情感的网等待着他。敌人还设置了无数别的绊脚石，企图引人陷入灭亡。\par

此外，他还能利用身体的情欲来助长其对灵魂的诡计。惟有心灵，若保持警醒，能胜过他，藉其向神圣事物的倾向挫败其计谋的攻势。如今，这位可敬的皇帝既有人性，亦未能免于情感，其义怒越过了适度的界限，导致了一件野蛮不法之事的实施。我必须讲述此事，以飨日后得见此书者；此事不仅涉及对可敬皇帝的责备，亦如此有助其声望，值得铭记。\par

% structure: p0089 source=h2 target=subsection reason=relative-heading-rank; base=h2

\subsection{第十七章　论\Place{得撒洛尼}的屠杀、\Person{安波罗削}主教的勇敢与皇帝的虔诚}

\Place{得撒洛尼}是一座广大而人口众多的城市，属\Place{马其顿}，但亦是\Place{帖撒利}、\Place{阿该亚}及由\Place{伊利里库姆}长官管辖的许多其他省的首府。此处爆发一场大骚乱，数名官员被投石并遭暴力对待。\par

皇帝闻讯怒火中烧，无法控制情感的冲动，甚至未用理性的缰绳勒住其发作，反而任由愤怒充当复仇的工具。当皇家的愤怒得到授权，仿佛本身即独立的君主，便挣断理性的束缚与轭，左右不分地拔出非正义的刀剑，不分青红皂白杀戮无辜与有罪者。判决前未经审判，对肇事者未予定罪。众民如收割时节的麦穗般被刈倒。据说有七千人丧生。\par

这悲惨灾祸的消息传到\Person{安波罗削}耳中。皇帝抵达\Place{米兰}，欲照常进入教堂。安波罗削在外门廊外迎住他，禁止他跨过神圣门坎。他说：\Quote{陛下，您似乎不知道所行流血之事的严重性。您的怒气已平息，但您的理性尚未认清此事的性质。或许您的皇权使您认不清罪过，权位遮蔽了理性。但我们应当知道我们的本性如何逝去，归于死亡；我们应当知道始祖的尘土，我们自彼而出，又将急速归回。我们不可因紫袍的光辉眩目，而看不见它所遮盖的身体的软弱。陛下，您统治的是与您同性的众人，他们实与您同为奴隶，因人类只有一个主子和君王，即宇宙的创造者。那么，您将用怎样的眼睛观看我们共同的主的圣殿？用怎样的脚踏这神圣门坎？您将如何伸出仍滴着不义屠杀鲜血的双手？您将如何用这双手接受主的至圣之体？您既在愤怒中不义地倾倒如此多血，又将如何举起那宝血到唇边？退去！勿在已犯的罪行上再加一罪。顺服于天主上主所同意对您施行的限制。祂将是您的医生，祂将赐您痊愈。}\par

\Person{德奥多西}既受教于神圣经卷，明确知道何为司铎之事，何为皇帝之事。因此他向\Person{安波罗削}的斥责低头，叹息流泪退回宫中。过了相当一段时间，八个月过去，救主诞生节来临，皇帝坐在宫中，泪如雨下。\par

当时，管理宫廷的鲁菲诺因与皇帝熟稔，能大胆进言，便上前问他为何哭泣。皇帝痛苦地叹息，泪水更多，说道：\Quote{您是在戏弄我，鲁菲诺，因为您没有感受到我的痛苦。我哀叹并悲泣，思及自己的灾祸；仆役和乞丐进入教会的道路是敞开的；他们可以毫无畏惧地进去，向他们的主呈上祈求。我却不敢踏入那里，此外，天堂之门也为我关闭，因为我记得主的话明确说：\BibleQuote{凡你在地上所束缚的，在天上也要被束缚}。}\par

鲁菲诺回答说：\Quote{若蒙允许，我立即前往主教那里，以恳求使他宽免您的惩罚。}皇帝说：\Quote{他不会让步。我知道盎博罗削所判的公正，他绝不会因畏惧我的皇权而违反天主的法律。}\par

鲁菲诺一再恳求，许诺能说服盎博罗削；最终皇帝命他速去。于是，德奥多西受虚假希望所骗，依赖鲁菲诺的承诺，亲自随后前往。神圣的盎博罗削一看见鲁菲诺，便喊道：\Quote{鲁菲诺，你的无耻如同狗一般，因为你是这场可怕屠杀的献策者；你已从脸上抹去羞耻，你犯下这疯狂暴行对抗天主的肖像，竟毫不脸红地站在这里。}鲁菲诺便开始哀求，并报告皇帝即将到来。神圣的盎博罗削怀着神圣的热忱喊道：\Quote{鲁菲诺，我预先告诉你：我将阻止他跨过圣门。若他要把他的主权变为暴政，我也欣然接受暴力之死。}于是鲁菲诺派使者告知皇帝总主教的心意，并劝他留在宫中。德奥多西已在广场中央时接到消息。他说：\Quote{我要去，接受我应得的羞辱。}他走到圣所前，但没有进入圣殿。总主教坐在接见厅里，皇帝前去求他解开束缚。\par

盎博罗削说：\Quote{你的到来是暴君的来临。你向天主发怒；你践踏他的法律。}德奥多西说：\Quote{不，我不攻击已定的法律，我不寻求不义地跨过圣门；而是求你解开我的束缚，顾及我们共同的主的慈悲，不要关闭那扇主为所有悔改者打开的门。}总主教回答说：\Quote{自从你犯下那重大罪行后，你显出怎样的悔改？你造成了难以愈合的创伤，你用了什么药膏？}皇帝说：\Quote{指示并调制药膏都是您的职责。我只需领受所给的。}神圣的盎博罗削于是说：\Quote{你让愤怒行正义，愤怒而非理性作判决。因此，颁布一道谕令，使你的愤怒判决无效；凡已发布的死刑或没收判决，暂停三十日，等候理性判决。当日期过去后，让写下判决的人展示命令，然后，待愤怒平息后，唯一的法官理性将审查判决，看它们是正确还是错误。若发现错误，就撤销行为；若正义，就确认；时间间隔不会给被正确定罪者带来不公。}\par

皇帝接受了这个建议，认为极好。他立即下令颁布谕令，并亲笔签署。于是神圣的盎博罗削解开了束缚。\par

这位非常虔诚的皇帝勇敢地进入圣殿，但没有站着或跪下向主祈祷，而是俯伏在地，发出大卫的呼喊：\BibleQuote{我的灵魂贴于尘土，求你照你的话使我生存。}\par

他拔头发，打头，以泪滴洒地，祈求宽恕。当他献上祭品到圣桌时，他哭泣着站起来，走近圣所。\par

献祭后，他照常留在祭台栏杆内，但伟大的盎博罗削再次没有保持沉默，教导他场所的区别。他先问皇帝是否需要什么；当皇帝说他在等候参与神圣奥迹时，盎博罗削通过首席执事传话：\Quote{先生，内室只对司铎开放；对其他所有人不可进入；出去，站在别人站的地方；紫红袍能造就皇帝，但不能造就司铎。}这位虔诚的皇帝也欣然接受这教导，并回复说，他并非出于鲁莽而留在栏杆内，而是因为他以为这是君士坦丁堡的习俗。他补充说：\Quote{我也为能纠正这个错误而感谢。}\par

于是，总主教和皇帝都显示了强大的美德光辉。两人在我眼中都令人钦佩：前者因他的勇敢言辞，后者因他的顺服；总主教因他热烈的热忱，君王因他纯洁的信德。\par

当狄奥多西返回君士坦丁堡时，他恪守从伟大总主教所习得的敬虔界限。因在节日的场合，他再次进入天主的圣殿，将祭品献给圣祭台后，便径直离去。当时的主教是奈克塔里乌斯，当主教询问皇帝为何不留在殿内时，狄奥多西叹息着回答：\Quote{我历经艰难才分清了皇帝与司铎的区别。找到一个能教导我真理的人并不容易。唯有安波罗修堪当主教的称号。}\par

一个光明而耀眼的善人将确信带回家中，这是何等大的收获啊。\par

% structure: p0105 source=h2 target=subsection reason=relative-heading-rank; base=h2

\subsection{第十八章 论皇后普拉奇拉。}

皇帝还有其他的成长机会，因为他的妻子时常提醒他天主的律法，她自己也首先仔细地学习了这些律法。她并未因皇后的身份而自高，反而因此激发出对属神之事的更大渴望。赐予她的这份伟大恩赐使她更加热爱赐予者。因此，她对残疾人和残废者关怀备至，拒绝所有仆从和卫兵的帮助，亲自探访病患的住所，为每人提供所需。她还在教堂的客房中走动，照料病人的需要，亲手拿起锅碗瓢盆，品尝肉汤，时而端上菜肴、掰开面包、递送食物，时而清洗杯子，并履行所有被认为是仆婢职责的工作。有人试图阻止她亲自做这些事，她便说：\Quote{君主应当分发黄金；而我，为了所赐予的皇权，正将自己的服侍献给赐予者。}她也常对她的丈夫说：\Quote{夫君，你应当常常回想你曾经是什么，如今成了什么；若常存此念，你便不会对恩主忘恩，而是依法治理所赐予的帝国，从而敬拜那赐予者。}她时常使用这样的言辞，如用美丽而有益的甘露，浇灌了植入丈夫心中的德行之种。\par

她先于丈夫去世，而就在她去世后不久，发生了表明丈夫何等爱她的事件。\par

% structure: p0108 source=h2 target=subsection reason=relative-heading-rank; base=h2

\subsection{第十九章 论安提约基雅的暴乱。}

由于连年征战，皇帝被迫向帝国各城加征重税。\par

安提约基雅城拒绝承受新税。当民众看到因苛税而遭受酷刑和凌辱的受害者时，除了暴民在借机骚乱时常做出的举动外，他们还推倒了著名的普拉奇拉（皇后之名）的铜像，并将其在全城拖行。皇帝得知此事后，如所预料地感到愤怒。他随后剥夺了该城的特权，将她的尊荣赐予邻城，以为这样可以给予她最大的侮辱，因为安提约基雅自远古以来便与劳迪基雅是竞争者。他还威胁要烧毁并毁灭该城，将其降为村庄。然而，地方官员在悲剧传到皇帝耳中前，已当场逮捕了一些人并处死了他们。这些命令都是皇帝下达的，但因伟大的安波罗修建言所颁布的法令的限制而未能执行。当带着皇帝威胁的钦差（其中包括当时的军事指挥官埃勒比库斯和宫廷长官凯撒里乌斯，罗马人称之为\OriginalTerm{御前大臣}{magister officiorum}）抵达时，全体民众在惊恐中战栗。然而，居住在山脚下的德行斗士们（当时有许多优秀之人）向帝国官员们做了许多恳求和祈祷。最神圣的玛刻多尼乌斯全然不谙世事，完全不知晓神圣的论说，居住在山顶，日夜向万有的救主献上纯洁的祈祷，他毫不畏惧皇帝的暴力，也不为钦差的权威所动。当他们骑马进入城中心时，他抓住其中一人的披风，命令他们二人下马。看到这个衣衫褴褛的小老头，他们起初感到愤怒，但陪同他们的一些人告知了玛刻多尼乌斯的崇高品格，于是他们一跃下马，抱住他的膝盖，请求他的宽恕。老人受上智的推动，对他们说了以下的话：\Quote{亲爱的先生们，请对皇帝说：您不仅是皇帝，您也是一个人。因此，请不仅想到您的统治权，也要想到您的本性。您是一个人，统治理应统治您的同胞。而人的本性是按照天主的肖像和模样受造的。因此，不要如此野蛮残忍地下令屠杀天主的肖像，因为惩罚祂的肖像，您会激怒造物主。请想，您现在这样因一个铜像而发怒行事。凡有理性的人都明白，一个无生命的肖像远逊于一个有灵魂和知觉的活人。同样，一个铜像，我们可以轻易地制造许多个；而您本人却连被杀者的一根头发也无法制造。}\par

善人听了这些话后，便报告给皇帝，并熄灭了他怒火的火焰。皇帝不再威胁，反而写了一份辩护，解释了他愤怒的原因。他说：\Quote{这是不应当的，因为我犯了错误，却在她死后将一个最值得最高赞美的妇人施加侮辱。那些受委屈的人本该将他们的愤怒对准我。}皇帝还补充说，当他听说有人被地方官员处决时，他感到悲伤和痛苦。叙述这些事件时，我有双重目的：我不认为应当遗忘那位杰出隐修士的勇敢，同时我也希望指出由伟大的\Person{盎博罗削}的建议所颁布的谕令的益处。\par

% structure: p0112 source=h2 target=subsection reason=relative-heading-rank; base=h2

\subsection{第二十章 关于全帝国境内庙宇的毁坏}

此时，真信皇帝将精力转向抵制异教，并颁布谕令，命令摧毁偶像的祠庙。首先以真宗教荣耀其帝国的是最值得一切赞颂的君士坦丁大帝；当他看到世界仍然沉迷于愚妄时，他普遍禁止向偶像献祭。但他并未摧毁庙宇，只下令将其关闭。他的儿子们继承了父亲的脚步。尤利安恢复了虚假信仰，重新点燃了古代欺骗的火焰。若维安继位后，再次禁止崇拜偶像；伟大的瓦伦提尼安以类似的法律统治欧洲。然而瓦伦斯允许所有人随意敬拜，并尊崇他们各自崇拜的对象；唯独对宗徒法令的护卫者他持续发动战争。因此，在他统治的整个时期，祭台之火被点燃，奠祭和祭献给偶像，公共宴会在广场上庆祝，狄奥尼索斯秘仪的入会者披着山羊皮奔跑，在酒神狂乱中撕裂猎犬，总之以种种行为展示其主人的邪恶。当真信皇帝特奥多西发现了这一切邪恶时，他将其连根拔起，并使之湮没无闻。\par

% structure: p0114 source=h2 target=subsection reason=relative-heading-rank; base=h2

\subsection{第二十一章 关于阿帕美亚主教玛尔塞禄，以及被他摧毁的偶像庙宇}

首位执行谕令、摧毁托付给他照管的城市中之庙宇的主教是玛尔塞禄，他倚靠天主而非众人的手。此事值得注意，我将叙述之。在我已提及的阿帕美亚主教若望去世后，心神热切、按照宗徒律法行事的神圣玛尔塞禄被任命接替他。\par

当时，东方总督带着两位统领及其部队到达了阿帕美亚。部队的威慑使民众保持安静。曾试图摧毁宏伟壮观的朱庇特庙，但建筑如此坚固坚实，以致于破碎其紧密砌合的石块似乎非人力所能及；因为石块巨大且垒砌精良，还用铁和铅牢牢固定。\par

当神圣的玛尔塞禄看到总督不敢开始进攻时，便打发他前往其余城镇；而他自己则祈祷天主帮助他进行破坏工作。次日清晨，有一个不请自来的人来到主教面前，他既非建筑匠、石匠，也非任何工匠，只是一个背负石头和木料的劳工。他说：\Quote{给我两个工人的工钱，我向你保证，我轻易就能摧毁这座庙。}圣主教照做了，那人的计谋如下：庙宇的四周有一排柱廊与之相连，庙的上层便支承在柱廊上。柱子体积巨大，与庙宇相称，每根柱子周长为十六肘尺。石材品质特别坚硬，对石匠工具造成很大阻力。那人在每根柱子上都凿了一圈缺口，并在继续处理下一根之前，用橄榄木支撑上部结构。在掏空了四根柱子中的三根之后，他点燃了木头。但出现了一个黑色的魔鬼，不肯让木头被火自然烧毁，阻止了火焰的力量。尝试了几次之后，发现计划无效，失败的消息传到了正在午睡的主教那里。玛尔塞禄立刻赶往教堂，命令将水倒入桶中，并将水置于神圣的祭台上。然后他俯首至地，恳求仁慈的主绝不要向魔鬼僭越的能力屈服，而要揭露其软弱，彰显祂自己的力量，免得不信者从此为更大的恶行找到借口。他以这些话及类似的话语，在水上划了十字圣号，并命令他的一位执事\Person{厄奎提乌斯}——这人满置信德与热忱——以信德取水洒之，然后点火。命令被服从了，魔鬼不能忍受水的接近，便逃走了。于是火被它的敌人水所影响，仿佛被油所激，立刻点燃了木头，瞬间将其焚尽。当支撑消失后，柱子本身倒塌，并拖倒了另外十二根。与柱子相连的庙宇一侧因倒塌的猛烈撞击而被拖垮，随之坍塌。巨大的轰鸣声响彻全城，所有人都跑来看热闹。群众一听说敌对魔鬼的逃跑，便齐声赞颂天主。\par

这位圣主教以同样的方式摧毁了其他祠庙。虽然我还有许多这位圣人的令人赞叹的事迹要叙述——因为他曾写信给凯旋的殉道者，并收到回信，并且他自己也赢得了殉道的荣冠——但此刻我犹豫不决，不敢叙述，惟恐过于冗长，使那些读到我的历史的人失去耐心。\par

因此我现在转向另一个主题。\par

% structure: p0120 source=h2 target=subsection reason=relative-heading-rank; base=h2

\subsection{第二十二章 关于亚历山大主教德奥斐禄，以及在该城摧毁偶像时发生的事}

杰出的亚大纳削由可敬的伯多禄继任，伯多禄后由弟茂德继任，弟茂德后由德奥斐禄继任，他是一位明智审慎、勇气卓越的人。亚历山大城藉他摆脱了偶像崇拜的谬误；因为他不满足于将偶像的庙宇夷为平地，还向受祭司诡计欺骗的人揭露了他们的伎俩。原来祭司们制造了内部中空的青铜和木制雕像，将背面固定在庙墙上，并在墙上留下了一些隐蔽的入口。然后他们从密室进入雕像内部，通过雕像随意发号施令，听者受骗上当，唯命是从。明智的德奥斐禄将这些伎俩公之于众。\par

此外，他进入了色辣彼斯的神庙，这座神庙曾被一些人描述为在规模和壮丽上胜过世上所有的庙宇。在那里他看到一尊巨大的雕像，其体积让观者心生恐惧，加上一个流传的谎言——有人说若有人靠近它，就会发生大地震，所有人将灭亡——更增加了恐怖。主教将这些传说视为醉酒老妇的胡言乱语，全然嘲笑那毫无生气的巨像的庞大，他吩咐一个手持斧头的人给色辣彼斯狠狠一击。那人刚一砍下，众人便大声呼喊，因为他们害怕所威胁的灾难。然而，遭受打击的色辣彼斯毫无感觉，因为他是由木头制成的，而且没有说一句话，因为他是一块没有生命的木头。他的头被砍下，立即涌出无数老鼠，因为这位埃及的神祇竟然是老鼠的居所。色辣彼斯被砸成碎片，其中一些被投入火焰，但他的头被抬着穿过全城，在他曾经的崇拜者面前示众，他们嘲笑他们曾向其屈膝叩拜的那位的软弱。\par

于是，全世界的偶像庙宇都被摧毁了。\par

% structure: p0124 source=h2 target=subsection reason=relative-heading-rank; base=h2

\subsection{第二十三章　论安提约基雅主教弗拉维亚诺及因保利诺在西方教会引发的骚乱}

在安提约基雅，伟大的墨利提乌斯由弗拉维亚诺继任，他同狄奥多禄一起，为拯救羊群经历了巨大的奋斗。保利诺确实曾希望获得主教之职，但受到神职人员的反对，理由是不应让一位与墨利提乌斯意见不合的人在他死后继任，而应当推举那位为羊群付出许多辛劳、冒过许多危险的人升任牧职。因此，在罗马人和埃及人与东方之间产生了持久的敌意，甚至保利诺死后，这种恶感仍未消除。在他之后，当厄瓦格里乌占据了他的主教座时，对伟大的弗拉维亚诺的敌意仍然存在，尽管厄瓦格里乌的晋升违反了教会的法律，因为他是由保利诺单独晋升的，无视了许多教规。因为一位临终的主教不得任命他人接替自己的位置，而是应召集所有省区的主教；此外，若无三位主教在场，不得举行主教的祝圣。然而，他们无视所有这些法律，接纳了厄瓦格里乌的共融，并向皇帝申诉反对弗拉维亚诺，致使皇帝屡受烦扰，将他召至君士坦丁堡，命令他前往罗马。\par

然而，弗拉维亚诺回答说，当时正值冬季，他答应在春天遵命。然后他返回了家乡。但是，罗马的主教们——不仅可敬的达玛稣，还有他的继任者西利西乌以及西利西乌的继任者阿纳斯大削——更加激烈地恳求皇帝，指出皇帝虽镇压了反对自己权威的竞争者，却容忍了那些胆敢违背基督法律的人维持其篡夺的权威。于是皇帝再次召见弗拉维亚诺，试图迫使他前往罗马。对此，弗拉维亚诺以极大的智慧非常勇敢地说：\Quote{陛下，若有人控告我信仰不正，或生活及行为不配司铎之职，我愿接受控告者本人作为审判官，并服从他们可能作出的任何判决。但若他们争论的是教座和首席权，我绝不争辩；我不反对那些想要取得它们的人；我愿退让，辞去主教之职。所以，陛下，您愿意将安提约基雅的主教座交给谁就交给谁吧。}\par

皇帝赞赏他的刚毅和明智，命他再次回家，并照料托付给他的教会。\par

过了相当长的时间，皇帝来到罗马，再次面对主教们提出的指控，说他并未试图铲除弗拉维亚诺的暴政。皇帝吩咐他们说明暴政的性质，并说自己就是弗拉维亚诺，成了他的保护者。主教们回答说，他们不可能与皇帝争辩。于是皇帝劝告他们今后要同心合一地联合教会，结束纷争，扑灭无益争论的烈火。他指出，保利诺早已离世；厄瓦格里乌的晋升不合规矩；东方教会接纳弗拉维亚诺为主教。不仅东方，而且整个亚细亚、本都和色雷斯都与他共融，整个伊利里亚也承认他对东方主教们的权威。听从这些劝告后，西方主教们承诺结束敌意，并接待将要派来的使节。\par

当弗拉维亚诺得知这一决定后，他派遣了几位可敬的主教，连同安提约基雅的长老和执事前往罗马，其中授予贝洛亚主教阿加丘最高职权，他在全世界享有盛名。阿加丘及其一行到达罗马后，结束了长期的分歧，在十七年的战争之后，将和平带给了教会。埃及人得知和解的消息后，也放弃了反对，欣然接受了所达成的协议。\par

那时，在罗马教会的首席权上，阿纳斯大削由依诺增爵继任，他是一位明智而机敏的人。此前我提到的德奥斐禄仍掌管亚历山大主教座。\par

% structure: p0131 source=h2 target=subsection reason=relative-heading-rank; base=h2

\subsection{第二十四章 论欧根尼的暴政以及德奥多西皇帝因信德而获得的胜利}

如此，最虔诚的皇帝确保了教会的和平。在和平建立之前，他已听闻瓦伦提尼安努斯之死及欧根尼的篡位，并已向欧洲进军。\par

此时，在埃及住着一位名叫若望的人，他度着隐修生活。他充满属神的恩宠，时常向前来请教的人预言许多未来之事。基督爱戴的皇帝心中焦虑，不知是否应当讨伐僭越者，便派人去请教他。对于前一次战争，他预言了一场不流血的胜利；对于第二次，他预言皇帝将在大屠杀之后才能获胜。皇帝怀着此期望出发，在列阵时，射杀了许多敌人，但也失去了许多蛮族盟友。\par

当他的将领们表示己方兵力不足，建议暂停战事，以便在初春集结大军，在人数上超过敌军时，德奥多西拒绝听从他们的建议。\Quote{因为把救恩之十字架说成如此软弱是错误的，}他说，\Quote{是十字架引领我们的军队，却将如此能力归于敌军阵前引导的赫丘利像。}他以正确的信德说了这话，尽管他手下的人寥寥无几，且十分沮丧。然后他在营地所在的山顶上找到一个小祈祷所，整夜向万有的天主祈祷。\par

大约鸡鸣时分，睡意征服了他，当他躺在地上时，他看见两个身穿白衣、骑着白马的人，他们吩咐他振作精神，驱散恐惧，在黎明时武装起来，排列人马准备战斗。\Quote{因为，}他们说，\Quote{我们被派来为你作战。}其中一个说：\Quote{我是若望福音作者，}另一个说：\Quote{我是斐理伯宗徒。}\par

皇帝看见这异象后，并没有停止恳求，反而更加热切地祈祷。这异象也被军中一个士兵看见，他向百夫长报告。百夫长带他去见千夫长，千夫长又带他去见将军。将军以为他讲的是新奇之事，便向皇帝报告。于是德奥多西说：\Quote{这异象被此人看见，并非为了我的缘故，因为我已经信靠那些应许我胜利的人。但为了不让人以为我因渴望战斗而编造了这异象，我帝国的护卫者也将这信息启示给这人，使他能为我说的话作证：当初是主首先将这异象赐予我。那么，让我们抛开恐惧。让我们跟随我们的前锋和将领。不要让任何人按参战人数衡量胜利的机会，而要人人思想领袖的能力。}\par

他向士兵们讲了类似的话，如此激励全军的崇高希望后，便带领他们从山顶下来。僭越者远远看见军队来攻击他，便武装他的军队，列阵准备战斗。他自己留在高处，说皇帝是求死，来战斗是因为想脱离今生；于是他命令将领们将皇帝活捉并锁链捆绑。当两军列阵时，敌军人数似乎远多于皇帝军队，皇帝的军队人数屈指可数。但当双方开始发射兵器时，前阵证明了他们的应许是真的。一阵狂风直吹敌军的面孔，将他们的箭矢、标枪和长矛都吹偏了，以致任何投射武器都对它们无用，无论是骑兵、弓箭手还是长矛手，都无法对皇帝的军队造成任何伤害。大团的尘土也吹到他们脸上，迫使他们闭上眼睛并保护自己免受攻击。另一方面，皇帝的军队没有受到风暴的丝毫伤害，他们奋勇攻击并杀死敌人。战败者于是认出了上天给予胜利者的帮助，扔掉武器，向皇帝求饶。德奥多西便应允他们的恳求，怜悯他们，命令他们立即将僭越者带到他面前。欧根尼不知道战况如何，当他看见他的士兵上气不接下气地跑到他所在的山丘上，气喘吁吁地显示出他们的急切时，他以为他们是胜利的信使，问他们是否如他所命令的那样将德奥多西锁链带来。\Quote{不，}他们说，\Quote{我们不是带他到你这里来，而是来把你带到他那里去，因为伟大的统治者如此安排。}他们一边说，一边把他从战车上拉下来，给他戴上锁链，就这样捆着带走，将片刻前还妄自尊大的人作为战俘押走了。\par

皇帝提醒他所犯下的对瓦伦提尼安努斯的恶行、他篡夺的权力以及他针对合法皇帝发动的战争。他还嘲笑了赫丘利像及其引发的愚蠢自信，最后宣布了公正合法的判决。\par

德奥多西在和平与战争中便是如此，总是祈求天主的帮助，从未被拒绝。\par

% structure: p0140 source=h2 target=subsection reason=relative-heading-rank; base=h2

\subsection{第二十五章 论德奥多西皇帝之死}

这次胜利后，德奥多西患病，将他的帝国分给他的两个儿子，将他自己曾经治理的东部分给长子，将欧洲的王位分给幼子。\par

他嘱咐两人持守真宗教，\Quote{因为藉着它，}他说，\Quote{和平得以保持，战争得以停止，敌人被击退，战利品被设立，胜利被宣告。}给他的儿子们这样嘱咐后，他便去世了，留下了不朽的名声。\par

他帝国中的继任者也继承了他的虔诚。\par

% structure: p0144 source=h2 target=subsection reason=relative-heading-rank; base=h2

\subsection{第二十六章 论霍诺留皇帝与隐修士特莱玛科斯}

欧诺留继承欧洲帝国，终止了长期在罗马举行的角斗比赛。他这样做的起因如下：有一位名叫德肋玛苛的人，度着苦修生活，从东方出发，为此原因来到罗马。在那里，当这可憎的表演进行时，他亲自进入运动场，走下竞技场，试图阻止那些互相挥舞武器的人。嗜血的观众被激怒，被那喜爱血腥行为的魔鬼的疯狂所驱使，用石头将这位和平缔造者砸死。\par

当这位可敬的皇帝得知此事后，他将德肋玛苛列入胜利的殉道者行列，并终止了那不敬神的表演。\par

% structure: p0147 source=h2 target=subsection reason=relative-heading-rank; base=h2

\subsection{第二十七章 论皇帝阿卡狄乌斯的虔诚与金口圣若望的祝圣}

在君士坦丁堡，该教区主教内克塔里乌斯去世后，已继承东帝国皇帝位的阿卡狄乌斯召来了若望——世界伟大的明灯。他听闻若望名列司铎团中，便下令召集的主教们将天主恩宠授予他，并任命他为那座大城市的牧者。\par

单是这一事实就足以显示皇帝对神圣事物的关心。同时，安提约基雅教区由弗拉维亚努斯主持，劳迪西亚教区由厄尔丕丢斯主持，他原是伟大美利提乌斯的同伴，从他的生活和言谈中所受的印记，比蜡从戒指上所取的印记更为清晰。\par

他继承了伟大的贝拉基；神圣的玛尔克路斯之后是著名的阿加丕图斯，我已描述过他以高超的苦修德行而卓越。在异端风暴时期，在陶鲁斯山麓的塞琉西亚\OriginalTerm{陶鲁斯山麓的塞琉西亚}{Seleucia ad Taurum}，伟大若望的同伴玛克西穆斯是主教，在摩普绥提亚则是戴奥多勒斯，两人都是杰出的教师。圣阿卡基乌斯，贝若亚主教，在智慧和品格上也同样卓越。\par

莱昂丢斯，许多德行的光辉榜样，牧养了加拉提雅人的羊群。\par

% structure: p0152 source=h2 target=subsection reason=relative-heading-rank; base=h2

\subsection{第二十八章 论若望为天主的勇敢}

当伟大的若望接掌教会的舵柄时，他勇敢地定了某些作恶者的罪，适时地向皇帝和皇后提出劝诫，并告诫圣职人员要按照所制定的法律生活。违反这些法律的人，他禁止他们接近教会，坚称那些不愿度真正司铎生活的人，不应享有司铎的尊荣。他如此关心教会，不仅在君士坦丁堡，而且遍及整个色雷斯——分为六省，同样也在亚细亚——由十一位总督治理。本都，也有与亚细亚数目相等的统治者，有福地被纳入同样的纪律之下。\par

% structure: p0154 source=h2 target=subsection reason=relative-heading-rank; base=h2

\subsection{第二十九章 论若望在腓尼基所摧毁的偶像庙宇}

若收到信息，得知腓尼基仍在遭受魔鬼礼仪的疯狂，若望召集了一些被神圣热忱所激励的隐修士，以帝国谕令武装他们，派遣他们去对付偶像的庙宇。支付从事破坏工作的工匠及其助手所需的金钱，若望并未取自帝国资源，而是说服某些富裕而忠信的女子慷慨捐献，向她们指出她们的慷慨将赢得多么大的祝福。\par

于是，剩下的魔鬼庙宇被彻底摧毁。\par

% structure: p0157 source=h2 target=subsection reason=relative-heading-rank; base=h2

\subsection{第三十章 论哥特人的教会}

若望察觉到司奇提亚人陷入了亚略派的网罗\OriginalTerm{亚略派}{Arian}；因此他设想了对策，发现了赢得他们的方法。他任命了能说司奇提亚语的司铎、执事和神圣训言诵读者，指定了一座教堂给他们，并藉着他们使许多人脱离了错误。他本人经常去那里讲道，使用一位精通两种语言的翻译，并让其他善于辞令者也这样做。这是他在城内的惯常做法，许多受骗的人，他通过向他们指出宗徒宣讲的真理而拯救了他们。\par

% structure: p0159 source=h2 target=subsection reason=relative-heading-rank; base=h2

\subsection{第三十一章 论他对司奇提亚人的关心及他对马西翁派\OriginalTerm{马西翁派}{Marcionists}的热忱}

若得知一些游牧民族驻扎在多瑙河沿岸，渴望救恩，却无人将活水带来，若望便寻找那些充满热爱辛劳之人，如同宗徒所突出的那般，并将工作托付给他们。我本人见过他写给安卡拉主教莱昂丢斯的信，信中描述了司奇提亚人的归化，并请求派遣适合教导他们的人。\par

若听闻在我们的地区有人感染了马西翁\OriginalTerm{马西翁}{Marcion}的瘟疫，他便写信给当时的主教，嘱咐他驱除瘟疫，并以帝国谕令提供援助。我已说得足够，表明他如何，借用神圣宗徒的话，在心中怀着\Quote{所有教会的挂虑}。\par

他的勇敢也可以从其他来源得知。\par

% structure: p0163 source=h2 target=subsection reason=relative-heading-rank; base=h2

\subsection{第三十二章 论盖纳斯的要求及若望·金口圣的答复}

有一位名叫盖纳斯的司奇提亚人，但性格更为野蛮，性情残酷暴力，当时是一位军事指挥官。他手下有许多自己的同胞，与他们一起指挥罗马骑兵和步兵。他不仅使所有人恐惧，甚至皇帝本人也恐惧他，皇帝怀疑他企图篡位。\par

他曾是亚略异端瘟疫的参与者，并请求皇帝准许他使用一座教堂。阿卡狄回答说，他将留意此事并予以办理。随后他派人请来金口若望，告知他那项请求，提醒他盖纳斯的势力，暗示其图谋篡位，恳求他通过这项让步来抑制那蛮族的愤怒。\Quote{但是，}那位高尚的人说，\Quote{大人，不要作出这样的许诺，也不要把圣物给狗。我绝不能容忍崇拜和歌颂圣言的人被驱逐，而将他们的教堂交给亵渎祂的人。大人，不要怕那个蛮族；将我和他一同召到你面前；你静听我们所说的话，我必能制住他的舌头，劝他不要要求那些不应允准的事。}\par

皇帝对金口若望的话甚为喜悦，次日便召见了主教和将军。盖纳斯开始要求履行诺言，但金口若望回答说，信奉真宗教的皇帝无权做任何反对真宗教的事。盖纳斯反驳说，他也必须有一个祈祷的地方。\Quote{为何？}金口若望说，\Quote{每座教堂都向你开放，你愿意祈祷时，无人阻止你。}\Quote{但是我，}盖纳斯说，\Quote{属于另一个派别，我请求与他们共用一座教堂。我为罗马人历尽战阵辛劳，提出这样的请求应是合理的。}\Quote{但是，}主教说，\Quote{你为你的劳苦获得了更大的报酬；你是一位将军，你身着执政官袍。你必须思想你先前如何，现今如何——你过去的贫困与今日的兴盛；你渡伊斯特河之前穿着怎样的衣服，如今又披着什么。我再说，要思想你劳苦之微小与报赏之宏大，不要对赐你荣耀者忘恩。}这位天下师表用这些话使盖纳斯哑口无言，迫使他呆立不语。然而，随着时间的推移，盖纳斯显露出他蓄谋已久的反叛之心，在色雷斯聚集兵力，四处劫掠抢掠。此消息传出，无论君主还是臣民都惶恐不安，无人愿意率军与他作战；无人认为遣使接近他是安全的，因为人人都怀疑其蛮族本性。\par

% structure: p0167 source=h2 target=subsection reason=relative-heading-rank; base=h2

\subsection{第三十三章　论金口若望出使盖纳斯}

当时，由于普遍的恐慌，其他人都被撇下，这位伟大的首领被说服承担了出使之职。他不计较所述之争及其所生的恶感，欣然前往色雷斯。盖纳斯一听说使者到达，便想起他为真宗教所作的勇敢言论。他急切地从远处赶来迎接，将右手放在若望的眼上，并将自己的儿女带到圣人膝前。善性就是如此，它能使那些最反对它的人羞愧降服。然而，嫉妒之心不能容忍他哲学的灿烂光芒。它施展惯用伎俩，使帝都——诚然，也使整个世界——失去了他的口才与智慧。\par

% structure: p0169 source=h2 target=subsection reason=relative-heading-rank; base=h2

\subsection{第三十四章　论因金口若望所发生的事}

在我史书的这一部分，我不知道该抱持何种情感；我既希望叙述对金口若望所行的不义，却又在另一方面顾念那些加害于他之人的崇高品德。因此，我将尽力甚至隐去他们的姓名。这些人敌视他的理由各不相同，且不愿正视他耀眼的德行。他们寻得一些卑劣之徒控告他，又因看穿这诽谤的露骨，便在远离城市的地方聚会并作出判决。\par

皇帝信任神职人员，便命令将他流放。于是金口若望未听取对他的指控，也未为自己辩护，就如已被定罪一般，被迫离开君士坦丁堡，前往欧克辛海口的希耶隆——海军驻地名即如此。\par

当夜发生了大地震，皇后惊惶不已。于是天亮时便派使者到被流放的主教那里，恳求他立即返回君士坦丁堡，以消除城中的危难。此后又派了另一批使者，又再派一批，博斯普鲁斯海峡上挤满了信使。当信友得知此事后，他们用船只覆盖了普罗庞提斯海的入口，全体民众点燃蜡炬，出来迎接他。那时，他的敌党一时四散。\par

然而过了数月，他们试图对其施加刑罚，不是因那伪造的指控，而是因他在被革职后仍参与神圣礼仪。主教申辩说，他未申诉，未听指控，未作辩护，在缺席中被判罪，被皇帝流放，又由皇帝召回。于是召开了另一次主教会议，反对者并未要求审讯，而是说服皇帝，认定判决合法且正确。金口若望于是不仅被流放，更被放逐到亚美尼亚一个偏远小镇，名为库库苏斯。甚至从那里，他又被转移并放逐到皮提乌斯，那是欧克辛海的尽头，地处罗马帝国边境，紧邻最野蛮的蛮族。但仁慈的主并未允许这位胜利的运动员被带至那个小岛；因为他到达科玛纳时，便被接引到那无老无痛的永生。\par

那英勇奋斗的身体被安葬在殉道者巴西利斯库斯的棺椁旁——如此殉道者曾在梦中预示。\par

我认为没有必要再延长叙述，说明有多少主教因\Person{金口圣若望}的缘故被逐出教会，被遣往地极居住，或有多少苦修哲人遭受同样的灾祸；我之所以如此，更是因为我认为有必要缩减这些可憎的细节，并对同我们一样信仰之人的恶行加以遮掩。然而，惩罚确实降在了大多数有罪者身上，他们的受苦对其他人而言是行善的途径。这件大错尤其为\Place{欧洲}的主教们所憎恶，他们与有罪者断绝了共融。在此行动中，所有\Place{依利里亚}的主教也都加入。在\Place{东方}，大多数城市虽不愿参与此错，却未在教会身体上造成分裂。\par

当世界伟大的导师去世后，西方的主教们拒绝与\Place{埃及}、\Place{东方}、\Place{博斯普鲁斯}和\Place{色雷斯}的主教们共融，直到那位圣人的名字被列入已故主教的名录中为止。他们拒绝承认他的直接继任者\Person{阿尔萨修}，但\Person{阿尔萨修}的继任者\Person{阿提古}多次恳求和平的恩赐后，在将名字加入名录时被接纳。\par

% structure: p0177 source=h2 target=subsection reason=relative-heading-rank; base=h2

\subsection{三十五、论\Place{安提约基雅}主教\Person{亚历山大}}

此时，\Place{亚历山大}主教座堂由\Person{济利禄}管理，他是他的叔父\Person{德奥斐洛}的继任者；同时，\Place{耶路撒冷}由\Person{若望}管理，他是我们先前提到过的\Person{济利禄}的继任者。\Place{安提约基雅}则受\Person{亚历山大}照管，他的生活和言行与其主教职务相符。在他祝圣之前，他过着苦修操练和严格身体锻炼的生活。他被称为高贵战士，以言教导，并以行证实其言。他的前任是\Person{波菲里}，他在\Person{弗拉维亚诺}之后领导那教会，并留下许多显示其仁爱品格的纪念。他也以智力卓著著称。圣\Person{亚历山大}尤其富有自我节制和修行；他过着贫穷和克己的生活；他口才流利，其他恩赐无数；通过他的建议和劝勉，伟大的\Person{欧斯达提}的追随者——\Person{保利诺}及其后继者\Person{厄瓦格利}曾不允许他们恢复——与教会身体的其他部分重新合一，举行了一场前所未见的庆典。主教聚集所有信友，包括圣职人员和教友，与他们一起前进到集会处。游行队伍伴有乐师；众人口唱同一首圣诗，和谐一致，就这样，他和他的队伍从西侧小门游行到大教堂，整个广场挤满了人，形成一股有灵的活人之流，如同\Place{奥龙特斯河}的奔流。\par

当这一切被犹太人、亚略瘟疫的受害者和那微不足道的外教人残余看见时，他们就哀叹痛哭，为看到其余的河流将水倾注于教会而苦恼。通过\Person{亚历山大}，伟大的\Person{若望}的名字首次被记入教会的记录中。\par

% structure: p0180 source=h2 target=subsection reason=relative-heading-rank; base=h2

\subsection{三十六、论\Person{若望}遗骸的迁移以及\Person{德奥多西}及其姊妹的信德}

后来，伟大的圣师的实际遗骸被运往帝都，信众再次以密集的船只将海化为陆地，用火把覆盖了\Place{博斯普鲁斯}朝向\Place{普罗庞提斯}的入口。这宝贵的所有物由当今皇帝带入\Place{君士坦丁堡}，他继承了他祖父的名字，并保持其虔诚无玷。他先瞻仰了棺架，然后将头靠在上面，为他的父母和那些无知犯罪的人祈求宽恕，因为他的父母早已去世，使他幼年成为孤儿；但他祖先的天主未允许他因孤儿身份而受苦，反为他提供虔敬的教养，保护他的帝国免受暴乱袭击，并勒住悖逆之心。他时常记念这些恩惠，以赞美颂歌尊崇他的恩人。与他一同参与这神圣敬拜的还有他的姊妹，她们终身保持童贞，认为研读神圣神谕是最大的乐趣，并看为穷人服事是盗贼无法夺走的财富。皇帝本人被许多恩宠装饰，尤其以他的仁慈和宽恕、灵魂的宁静与纯洁昭著的信德为最。对此我将给出一个无可辩驳的证据。\par

一位脾气有些暴躁的苦修者带着请愿书来到皇帝面前。他来了好几次都未能达到目的，最后他绝罚了皇帝，并使皇帝处于禁令之下。这位信德的皇帝回到宫中，正值宴饮之时，宾客齐聚，他说在禁令解除之前他不能享用宴席。因此，他派出自己最亲信的随从去见主教，恳求主教命令施加禁令者解除禁令。主教回答说，不应接受每个人的禁令，并宣布该禁令无效，但皇帝不接受这一宽免，直到经过很大困难找到施加禁令者，并恢复了被撤销的共融为止。他对神圣法律的服从如此之甚。\par

依照同样的原则，他下令彻底摧毁偶像神庙的遗迹，以使我们的后代甚至避免看到一丝古代错误的痕迹，这是他为此颁布的敕令中所表达的动机。他不断收获所播下的良种果实，因为万有的主与他同在。因此，当斯基泰游牧民族的首领\Person{罗伊拉斯}率领大军渡过\Place{多瑙河}，蹂躏掠夺\Place{色雷斯}，并威胁要围攻帝都，迅速占领并将其交付毁灭时，天主从高处用雷电和暴风雨打击了他，烧毁了入侵者并毁灭了他的全军。同样，在波斯战争中也显示了类似的天主眷顾。波斯人得知罗马人正忙于他处，便违反和约，向他们的邻国进军，邻国在进攻中找不到援助者，因为皇帝信赖和约，已派遣将领和士兵去其他战争。随后，一场非常猛烈的暴风雨和冰雹阻止了波斯人的进一步前进；他们的马匹拒绝前进；二十天里他们未能前进二十弗隆。与此同时，将领们返回并集结了军队。\par

在先前的战争中，同样是这些波斯人，在围攻与皇帝同名的城市时，蒙天主眷顾而显得可笑。因为\Person{瓦拉兰}率领全军围攻上述城市超过三十天，调动了许多攻城塔，使用了无数器械，并在墙外筑起了高塔，但抵抗仍在进行，攻击器械的冲锋被\Person{欧诺弥乌斯}主教一人击退。我方士兵拒绝与敌人作战，不愿前来援助被围困者，而主教却独自挺身而出，保护城市免于沦陷。当一位蛮族酋长胆敢一如既往地亵渎，用类似\Person{拉波沙刻}和\Person{散乃黑黎布}的言语疯狂威胁要焚烧天主的圣殿时，圣主教不能容忍他的暴怒，亲自命令将一架称为宗徒\Person{多默}的弩炮架在城垛上，并装上大石。然后，以被亵渎之主的名号，他下令发射——石头砸在那邪恶酋长身上，击中他的邪恶之口，砸碎他的面容，打碎他的头颅，将他的脑浆洒在地上。当期望攻取该城的军队统帅看到所发生的事，承认失败并撤退，惊恐中缔结了和平。\par

因此，宇宙的君王保护忠信皇帝，因他清楚承认自己是谁的奴仆，并向他的主人履行合适的服务。\par

% structure: p0186 source=h2 target=subsection reason=relative-heading-rank; base=h2

\subsection{第37章 论安提约基雅主教\Person{德奥多图}}

\Person{德奥多西}将世界大光的遗骸归还给深深哀悼其损失的城。然而这些事发生在后来。\par

\Person{依诺森}，罗马卓越的主教，由\Person{波尼法修}继任，\Person{波尼法修}由\Person{佐西莫}继任，\Person{佐西莫}由\Person{塞勒斯丁}继任。\par

在耶路撒冷，令人钦佩的\Person{若望}之后，教会的管理托付给了\Person{普拉依略}，一位名实相符的人。\par

在安提约基雅，神圣的\Person{亚历山大}之后，纯洁的珍珠\Person{德奥多图}继任教会的最高权力，他是一位异常温良、生活严谨的人。通过他，\OriginalTerm{阿波利纳里}{Apollinarius}教派被允许与其他羊群共融，因其成员恳切要求与羊群合一。然而其中许多人仍留下了先前不健全的印记。\par

% structure: p0191 source=h2 target=subsection reason=relative-heading-rank; base=h2

\subsection{第38章 论波斯的迫害及在那里殉道者}

此时，波斯王\Person{伊斯迪吉尔德}开始向教会发动战争，其原因如下：一位名叫\Person{阿布达斯}的主教，德才兼备，却因不当的热心而摧毁了一座\OriginalTerm{火庙}{Pyreum}——Pyreum是波斯人对他们视为天神的火之寺庙的称呼。\par

\OriginalTerm{祆教徒}{Magi}将此禀告给\Person{伊斯迪吉尔德}后，他召来\Person{阿布达斯}，首先以温和的语气抱怨所发生的事，并命令他重建那火庙。\par

主教坚决拒绝如此做，于是国王威胁要摧毁所有教堂，最终实现了所有威胁：他先下令处死那位圣人，然后下令摧毁教堂。现在我认为，摧毁火庙是错误且不明智的，因为即使是神圣的宗徒，当他来到\Place{雅典}，看见全城充满偶像崇拜时，也没有摧毁雅典人所敬重的任何祭坛，而是通过论证说服他们的无知，并揭示了真理。但拒绝重建倒塌的神庙，并决心选择死亡也不如此行，我大大称赞并尊敬，认为这是值得殉道者冠冕的事；因为建造一座供奉火的圣殿，在我看来等同于崇拜它。\par

从这开始，掀起了一场风暴，向真信仰的哺育者激起凶猛残酷的波浪，三十年过去，由\OriginalTerm{祆教徒}{Magi}煽动的骚动仍未平息，如同大海被狂风激起的巨浪。\OriginalTerm{祆教徒}{Magi}是波斯人对太阳和月亮的崇拜者的称呼，但我在另一篇论文中揭示了他们神话般的体系，并提出了对其难题的解答。\par

当\Person{伊斯迪吉尔德斯}去世后，其子\Person{瓦拉兰}即继承了王国，也继承了针对信德的战争；待他离世时，又两者同时留给了他的儿子。要详述对圣人们施加的各种酷刑及暴行，实非易事。有时，他们的双手被剥皮，有时是背部；另一些人从头至须的皮肤被剥去；还有一些人被劈开的芦苇包裹，切口朝内，从头到脚被紧紧缠绕的绷带束缚；然后每一根芦苇被用力拉出，扯下邻近的皮肤，造成剧痛。挖好了坑，仔细涂上油脂，放入大量老鼠，然后放下手脚被绑的殉道者，使他们无法抵御这些动物，成为老鼠的食物；老鼠在饥饿逼迫下，一点点啃食受害者的肉，造成漫长而可怕的痛苦。另一些人则遭受比这些更可怕的折磨，那是由人性之敌、真理之敌所发明的；但殉道者的勇气并未消减，他们为赢得那引入不朽生命的死亡而迫不及待地奔赴。\par

在这些殉道者中，我将引述一两个事例，以作为其余勇气之榜样。波斯贵族中有一位名叫\Person{霍尔米斯达斯}，出自阿契美尼德家族，是一位总督之子。有人报告他是基督徒，国王召见他，命他背弃他的救主天主。他回答说，王命既不义又不合理，\Quote{因为，}他继续说道，\Quote{一个被教导认为轻蔑和否认万有之天主毫无困难的人，或许会更容易轻视一个凡人有死的君王；并且，陛下，如果否认你的主权应受最严厉的惩罚，那么否认创造世界的那一位，岂不是更该受更可怕的刑罚？}国王本应赞赏这番话的智慧，但他却剥去了这位高贵斗士的财富和尊位，命他除了一块腰布外赤身露体，去赶军队的骆驼。几天后，国王从寝宫向外望，看见那位杰出的人被阳光灼烤，满身尘土，便想起他父亲的显赫地位，派人召他来，命他穿上一件细麻布长衣。国王以为他受过的劳苦和所施的恩惠已软化了他的心，便说：\Quote{现在至少，}他说，\Quote{放弃你的反对，否认那木匠的儿子。}满怀神圣热忱的\Person{霍尔米斯达斯}撕碎了长衣，扔开说道：\Quote{如果你以为这会使我放弃真信仰，那你就留着你的礼物和你的伪信吧。}国王见他如此大胆，便把他赤身逐出王宫。\par

有一位名叫\Person{苏埃内斯}的人，拥有上千名奴隶，他抗拒国王，不肯否认他的主人。于是国王问他哪一个奴隶最卑贱，便将其他所有奴隶的所有权交给了他，并把\Person{苏埃内斯}赐给他为奴。国王还把他的妻子嫁给那奴隶，以为如此便能折服这位真理捍卫者的意志。但他失望了，因为\Person{苏埃内斯}把他的房子建在了磐石上。\par

国王还逮捕并监禁了一位名叫\Person{本雅明}的执事。两年后，有一位来自\Place{罗马}的使节前来处理其他事务，他得知此事后，请求国王释放这位执事。国王命\Person{本雅明}承诺不向任何玛吉传讲基督信仰，使节也劝\Person{本雅明}顺从；但\Person{本雅明}听了使节的话后，回答说：\Quote{我不能不把自己所得的光明传给别人；因为隐藏塔冷通该受多么大的惩罚，神圣福音的记载已经教导了我们。}直到此时，国王尚未得知这一拒绝，便命释放了他。\Person{本雅明}一如既往，设法捕捉那些被无知黑暗所压制的人，把他们领向知识的光明。一年后，有人向国王报告了他的行为，国王便召见他，命他否认他所敬拜的那一位。他问国王：\Quote{一个人若背弃自己的忠诚而转向另一位，应受什么惩罚？}\Quote{死刑和酷刑，}国王说。\Quote{那么，}这位明智的执事继续说道，\Quote{一个人若背弃他的创造者和造主，反而以同是受造的人为神，并把应归于他主的荣耀献给那一位，又该如何处置？}国王于是满怀愤怒，命人将二十根削尖的芦苇钉入他的指甲和脚趾中。当他看到\Person{本雅明}视此折磨如同儿戏时，又拿起一根芦苇，刺入他的私处，来回搅动，造成难以言喻的痛苦。在此酷刑之后，这邪恶残暴的暴君命人把他插在一根粗大的多节木桩上，于是这位高贵的受苦者便断了气。\par

这些不虔敬的人还犯下了无数类似的其他暴行；但我们不必惊异于万有的主竟容忍了他们的残暴和不虔敬，因为实际上，在\Person{君士坦丁大帝}统治之前，所有罗马皇帝都曾向真理的朋友们发泄愤怒；并在救主受难之日，\Person{戴克里先}曾摧毁了整个罗马帝国境内的教会；但九年过去后，教会重新兴起，更加繁荣美丽，规模比先前大了许多倍，更加辉煌；而他本人及其邪恶却一同灭亡了。\par

这些战争和教会的胜利，主早已预言；事件教导我们，战争比和平带来更多的祝福。和平使我们变得娇弱、安逸和怯懦；战争则磨砺我们的勇气，使我们轻视这转瞬即逝的现世。但我们在其他著作中已多次作过这些观察。\par

% structure: p0202 source=h2 target=subsection reason=relative-heading-rank; base=h2

\subsection{第三十九章　论\Place{莫普苏埃斯提亚}主教\Person{戴奥多若}}

当\OriginalTerm{神圣的狄奥多鲁斯}{Theodorus}治理\OriginalTerm{安提约基亚}{Antioch}教会时，\OriginalTerm{摩普绥提亚}{Mopsuestia}的主教\OriginalTerm{狄奥多鲁斯}{Theodorus}——他是整个教会的导师，也是成功抗击一切异端方阵的斗士——结束了此生。他曾受教于伟大的\OriginalTerm{狄奥多罗斯}{Diodorus}，并且是\OriginalTerm{圣若望}{John}的朋友和同工，因为他们二人共同从\OriginalTerm{狄奥多罗斯}{Diodorus}所赐的灵性甘露中受益。他在主教职位上度过了三十六年，与\OriginalTerm{亚略}{Arius}和\OriginalTerm{欧诺弥}{Eunomius}的势力争战，与\OriginalTerm{阿波利纳里}{Apollinarius}的海盗团伙搏斗，并为天主的羊群找到了最好的牧场。他的兄弟\OriginalTerm{波里赫罗尼乌斯}{Polychronius}是\OriginalTerm{阿帕米亚}{Apamea}的杰出主教，一位口才极佳、品格卓越的人。\par

我将在此结束我的历史，并恳请那些读到它的人以祈祷回报我的辛劳。本叙事涵盖105年，从\OriginalTerm{亚略派}{Arian}的疯狂开始，到令人敬仰的\OriginalTerm{狄奥多鲁斯}{Theodorus}和\OriginalTerm{狄奥多图斯}{Theodotus}的去世为止。我将列出一份大城主教名单，在迫害之后。\par

\Emph{大城主教名单。}\par

\begin{itemize}
\item \OriginalTerm{罗马}{Rome}：—
\end{itemize}

\begin{itemize}
\item \OriginalTerm{米提亚德}{Miltiades} [\OriginalTerm{梅尔基亚德斯}{Melchiades}. 311-314]
\end{itemize}

\begin{itemize}
\item \OriginalTerm{西尔维斯特}{Silvester} [314-335]
\end{itemize}

\begin{itemize}
\item \OriginalTerm{尤利乌斯}{Julius} [337-352. \OriginalTerm{马尔谷}{Mark} Jan. to Oct., 336]
\end{itemize}

\begin{itemize}
\item \OriginalTerm{利贝里乌斯}{Liberius} [352-366]
\end{itemize}

\begin{itemize}
\item \OriginalTerm{达玛稣}{Damasus} [366-384]
\end{itemize}

\begin{itemize}
\item \OriginalTerm{西里修}{Siricius} [384-398]
\end{itemize}

\begin{itemize}
\item \OriginalTerm{阿纳斯塔修斯}{Anastasius} [398-401]
\end{itemize}

\begin{itemize}
\item \OriginalTerm{依诺增爵}{Innocentius} [402-417]
\end{itemize}

\begin{itemize}
\item \OriginalTerm{博尼法爵}{Bonifacius} [418-422]
\end{itemize}

\begin{itemize}
\item \OriginalTerm{佐西默}{Zosimus} [417-418]
\end{itemize}

\begin{itemize}
\item \OriginalTerm{策肋定}{Cælestinu} [422-432]
\end{itemize}

\begin{itemize}
\item \OriginalTerm{安提约基亚}{Antioch}：—
\end{itemize}

\begin{itemize}
\item \OriginalTerm{维塔利}{Vitalius}（\OriginalTerm{正统派}{Orthodox}）[312-318]
\end{itemize}

\begin{itemize}
\item \OriginalTerm{斐洛格尼}{Philogonius}（\OriginalTerm{正统派}{Orthodox}）[318-323]
\end{itemize}

\begin{itemize}
\item \OriginalTerm{厄斯塔提}{Eustathius}（\OriginalTerm{正统派}{Orthodox}）[325-328]
\end{itemize}

\begin{itemize}
\item \OriginalTerm{厄乌拉里}{Eulalius}（\OriginalTerm{亚略派}{Arians}）[328-330]
\end{itemize}

\begin{itemize}
\item \OriginalTerm{欧弗罗尼}{Euphronius}（\OriginalTerm{亚略派}{Arians}）[330-332]
\end{itemize}

\begin{itemize}
\item \OriginalTerm{普拉西多}{Placidus}（\OriginalTerm{亚略派}{Arians}）[332-342]
\end{itemize}

\begin{itemize}
\item \OriginalTerm{斯德望}{Stephanus}（\OriginalTerm{亚略派}{Arians}）[342-348]
\end{itemize}

\begin{itemize}
\item \OriginalTerm{良提}{Leontius}（\OriginalTerm{亚略派}{Arians}）[348-357]
\end{itemize}

\begin{itemize}
\item \OriginalTerm{欧多克修}{Eudoxius}（\OriginalTerm{亚略派}{Arians}）[357-359]
\end{itemize}

\begin{itemize}
\item \OriginalTerm{梅里提}{Meletius}（\OriginalTerm{正统派}{Orthodox}）[360 (died) 381]
\end{itemize}

\begin{itemize}
\item \OriginalTerm{弗拉维雅诺}{Flavianus}（\OriginalTerm{正统派}{Orthodox}）[381-404]
\end{itemize}

\begin{itemize}
\item \OriginalTerm{波菲里}{Porphyrius}（\OriginalTerm{正统派}{Orthodox}）[404-413]
\end{itemize}

\begin{itemize}
\item \OriginalTerm{亚历山大}{Alexander}（\OriginalTerm{正统派}{Orthodox}）[413-419]
\end{itemize}

\begin{itemize}
\item \OriginalTerm{狄奥多图斯}{Theodotus}（\OriginalTerm{正统派}{Orthodox}）[419-429]
\end{itemize}

\begin{itemize}
\item \OriginalTerm{保利诺三世}{Paulinus III.}（\OriginalTerm{厄斯塔提派}{Eustathians}）[362-388]
\end{itemize}

\begin{itemize}
\item \OriginalTerm{厄瓦格里}{Evagrius}（\OriginalTerm{厄斯塔提派}{Eustathians}）[388-]
\end{itemize}

\begin{itemize}
\item \OriginalTerm{亚历山大}{Alexandria}：—
\end{itemize}

\begin{itemize}
\item \OriginalTerm{伯多禄}{Peter} [301-312]
\end{itemize}

\begin{itemize}
\item \OriginalTerm{阿基拉斯}{Achillas} [312-313]
\end{itemize}

\begin{itemize}
\item \OriginalTerm{亚历山大}{Alexander} [313-326]
\end{itemize}

\begin{itemize}
\item \OriginalTerm{亚大纳削}{Athanasius} [326-341]
\end{itemize}

\begin{itemize}
\item \OriginalTerm{额我略}{Gregory}（\OriginalTerm{亚略派}{Arian}）[341-347]
\end{itemize}

\begin{itemize}
\item \OriginalTerm{亚大纳削}{Athanasius} [347-356]
\end{itemize}

\begin{itemize}
\item \OriginalTerm{乔治}{George}（\OriginalTerm{异端者}{heretic}）[356-362]
\end{itemize}

\begin{itemize}
\item \OriginalTerm{亚大纳削}{Athanasius} [363-373]
\end{itemize}

\begin{itemize}
\item \OriginalTerm{伯多禄}{Peter}（\OriginalTerm{亚大纳削}{Athanasius}的门徒）[373-373]
\end{itemize}

\begin{itemize}
\item \OriginalTerm{路济}{Lucius}（\OriginalTerm{亚略派}{Arian}）[373-377]
\end{itemize}

\begin{itemize}
\item \OriginalTerm{伯多禄}{Peter} [377-378]
\end{itemize}

\begin{itemize}
\item \OriginalTerm{弟茂德}{Timothy} [378-385]
\end{itemize}

\begin{itemize}
\item \OriginalTerm{德奥斐禄}{Theophilus} [385-412]
\end{itemize}

\begin{itemize}
\item \OriginalTerm{济利禄}{Cyril} [412-444]
\end{itemize}

\begin{itemize}
\item \OriginalTerm{耶路撒冷}{Jerusalem}：—
\end{itemize}

\begin{itemize}
\item \OriginalTerm{玛加略}{Macarius} [324-336]
\end{itemize}

\begin{itemize}
\item \OriginalTerm{马克西默}{Maximus} [336-350]
\end{itemize}

\begin{itemize}
\item \OriginalTerm{济利禄}{Cyril} [350-388]
\end{itemize}

\begin{itemize}
\item \OriginalTerm{若望}{John} [388-416]
\end{itemize}

\begin{itemize}
\item \OriginalTerm{普莱伊略}{Praylius} [416-425]
\end{itemize}

\begin{itemize}
\item \OriginalTerm{犹文纳里}{Juvenalius} [425-458]
\end{itemize}

\begin{itemize}
\item \OriginalTerm{君士坦丁堡}{Constantinopole}:
\end{itemize}

\begin{itemize}
\item \OriginalTerm{亚历山大}{Alexander} [326-340]
\end{itemize}

\begin{itemize}
\item \OriginalTerm{尼科米底亚的欧瑟伯}{Eusebius of Nicomedia}（\OriginalTerm{亚略派}{Arian}）[340-342]
\end{itemize}

\begin{itemize}
\item \OriginalTerm{宣信者保禄}{Paul the Confessor} [342-342]
\end{itemize}

\begin{itemize}
\item \OriginalTerm{马其顿尼}{Macedonius}圣神的仇敌[342-360]
\end{itemize}

\begin{itemize}
\item 不虔的\OriginalTerm{欧多克修}{Eudoxius} [360-370]
\end{itemize}

\begin{itemize}
\item 色雷斯地\OriginalTerm{贝罗亚}{Berœa}的\OriginalTerm{德摩菲利}{Demophilus}（\OriginalTerm{异端者}{heretic}）[370-]
\end{itemize}

\begin{itemize}
\item \OriginalTerm{纳齐盎的额我略}{Gregory of Nazianzus} [380-381]
\end{itemize}

\begin{itemize}
\item \OriginalTerm{纳克塔里}{Nectarius} [381-398]
\end{itemize}

\begin{itemize}
\item \OriginalTerm{金口若望}{John Chrysostom} [398-404]
\end{itemize}

\begin{itemize}
\item \OriginalTerm{阿尔萨修}{Arsacius} [404-406]
\end{itemize}

\begin{itemize}
\item \OriginalTerm{亚提古}{Atticus} [406-426]
\end{itemize}

\begin{itemize}
\item \OriginalTerm{西西尼乌}{Sissinnius} [426-428]
\end{itemize}

\SourceRecord{Translated by Blomfield Jackson.；Nicene and Post-Nicene Fathers, Second Series；Vol. 3.；Edited by Philip Schaff and Henry Wace.；Buffalo, NY: Christian Literature Publishing Co.,；1892.；<http://www.newadvent.org/fathers/27025.htm>.}
